% Supplementary material to the paper (Online Resource 1): the algebraic steps of every derivation
% the paper marks derived, in Springer Nature's class like the article. Standalone: pdflatex Universe24_Supplement.tex with TEXINPUTS=sn:, three times.
\pdfoutput=1
\documentclass[pdflatex,sn-mathphys-num]{sn-jnl}
\usepackage{amsmath,amssymb}
\usepackage{booktabs,multicol,enumitem,array,microtype,longtable}
\usepackage{orcidlink}
\setlength{\columnsep}{14pt}\setlist{nosep}
\setlength{\emergencystretch}{1.5em}
\renewcommand{\labelitemi}{$\bullet$}
\hypersetup{colorlinks=true,allcolors=black}
\newcommand{\num}{\mathrm{num}}
\newcommand{\den}{\mathrm{den}}
\newcommand{\anow}{a_{\mathrm{now}}}
\newcommand{\abefore}{a_{\mathrm{before}}}
\newcommand{\dmod}{\ \mathrm{mod}\ }
\newcommand{\bM}{\mathbf M}
\newcommand{\bD}{\mathbf D}
\newcommand{\bK}{\mathbf K}
\newcommand{\claimmark}[1]{\textup{\textsc{#1}}}
\newcommand{\fence}[1]{\textup{\textsc{#1}}}
\newcommand{\anext}{a_{\mathrm{next}}}
\newcommand{\ddiv}{\ \mathrm{div}\ }
\newcommand{\Wc}{W_c}
\newcommand{\nom}[2]{\item[{#1}] #2}
\newcolumntype{L}[1]{>{\raggedright\arraybackslash}p{#1}}
\theoremstyle{thmstyleone}
\newtheorem{derivation}{Derivation}
\renewcommand{\thederivation}{S.\arabic{derivation}}
\raggedbottom
\begin{document}
\title{Supplementary material (Online Resource 1): the derivations\\ \large to the article ``Universe24: one simple rule, the formulas it shares with nature and adds''}
\author*{\fnm{Alon} \sur{Gonen}\,\orcidlink{0009-0008-2701-4599}}\email{alon@defounder.ai}
\affil{\orgname{Independent researcher}, \orgaddress{\city{Haifa}, \country{Israel}}}
\maketitle

\noindent Journal: Foundations of Physics.


\noindent This supplement gives the algebraic steps of every result the paper marks \textsc{derived} or \textsc{computed} without a proof in its text, Derivations S.1 to S.56, numbered in their own order, each naming the paper's section it serves; the glossary of the derivations' words stands first, then the paper's claims list, every claim with its status, its fence, what it rests on and its derivation, then the Nomenclature. The derivations keep their own words, mapped to the paper's in the glossary below; the derivations and the passages the paper's statuses cite beyond S.1 to S.56 stand in Section C at the end, C.1 to C.9, cited as (C.$n$). Theorems 1 to 3 of the paper are proved in S.45 and Theorem 4 in S.5 and S.6; the one-line theorems carry their proof in their sentence. Each derivation, S.48 to S.50 excepted, names its inputs from the paper, its steps, and, where a number was computed, the computation and its check. The notation is the paper's: Rule3's coefficients at a Node of clock pace $p_0$ and Link paces $p_a$ are $R_a = 2\,\num\,p_a^2$ on the two arrivals of the axis $a$ where the axis's two Links share one pace (the six Port coefficients $R_{ij}$ of Eq.~(1) of the paper in general), $S = 12\,\den\,\Gamma^2 - 12(\den - \num)\,p_0^2 - 4\,\num\sum_a p_a^2$ on the Node's own level, and the wall $w = 6\,\den\,\Gamma^2$; the content read at the Node is $x = c / \Gamma$, also written $U$; $g = (\den - \num) / \den = 1 - \cos\omega_0$. The lattice's own order: every derivation below is exact algebra of the line where it says so; where a continuum method is used on the lattice (the far field $3 / (4\pi r)$ of a holder's rest, S.25; the retarded solution of the wave equation, S.40; the shell theorem, S.29; the ray integrals along the straight path, S.20 and S.30; the radial mesh, S.52), it holds at long wavelength and large $r$ with the lattice's correction of relative order $(\mathrm{Link} / r)^2$ or $k^2 / 12$ along an axis, the order the law's lines name; and the theorems of mathematics cited by name (the discrete Fourier transform's resolution $2\pi / n_w$ in S.8, Watson's constant $G_0(0) = 0.252731$ \cite{watson} in S.36, the post-Newtonian periapsis formula $6\pi U_1 / (1 - e^2)$ in S.19 \cite{will}) are standard results used without proof.

\paragraph{The use of a large language model.} A large language model was used as a tool under the author's direction, as the article's Statements and Declarations state.

\section*{The words: the derivations' and the paper's}

\noindent The derivations keep their own words; the paper's word stands beside each of them here.
\begin{multicols}{2}\small\begin{description}[style=unboxed,leftmargin=0pt,itemsep=0pt,parsep=0pt,font=\normalfont\itshape]
\item[family] a field, one kind of physical value with its pair and its shape (the paper's Sections 2 and 4.1)
\item[record] a configuration, a field's levels over the lattice
\item[holder] a potential field, a field that holds a value for others to read: the sign holder is the charge potential, the massless holder gravity's potential field, the binding holder the binding potential, a gapped holder a potential field with a gap
\item[content] the potential level, what a holder holds at a Node, $x = c / \Gamma$, also written $U$
\item[share] the quantum share, a configuration's count in a region
\item[credit, the credit rule] the counting rule, which credits one quantum to one NodeDetector at a click
\item[books] the count register of the click layer, one integer per configuration
\item[draw] the NodeDetector's random draw
\item[window] the counting window
\item[lay, laid] the initial data, the configurations' initial writing
\item[root] the source Node of a configuration's initial data
\item[blind, the blind row] the prediction written before the computation
\item[front] a region's boundary, the Ports through which its inflow is read; light's front and the erasing front keep their names
\item[tension] the stress, the momentum flux of a wave
\item[body] a bound state, bound or free, a configuration standing in a potential well or falling free
\item[the taker] the configuration whose transition the click reads
\item[the wall] Rule3's positive integer divisor $w = 6\,\den\,\Gamma^2$, $6\Gamma^3$ at $\den = \Gamma$; it bounds the remainder, not the field level
\item[the hold] the potential's write, the derivations' name of the act
\item[the event lattice] the lattice
\item[packet] a wave packet
\item[the declared universe] the declaration of a world: its fields with their pairs and shapes, its holders, its NodeDetectors and its initial data
\end{description}\end{multicols}

\section*{The claims of the paper}

\noindent Every claim of the paper, with its status (the key of Section 1.3 of the paper), its fence, what it rests on and its derivation; the paper carries the same status in the parenthesis that closes each claim's sentence. A calibration is a number taken from nature; experiment, a number the model does not fix.

\begin{footnotesize}\setlength{\tabcolsep}{3pt}\begin{longtable}{L{4.3cm}L{2.4cm}L{1.2cm}L{2.0cm}L{2.0cm}}
\caption{\label{tab:claims}The claims of the paper, each with its status, its fence, what it rests on and its derivation or section}\\
\toprule
The claim & Status & Fence & Rests on & Where\\
\midrule
\endfirsthead
\toprule
The claim & Status & Fence & Rests on & Where\\
\midrule
\endhead
\bottomrule
\endfoot
Universe24 is the meeting of the past with the future; a click is that meeting, the one measurement; only clicks are seen & assumption, the paper's first statement & clicks & the model's statement & Section 1.1\\
Rule3's step is reversible exactly, integer for integer; the whole evolution goes back under the condition on the reads & theorem for the step and where the writes change no level; the write an assumption & lattice & Rule3 & Section 3.1; S.45\\
The two slits' run forward and back returns every level over $40$ intervals, before the counting window opens at $45$ & computed, a run & lattice & Rule3's integer line on the declared universe & Section 3.1; S.23\\
Rule3 is covariant under the cube's group of $48$; $D$ a scalar, $F$ a vector, $W$ a scalar under the $48$, odd under the time reversal and the plane's reflection; no boost in the group; The rule's acts form no product of levels, on the infinite lattice, for one field at fixed coefficients; Three exact massive rotations, $[1, 2]$, $[0, 1]$, $[-1, 2]$ & theorem & lattice & Rule3 & Sections 2.1, 3.2 and 3.6; S.4, S.45\\
A conserved form at static paces with the weights $1 / p_i^2$, its current with the Link's factor, and the momentum flux as the wave's stress & theorem at static paces, weighted; derived for the momentum flux & lattice & Rule3 & Section 3.3; S.5, S.7, S.32, S.45\\
The count is the configuration's quantum share; the rational total is non-negative under the two conditions, the floors' total within the number of Nodes of it & theorem under the two conditions & lattice & Rule3 & Section 3.4; S.5, S.6\\
The Wronskian is conserved exactly in the rational line under the condition on the paces; the integer step adds the remainders' term; antimatter is the opposite sense & theorem for the conservation; derived for the opposite sense, the sign declared & lattice; clicks for antimatter & Rule3, the rotation & Section 3.5; S.9, S.33\\
An event spreads by Rule3 at one Link per interval, the causal bound; light's speed $1 / \sqrt 3$ Link per interval, its dispersion, none along the diagonal; the two fronts, the octahedron and the inscribed sphere; the lattice Klein-Gordon equation with Schr\"odinger's slow limit; $E = m^* c_m^2 = \omega_0$ & the bound postulate 4, an assumption, the step's reach (the line reads the six Ports and nothing farther); derived for the speed, the fronts and the rest, $E = m^* c_m^2$ by construction and no test & lattice; clicks for $E$ & Rule3 for the pairs & Sections 2.2 and 2.5.1; S.1, S.18, S.21, S.53\\
The paces compose, $p_0 = \Gamma(1 - 1 / \Gamma)^c$; the Link clocked twice by its reading Node's clock, $hN = 1$; the odd Link part, gravity's holder's three odd lines sourced by the count's flux and read at the Port as the phase $\theta_{ij} = 4V_{ij}\omega_0 / c^2 = 12\,V_{ij}\omega_0$ & assumption, fitted; the composition derived from it; assumption for a moving source, the coefficient derived under it & lattice & fixed by the bending's measurement, the perihelion choosing the reading and testing the choice & Section 2.3; S.34; S.56\\
Frame dragging's form $\delta\omega / \omega_0 = +4\,\mathbf V\cdot\mathbf v / c^2 = 12\,\mathbf V\cdot\mathbf v$ for a plane configuration of one rotation sense, the opposite sense opposite, a real line undragged; a body's drag its plane configurations' signed share of its inertia, about one half for nature's matter, a finding against nature; $\alpha_1 = \alpha_2 = 0$ and $\eta_N = 0$ for one sense alone ($\alpha_2$ and $\eta_N$ computed under the relabelling of S.56 (g) and not derived, $g_{00}$'s velocity terms not fixed, S.35), $\alpha_1 = 8f - 8$ for a body of share $f$, $-4$ at one half & derived under the one assumption, to first order in $\mathbf V$, slow matter; the body's share and its $\alpha_1$ under the proton hypothesis of the paper's Section 4.1 as well; the comparisons supporting findings & clicks & the one assumption for a moving source; S.35 (c), S.48, S.56 & Sections 5.1 and 8.1; S.56 (e), (f), (g)\\
Light undragged by a rotating mass; a fast configuration's coupling falling as $1 / \gamma_L$ & derived under the one assumption; untested & lattice & the one assumption for a moving source; S.56 & Sections 4.1 and 5.2; S.56 (h)\\
The exponential metric, $\gamma = \beta = 1$ in $U$; the bending $4U_b$, Shapiro's delay, the perihelion, the redshift's factor, the shadow's $2e\,m$; $\gamma_K$, $\beta_K$ and $\alpha_K$ from the pair alone, closing with the gap & derived under the one assumption; the bending's measurement its calibration, the perihelion the test of the end's choice, the rest its consequences, the shadow bounded by measurement and not confirmed, conditional on the interior's fate & clicks for $\gamma$, $\beta$ and the redshift's factor, lattice for the metric's form at the computed pair & $hN = 1$, Eq.~(2) & Sections 5.1 and 5.2; S.15, S.19, S.20, S.30, S.31, S.51\\
The guard is the Courant-Friedrichs-Lewy condition \cite{cfl}; under the composed paces no Link closes in the rational paces; a frozen Node is dark & theorem for the edge; derived for the closing; assumption for the frozen Node's energy & lattice & Rule3's coefficients & Section 2.3; S.3, S.39\\
A click is the meeting of two quanta, decided by the counting rule's one draw, written at one Node; one click per quantum (antibunching); no signal through the amplitudes nor through a bound state's clicks, a Nodes-only region's count coupled to the others' inflows by the credit's rounding by at most one quantum per window, $1 / N$ relative; the write is the NodeDetector's act; the direction of time belongs to the clicks & assumption for the draw, the write and the counting rule; derived under the counting rule's one draw, the NodeDetector's declaration, for exclusivity; derived for antibunching and the one form; theorem for the step; the arrow the NodeDetector's declaration & clicks; lattice for the write and the arrow & the meeting, the NodeDetector's declaration & Section 2.5, the steps paragraph; S.46\\
Born's rule is the NodeDetector's declaration; the uncertainty relation is its declarations, the region, the counting window and the random draw; the draw's frequencies within $2^{-b}$, $b$ the declared bound's exponent & declaration; derived under it & clicks & the counting rule & Section 2.5; S.8, S.46\\
Only a click moves a quantum between fields; between clicks every count is constant & derived for the count constant between clicks (S.46); the exchange the NodeDetector's declaration, an assumption; for nature a hypothesis under its own name & clicks & Rule3's lines & Section 2.5; S.46\\
Light is the charge potential's own configuration; a steady rotation writes no light, a breathing bound state at its beat; a uniformly moving one none in the long-wavelength band & derived & lattice (the Wronskian and the source); clicks (the far NodeDetector's rate) & the write & Sections 2.5 and 4.3; S.40\\
A potential field's reach from its gap, Yukawa's form; Poisson's rest $3s / (4\pi r)$; gravity enters as a slowed clock & derived & lattice & Rule3 for a gapped pair & Section 4.1; S.11, S.37; Online Resource 1, B.2\\
The vacuum level $60$ at $\Gamma = 6{,}000$, above the lattice's floor; no click bounds it above; gravity's events at light's speed in every level & declaration above the floor; derived for the speeds & lattice for the level and its floors; clicks for the speeds & every field reads the potential level & Section 4.1; S.12, S.48; Online Resource 1, B.11\\
The binding potential's pair and range and gravity's divisor, chosen where nature's scale forced them & the pair and the divisor an assumption, fitted; the range derived from the pair (S.11) & lattice & the declared universe & Section 4.1; S.11\\
The bound state's fixed point and functional; the stable bound state along the dilation to first order in the well; the two branches & theorem for the equation; derived for the stability, the integer version open; the branches a lattice reading, named & lattice & the functional & Section 4.2; S.13, S.45; Online Resource 1, B.3\\
No standing bound state is claimed from a run; the compact branch's budget $(\mathrm{Link} / R)^2$, $\sqrt n / A$ & a lattice reading, named & lattice & the lattice's readings & Online Resource 1, B.3; C.3\\
Bohr's form with the nucleus's angle alone; the self-read gives the law's form at large $Z$ only; hydrogen's lines under line (1); the criterion, what nature splits is a bound state & derived in the slow limit, for the levels' differences read by the absorber's resonance; the emitted spectrum of the one-Node write not derived; line (1) of S.52 an assumption, fitted; line (2) an assumption; the criterion the model's decision, an assumption & clicks & the charge potential & Sections 4.2 and 5.1; S.43, S.52; Online Resource 1, B.3\\
The adiabatic invariant and Planck's relation; the drift under a moving well; the clock's rate $1 - f(\omega_0)U$; the resonance at $2\omega_b$ under the pace read alone (none under the rotation, S.41) and Rabi's rate, two derivations; a bound state counts photons at Millikan's slope $T$; no recoil in the law's write & derived (the clock's rate under the one assumption; the invariant computed) & clicks; the drift lattice & Rule3 at the bound state's Nodes & Section 4.3; S.14, S.15, S.16, S.41; Online Resource 1, B.4\\
A cluster: the three levels at a Node as the nesting, three kinds of binding, the click the only thing between clusters; dark matter in kind & definition, derived; the source a hypothesis & lattice, clicks & every bound state is bound & Section 4.4; Online Resource 1, B.5\\
Friedmann's dust equations and Hubble's law at $z \ll 1$; the deceleration $q_0 = \Omega_m / 2 > 0$ & derived, two derivations, under the declared homologous start & clicks & Kepler's $G_K$ & Section 4.4; S.29\\
The two potentials and the Kepler map; the moving clock at $c_m$ with its fourth order; de Broglie; the fall and the equivalence between fields, $\eta = 32.3\,\omega_e^2$; the lens with its chromatic term; Coulomb over $\gamma_L$ (the Lorentz factor) and magnetism & derived; for the equivalence between fields derived under the proton's and the neutron's gaps read as two fields, a hypothesis under its own name, a free-packet upper estimate; for Coulomb and magnetism theorem for the conservation, derived for the force's form under the rotation, not under the plain read & clicks; the index lattice; for Coulomb and magnetism lattice for the identity, clicks for the force & the band, the paces, the rotation & Section 5.1; S.18, S.21, S.22, S.28, S.33, S.44\\
$G$ as a reading, Eq.~(11); $\alpha_{\mathrm{law}} = (3\sqrt 3 / 8\pi)\,k / (\Gamma E_s)$; charge universality by construction; the two forces' ratio $(m_p / m_e)^2$; the units; the Link's and the gap's bounds & derived; $k$ declared; charge universality and the forces' ratio in the small-gap limit, by construction, no test; the bounds comparisons & clicks & the rotation, Poisson's rest, the Large High Altitude Air Shower Observatory's (LHAASO), the Major Atmospheric Gamma Imaging Cherenkov telescopes' (MAGIC) & Section 5.1; S.24, S.25, S.26, S.27, S.42\\
The gravity of light at $1 \times E / c^2$ & derived under the static kernel, an assumption named, the packet's field open; untested & clicks & the static kernel & Section 5.2; S.39; Online Resource 1, B.6\\
The two-path check is $0.989$ for a strontium clock at one metre for one second & derived under the one assumption, S.31's scaling; untested & clicks & the one assumption, S.31's scaling of the beat & Online Resource 1, B.6; S.39 (1)\\
The two slits: the two slits' $284.7$ by the law's line, the design's blind row $273$, the Huygens sum, light's line propagated from the two gaps on the design's lattice; the run $N = 285$ at both seeds within the integer step's walk, the seed-$24$ clicks $28$ of $275$ from the prediction's counts ($275$ the blind row's sum by its per-Node rounding, $273$ its $N$), a comparison and not a statistical test; the inflow's peak at $69.0$ against $62.0$, the flight time $62.2$ & computed; the run a check & clicks; $674$ photons at the light's own rotation, lattice & Rule3 on the declared universe & Section 6; S.23; Online Resource 1, B.7\\
Bell's $E = \cos 2(a - b)$, $S = 478 / 169$, Tsirelson's $2\sqrt{1 + \rho^2}$, Greenberger-Horne-Zeilinger's (GHZ) $M = -4$; every local counting rule gives at most $2$; no signalling & the click's one form an assumption, fitted; the square declared; derived under them; computed for the numbers; theorem for the local bound & clicks & the counting rule, the meeting's square & Section 7; S.10, S.38, S.47; Online Resource 1, B.8\\
The misses, thirteen by name: ten findings against nature, frame dragging for bodies, $q_0 > 0$, the nuclear binding and its size law, no Fermi constant $G_F$, the $Z$ and parity (one root), the reduced mass and the isotope shifts (the record's own level excluded by line (1), S.43 (c), S.52), the mass ratios and the cluster's defect, the fine structure, helium's correlation energy and positronium's binding, and the $\times$ polarisation, derived under the static read, which the odd write keeps; three places with no line, the neutron's lifetime, Pauli's exclusion and the loops ($g - 2$ and the Lamb shift) & derived, the comparisons supporting findings; no line for the three & clicks; lattice for helium and positronium & Rule3 and the declared entries & Section 8.1; S.35, S.50, S.52; Online Resource 1, B.9\\
Neither layer keeps a total energy & theorem for the step's conservations; assumption for the clicks' & lattice, clicks & Rule3, the counting rule & Section 2.4; S.49\\
The universe's physical coefficients $E_g$, $k$ and the gaps are each read from one measurement of nature; the binding pair, $T$ and the integer bound are declared and the Link unit read from one measurement; no number comes from the lattice itself & calibration, experiment for $E_g$, $k$, the gap and the Link's length; hypothesis for any such number & clicks & nature's measurements & Section 8.2; S.48; Online Resource 1, B.10\\
The declarations: the entries, the NodeDetectors, the rates, the conversion's table, the initial data, the receding face, the one draw, the erasing front and the empty wave beyond it & assumption for every declaration & lattice, clicks & Online Resource 1, B.10's list, B.1 for the empty wave & Section 8.2; Online Resource 1, B.1, B.10\\
The hypotheses under their own names: the conversion click, the dimension-$3$ field with its confinement in the model's words, the nuclear potentials, the parts as particles, the masses as fields, only a click moves a quantum in nature (the single atom's lifetime its test) & hypothesis; no number of the paper rests on the last & clicks & stated under their names & Section 9; Online Resource 1, B.12, C.5, C.6, C.7\\
The click's algebra beyond the law: the count's change and the booking identity; the draw of whether and which future continues; the momentum's chain; under the twist proposition the wave vector conserved, the recoil $\propto 1 / M$, Compton with the factor $\omega_0 / \tan\omega_0$, nature's at $1$ & derived for the count, the booking and the chain; assumption, the counting rule, for the draw; hypothesis, fitted, for the twist, derived under it for the rest & lattice; clicks for the booking; Compton's factor lattice at the computed pair, nature's by construction, no test & the law's definitions; no run & Section 2.5.2; S.55\\
\end{longtable}\end{footnotesize}

\section*{Nomenclature}

Every letter of the paper and this supplement, with each of its senses; a letter that serves twice is named at each first use as well; the glossary of the derivations' words, with the paper's beside them, stands above.
\begin{multicols}{2}\small\begin{description}[style=unboxed,leftmargin=0pt,itemsep=0pt,parsep=0pt,font=\normalfont]
\nom{$A$}{an amplitude, a mass number and Bell's side; $A_0$, $A_a$ the vector potential on the time step and the Link (S.33, S.41)}
\nom{$a$}{a Node's level ($a_t$, $a_{\mathrm{now}}$); $\mathrm{arr}_j$ the level of the neighbour $j$ in Rule3's line, the six Ports summed (Eq.~(1)); a radius (the ball's $a(t)$, Bohr's $a_0$), Bell's setting, the pressure's coefficient and the axis index where so named; the fall's acceleration (Section~5.1, S.18) and the semi-major axis}
\nom{$B$}{a standing record's peak amplitude, Bell's side and the bound mode's binding where so named}
\nom{$b$}{the impact parameter, the leaving level in the booking identity, Bell's setting, the hollow's coefficient and the axis index where so named; the declared bound's exponent $b$ (Section~2.5); $b_L$ the light level read; $b = 1{,}241$ the laid packet's peak level (Section~6); $\mathbf b$ the vector of before levels in S.6, with $\mathbf n$ the now levels; $b_t$ a second standing record's level (S.17)}
\nom{$C$}{a body's count in quanta; an amplitude or a coefficient in the supplement where so named; $C_i$ the counts at the Nodes; the third Greenberger-Horne-Zeilinger (GHZ) side; $C_1(k)$, $C_0(k)$ the step symbol's coefficients in S.56 (d)}
\nom{$c$, $c_m$}{light's speed, and the content in level units where so written ($c$, $c_{\mathrm{vac}}$); the band's kinetic scale (Section~2.2, used in 5.1); $c_{\mathrm{local}} = cN^2$ the local light speed at the composed paces (Eq.~(9)); $c_K$ the lattice Klein-Gordon equation's coefficient, $c_s$ a family's own light speed ($c_s^2 = \num_s / (3\,\den_s)$), $c_k$ a part's laid weight (Eq.~(12)), the third GHZ setting and the well's coefficient where so named; $c_n$, $c_a$ cosines in S.7 and S.10}
\nom{$D$}{the form of a record}
\nom{$d$}{a radius, the slits' distance and a dimension where so named; $d_i$ a Node's weight, $d_{jm}$ the dipole element, $d = \den - \num$ (S.27); $d_L$ the luminosity distance (S.29); $d_1$, $d_2$ two bodies' distances from a Link (S.17)}
\nom{$E$, $E(a, b)$}{the energy, the conserved form of Eq.~(4), the part's wall in Eq.~(1) and the level weights $E_s$, $E_g$ (gravity's) and $E_r$, the divisor of the potential field $r$ (S.13); $E_{\mathrm{QG},2}$ the quadratic Lorentz-violation scale (S.54; B.9); Bell's correlation, $E_3$ the three-quantum Greenberger-Horne-Zeilinger correlation $E_3 = \cos 2(a + b + c)$ (Section~7.2); $E_0$ the field's amplitude (S.41)}
\nom{$e$}{Euler's number, the share $e_i$, a port's coefficient $e_k$, the unit charge (one quantum's Wronskian $T / 2$, Section~5.1), an orbit's eccentricity, a record's ellipticity and the part label $e$}
\nom{$F$}{the current, the functional, a force, each where so named}
\nom{$f$}{the redshift's factor $f(\omega_0) = 2(1 - \cos\omega_0) / (\omega_0\sin\omega_0)$, $f_v$ the moving clock's factor, the sourcing family index in $w_f q_f$, a generic function, and the subscript of the frozen content $U_f$}
\nom{$G$}{Newton's constant, $G_K$ as the orbits read it, $\mathbf G_r$ the static kernels of the potential fields (S.13; B.3), $G_{\mathrm{clock}} = 1 / (4\pi\Gamma E_g)$ as the clock's potential reads it, before Kepler's factor (Eq.~(11)), and $G_0(r)$ the lattice's Green's function (S.37), and $G_{ab}$ the stress reading; $G_F$ the Fermi constant}
\nom{$g$}{the rest term's coefficient, the metric's components, the ground part $g$ and Land\'e's $g_L$; $g_{ij}$ the sign holder's Link current (S.33); Earth's acceleration in $gh / c^2$ (S.39), Yukawa's coupling $g^2$ (S.37), the transfer generator $g = \Omega_R$ (S.41) and the anomaly $g - 2$}
\nom{$H$}{Hubble's rate ($H_0$); $T_H$ Hawking's temperature}
\nom{$h$, $h_a(i)$}{the ruler of Section~2.3 and Planck's constant where nature's is meant; the Link's stress reading (S.7); a hazard per window, $(1 - h)^n$ (C.9 (c)); the height in $gh / c^2$ (S.39); $h_{00}$, $h_{aa}$, $h_{xy}$ the metric perturbation's components (S.50 (1)), beside $h = p_0 / p_a$ the Link factor}
\nom{$I$}{the adiabatic invariant $A^2\sin\omega / T$, the number of quanta; the integral $I(z)$ of the running exponent (S.13)}
\nom{$i$}{the Node index ($i$, $j$ its neighbour) and the imaginary unit in $e^{i\theta}$; a bound mode's index where paired with $j$}
\nom{$J$}{the joint share $J(o_A, o_B)$ of Bell's pair over the two sides' ports $o_A$, $o_B$, $J_a$ the sign's Node current on the axis $a$, and the atom's Hartree self-term where so named}
\nom{$j$}{the neighbour Node index, a bound mode's index, a spatial index ($g_{0j}$, $V_j$), nature's current density where so named}
\nom{$K$}{Friedmann's constant, the start's kick and a pressure's weight where so named; $K_{jm}$ the two modes' coupling (S.41)}
\nom{$k$}{$k_L$ light's wave number and $\mathbf k_L$ its wave vector (S.55); the wave number and the sign holder's write weight; $k_{\mathrm{in}}$ the wave number a configuration takes on entering a potential level $U$ at fixed rotation (Eq.~(10)); $k_w$ the write's and $k_r$ the read's weight, $k = k_w k_r$ (S.26); the part's label in Bell's sums (Eq.~(12)); $k_B$ Boltzmann's constant; the $1s$ cloud's curvature in atomic units in S.52 (b)}
\nom{$L$ and $\ell$}{$L_{\mathrm{box}}$ the declared box's side in Links (S.55); both a holder's level in levels, $L$ also a length in Links where so named, $L'$ the packet's flight to the screen (Online Resource 1, Section B.7); $\ell_{\mathrm{Link}}$ the Link's length in metres (S.54)}
\nom{$l$}{the ray's path element $dl$ and a level increment in Cauchy's equation $P(U + l) = P(U)P(l)$}
\nom{$M$}{the taker's count at a click (Section 2.5.2, S.55); the mass in $GM$ and Mermin's $M$; $\mathbf M$ Rule3's step matrix, with $\mathbf D$ the weights' diagonal and $\mathbf K = \mathbf D\mathbf M$ (S.56 (a)); $M_1$, $M_2$ the parts' products (S.10)}
\nom{$m$}{the rest term of the lattice Klein-Gordon equation, $m^*$ the inertia, $m = GM / c^2$ in Links in the shadow's radius; the mass, Table 1's free number of the matter field; $m_W$ the W boson's mass, nature's (S.50 (17); B.9); a bound mode's index ($\phi_m$, $\omega_m$) and $m_\ell$ the orbital quantum number (S.43)}
\nom{$N$}{the books' count and the lapse $N = p_0 / \Gamma$, $e^{-U}$ to $O(U / \Gamma)$; a number of Nodes ($N_x \times N_y \times N_z$), intervals or bodies where so named, $N_R$ a region's count}
\nom{$n$}{a count of intervals or quanta where so named; light's index $n = e^{2U}$, $n_m$ slow matter's (Eq.~(10)); $\mathbf n$ a unit direction, $n_a$ its components; $n = 1$ nature's survival exponent (C.9 (c)); a mode's index ($E_n$, $\omega_n$ of the modes)}
\nom{$O$, $o_A$, $o_B$}{$O_h$ the cube's symmetry group of order $48$, $O(3)$ the rotation group of a three-line record, $O(\cdot)$ the order symbol; $o_A$, $o_B$ the two sides' ports in Bell's joint share (Eq.~(12))}
\nom{$P$}{the momentum density, a probability, the guard's edge on a pace and the momentum $P = cN^2 k$ at the composed paces; $\mathbf P$ a signed permutation (S.4), $P(U)$ the pace as a function of the content (S.34)}
\nom{$p_0$, $p_i$, $p_a$}{the paces and $q_{ij}$ the Link's factor, with $(p, q)$ Bell's setting pair; $\mathbf p$ the momentum (Section~5.1), $p$ an orbit's semi-latus rectum (S.19) and nature's pressure (S.39)}
\nom{$Q$, $Q_R$}{the two-level form and $q_{ij}$ the Link's factor; a region's inflow; $Q(t)$ a second moment (S.40)}
\nom{$q$}{$q_{ij}$ the Link's factor, a body's charge $q = \pm Ce$, Bell's setting in $(p, q)$, the momentum transfer of a hard click, the second argument of the two-level form $Q(p, q)$, and $q_0$ the deceleration parameter; $q_f$ the sourcing field $f$'s quantity at the Node in the one write, a form $D_f$, a Wronskian $W_f$ or a stress $T_{aa,f}$ (Eq.~(1))}
\nom{$R$}{the read coefficient of Rule3, a body's radius in Links and the binding range; a NodeDetector's region where it carries no index letter and no number; Rydberg's unit in $E_{\mathrm{tot}} / (Z^2 R)$ (S.43, S.52), a scattering rate where so named}
\nom{$r$}{Rule3's remainder and a radius; the parts' amplitude ratio in Bell's $\rho = 2r / (1 + r^2)$, $r_P$ Pekar's constant, $r_p$ the proton's radius; $r_a$ the taker's books' remainder (S.55)}
\nom{$S$}{Rule3's self coefficient in Eq.~(1), Bell's $S$ in Section~7 and $S_t$ the step map; $S_6$ the six-neighbour sum (S.5), $S_4 = \sum_a n_a^4$ (S.54); a source's total (S.40)}
\nom{$s$}{a source's write per interval over the holder's divisor (the far level $3s / (4\pi r)$, Section~5.1, S.20, S.25), a share, the refinement scale, the family index and a scaling power where so named; $s(x) = 2\sin(x / 2)$ (S.28, S.31, S.34), $s_\parallel$, $s_\perp$ a separation's components (S.39)}
\nom{$T$}{the invariant's unit and $T_{aa}$ the tension; $T_b$ a beat's length (S.17)}
\nom{$t$}{the interval (the time index): $a_t$, $S_t$, the click $(R, t)$; $t_a(i, j)$ the Link's tension read into $q_{ij}$}
\nom{$U$}{the potential, the clock's fraction $c / \Gamma$; $U_b$ its value at the impact parameter $b$ (Eq.~(9)); $U_K = G_K M / (c^2 r)$ the potential the orbits of matter read, Kepler's (Eq.~(8))}
\nom{$\mathcal U$}{$\mathcal U_t = e^{-i\theta_t}$ the time turn (S.41)}
\nom{$u$}{the direction call's working integer, the turned level $u_i$, the separation in units of the width (S.13's integration variable), the expansion variable $1 - \cos k$, the band-edge solution $u_t$ and a mode's projection $u_n$ where so named; the gauge-turned record $u_t$ in S.41 (a)}
\nom{$V$}{$V_r$ a holder's well energy at the dilation $1$, $V_j$ a post-Newtonian vector potential, $V_{ij} = -V_{ji}$ the level of gravity's odd line along a Link, the mean of its two ends' levels, sourced by the count's flux (S.56), the slow limit's potential term and the Link-phase perturbation operator where so named}
\nom{$v$}{a velocity (the group velocity, a source's speed, $\mathbf v$), the level Rule3 would write at the Node in $\Delta(wQ) = w\,(x - v)(x - b)$, the turned level $v_i$ and a mode's projection $v_n$ where so named}
\nom{$W$}{the Wronskian, $\Wc$ and $W_{\mathrm{rec}}$ the credit's units; a holder's declared write weight ($W$, $W_1$, $W_2$ of the nuclear holders) where so named; $W_j$ a post-Newtonian vector potential (S.35)}
\nom{$w$}{the wall $w = 6\,\den\,\Gamma^2$, Rule3's divisor; $w_f$ a sourcing family's weight}
\nom{$X$}{the state space}
\nom{$x$}{a content fraction, a level or a coordinate where so named; $x_{\mathrm{self}}$ the self-read's content fraction (S.36; B.9); $x_t$ the lattice's state at the interval $t$ (Section~2.5); the draw's state in S.46}
\nom{$Y$}{$Y$ the line's read $\sum_b R_b(\mathrm{now}_{i+b} + \mathrm{now}_{i-b}) + S_i\,\mathrm{now}_i$}
\nom{$y$}{a lattice coordinate or row, the before level in the direction call $(x, y, \rho)$, and the second vector of levels in $Q(\mathbf x, \mathbf y)$}
\nom{$Z$}{the nuclear charge number; the $Z$ boson ($M_Z$)}
\nom{$z$}{the redshift and the complex level where so named; the direction call's quotient (Theorem~1); $z = \kappa R$ in S.13}
\nom{$\alpha$}{the fine-structure constant, the redshift's $\alpha_U$, $\hat\alpha$ the bending angle (Eq.~(9)) and the frame's $\alpha_1$, $\alpha_2$; $\alpha_k$ a part's amplitude at one side's port (Eq.~(12)), $\alpha_K$ the post-Newtonian parameter against Kepler's potential, $\alpha_g(s)$ gravity's coupling, $\alpha_{\mathrm{law}}$ the law's coupling between two quanta (Section~5.1, S.26); a phase in S.17 with $\beta$}
\nom{$\beta$}{the post-Newtonian parameter, $\beta_K$ its value against Kepler's potential (Section~5.1), and a speed over light's; $\beta_k$ a part's amplitude at the other side's port (Eq.~(12)); a standing record's phase in S.17}
\nom{$\chi$}{the ratio of the two energy changes under the write (S.26); the post-Newtonian superpotential where so named}
\nom{$\Delta$}{the six-neighbour second difference and, as $1 / m^*$, the band's curvature; a detuning ($\Delta = \omega_L - 2\omega_b$), a spread ($\Delta\omega$, $\Delta y$, $\Delta U$) and the centred difference $\Delta_a$ where so written}
\nom{$\delta$}{the remainders' change $r - r'$, the settings' difference $a - b$ (S.10), a rotation's shift from rest, a small change, the breathing's depth, the detuning per interval, Kronecker's $\delta_{ab}$, each where so named}
\nom{$\epsilon$}{the step's remainder term $(r - r') / w$ and a resonance's modulation depth where so named; $\epsilon_4 = (3\sum_a n_a^4 - 1) / 36$ the band's quartic coefficient (S.54), $\epsilon_{ij}$ the two modes' modulation, $\epsilon_a(i)$ the light's profile (S.41); the mode's frequency over $Z^2 R$ in S.43 (c); the odd real coefficient of a one-part line in S.56 (d)}
\nom{$\eta$}{the E\"otv\"os parameter; $\eta_N$ the Nordtvedt parameter (Section~5.1)}
\nom{$\Gamma$}{the Node clock, the lattice's clock ($\Gamma = 6{,}000$), the clock integer of the pace's base $1 - 1 / \Gamma$; $\Gamma_\theta$ the Port phase's common divisor (S.56 (a))}
\nom{$\gamma$}{the post-Newtonian parameter, $\gamma_K$ against Kepler's potential, Lorentz's factor ($\gamma_L$, the coupling's fall at speed in S.56; $\gamma_s$ a family's own), and $\gamma_{\mathrm{GW}}$, $\gamma_{\mathrm{EM}}$ the gravitational wave's and the electromagnetic wave's Shapiro-delay parameters}
\nom{$\hbar$}{Planck's reduced constant, nature's; $T$ stands for it in the lattice's units}
\nom{$\kappa$}{a holder's inverse reach; $\kappa_b$ the bound states' own inverse reach at $\omega_b$ (S.17)}
\nom{$\Lambda$}{the fringe's period in rows (S.8, S.46); the cosmological constant in the $\Lambda$CDM model (cold dark matter with a cosmological constant)}
\nom{$\lambda$}{the wavelength and a Lagrange multiplier; the ratio $U / U_K$ (S.19), the tail's level ratio (S.11), a gauge parameter (S.35); a root of the step's symbol (S.56 (d)); $\bar\lambda_e$ the electron's reduced Compton wavelength (Section~5.1)}
\nom{$\mu$}{a growth rate and $GM / c^2$ where so written; the reduced inertia of a pinned pair where so named}
\nom{$\nabla$}{the gradient; $\nabla^2$ the Laplacian}
\nom{$\nu$}{the pair's ratio $\num / \den = \cos\omega_0$; the neutrino in the pair $(\nu, e)$}
\nom{$\Omega$}{a transition's rotation and the moving centre's rotation $\Omega(k)$; $\Omega_R$ Rabi's rate, $\Omega_w$ the weak drive's, $\Omega_m$ and $\Omega_\Lambda$ the density parameters}
\nom{$\omega$}{a record's rotation per interval, the frequency: $\omega(k)$ the band, $\omega_0$ the rest rotation (the gap), $\omega_b$ a bound body's, $\omega_{\mathrm{part}}$ a laid part's own (Section~2.5), $\omega_s$ the family $s$'s rest rotation (its mass), $\omega_e$ the electron's and, in a transition, the excited part's against the ground part's $\omega_g$, $\omega_L$ the light's, $\omega_h$ the holder's own events', $\omega_{\mathrm{zp}}$ the nucleus's zero-point, $\omega_i$, $\omega_j$, $\omega_n$ the bound modes'; $\omega_n$ also the neutron's gap, $\omega_p$ the proton's}
\nom{$\Phi$}{Newton's potential $\Phi = -c^2 U$; the plaquette flux of the orbital moment; $\Phi_t$ the time turn's accumulated phase}
\nom{$\phi$}{the body's normalised record or mode ($\phi_i$, $\phi_m$, $\phi_\lambda$); a phase in $\cos(kx + \phi)$; a retarded potential in S.39 (2)}
\nom{$\Pi$}{$\boldsymbol\Pi = E\mathbf v$ the covariant momentum variable (S.51)}
\nom{$\pi$}{the number $\pi$}
\nom{$\psi$}{the ray's angle to the radius, $nr\sin\psi = b$; the slowly varying envelope of the slow limit}
\nom{$\rho$}{Bell's $2r / (1 + r^2)$ and the dust density; the direction call's carried remainder (Theorem~1); $\rho_i = w\,W_i$ the continuity equation's density (S.33); $\phi^2$ of unit mass (S.13)}
\nom{$\sigma$}{a holder's source per interval, a sign and a spread where so named; a cube map (S.4)}
\nom{$\tau$}{the lifetime, $\tau_{\mathrm{hold}}$ the hold of the two paths, $\tau_D$ the dark periods' declared mean, $\tau_{\mathrm{int}}$ the interval in seconds}
\nom{$\Theta$}{$\Theta_t$ the midpoint phase of the adiabatic ansatz (S.14)}
\nom{$\theta$}{the sign holder's turn angle ($\theta_t$, $\theta(r)$, $\theta_0$ the nucleus's), $\theta_L$ the light's angle on the Link, $\theta_a$ the Link phases, $\theta_{ij} = -\theta_{ji}$ the Port's phase from the odd line's level along the Link (S.56), the pairs' common phase and nature's GHZ angles where so named; $\theta_W$ the weak mixing angle (S.50); Compton's scattering angle in S.55}
\nom{$\vartheta$}{the modulation's phase in $\cos(\omega_L t + \vartheta)$ (S.16)}
\nom{$\xi$}{the post-Newtonian (Whitehead) parameter}
\nom{$\zeta$}{the detuned slow amplitude (S.16); $\zeta_1$ a post-Newtonian parameter}
\end{description}\end{multicols}

\section*{The sections moved from the main text}\label{sec:moved}

\noindent The paragraphs of this section belong to the paper's Sections 1.2, 2.5.1, 4.1, 4.2, 4.3, 4.4, 5.2, 6.1, 6.2, 7.1, 7.2, 8.1, 8.2, 8.3 and 9 and stand here whole; their pointers to the paper are written as the paper's Section, Eq., Fig. and Theorem numbers.

\subsection*{B.0 The precursors (the paper's Section 1.2)}

The world has been read as a cellular automaton \cite{wolfram}\cite{fredkin}, as unitary lattice automata \cite{arrighi}, in event-based simulations with counting units that click \cite{deraedt}, the closest prior work to the click, and as a deterministic foundation beneath quantum theory \cite{thooft}. The time-symmetric accounts of measurement are Wheeler and Feynman's absorber theory \cite{wheeler}, Cramer's transaction \cite{cramer}, the two-state vector formalism \cite{vaidman} and Wharton and Argaman's locally mediated reformulations \cite{wharton}; the click as a meeting is their lattice form, both directions one integer line exact to the integer, and the choice among the meetings the only non-local element. The exponential metric is Yilmaz's scalar-field theory \cite{yilmaz} and the position-dependent index on flat space \cite{dicke}\cite{puthoff}; the paper adds the route to that metric from Rule3 under the one assumption and the finite-gap factors, and its photon sphere at $r = 2m$ and critical impact parameter $2e\,m$ \cite[Section 7]{boonserm} are the shadow's form here. The Link phases have their precursor in the lattice gauge field \cite{wilson}. An integer random walk has reproduced Schr\"odinger's statistics \cite{sciarretta}; here the lattice is deterministic and the one draw is the click's \cite{bohm}\cite{philippidis}. The derivations of Born's rule \cite{gleason}\cite{deutsch}\cite{zurek} the paper declines: Born's rule is the NodeDetector's declaration (the paper's Section 2.5). Meyer's one-dimensional quantum lattice gas \cite{meyer1997} has the dispersion relation $\cos\omega = \cos k\cos\theta\cos\rho + \sin\theta\sin\rho$ with $\tan\theta$ in the role of the mass; along one axis the band of the paper's Eq. (3) is that relation at $\theta - \rho = \omega_0$ and $\tan\theta\tan\rho = 2$, so the inertia $3\tan\omega_0$ is the lattice's and not this paper's; what the line adds is the three axes' equal share, which fixes both coefficients from the one pair and gives the band's $k^4$ anisotropy, and the same gap's factors under the potential (the paper's Section 5.2).

\subsection*{B.1 How an event spreads (the paper's Section 2.5.1)}

On the lattice there are events and amplitudes. An event is a level split by Rule3 to the six neighbours. An amplitude is the level Rule3 produces at a Node, the sum of the split events there. The configuration, its levels over the lattice, is the amplitudes. In Universe24 there are clicks, the NodeDetectors' reports, and nothing else.

In the two slits the initial data's events pass both gaps, every path of the future at once; the fringes are where split events from the two gaps meet at the screen with the same sign. The past is made of the same split events: read backward from the realised click by the same line, Rule3 at $-1$, they form the past's tree, which also passes both gaps back to the source Node. What stays on the lattice, the amplitudes no click realised, is not in Universe24. After a click they are a cancelled future at three levels, the counting rule's, the NodeDetectors' and the lattice's: in the count register the configuration's count is at $0$ and no NodeDetector credits it again. Where they reach a NodeDetector backed by a declared face, their quantum share leaves the lattice there. On the lattice the absorbed configuration is erased from the click's Node outward at the causal bound, one Link per interval, inside the click's causal cone ($\sqrt 3$ times light's speed), the ball of radius $d$ holding $(2d + 1)(2d^2 + 2d + 3) / 3$ Nodes of zeros, never faster. Beyond the front its levels remain, an empty wave with count $0$, stepped on and never credited, that still sources the potential fields it writes to until the front arrives (an assumption, the NodeDetector's declaration; lattice, clicks; S.46).

\subsection*{B.2 Fields, potential fields and the declared universe (the paper's Section 4.1)}

Gravity enters in four steps, with no metric, no force and no law of motion (derived under the one assumption; lattice; the paper's Section 2.3; S.25, S.34, S.37, S.48, S.56).
\begin{enumerate}
\item Every field of quanta writes its form over the action, $D_i \ddiv T$, into gravity's potential over the divisor $E_g$.
\item Gravity's potential, $[1, 1]$, steps by Rule3 at its pair, resting around a static source at Poisson's solution, $3\,G_0(r)\,s$, $s$ the bound state's write per interval over the divisor and $G_0$ the lattice's Green's function (S.37), Poisson's rest far away (the paper's Section 5.1).
\item Every field reads the potential's level into its own coefficients, the paper's Eq. (2), the clock's pace once and each Link's pace twice.
\item Gravity's potential field carries three odd lines beside its time line, sourced at the matter's Links by the count's flux, $V_{ij} = -V_{ji}$, and resting around a moving source at $\mathbf V = U\,\mathbf v_s$ as the time line rests at Poisson's, which a plane configuration reads at the Port as the phase $\theta_{ij} = 4V_{ij}\omega_0 / c^2 = 12\,V_{ij}\omega_0$, $c^2 = 1/3$ light's own speed squared on the lattice, the three axes' $3$ and no number of nature, the delay of a rotating configuration along a Link $4V / c^2$ intervals, and a one-part real line does not read.
\end{enumerate}

\subsection*{B.3 The bound state: the fixed point and its stability (the paper's Section 4.2)}

A bound state is its count and its configuration over its Nodes, a configuration held at its Nodes by the potential wells; its count is the number of its quanta, it clicks when a quantum from outside meets its own transition, and what a NodeDetector over all its Nodes credits of it is the energy of its configuration in quanta, a reading and not a state. No law is written for a bound state: it is the support of a standing mode of the one rule, the fixed point of its binding potential, standing where the sources return themselves, its quanta its configuration's quantum share, and a free packet falls as $a = \Delta\,|d\omega_0 / d\ell|\,\nabla \ell$, $\Delta = 1 / m^*$ the band's curvature at rest and $\ell$ the potential level (derived from the line, with the decisions named; lattice, clicks for the fall; the paper's Theorem 4; S.13, S.18). The free packet's fall coefficient is $2\nu / (3(1 + \nu))$, $\nu = \num / \den$, differing between fields at finite gaps, $4 / 15$ at $[2, 3]$ against $2 / 7$ at $[3, 4]$ (derived; clicks; S.18, S.44 and C.7). A bound state in the dark is only the future and still gravitates. It is computed by the rule's own iteration, the wells a spread of the declared count lays, the configuration's top mode in those paces by Rule3's own read act, the levels scaled to the declared count and the wells rewritten until the levels repeat, in integers, at the fixed point of the division act; the branch is chosen by the initial data's width alone (a lattice reading; lattice; B.3 and C.3).

For a bound state of $C$ quanta with the configuration $\phi$ over its region, $\sum\phi_i^2 = 1$, sourcing potentials with the static kernels $\mathbf G_r$, the fixed point is to first order in the well a critical point on the sphere of the functional $F[\phi]$ of S.13, the band's pressure less the wells' energy, each well counted once because the configuration is its own source (Hartree's form), with $2\cos\omega_b = 2\cos\omega_0 - \lambda$ at $dF / d\phi = 2\lambda\phi$ (a theorem; lattice).

A hollow beyond its reach adds a barrier with no critical point where $3bc > a^2$ ($a$, $b$, $c$ the pressure's, the hollow's and the well's coefficients at unit size), the collapse, which ends in a frozen clock under the one assumption (derived; lattice; S.13; the frozen clock derived under the one assumption, S.39). And a screened kernel's exponent runs from the massless kernel's $1$ to the contact kernel's $3$ (computed; lattice; S.13). The binding potential's fixed point has two branches, a wide cloud and a compact branch one Node wide, and the initial data's width alone chooses between them (theorem for the equation; derived for the stability, the integer version open; the branches a lattice reading, named, in C.3; lattice; S.13, S.45). No standing bound state is claimed from a run (a lattice reading, named; lattice). The compact branch is coarse at the Node's width by construction, its drift and breathing bounded by $(\mathrm{Link} / R)^2$ and $\sqrt n / A$ (C.3), and whether a laid bound state stands under Rule3 within that walk is not this paper's claim.

The atom is known from nature's measurements and nothing of it is seen on the lattice; what is tested is whether the way from the clicks to the lattice and back carries it (S.43, S.52). In the model an atom is a bound state of the charge potential. A hydrogen nucleus, one quantum of the proton's field, writes the time angle $\theta_0 = Z\alpha c / r$, with $c = 1 / \sqrt 3$, into the charge potential by its Wronskian, and electrons, quanta of the matter field, read that angle into their rotation. The standing configurations Rule3 makes in that angle are the atom's levels, with Bohr's form $E_n \propto 1 / n^2$ and $a_0 = \sqrt 3 / (m^* Z\alpha)$ when the nucleus's angle alone is read (derived in the slow limit under line (1) of S.52, an assumption, fitted, made with hydrogen's levels in view, its ground the click's line (a configuration reads every charge potential's row but its own, a meeting being of two quanta and never one; S.52); line (2) an assumption; the criterion the model's decision, an assumption; clicks; S.43, S.52). A configuration that reads its own charge field has the law's form at large $Z$ only, $1 - 1.25 / Z$, and not hydrogen, so under line (1) of S.52 a configuration reads every charge field but its own, and the level's potential fields keep the self-read (line (2), an assumption). The criterion is a rule of reading: what nature splits with energy is read as a bound state, and what no energy splits as a quantum of a field, a quantum itself never breaking; the atom, the nucleus and the molecule are bound states by the one mechanism the law has, a configuration standing in a potential's well (the model's decision, an assumption; the one mechanism an assumption; C.3).

\subsection*{B.4 A bound state's clicks (the paper's Section 4.3)}

For one Node of a bound state, $\anext + \abefore = 2\cos\omega(t)\,\anow$ with $\omega$ changing slowly keeps $A^2\sin\omega_b$ invariant, the discrete adiabatic invariant of Rule3's line, while the form $D = A^2\sin^2\omega_b$ and the energy are not invariant, the recurrence stepped over $20{,}000$ intervals with $\omega$ lowered from $0.841$ to $0.600$ keeping $D / \sin\omega$ within $4 \times 10^{-5}$ while $D$ falls by a quarter (computed; lattice; S.14). A bound state's rotation reads the level once, so at rest its rate falls by $f(\omega_0)\,U$ with $f(\omega_0) = 2(1 - \cos\omega_0) / (\omega_0\sin\omega_0)$, $1.063$ at $[2, 3]$ and $1$ as the gap closes, $f = 1 + \omega_0^2 / 12 + \dots$ in series, the finite-rotation correction in $\omega$ of a rate that is exactly $N$ in $2\sin(\omega / 2)$ (S.15, S.31); a light clock between two faces falls by $2U$, the Link's pace read twice.

A bound configuration at the rotation $\omega_b$ that reads a light level into its potential level, the pace read, is modulated by $\epsilon = 4g(1 - 2x_0)\,b_L / \Gamma$, with $g = 1 - \num / \den$ the rest term's coefficient (half the mass term $m^2$ of the paper's Section 2.2), $x_0$ the bound state's potential level in units of $\Gamma$ and $b_L$ the light level's amplitude. The resonant part at $\omega_L = 2\omega_b$ grows the amplitude at $\mu = \epsilon / (4\sin\omega_b)$ per interval, under that read alone, none under the rotation (S.41), and a bound state with two modes $\omega_i < \omega_j$ transfers between them at the difference, at Rabi's rate $\Omega_R = \epsilon_{ij} / (4\sqrt{\sin\omega_i\sin\omega_j})$ (derived, the clock's rate under the one assumption, the invariant computed; clicks; the drift lattice; S.14, S.15, S.16, S.41). Under the charge potential's rotation the rate has the small-line form $\Omega_R = \theta_L\omega_L|d| / 2$, $\theta_L$ the light's angle on the Link and $d$ the transition's dipole element (S.41). The fraction absorbed, with the bound state's write back, is open. A bound state's count in clicks is its energy summed over its Nodes, a gapped quantum's energy in a well differing from $W_{\mathrm{rec}}$ by the rest term's shift, $m^2 = 2(1 - \nu)(1 - 2U)$ changing by $-4U(1 - \nu)$ in $m^2$'s unit, $\nu = \num / \den$, the quantum's energy by $-2\nu U / (1 + \nu)$ relative (S.18), the mass defect read as the count.

At the atom the born light is at the resonance $\omega_e - \omega_g$ and not at the energy-conserving frequency (the NodeDetector's declaration, an assumption; S.46). Absorption erases the photon's empty wave inside the click's causal cone. There is no recoil in the law's write, since the write conserves no momentum (under the proposition of the paper's Section 2.5.2 the twist is $\mathrm{sgn}(k)(|k|\,\mathrm{div}\,M)$ per Link and the recoil $k^2 / (2Mm^*)$ in rotation; the atom's needs the nucleus's finite inertia, which line (2) of S.52 pins, so it stays outside the claim). The orbital Zeeman splitting has Land\'e's $g_L = 1$ (derived; clicks for the levels' ratios, lattice for the configurations; S.43). Two standing bound states at the two ends of a Link write a cross current at their beat, which falls under one quantum's measure beyond the reach $d = \ln(\num\,a_1 a_2 / (3\,\den\,T\sin\omega_0)) / \kappa_b$, $\kappa_b$ the bound states' own inverse reach at $\omega_b$, $\cosh\kappa_b = 3\cos\omega_b / \cos\omega_0 - 2$ (derived; lattice; S.17).

\subsection*{B.5 Clusters: the three kinds of binding, and the cosmology at small redshift (the paper's Section 4.4)}

The three potential levels at a Node are the nesting, the binding potential's level the matter's tail, the charge potential's the charge's well and the massless potential's gravity's well, three kinds of binding by rank, quanta into a bound state, bound states into a charged cluster, charged clusters into a gravitational cluster (derived from the reads of the paper's Section 4.1; lattice). The couplings are at the edge, every divisor $1$ but gravity's $E_g$ (assumption, the divisors; lattice). A bound state declared a NodeDetector clicks in the meeting's form, and one with no declaration does the Nodes' laws alone. A click is not the NodeDetector's privilege: every bound state, in the model's intent and in the hypothesis of the paper's Section 9, absorbs, emits and converts by the same meeting, and the NodeDetector is the one bound state whose clicks are reported. What it credits of a bound state moves only by the clicks at its boundary, up to the drift of the paper's Section 4.3 (derived from the definitions; clicks, the drift lattice; S.14).

What nature defines as dark matter, gravitating and unseen, is in Universe24 every configuration that meets no NodeDetector's transition, the free quanta of the gapped scalar matter field, with $W = 0$, reading no charge potential and scattering no light, and the bound states no declared NodeDetector meets, each writing the potential level by its form and clicking nowhere (definition, derived under the reads of the paper's Section 4.1; the source a hypothesis; lattice, clicks). Its kind is derived under the universe's table as declared, its source is a hypothesis, its amount is not derived, and the paper claims no identification with nature's dark matter.

\subsection*{B.6 The formulas the paper adds, and the criterion of refutation (the paper's Section 5.2)}

The post-Newtonian column at the computed pair departs with the gap, so at the rule's own matter pair the model does not describe the gravity of the solar system, and the closing gap is nature's pair, where $\gamma_K - 1$ is below the MAGIC bound of the paper's Section 5.1 and $6.7 \times 10^{-20}$ at LHAASO's, whether the model's bound states stand at that pair not shown. Read as a metric, the composed clock gives $N^2 = e^{-2U}$ and the ruler $h^2 = e^{2U}$, against the isotropic Schwarzschild chart's $1 - 2U + 2U^2$ and $1 + 2U + 1.5U^2$: the spatial second order is $2$ against $1.5$, a prediction no measurement has read, which one click in the deep level would read (S.34). Light's index $e^{2U}$ places the photon sphere at $r = 2m$ with $U = m / r$, $m = GM / c^2$ in Links, so the capture radius of a black hole's shadow is $b = 2e\,m = 5.44\,m$, $e$ Euler's number, against Schwarzschild's $3\sqrt 3\,m = 5.20\,m$, $4.6$ percent larger. This is a conditional geometric result under the one assumption, known for the exponential metric \cite{yilmaz}\cite{boonserm} and reached here from Rule3: the ray's capture follows from $n(r) = e^{2m / r}$ alone, while the fate of the rays inside is open, the law having no horizon, so that a shadow needs absorption by a bound state's clicks, not computed (S.30). The Event Horizon Telescope's Sagittarius A$^*$ (Sgr A$^*$) has the fractional deviation from Schwarzschild $-0.08 \pm 0.09$ \cite{eht_sgra}, $1.4$ standard deviations from the prediction under the Very Large Telescope Interferometer mass prior, and M87$^*$'s is $-0.01 \pm 0.17$, about $0.3$ standard deviations \cite{psaltis}.

One number follows from the law's lines where general relativity and quantum theory do not meet, entered as untested. In the model a light packet's integrated write is a static bound state's of the same count, $1 \times E / c^2$ under the static kernel, the packet's own field open, where general relativity's pencil pulls at $2 \times E / c^2$ \cite{tolman} (derived under the static kernel, an assumption named, the packet's field open; untested; clicks; S.39). One check stands beside it: a bound state laid as one configuration on two paths at two heights has its beat scale as $N = 1 - U$, so that the visibility of the fringes falls as $|\cos(\Omega\,\Delta U\,\tau_{\mathrm{hold}} / 2)|$, $0.989$ for a strontium clock at one metre for one second, the magnitude of the semiclassical expectation \cite{zych}, unmeasured (derived under the one assumption, S.31's scaling; untested; clicks; S.39 (1)). None of the added formulas is confirmed by nature, each following from the line and the assumptions named beside it, bounded where a measurement exists and a prediction where none does.

\subsection*{B.7 The two slits: the declared universe; what the model says it must give, and the run (the paper's Sections 6.1 and 6.2)}

The wall at the column $20$ has two gaps of three Nodes, $12$ Links apart. The screen at the column $44$, $24$ Links from the wall, is twelve declared NodeDetectors of four rows each. A thirteenth NodeDetector of sixteen Nodes stands behind the first gap, backed by no face; its net inflow over the window is nothing ($-0.58$ by the real line), and it is never drawn. The counting window is the intervals $45$ to $130$. The run is $130$ intervals of Rule3 at every Node from these initial data, with the wall's face and the screen's NodeDetectors and nothing else (the run's declaration; S.23).

The integer step's deviation from the line is bounded exactly per step, below one level per Node per interval. The error per laid mode is bounded by $(n - 1) / \cos(\omega / 2) + 1 / \sin\omega + 1$, $135.6$ at the carrier at $n = 130$, while under the model of independent uniform remainders it walks as $\sqrt{n / (12(1 + \cos\omega_k))}$ per band mode, $2.4$ levels at the carrier and $3.0$ per Node over the lattice's zone at $n = 130$, against $b = 1{,}241$, which is $0.4$ to $0.5$ percent of the energy and $1.1$ to $1.4$ units of $N$ (derived for the bound, computed for the walk; lattice; S.53 (i)).

Huygens's construction and Young's spacing give the far field's spacing $\lambda L / d = 16$ Links and, at $L = 2d$, the near field's minima near $y = 15.3$ and $32.7$, in the regions $3$ and $8$ of the twelve four-row blocks counted from $0$ (derived; clicks; S.23). The blind row's dips are at the regions $4$ and $8$ (computed; clicks; S.23, the paper's Fig. 2(b)). The pattern's centre at $y = 24$ is no mirror of the blocks' $23.5$, which is why the two readings can differ by one region.

The run's $N = 285$ is $674$ photons of energy $T\sin\omega$ at the run's rotation $\omega = 0.4366$ on the level-$60$ paces ($285 / \sin\omega$). \emph{The two rotations.} Light's one-axis band at the carrier $k = \pi / 4$ is $\cos\omega = 1 - (p / \Gamma)^2 (1 - \cos k) / 3$: at the level $0$'s paces, $p / \Gamma = 1$, $\omega = 0.4456$, $\sin\omega = 0.4310$, $v_g = \sin k / (3\sin\omega) = 0.5469$ Link per interval, the flight $34 / v_g = 62.2$ intervals and the remainders' bound $(n - 1) / \cos(\omega / 2) + 1 / \sin\omega + 1 = 135.6$ at $n = 130$; at the level $60$'s paces, $p / \Gamma = (1 - 1 / \Gamma)^{120} = 0.9802$, $(p / \Gamma)^2 = 0.9608$, $\omega = 0.4366$ and $285 / \sin\omega = 674.0$ photons, $661.3$ at the level $0$'s rotation (computed; lattice). \emph{The counting unit.} The run's $N = 285$ is counted in $W_c = 3\,\den\,T$, and a quantum of count $1$ is laid at $B^2\sin\omega = T$ on light's one real line, the unit of C.9 (b), so the photon count is $N / \sin\omega$ and no other normalisation enters (computed; lattice; C.9 (b)). The line stepped on the laid integers themselves, laid at the level $60$'s paces and stepped at the level $0$'s, gives $278.9$, $6.1$ below the run and outside the walk: the before level's mismatch carries an opposite-branch admixture that the receding face loses, so $278.9$ is no comparator of the run (computed for $278.9$, derived for the admixture; clicks for the total, lattice for the admixture; S.23). The run's $285$ and the line's $284.7$ are quoted at the scale $s = 1$; the lattice's transmission per laid quantum sits above its continuum limit, and the line's $N$ moves with the scale. The scaling: with the lattice refined $s$-fold about the same geometry, the transmission per laid quantum at $\lambda = 8$ Links runs $0.157$, $0.148$, $0.142$, $0.139$ at $s = 1, 2, 4, 8$, converging as $1 / s$ to the continuum's $0.136$ (the fit from $s = 4$ and $8$ checked against $2$ and $4$; the one-Node wall and screen the first-order term), and the line's $N$ runs $284.7$, $308.3$, $308.1$, $304.9$ with the limit $303.8$, so the run's $285$ and the line's $284.7$ are the lay's own numbers at $s = 1$ and not the continuum's (computed; lattice).

The arrival is $t = L' / v_g$ with $v_g = (\num / (3\,\den))\sin k / \sin\omega = 0.547$ Link per interval and $L' = 34$ Links from the packet's centre to the screen, $62.2$, beside $62.0$ on the same universe with the lattice refined fourfold ($62.1$ eightfold); the peak of the screen's inflow is at $69.0$ on the lattice as declared (derived, computed; lattice; S.1, S.23). The model's opaque NodeDetector in one gap loses the fringes and reads one gap's envelope, as nature's absorbing which-way markers do, by the continuity of the one inflow through the two boundaries (derived, the continuity of the one inflow through the two boundaries; clicks; S.23). A which-way measurement without absorption is not computed here.

\subsection*{B.8 Bell and GHZ: the declared universe; what the model says it must give (the paper's Sections 7.1 and 7.2)}

The entanglement the model computes is one event with two trees: a pair is two configurations from one source Node, each read forward and backward, and the pair's correlation is the partner's configuration read through the same source, re-weighted by what was realised. The pair is laid as the pair field, two real lines of light never summed at a Node, with equal parts. Each side's NodeDetector is a region at least half a wavelength across, reporting the signed level sums of the two parts over its region, and it applies the declared setting, a pair of integers $(p, q)$, as the coefficients of the analyser's two channels $+$ and $-$, $e(+) = (p, q)$ and $e(-) = (-q, p)$, orthogonal, so each side's marginal is one half for any setting. The settings of the Clauser-Horne-Shimony-Holt (CHSH) test are $(1, 0)$, $(1, 1)$, $(12, 5)$ and $(5, 12)$, at $0$, $45$, $22.62$ and $67.38$ degrees.

(1) Rule3 gives each part's amplitudes at the ports; the two parts of one beam, laid equal, are two configurations of one field stepped by the same line along the same path, and they are identical exactly, integer for integer (a theorem: one integer map on equal integers; lattice). (2) The meeting, with the linearity of the real line that Rule3 follows within its bounded remainder (the paper's Theorem 2; the bound per laid mode as Section B.7 states it for the slits), gives the joint amplitude. The realised click at A's port weights the parts by A's analyser, is read back to the source Node, where the two parts were one event, and forward to B, where B's analyser weights them; the source sums the two branches, $\sum_k c_k\,e_k(a)\,e_k(b) = \cos(a - b)$ for the normalised vectors $e(a)$ and $e(b)$ (derived; clicks; S.10). (3) That the energy is the density of the meeting, its square, is Born's rule for a pair \cite{born}; nothing in Rule3 forces it, and everything after it is derived in the algebra (an assumption of the law, fitted, adopted with nature's correlations in view; the square Born's declaration: the counting rule was chosen to reach nature's correlations, so Bell's $2\sqrt 2$ follows from it and does not test it; clicks; S.47).

For unequal parts, with $r$ the ratio of their products at the ports, the same counting rule gives $E = \cos 2a\cos 2b + \rho\sin 2a\sin 2b$ with $\rho = 2r / (1 + r^2)$ and $S = \sqrt 2\,(1 + \rho)$ at the ideal angles, the maximum over the four settings $2\sqrt{1 + \rho^2}$, Tsirelson's bound \cite{cirelson} at $\rho = 1$ and never the algebraic $4$; the declared pairs' $478 / 169$ sits below $2\sqrt 2$ because the angles are rational (derived; clicks; S.10, S.47). No signal passes, $P(A+) = (\cos^2 a + r^2\sin^2 a) / (1 + r^2)$ being independent of $b$ for parts of any amplitude ratio (derived; clicks; S.10).

Bell's value sits in the counting rule, a declaration, no independent evidence that the rule accounts for entanglement. A counting rule by one side's contrast (S.10) reaches the quantum correlation only by discarding pairs \cite{pearle}\cite{eberhard}\cite{gargmermin}, the event-based simulations with a time-window rule the named instance \cite{deraedt2007}\cite{michielsen2014}, which the loophole-free tests of 2015 close for nature \cite{hensen}\cite{giustina}\cite{shalm}. No entanglement other than the pair's and the entanglement swap's is named as met, the swap's joint energy a hypothesis under its own name with the criterion that separates a transferred quantum from a coupling, the counting rule's bilinear form after the exchange (S.47); the gate-joined ions and transmons that exceed the law's $2$ are a finding against nature named in S.50 (14), not counted among the thirteen of the paper's Section 8.1 (derived for the bound, the comparison a named finding; clicks; S.50 (14)). No settings are chosen by the lattice's own past, so a superdeterministic world \cite{hossenfelder} is not this one.

Keeping the lattice local has three prices. Born's rule is the NodeDetector's declaration. A clicked configuration's empty wave still writes its form into the potentials it sources until the front reaches it. The meeting through the source Node gives $S = 2\sqrt 2$ with nothing signalled, resting on the counting rule, a declaration.

\subsection*{B.9 The misses, with their numbers (the paper's Section 8.1)}

The massless potential's travelling events carry the breathing polarisation and the lattice-aligned $+$ and no $\times$ (derived under the static read, which the odd write keeps; lattice for the form; S.50 (1); S.56 (d)). GW170814, read by LIGO's two interferometers and Virgo's, favours pure tensor polarisation \cite{gw170814} (a comparison; clicks). The write without the odd Link part gives $\alpha_1 = -8$, $\alpha_2 = -1$, $\eta_N = 7.3$ and no frame dragging (S.35, S.50 (1)). These four values are the write without the odd Link part's. Under the law's write a plane configuration of one rotation sense has S.56 (e) to (g)'s values and a body its signed share, $\alpha_1 = 8f - 8$ (S.56 (h), the paper's Section 8.1). On Rule3's own clock stands this: of the five forms of C.4, none is a real, local, explicit, bijective lattice line carrying $g_{0a}$. That is why the odd part is a phase on a plane configuration and no pace (a theorem for one-part real lines with a $V$-proportional odd part, the staggered step's whole-Node shift the one exception, S.56 (d); lattice; C.4 its five instances).

At the second order in $z$ the model's finite ball of dust, bound or free, decelerates, $q_0 = \Omega_m / 2 > 0$ (derived under the declared homologous start; clicks; S.29). The supernovae read $q_0 \approx -0.5$, the flat $\Lambda$CDM model's (cold dark matter with a cosmological constant) $\Omega_m / 2 - \Omega_\Lambda$, $\Omega_\Lambda$ the dark energy density parameter \cite{riess}\cite{perlmutter}\cite{pantheon}. Rule3 has no repulsive sector, every potential level held being a hollow, so an accelerating Hubble diagram is a finding against nature by name. The nuclear binding under the declared potentials misses in two directions, the deuteron above and iron below, by factors of order ten and two under the root-mean-square (rms) convention of the quoted radii ($12$ and $2$; $4$ and $6$ under a Bohr-like radius), the profile and the screening entering at the same order (C.6 (b)), and the size law gives $R \propto 1 / A$ against $A^{1/3}$, one well with no saturating core (derived under the nuclear hypothesis; clicks; C.6 (b), (e)). The exclusion's part in the size law has no line (S.50 (2), (5)). A bound state's binding per quantum is bounded by $0.76\,CW / (\Gamma E)$ at one Node, $W$ the declared write weight and $E$ the divisor (derived, a bound under the declared $W$, which no click bounds; lattice for the well, clicks for the defect; S.36). The shells and the magic numbers, the exclusion's, and the neutron's lifetime are not derived (no line; clicks; S.50 (3), (5)).

The model has no electroweak theory; its weak interaction is the conversion click (the declarations of the paper's Section 8.2). Every rotation potential couples at $\alpha$ at the Link, so no Fermi constant $G_F$ follows: at $\alpha$ and nature's $m_W$, the W boson's mass \cite{pdg}, $G_F$ falls short of nature's by a factor $4.6$ ($2.51 \times 10^{-6}$ GeV$^{-2}$ against $1.166 \times 10^{-5}$; or $m_W$ from $G_F$ by $2.2$), the mixing angle's factor being absent (S.50 (17); S.37 (c)). A three-line configuration read as $(\nu, e)$ under one $O(3)$ has no fourth rotation, so there is no $Z$ boson, where nature has the neutral current \cite{hasert} and the $Z$ \cite{lep}. Parity is conserved, where nature's weak interaction violates it \cite{wu}; the two misses have one root, a step that commutes with the cube's reflections. The law has no spontaneous click, so the conversion has no derived rate and the neutron's $878.4 \pm 0.5$ s \cite{pdg} has no source.

In the atom, the self-read gives hydrogen's $1s$ at $-0.043$ of Rydberg's unit (S.43), which line (1) of S.52 mends. The fine structure under the spin-0 configuration has the Klein-Gordon form $n / (l + 1 / 2)$, Sommerfeld's $n / k$ \cite{sommerfeld} with $k$ read as $l + 1 / 2$. It splits hydrogen's $2p$ from $2s$ by $29$ GHz, where Dirac's $n / (j + 1 / 2)$ gives $2p$'s splitting $10.95$ GHz (S.50). The law holds no amplitude over two positions on the lattice. So the reduced mass, the isotope shifts and the muonic atoms are outside what it derives, and the many-body correlation is not reproduced by a product of configurations, helium $1.4$ percent short of nature's binding energy and positronium at $43$ percent of it (computed by hand from the modes, a comparison and no click claimed; lattice; S.52). Pauli's exclusion has no line. Universe24 is a mean-field theory of one quantum. It computes helium at the Hartree-Fock value, $-2.8617$ Ha (hartree) \cite{roothaan}, and lacks the correlation energy and the radiative corrections in the amounts Hartree-Fock lacks them in nature. Its many-body misses (the reduced mass, positronium's binding, Pauli's two per orbital, the nuclear size law in part) are one absence, the parts' pairing, which the counting rule reads through the source Node and the lattice's step does not read, antisymmetric with spin for the exclusion; its radiative misses are a vacuum that does not fluctuate (derived under the two lines where the numbers are, a comparison for the Hartree-Fock value, no line where the key says so; lattice; S.52).

Nature's mass ratios are not whole numbers ($m_p / m_e = 1836.153$ \cite{codata}), and a cluster's defect is bounded by the law at $C\omega_0 x_{\mathrm{self}}$, of the order of $10^{-22}$ at the write weight $1$ and the illustration's $\Gamma = 10^{14}$ (S.36 (d)), smaller still at nature's $2^{94}$, the one denominator of S.54 declared as $\Gamma$ (C.7) (derived, conditional on the declared $W$, which no click bounds; lattice for the well, clicks for the defect; S.36 (d)). So no particle of nature can be read as an aggregate of the electron's quanta. The rule's own universe at $[2, 3]$ is no finding against nature, because nature measures at its own pair. The gaps have their causes in the findings above, and none is the fault of a number: no loops ($\alpha$'s running, the Lamb shift, $g - 2$), no spontaneous click (the lifetimes), no vector for the exclusion, and a step that commutes with the cube's reflections (the $Z$ and parity). Between separately prepared bound states the law's counting rules give at most $2$, where gate-joined ions \cite{rowe} and transmons \cite{storz} exceed it (S.50 (14); B.8).

The lattice's scale follows from the massless band's quartic term, which gives the quadratic Lorentz-violation scale $E_{\mathrm{QG},2} = \hbar c / (\ell_{\mathrm{Link}}\sqrt{\epsilon_4})$, $\ell_{\mathrm{Link}}$ the Link's length and $\epsilon_4$ the band's quartic coefficient. The Fermi Large Area Telescope's (Fermi-LAT) bound \cite{vasileiou} puts the Link at $\ell_{\mathrm{Link}} < 6.4 \times 10^{-27}$ m, so that the electron's pair needs $\den \ge 2^{94}$ (derived, a constraint on any declared universe meant for nature, the law fixing no number; lattice; S.54). The paper uses two Link bounds: MAGIC's $1.3 \times 10^{-26}$ m in the paper's Section 5's rows (S.24; S.36 (d) and S.37 (c)) and Fermi-LAT's $6.4 \times 10^{-27}$ m, the tighter, in S.54's $\den \ge 2^{94}$; each row names the one it uses. The click's write conserves no momentum, so the recoil is the NodeDetector's, where nature's counter recoils, as in the M\"ossbauer line \cite{mossbauer} and the Compton shift (its form under the proposition is the paper's Section 2.5.2's). The model declares a NodeDetector and derives what the click does once declared; that nature's bound states jump \cite{cook}\cite{bergquist} is the method's widest open gap. Neither layer keeps a total energy (a theorem for the step's conservations, an assumption for the clicks'). No standing bound state is claimed from a run.

\subsection*{B.10 What is put in (the paper's Section 8.2)}

\begin{itemize}
\item a field's entry and the integers $\Gamma = 6{,}000$, a multiple of $6$ so that light's and matter's pairs are exact over it, the binding potential's pair $[2400, 2401]$ the one pair not over it, and $T = 2^{15}$;
\item a NodeDetector with Nodes alone, its region, the universe's counting window and the seed;
\item a NodeDetector with a configuration of its own;
\item the emission's rate constant;
\item the conversion's table and rate, a hypothesis under its own name;
\item the initial data;
\item the receding face;
\item the counting rule's one draw through the source Node, the one declared non-local element;
\item the erasing front.
\end{itemize}
The two slits' run carries the law's lines without the write of the odd Link part: its universe has gravity off in all but name, and Bell's derivation reads no current (a declaration; the paper's Section 6).

The numbers of the declared universe are nature's or declarations, none the rule's. The vacuum level is a declaration. The binding range follows from the pair (S.11). Gravity's weight $E_g$ and the charge potential's $k$ are read from $G$ and from $\alpha$, and the gaps are read from the masses. No number comes from the lattice itself (calibration, experiment for $E_g$, $k$, the gap and the Link's length; hypothesis for any such number; clicks; S.48). That the declared universe is the one the paper tests is a decision and not a theorem (an assumption). The selective read of the nuclear potentials by the nucleon fields and the energy line of S.26 are declared lines of the law. Three things are held to be computable and not computed: the bound states at the closing gap, the moving fixed point of a bound state, and the absorbed fraction with the bound state's write. Four things are not claimed: that nature is a lattice, that any number of the computed pair $[2, 3]$ is nature's, that a comparison with nature is more than a consistency statement at the stated order, or that a formula known in the continuum is new where the paper gives its lattice form.

\subsection*{B.11 What a world closer to nature needs (the paper's Section 8.3)}

The first is the box's: a finite box returns the potentials' radiation, and a box larger than the light-crossing of its evolution ($10^{12}$ Nodes, about $17{,}300$ intervals at $1 / \sqrt 3$ Link per interval) behaves as open (the return read on the lattice of a closed box and not derived, the crossing arithmetic from light's $1 / \sqrt 3$; lattice). The second is $\Gamma$'s, the integers' fineness being $\Gamma$'s and not the box's: the vacuum level's floor, the least level $c$ with $c - K / \sin\omega(\pi) \ge 1$ for the start's kick $K$, $(K^2\Gamma / 16)^{1 / 3}$ at large $\Gamma$, is $8$ to $13$ levels at $\Gamma = 6{,}000$ for the start's kick $K$ this world's lattice reads, the kick read on this world's lattice alone and the step not taken at a second $\Gamma$, and below $10^{-19}$ of $\Gamma$ at nature's $2^{94}$, the one denominator of S.54 declared as $\Gamma$ (C.7), so the discreteness' marks shrink as $\Gamma^{-2 / 3}$ whatever the box, and a world closer to nature is a larger $\Gamma$ with the one-way face, not a larger count of Nodes. The floor's derivation: at the level $0$ the massless band's corner mode sits on the step's edge, $2\cos\omega(\pi) = 2 - 4(p / \Gamma)^2 = -2$, a level below $0$ grows, the margin at the level $c$ is $4 - 4(1 - 1 / \Gamma)^{4c} = 16c / \Gamma$ to first order, the start's kick $K$ rings the mode at $K / \sin\omega(\pi)$ with $\sin\omega(\pi) = 4\sqrt{c / \Gamma}$ to first order, and the least $c$ with $c - K / \sin\omega(\pi) \ge 1$ solves $c^{3 / 2} = K\sqrt\Gamma / 4$, $(K^2\Gamma / 16)^{1 / 3}$: $8$ for $K = 1$ and $13$ for $K = 2$ at $\Gamma = 6{,}000$ by the exact margin ($c - K / \sin\omega(\pi)$ is $-0.35$ at $7$ and $1.13$ at $8$, $0.75$ at $12$ and $2.19$ at $13$), $5.4 \times 10^{-20}$ of $\Gamma$ at $2^{94}$ (derived, with the kick read on the lattice; lattice). The calculations such a world needs are named and not done here, none a claim of this paper. Universe24 stays as it is, since it is not nature. The lines the algebra reached beyond it (the stresses read into the clock, a write of the click's momentum) are an option for building another world, which would not be Universe24, and no claim of this paper.
\subsection*{B.12 The conjectures (the paper's Section 9)}

The conjectures, one line each (inspiration):
\begin{itemize}
\item the weak force as the conversion click and the strong force as one three-line configuration read whole;
\item the nuclear potentials;
\item the configuration's parts as nature's particles;
\item the masses as fields with one denominator;
\item dark matter's source;
\item the hypothesis that in nature only a click moves a quantum between fields, with the single atom's lifetime as its test;
\item this paper's own, the click's twist of the paper's Section 2.5.2, Compton's factor its lattice term, $1$ at nature's gap, no test;
\item the frozen Node as the law's one-Link black hole.
\end{itemize}

\section*{The model (Section 2 of the paper)}

\begin{derivation}[the band at a pace, Eq.~(3) of the paper]
\emph{Inputs:} Rule3's line at paces where the two Links of an axis share one $R_a$, $2\,\num\,\Gamma^2$ at the vacuum's and $2\,\num\,p_a^2$ in general, $w\,\anext = \sum_a R_a(\mathrm{arr}_{+a} + \mathrm{arr}_{-a}) + S\,\anow - w\,\abefore$, with the remainders dropped (they are bounded by $w$ and, below one level per Node per interval, leave a plane wave's rotation unmoved at amplitudes large against a level, the clause of Online Resource 1, Section B.7, below one level per Node per interval; at an amplitude of one level the integer orbit is its own cycle). \emph{Steps:} for a plane wave $a_t(\mathbf x) = A\cos(\mathbf k \cdot \mathbf x - \omega t)$ the two arrivals of the axis $a$ sum to $2\cos k_a\,\anow$ and $\anext + \abefore = 2\cos\omega\,\anow$, so
\begin{align*}
2\cos\omega &= \frac{2\sum_a R_a\cos k_a + S}{w}\\
&= \frac{4\,\num\sum_a p_a^2\cos k_a + 12\,\den\,\Gamma^2 - 12(\den - \num)\,p_0^2 - 4\,\num\sum_a p_a^2}{6\,\den\,\Gamma^2},
\end{align*}
that is $\cos\omega = 1 - g\,(p_0 / \Gamma)^2 - \frac{\num}{3\,\den}\sum_a (p_a / \Gamma)^2(1 - \cos k_a)$, which is Eq.~(3). At the vacuum's pace, $p_0 = p_a = \Gamma$, this is $\cos\omega = \num / \den - \frac{\num}{3\,\den}\sum_a(1 - \cos k_a) = \frac{\num}{3\,\den}\sum_a\cos k_a$; at $k = 0$, $\cos\omega_0 = \num / \den$. Written as $2(1 - \cos\omega) = 2(1 - \num / \den) + \frac{\num}{3\,\den}\sum_a 2(1 - \cos k_a)$, the band is the discrete Klein-Gordon form with $c_K^2 = \num / (3\,\den)$ and $m^2 = 2(1 - \num / \den)$, an identity of Eq.~(3). \emph{Light's speed and dispersion:} for $[1, 1]$ along one axis $\cos\omega = (2 + \cos k) / 3$; expanding, $1 - \omega^2 / 2 + \omega^4 / 24 = 1 - k^2 / 6 + k^4 / 72$, so $\omega = (k / \sqrt 3)(1 - k^2 / 36 + \dots)$ and the group velocity $d\omega / dk = \sin k / (3\sin\omega) = (1 / \sqrt 3)(1 - k^2 / 12 + \dots)$: the speed $1 / \sqrt 3$ at long wavelength and the group speed falling as $1 - 0.083\,k^2$. Along the diagonal $k_x = k_y = k_z = k$ the band is $\cos\omega = \tfrac13 \cdot 3\cos k = \cos k$ exactly, with no dispersion; one Link per interval along the axes is the octahedron $|x| + |y| + |z| = t$, whose face $x + y + z = t$ stands at the distance $t / \sqrt 3$, light's $1 / \sqrt 3$. \emph{Check:} at $k = 0.05$, $0.1$ and $0.2$ the exact $\sqrt 3\,d\omega / dk - 1$ over $k^2$ is $-0.0833$, $-0.0834$ and $-0.0835$.
 \emph{Status:} derived; lattice.
\end{derivation}

\begin{derivation}[light's index at a content, Section 2.3]
\emph{Inputs:} Eq.~(3) with the paces of Eq.~(2), $p_0 = \Gamma(1 - 1 / \Gamma)^L = \Gamma N$, $N = e^{-x} + O(1 / \Gamma)$ at the content $x = L / \Gamma$ ($L$ the holder's level, the $c$ of the preamble), and $p_a = p_0^2 / \Gamma = \Gamma N^2$ for a reader at the weight $1$ and no tension. \emph{Steps:} for light, $g = 0$, so $\cos\omega = 1 - \frac{1}{3}N^4\sum_a(1 - \cos k_a)$, and at long wavelength $\omega = N^2\,|\mathbf k| / \sqrt 3$. Light's speed at the content $x$ is $N^2 = e^{-2x}$ times its vacuum speed, so its index is $n = e^{2x} = 1 / (1 - 2x) + O(x^2)$ and $n - 1 = 2x + O(x^2)$: the Link's pace is the clock twice, and so is light's slowing.
 \emph{Status:} derived; lattice.
\end{derivation}

\begin{derivation}[the guard's edge and the horizon, Section 2.3]
\emph{Inputs:} Eq.~(3); Eq.~(2)'s composed paces $p_0 = \Gamma(1 - 1 / \Gamma)^c$, $p_a = p_0^2 / \Gamma$, with the rounding to the unit; S.39's frozen content; S.16's parametric line for the guard's limit. \emph{Steps:} the lowest value of $2\cos\omega$ on the band, its top $\omega$, is at $k_a = \pi$ on every axis, where $2\cos\omega = 2 - 2g\,(p_0 / \Gamma)^2 - \frac{4\,\num}{3\,\den}\sum_a(p_a / \Gamma)^2$; the recurrence $\anext + \abefore = 2\cos\omega\,\anow$ has no exponentially growing solution iff $|2\cos\omega| \le 2$; at the edges $\pm 2$ a linearly growing solution exists ($u_t = t$ at $+2$ and $u_t = (-1)^{t + 1}\,t$ at $-2$, from $u_0 = 0$, $u_1 = 1$), which the guard admits, light's checkerboard mode sitting at $-2$; for a negative numerator the lowest $2\cos\omega$ is at $k = 0$ and the same bound holds with $|\num|$. The guard is the fixed-coefficient one: with the coefficients changing in time the product of the one-step transfers can have a trace below $-2$ while every instant's band obeys the bound, and a solution grows (S.16's parametric line is the case), so the guard bounds each instant's band and not a sequence of paces in time. With $p_0 = p_a = p$ the condition $2\cos\omega \ge -2$ reads $(p / \Gamma)^2\,(2g + 4\,\num / \den) \le 4$ at wave number $\pi$ on every axis and $2g\,(p / \Gamma)^2 \le 4$ at $k = 0$, that is $(p / \Gamma)^2 \le 2\,\den / (\den + \num)$ for $\num \ge 0$ and $2\,\den / (\den - \num)$ for a negative numerator, so the edge is $P^2 = 2\,\den\,\Gamma^2 \ddiv (\den + |\num|)$, which is $\Gamma^2$ for light (at $[-1, 2]$, $\Gamma = 3$ and $p = 4$: $P^2 = 12 < 16$, the uniform mode's $2\cos\omega = -10 / 3$ with the roots $-3$ and $-1 / 3$ growing while the checkerboard's $\cos\omega = 1 / 9$ stands). \emph{No horizon:} under the composed paces $p_0 = \Gamma(1 - 1 / \Gamma)^c > 0$ at every content and $p_a = p_0^2 / \Gamma > 0$, so $R_a = 2\,\num\,p_a^2$ never vanishes and no Link closes in the rational paces, on the domain $0 < \num \le \den$; the quadratic truncation $p_a = \Gamma(1 - 2x)$ vanished at $x = 1 / 2$, a horizon of the truncation. \emph{The frozen clock:} at a content $c$ the clock runs at $N = (1 - 1 / \Gamma)^c$ of the vacuum's, $e^{-c / \Gamma}$ to the unit, so a deep well freezes the Node's Link pace $p_i = p_0^2 / \Gamma$ to $0$ at the content $(\Gamma / 2)\ln 2\Gamma$ (S.39), the clock pace there $\sqrt{\Gamma / 2}$, and the Node's share is undefined there.
 \emph{Status:} theorem for the edge and the extreme mode; derived for the horizon of the quadratic form, a feature of the truncation; lattice.
\end{derivation}

\begin{derivation}[the cube's group of 48 and its action, Section 2.1]
\emph{Inputs:} Rule3's line, whose six arrivals enter through the six Port coefficients $R_{ij}$, the two of an axis one $R_a = 2\,\num\,p_a^2$ with no coefficient naming an axis beyond its own pace; the readings $D$, $F$ and $W$ of Eq.~(1). \emph{The group:} a map of the six Ports onto themselves that keeps opposite Ports opposite is a signed permutation of the three axes, $\mathbf x \mapsto \mathbf P\mathbf x$ with $\mathbf P$ a permutation matrix with signs; there are $3! \cdot 2^3 = 48$ of them, the group $O_h$ of the cube, of which the $24$ with $\det\mathbf P = +1$ are the rotations and the other $24$ reflections and rotatory reflections; with the translations of the lattice they form the space group of the simple cubic lattice. \emph{Covariance:} let $\sigma$ be such a map and let the turned lattice carry the levels $a^\sigma_{\sigma(i)} = a_i$ and the paces $p^\sigma_{\sigma(a)}(\sigma(i)) = p_a(i)$. The arrivals of the axis $\sigma(a)$ at $\sigma(i)$ in the turned lattice are the arrivals of the axis $a$ at $i$ in the original, in either order, and $R_a$ is even in the axis's sign; so the sum $\sum_a R_a(\mathrm{arr}_{+a} + \mathrm{arr}_{-a})$ is the same number at $\sigma(i)$ as at $i$, $S$ and $w$ are the same, and the division gives the same quotient and remainder: the step commutes with $\sigma$, and a lattice turned by $\sigma$ runs as the turned lattice. A translation commutes for the same reason. \emph{The readings' kinds:} $D_i$ is one number at a Node with no axis in it, carried to $\sigma(i)$ unchanged, a scalar. The currents $F_{ij}$ are attached to the Links and carried with them, and the momentum density $P_a(i) = (F_{i, i-a} - F_{i, i+a}) / \num$ (this derivation's orientation, positive for a record moving toward $+a$) changes sign where $\sigma$ reverses the axis $a$ and is carried to the image axis otherwise, a vector; $W_i$ carries no axis and is carried unchanged, a scalar under the cube's $48$, odd under the plane's internal reflection $\mathrm{im} \mapsto -\mathrm{im}$ and under the time reversal. The click $(R, t)$ names a region and an interval and is carried to the turned region, a scalar. \emph{The time reversal:} it exchanges $\mathrm{next}$ and $\mathrm{before}$ and keeps $\mathrm{now}$; $D_i = \mathrm{now}_i^2 - \mathrm{next}_i\,\mathrm{before}_i$ is symmetric in that pair and is unchanged; $F_{ij}$ and $W_i$, antisymmetric in the two adjacent levels they read, are carried to minus the next interval's readings, $F_{ij}(t) \mapsto -F_{ij}(t + 1)$ and $W_i(t) \mapsto -W_i(t + 1)$ (the reversed pair at $t$ is $(\mathrm{now}, \mathrm{next})$), the reversed current being the forward one read backward; and it carries the step to its inverse by Theorem 1. \emph{The rational rotation:} the two parts of a plane record obey the same linear line with the same coefficients, so any linear combination of them does; the combination, with $p$, $q$, $r$ here the integers of the rotation and not the paces, $(\mathrm{re}', \mathrm{im}') = ((p\,\mathrm{re} - q\,\mathrm{im}) / r, (q\,\mathrm{re} + p\,\mathrm{im}) / r)$ with $p^2 + q^2 = r^2$ carries solutions to solutions, and its Wronskian is $(p^2 + q^2) / r^2$ times the old one, unchanged. \emph{What is not in the group:} a boost changes the record's wave number and rotation together, which no permutation of Ports does, and a rotation by an angle other than the cube's carries a Node off the lattice; so neither is a symmetry, and Lorentz's factor, where the paper finds it, is a property of the clicks (Derivation S.21).
 \emph{Status:} theorem; lattice.
\end{derivation}

\section*{Exact properties (Section 3)}

\begin{derivation}[the share's change is the currents, Eq.~(5)]
\emph{Inputs:} the share at the vacuum's paces, $e_i = 3\,\den\,(\mathrm{now}_i^2 + \mathrm{before}_i^2) - \num\,\mathrm{now}_i\,S_6(\mathrm{before})_i$, with $S_6$ the six-neighbour sum; Rule3 at the vacuum's paces, where $S = 0$, $R_a = 2\,\num\,\Gamma^2$ and $w = 6\,\den\,\Gamma^2$, so that $3\,\den\,(\mathrm{next}_i + \mathrm{before}_i) = \num\,S_6(\mathrm{now})_i + 3\,\den\,\epsilon_i$ with $\epsilon_i = (r_i - r_i') / w$ the remainder term of the share's change (the paper's Section 3.4) and $w = 6\,\den\,\Gamma^2$ the wall at these paces, so that the remainders' term of this line, $(r_i - r_i') / (2\Gamma^2)$, is $3\,\den\,\epsilon_i$ (from $6\,\den\,\Gamma^2(\mathrm{next} + \mathrm{before}) = 2\,\num\,\Gamma^2 S_6(\mathrm{now}) + r - r'$). \emph{Steps:} after the step the pair $(\mathrm{now}, \mathrm{before})$ becomes $(\mathrm{next}, \mathrm{now})$, so
\[
e_i' - e_i = 3\,\den\,(\mathrm{next}_i^2 - \mathrm{before}_i^2) - \num\,[\mathrm{next}_i\,S_6(\mathrm{now})_i - \mathrm{now}_i\,S_6(\mathrm{before})_i].
\]
With $3\,\den\,(\mathrm{next}_i + \mathrm{before}_i) = \num\,S_6(\mathrm{now})_i + 3\,\den\,\epsilon_i$ the first term is $(\mathrm{next}_i - \mathrm{before}_i)(\num\,S_6(\mathrm{now})_i + 3\,\den\,\epsilon_i)$, and the bracket cancels its $\mathrm{next}_i$ part:
\begin{align*}
e_i' - e_i &= \num\,[\mathrm{now}_i\,S_6(\mathrm{before})_i - \mathrm{before}_i\,S_6(\mathrm{now})_i] + (\mathrm{next}_i - \mathrm{before}_i)\,3\,\den\,\epsilon_i\\
&= \sum_{j \sim i} \num\,(\mathrm{now}_i\,\mathrm{before}_j - \mathrm{before}_i\,\mathrm{now}_j) + (\mathrm{next}_i - \mathrm{before}_i)\,3\,\den\,\epsilon_i,
\end{align*}
the six currents $F_{ij}$ into the Node plus the remainder term. \emph{At any paces:} with the weighted share $\tilde e_i = [w(\mathrm{now}_i^2 + \mathrm{before}_i^2) - S_i\,\mathrm{now}_i\,\mathrm{before}_i - \mathrm{now}_i\sum_j R_{ij}\,\mathrm{before}_j] / (2p_i^2)$ and the line without its division, $w\,\mathrm{next}^{*}_i = \sum_j R_{ij}\,\mathrm{now}_j + S_i\,\mathrm{now}_i - w\,\mathrm{before}_i$, the difference $e_i(\mathrm{next}^{*}, \mathrm{now}) - e_i(\mathrm{now}, \mathrm{before})$ is $\sum_j (q_{ij}^2 / \Gamma^2)\,\num\,(\mathrm{now}_i\,\mathrm{before}_j - \mathrm{before}_i\,\mathrm{now}_j)$ with no anisotropy term, the $S_i$ terms and the Node's terms cancelling as in Theorem 3's proof, since $R_{ij} / (2p_i^2) = \num\,q_{ij}^2 / \Gamma^2$ is the same from both ends. The integer step's $\mathrm{next} = \mathrm{next}^{*} + (r - r') / w$ adds $-(\mathrm{next}_i - \mathrm{before}_i)(r'_i - r_i) / (2p_i^2)$, because $w(\mathrm{next}_i + \mathrm{next}^{*}_i) - S_i\,\mathrm{now}_i - \sum_j R_{ij}\,\mathrm{now}_j = w(\mathrm{next}_i - \mathrm{before}_i)$. \emph{Check:} exact fractions on a periodic chain of five Nodes at $[2, 3]$, the two other axes folded (four self-Ports per Node), $\Gamma = 10$, with contents $0$ to $3$, Link factors $5$ to $10$ and random levels: the identity holds at every Node, with and without the division; and a second witness, a uniform $\mathrm{now} = 3$ over $\mathrm{before} = 0$ at $\Gamma = 10$, $q = 5$, reads $81$ per Node, half of Eq.~(4)'s $162$, since $e_i$ carries $1 / (2p_i^2)$ where Eq.~(4) carries $1 / p_i^2$.
 \emph{Status:} two derivations; theorem; lattice.
\end{derivation}

\begin{derivation}[the total share is non-negative, Section 3.4]
\emph{Inputs:} $e_i$ as above, at the vacuum's coefficients, on any lattice whose Links read the same from both ends and whose every beyond reads $0$ (a closed, a periodic or an open-faced lattice); $S_6$ symmetric with row sums at most $6$ (six Ports, a Port onto a face reading $0$ contributing nothing on either side), so its spectrum lies in $[-6, 6]$; a declared beyond that is not $0$ adds linear terms and is outside the theorem (a universe whose massless potential rests at a declared level beyond the open faces among them). \emph{Steps:} $\sum_i e_i = 3\,\den\,(|\mathbf n|^2 + |\mathbf b|^2) - \num\,\mathbf n \cdot S_6\mathbf b$, and $|\mathbf n \cdot S_6\mathbf b| \le \|S_6\|\,|\mathbf n|\,|\mathbf b| \le 6\,|\mathbf n|\,|\mathbf b| \le 3(|\mathbf n|^2 + |\mathbf b|^2)$, so $\sum_i e_i \ge 3(\den - |\num|)(|\mathbf n|^2 + |\mathbf b|^2) \ge 0$ for $|\num| \le \den$. Equality needs $|\num| = \den$ and $\mathbf b$ an eigenvector of $S_6$ at $\pm 6$ with $\mathbf n = \pm\mathbf b$ ($\mathbf n$, $\mathbf b$ the vectors of the now and before levels): on an even periodic lattice the static uniform level of light or its checkerboard, $\mathrm{now} = (-1)^{x + y + z}$ with $\mathrm{before} = -\mathrm{now}$, whose shares are $0$. \emph{A Node's share:} at fixed $\mathrm{before}_i$ and $S_6(\mathrm{before})_i = s$, $e_i$ is a quadratic in $\mathrm{now}_i$ with the minimum $3\,\den\,\mathrm{before}_i^2 - \num^2 s^2 / (12\,\den)$, negative iff $|s| > 6(\den / |\num|)\,|\mathrm{before}_i|$ for $\num \ne 0$ (at $\num = 0$ the share is $3\,\den\,(\mathrm{now}_i^2 + \mathrm{before}_i^2)$ and never negative); a plane record has $|s| / |\mathrm{before}_i| = |2\sum_a\cos k_a| \le 6$ and is never negative, a lay's edge can be. \emph{A region:} pairing each boundary Link's cross term and bounding $2|\mathrm{now}_i\,\mathrm{before}_j| \le \mathrm{now}_i^2 + \mathrm{before}_j^2$ (the inside $\mathrm{now}_i^2$ absorbed by the interior row-sum deficit $(6 - d_i) / 2$, the outside $\mathrm{before}_j^2$ kept), $\sum_{i \in R} e_i \ge 3(\den - |\num|)\sum_R(\mathrm{now}_i^2 + \mathrm{before}_i^2) - \frac{|\num|}{2}\sum_{\text{boundary Ports}}(\mathrm{now}_j^2 + \mathrm{before}_j^2)$ over the outside neighbours $j$ (the numerator's modulus, so that the bound holds at a negative numerator too): a region's share is negative at most by half a level-square per boundary Port in the current's units.
 \emph{Status:} theorem; lattice.
\end{derivation}

\begin{derivation}[the tension of a plane wave, Section 3.3]
\emph{Inputs:} the stress reading $G_{aa}(i) = \mathrm{now}_i(\mathrm{now}_{i+2a} - \mathrm{now}_i) - \mathrm{now}_{i+a}(\mathrm{now}_{i+a} - \mathrm{now}_{i-a})$ and the cross reading $G_{ab}(i) = \mathrm{now}_i\,(\mathrm{now}_{i+a+b} - \mathrm{now}_{i-a+b}) - \mathrm{now}_{i+b}\,(\mathrm{now}_{i+a} - \mathrm{now}_{i-a})$, the stress booking on the $b$-Link $(i, i + b)$, $G_{aa}$ at $b = a$. \emph{Steps:} for $\mathrm{now}_i = A\cos(kx_i + \phi)$ along the axis, with $c_n = \cos(kx_i + nk + \phi)$, $G = A^2[c_0 c_2 - c_0^2 - c_1^2 + c_1 c_{-1}]$. The identities $c_0 c_2 = \frac{1}{2}[\cos(2kx + 2k + 2\phi) + \cos 2k]$, $c_1 c_{-1} = \frac{1}{2}[\cos(2kx + 2\phi) + \cos 2k]$, $c_0^2 = \frac{1}{2}[1 + \cos(2kx + 2\phi)]$ and $c_1^2 = \frac{1}{2}[1 + \cos(2kx + 2k + 2\phi)]$ cancel the phase-dependent parts, leaving $G = A^2(\cos 2k - 1) = -2\,A^2\sin^2 k$ at every Node and interval. So the Link's booking is $-\num\,G = 2\,\num\,A^2\sin^2 k$ and the Node's tension $T_{aa} = -\num\,h_a = \num\,A^2\sin^2 k$ with $h_a = G / 2$ on a plane wave, where the two ends' tensions are equal, and for a wave with the components $k_a$, $k_b$ the cross reading is $G_{ab} = -2\,A^2\sin k_a\sin k_b$.
 \emph{Status:} derived; a theorem for the line's identity (the plane wave's stress reading phase-free at every Node), derived for the momentum flux (S.45); lattice.
\end{derivation}

\section*{Measurement (Section 2.5)}

\begin{derivation}[the region's reading factor and the draw's scatter, Section 2.5]
\emph{Inputs:} the screen's reading factor over $n$ consecutive rows of a fringe of period $\Lambda$ (the paper's Section 2.5); the draw's binomial scatter (Section 2.5). \emph{Steps:} a region of $n$ consecutive rows sums a fringe $\cos(2\pi y / \Lambda)$, $\Lambda$ here the fringe's period, to $\mathrm{Re}\sum_{y = 0}^{n - 1} e^{2\pi i y / \Lambda} = \mathrm{Re}\,\frac{1 - e^{2\pi i n / \Lambda}}{1 - e^{2\pi i / \Lambda}}$, whose modulus is $|\sin(\pi n / \Lambda) / \sin(\pi / \Lambda)|$, the sum being $e^{i\pi(n - 1) / \Lambda}\sin(\pi n / \Lambda) / \sin(\pi / \Lambda)$; over $n$ rows the fringe's modulation is kept with the signed factor $\sin(\pi n / \Lambda) / (n\sin(\pi / \Lambda))$ (negative, the fringe inverted, for $\Lambda < n < 2\Lambda$), $0.91$ at $n = 4$, $\Lambda = 16$, and $0$ at $\Lambda = n$. A frequency read as a click rate over a window of $n_w$ intervals resolves $\Delta\omega \ge 2\pi / n_w$, the window's own resolution. $N$ quanta drawn independently by the shares $P_i$ give each region a binomial count, $NP_i \pm \sqrt{NP_i(1 - P_i)}$. \emph{Check:} $\sin(\pi / 4) / (4\sin(\pi / 16)) = 0.70711 / (4 \times 0.19509) = 0.9061$.
 \emph{Status:} derived; clicks.
\end{derivation}

\begin{derivation}[the Wronskian's conservation and its current, Section 3.5]
\emph{Inputs:} the two parts of a plane record step by the same read matrix $\mathbf M$, with $\mathbf D\mathbf M$ symmetric for $\mathbf D = \mathrm{diag}(1 / p_i^2)$ (the paper's Theorem 3), $\mathbf{re}_{\mathrm{next}} = \mathbf M\,\mathbf{re}_{\mathrm{now}} - \mathbf{re}_{\mathrm{before}}$ and the same for $\mathbf{im}$. \emph{Steps:} let $W = \mathbf{re}_{\mathrm{now}}^\top\mathbf D\,\mathbf{im}_{\mathrm{before}} - \mathbf{im}_{\mathrm{now}}^\top\mathbf D\,\mathbf{re}_{\mathrm{before}}$. After the step it is $\mathbf{re}_{\mathrm{next}}^\top\mathbf D\,\mathbf{im}_{\mathrm{now}} - \mathbf{im}_{\mathrm{next}}^\top\mathbf D\,\mathbf{re}_{\mathrm{now}} = (\mathbf M\mathbf{re}_{\mathrm{now}} - \mathbf{re}_{\mathrm{before}})^\top\mathbf D\,\mathbf{im}_{\mathrm{now}} - (\mathbf M\mathbf{im}_{\mathrm{now}} - \mathbf{im}_{\mathrm{before}})^\top\mathbf D\,\mathbf{re}_{\mathrm{now}}$; the $\mathbf M$ terms are $\mathbf{re}_{\mathrm{now}}^\top\mathbf M^\top\mathbf D\,\mathbf{im}_{\mathrm{now}} - \mathbf{im}_{\mathrm{now}}^\top\mathbf M^\top\mathbf D\,\mathbf{re}_{\mathrm{now}}$, which vanish since $\mathbf M^\top\mathbf D = \mathbf D\mathbf M$ is symmetric, and the rest is $W$ again. Per Node the same cancellation leaves the Link terms: the weighted Wronskian $W_i / p_i^2$ changes by $(2\,\num / w)\,(q_{ij}^2 / \Gamma^2)\,(\mathrm{im}_i\,\mathrm{re}_j - \mathrm{re}_i\,\mathrm{im}_j)$ summed over the Node's six Ports at the levels now, the Wronskian's current, which for the integer $W_i$ at the vacuum's paces is $(\num / 3\,\den)\,(\mathrm{im}_i\,\mathrm{re}_j - \mathrm{re}_i\,\mathrm{im}_j)$ per Port. \emph{Check:} a periodic chain of three Nodes at $\Gamma = 2$ with the two other axes folded, $[1, 1]$, $\mathrm{re}_{\mathrm{now}} = (0, 3, 0)$, $\mathrm{im}_{\mathrm{now}} = (3, 0, 0)$ and zeros before steps to $(1, 4, 1)$ and $(4, 1, 1)$ with no remainder; $W_0$ goes from $0$ to $3$, $W_0 / p_0^2$ by $3 / 4$, and the current is $(2 / 24)\cdot 9 = 3 / 4$; exact fractions on a chain of five Nodes with contents and Link factors confirm the identity at every Node. For the circular record $\mathrm{re} = A\cos\omega t$, $\mathrm{im} = eA\sin\omega t$ at one Node, with $e$ here the ratio of the two parts' amplitudes, $W = A\cos\omega t\cdot eA\sin\omega(t - 1) - eA\sin\omega t\cdot A\cos\omega(t - 1) = -eA^2\sin\omega$, constant. \emph{Domain:} the identity is the line's, $\mathbf{re}_{\mathrm{next}} = \mathbf M\,\mathbf{re}_{\mathrm{now}} - \mathbf{re}_{\mathrm{before}}$ exactly. The integer step's division adds $\epsilon = (r - r') / w$ to each part, and the remainder term $\epsilon_{\mathrm{re}}\,\mathrm{im}_{\mathrm{now}} - \epsilon_{\mathrm{im}}\,\mathrm{re}_{\mathrm{now}}$ is not zero in general: at the vacuum's pace $S = 0$ and $w\,\anext = 12\,\num\,\Gamma^2\anow - w\,\abefore$ for a uniform level, so at $[2, 3]$ the parts $(1, 3)$ with zeros before step to $(\lfloor 4 / 3\rfloor, \lfloor 4\rfloor) = (1, 4)$ and $W$ goes from $0$ to $1\cdot 3 - 4\cdot 1 = -1$. An exactly conserved integer charge is open.
 \emph{Status:} theorem, its witnesses computed; lattice.
\end{derivation}

\begin{derivation}[Bell's correlation under the declared credit, Section 7.2]
\emph{Inputs:} two parts with the products $M_1$, $M_2$ of their amplitudes at the two sides' ports; the port coefficients $e(+) = (\cos a, \sin a)$, $e(-) = (-\sin a, \cos a)$ on side A and the same with $b$ on side B; the joint share $J(o_A, o_B) = \bigl(\sum_k e_k(o_A)\,e_k(o_B)\,M_k\bigr)^2$ over the two sides' ports $o_A$, $o_B$; $E = \sum_{o_A, o_B} s(o_A)s(o_B)\,J(o_A, o_B) / \sum_{o_A, o_B} J(o_A, o_B)$ with $s(\pm) = \pm 1$. \emph{Steps:} write $c_a = \cos a$, $s_a = \sin a$ and the same for $b$. The four sums are $M_1 c_a c_b + M_2 s_a s_b$, $M_1 s_a s_b + M_2 c_a c_b$, $-M_1 c_a s_b + M_2 s_a c_b$ and $-M_1 s_a c_b + M_2 c_a s_b$ for $(+, +)$, $(-, -)$, $(+, -)$ and $(-, +)$. Their squares sum to $(M_1^2 + M_2^2)(c_a^2 c_b^2 + s_a^2 s_b^2 + c_a^2 s_b^2 + s_a^2 c_b^2) = M_1^2 + M_2^2$, the cross terms cancelling in pairs. The signed sum is $(M_1^2 + M_2^2)(c_a^2 - s_a^2)(c_b^2 - s_b^2) + 8M_1M_2\,c_a s_a c_b s_b$, so
\[
E(a, b) = \cos 2a\cos 2b + \rho\sin 2a\sin 2b, \qquad \rho = \frac{2M_1M_2}{M_1^2 + M_2^2},
\]
and the marginal $P(A+) = [(M_1 c_a c_b + M_2 s_a s_b)^2 + (-M_1 c_a s_b + M_2 s_a c_b)^2] / (M_1^2 + M_2^2) = (M_1^2 c_a^2 + M_2^2 s_a^2) / (M_1^2 + M_2^2)$, independent of $b$. For equal parts $\rho = 1$ and $E = \cos 2(a - b)$, the known form for parallel polarisations, derived here under the declared credit by Lagrange's identity. \emph{The integer settings:} with $e(+) = (p, q)$ and $e(-) = (-q, p)$ unnormalised and equal parts, $J(+, +) = J(-, -) = (p_Ap_B + q_Aq_B)^2$ and $J(+, -) = J(-, +) = (p_Aq_B - q_Ap_B)^2$, whose sum is $2(p_A^2 + q_A^2)(p_B^2 + q_B^2)$ by Lagrange's identity, so $E = [(p_Ap_B + q_Aq_B)^2 - (p_Aq_B - q_Ap_B)^2] / [(p_A^2 + q_A^2)(p_B^2 + q_B^2)]$, a rational number. At $a = (1, 0)$, $a' = (1, 1)$, $b = (12, 5)$, $b' = (5, 12)$: $E(a, b) = (144 - 25) / 169 = 119 / 169$, $E(a, b') = (25 - 144) / 169 = -119 / 169$, $E(a', b) = (17^2 - 7^2) / 338 = 120 / 169$, $E(a', b') = 120 / 169$, and $S = E(a, b) - E(a, b') + E(a', b) + E(a', b') = 478 / 169 = 2.8284$. At the ideal angles $S = \sqrt 2(1 + \rho)$; at the integer settings $S = (238 + 240\rho) / 169$. \emph{The two local credits:} with the pair's phase $\theta$ uniform, a credit of each side by the shares of its ports gives $E = \langle\cos 2(a - \theta)\cos 2(b - \theta)\rangle_\theta = \frac{1}{2}\cos 2(a - b)$ and $S = \sqrt 2$; a credit by the sign of the larger port gives $E = \langle\mathrm{sgn}\cos 2(a - \theta)\,\mathrm{sgn}\cos 2(b - \theta)\rangle_\theta = 1 - 4|a - b| / \pi$, the sawtooth, whose sum at the four settings is exactly $2$. \emph{Check:} the averages over $\theta$ on a mesh of $2 \times 10^5$ points give $S = 1.4142$ and $2.0000$, and the contrast credit, one side's count weighted by its contrast $|\cos 2(a - \theta)|$ and the other side's sign always counted, gives $E = \cos 2(a - b)$ exactly once divided by one side's efficiency $2 / \pi$, by $\langle\cos x\,\mathrm{sgn}\cos(x - 2\delta)\rangle = (2 / \pi)\cos 2\delta$, $S = 2.8284$ ($478 / 169 = 2.828402$ against $2\sqrt 2 = 2.828427$); the symmetric readings do not ($|\cos|$ on both sides: $E(\pi / 8) = 0.88$, $S = 3.52$).
 \emph{Status:} derived under the declared credit, computed for the numbers ($478 / 169$); clicks.
\end{derivation}

\section*{The families (Section 4.1)}

\begin{derivation}[the reach of a holder with a gap, Section 4.1]
\emph{Inputs:} a holder's rest outside its source, the line with $\anext = \abefore = \anow = a$: $6\,\den\,a = \num\,S_6(a)$ at the vacuum's pace, that is $\num\,(S_6 - 6)\,a = 6(\den - \num)\,a$. \emph{Steps:} for a rest falling along one axis as $e^{-\kappa r}$ and uniform across, $S_6(a) = (2\cosh\kappa + 4)\,a$, so $6\,\den = \num\,(2\cosh\kappa + 4)$ and $\cosh\kappa = 3\,\den / \num - 2$, for $0 < \num \le \den$ (the positive tail; the axial characteristic equation is $\lambda + \lambda^{-1} = 6\,\den / \num - 4$ for the ratio $\lambda$ of neighbouring levels, and at a negative numerator its decaying root is negative, an alternating tail, $\lambda = -8 + \sqrt{63} = -0.0627$ at $[-1, 2]$ with $\cosh\kappa = 8$ for $\kappa = -\ln|\lambda|$; at $\num = 0$ the static line forces $a = 0$ outside the source, no tail). For $\num$ just under $\den$, $\cosh\kappa = 1 + \kappa^2 / 2 + \dots$, so $\kappa^2 = 6\,\den / \num - 6 = 6(\den - \num) / \num$ and the range is $R = 1 / \kappa = \sqrt{\num / (6(\den - \num))}$ Links: $20.00$ at $[2400, 2401]$; $[1, 2]$ has $\cosh\kappa = 4$, $e^{-\kappa} = 0.127$. The gap is also the rest rotation of the holder's own travelling events, $\cos\omega_h = \num / \den$, $\omega_h^2 = 2(\den - \num) / \den + \dots$, whose Compton length at light's speed, $c / \omega_h = \sqrt{\den / (6(\den - \num))} = R\sqrt{\den / \num}$, is the range to the gap's first order: Yukawa's identity with no input.
 \emph{Status:} derived; lattice.
\end{derivation}

\begin{derivation}[the two speeds, Section 4.1]
\emph{Inputs:} S.2's $N^2 = e^{-2x}$ and the law's read of the content by every family. \emph{Steps:} light at a content $x$ travels at $N^2 / \sqrt 3 = (1 - 2x) / \sqrt 3 + O(x^2)$, $N^2 = e^{-2x}$ (Derivation S.2). Under the law's line the massless holder reads the content as every family does, its own level among it, so its travelling events have the same paces and the two arrive together at every content, the difference of their Shapiro delays $0$. Were the massless holder to read no content, its events travel at $1 / \sqrt 3$ and over $L$ Links the two arrive $L\sqrt 3\,(e^{2x} - 1) = 2xL\sqrt 3 + O(x^2)$ intervals apart at long wavelength, against GW170817's bound \cite{gw170817} on $\gamma_{\mathrm{GW}} - \gamma_{\mathrm{EM}}$ (the gravitational wave's against the electromagnetic wave's), between $-2.6 \times 10^{-7}$ and $1.2 \times 10^{-6}$; under the line the measurement bounds no vacuum content.
 \emph{Status:} derived; clicks.
\end{derivation}

\section*{Bodies (Section 4)}

\begin{derivation}[the bound body: the fixed point and the functional, Section 4.2]
\emph{Inputs:} Rule3's line at the paces of the wells the record sources; a standing record $a_{i, t} = \phi_i\cos(\omega_b t)$ with one rotation at every Node; the clock's pace $(p_0 / \Gamma)^2 = 1 - 2\ell_i / \Gamma + \dots$ at a well $\ell_i$ (in levels, $\ell_i = \Gamma U_i$; Newton's potential is $\Phi = -c^2 U$, S.28). \emph{The fixed point's equation:} with the remainders dropped the line gives at every Node
\[
2\cos\omega_b\,\phi_i = \frac{1}{w}\Bigl[\sum_a R_a(i)\,(\phi_{i+a} + \phi_{i-a}) + S_i\,\phi_i\Bigr] = (\mathbf M\phi)_i,
\]
so $\phi$ is an eigenvector of the read matrix $\mathbf M$ and the standing record is its top mode, the maximiser of $\phi^\top\mathbf D\mathbf M\phi$ on the weighted sphere $\phi^\top\mathbf D\phi = 1$, $\mathbf D\mathbf M$ symmetric under the law's read coefficient (the paper's Theorem 3) (the rule's own iteration of Online Resource 1, Section B.3 repeats this read act). To the first order in $\ell / \Gamma$, $S_i / w = 2 - 2g(1 - 2\ell_i / \Gamma) - \frac{2\,\num}{3\,\den}\sum_a (p_a / \Gamma)^2 = 2\cos\omega_0 + 4g\ell_i / \Gamma - \frac{2\,\num}{3\,\den}\sum_a(p_a / \Gamma)^2$ and $R_a / w = \frac{\num}{3\,\den}(p_a / \Gamma)^2$, so, keeping the band's term at the vacuum's Link pace (to first order in $U$ for a body whose band term balances its well, $\Delta\phi = O(U\phi)$, so that the product of the two is of the second order; a free packet in a well, with $\Delta\phi$ of the packet's own size, is outside this reduction),
\[
(\mathbf M\phi)_i = 2\cos\omega_0\,\phi_i + \frac{4g}{\Gamma}\,\ell_i\,\phi_i - \frac{\num}{3\,\den}\,(-\Delta\phi)_i, \qquad (-\Delta\phi)_i = 6\phi_i - \sum_a(\phi_{i+a} + \phi_{i-a}).
\]
\emph{The functional:} the well is sourced by the record itself, $\ell_i = \sum_r \sigma_r (C / E_r)\sum_j G_{r, ij}\,\phi_j^2$ for a body of $C$ quanta with the counts $C_j = C\phi_j^2$. Let
\[
F[\phi] = \frac{\num}{3\,\den}\sum_{\text{Links}}(\phi_i - \phi_j)^2 - \frac{2g}{\Gamma}\sum_r \sigma_r\frac{C}{E_r}\sum_{ij}\phi_i^2\,G_{r, ij}\,\phi_j^2,
\]
the functional of C.2. Its gradient is $\partial F / \partial\phi_i = \frac{2\,\num}{3\,\den}(-\Delta\phi)_i - \frac{8g}{\Gamma}\ell_i\phi_i$, using the symmetry of $G_r$ (each well counted once, Hartree's form, is what makes the factor $2g$ in $F$ give $4g$ in the gradient). So $\partial F / \partial\phi = 2\lambda\phi$ on the sphere is exactly the fixed point's equation with $\lambda = 2\cos\omega_0 - 2\cos\omega_b$: a body is a critical point of $F$ on the sphere, and the virial and the second variation along the dilation are the paper's proposition (Section 4.2). \emph{The scaling powers:} under $\phi_R(\mathbf x) = R^{-d / 2}\phi(\mathbf x / R)$ the sum over Links of $(\phi_i - \phi_j)^2$ is the lattice form of $\int|\nabla\phi|^2$ and scales as $R^{-2}$, the pressure $K / R^2$. The well energy $\sum_{ij}\phi_i^2 G_{ij}\phi_j^2$ is $\int\!\!\int\rho(\mathbf x)G(\mathbf x - \mathbf y)\rho(\mathbf y)$ with $\rho = \phi^2$ of unit mass, which scales as $R^{-(d - 2)}$ for the massless kernel $G \sim 1 / r$ and for a screened kernel within its reach, and as $R^{-d}$ for a screened kernel beyond its reach, where $G$ acts as a point kernel and the energy is $\int\rho^2$: the powers $s_r$ of the paper's proposition. \emph{The cloud's size:} for the massless kernel in three dimensions $F(R) = a / R^2 - c / R$ with $a = \frac{\num}{3\,\den}\int|\nabla\phi|^2$ and $c = \frac{2g}{\Gamma}\frac{C}{E}\int\!\!\int\frac{\rho\rho'}{4\pi|\mathbf x - \mathbf y|}\cdot 3$ at the unit size, so the minimum is at $R = 2a / c \propto \Gamma E_g / (gC)$: the radius falls as $1 / C$ and grows with $\Gamma$; the continuum constant of the root-mean-square (rms) radius is $(2\pi\cos\omega_0 / (9g))\,r_P\,\Gamma E_g / C$ with Pekar's $r_P = 4.635$ \cite{pekar} (the Choquard minimiser's rms radius, at Lieb's energy $-0.1085$ \cite{lieb} and the virial $1 / 2$), $(4\pi / 9)\,r_P = 6.47$ at $[2, 3]$ where $\cos\omega_0 / (2g) = 1$, the factor $\cos\omega_0 / (2g)$ tending to $1 / \omega_0^2$ as the gap closes ($4.5$ at $[9, 10]$, $500$ at $[999, 1000]$). \emph{The hollow's barrier (Section 4.2):} a hollow beyond the body's reach adds the point-kernel term $-b / R^3$ (the power $R^{-d}$ above at $d = 3$), $F(R) = a / R^2 - b / R^3 - c / R$; $F'(R) = (cR^2 - 2aR + 3b) / R^4$ vanishes at $R_\pm = (a \pm \sqrt{a^2 - 3bc}) / c$, two critical points where $a^2 > 3bc$ and none where $3bc > a^2$; at a root $F'' = 2(aR - 3b) / R^5$, positive at $R_+ \ge a / c > 3b / a$ and negative at $R_- = 3b / (a + \sqrt{a^2 - 3bc}) < 3b / a$, so the wide critical point is the minimum and the narrow one the saddle. \emph{The running exponent of a screened kernel (the proposition's commentary):} for a Gaussian body of width $R$ the integrated well energy is $(C / R)\,I(\kappa R)$ with $I(z) = \int_0^\infty u\,e^{-u^2 / 4 - zu}\,du$, so its exponent is $s = 1 + z\,\langle u^2\rangle / \langle u\rangle$ in that weight: $1 + \sqrt\pi\,\kappa R + O((\kappa R)^2)$ at a small $\kappa R$, $3 - 3 / (\kappa R)^2 + O((\kappa R)^{-4})$ at a large one, $2.130$ at $\kappa R = 1$ and $2.971$ at $10$, crossing $2$ at $\kappa R = 0.82$; the pair kernel's log-derivative at one separation, $1 + \kappa R$, is another number (three derivations). 
The equation is a theorem (the functional's variational derivative); the two branches are a lattice reading named in C.3 and not computed here, the remainder's budget clause in S.53 (i); the fence is lattice, the integer version open.
 \emph{Status:} theorem for the critical point, the fixed point computed and the continuum constants; the stability along the dilation derived; lattice.
\end{derivation}

\begin{derivation}[the invariant and the count's drift, Eq.~(7)]
\emph{Inputs:} the one-Node line $a_{t+1} + a_{t-1} = 2\cos\omega_t\,a_t$ with $\omega_t$ changing slowly. \emph{Steps:} at fixed $\omega$ the form $D_t = a_t^2 - a_{t+1}a_{t-1} = a_t^2 + a_{t-1}^2 - 2\cos\omega\,a_t a_{t-1}$ is conserved and equals $A^2\sin^2\omega$ for $a_t = A\cos(\omega t + \phi)$; in the plane $(a_t, a_{t-1})$ the step is a linear map of determinant $1$ whose orbits are the ellipses $D = \text{const}$, of area $\pi D / \sin\omega$. A slow change of $\omega$ keeps the area enclosed by the orbit, the action of the map, so $D / \sin\omega = A^2\sin\omega$ is the adiabatic invariant to the first order in the slow variation: the discrete form of $E / \omega$. Equivalently, the ansatz $a_t = C\cos(\Theta_t + \phi) / \sqrt{\sin\omega_t}$ with the midpoint phase $\Theta_t = \sum_{s < t}(\omega_s + \omega_{s+1}) / 2$, advancing by $(\omega_t + \omega_{t+1}) / 2$ between $t$ and $t + 1$, satisfies the line to the second order in $\omega_{t+1} - \omega_t$ (the phase $\sum_{s \le t}\omega_s$ fails at the first order; two derivations). \emph{Check:} the recurrence stepped $20{,}000$ intervals with $\omega$ lowered from $0.841$ to $0.600$ keeps $D / \sin\omega$ within $4 \times 10^{-5}$ of its value while $D$ falls by a quarter. \emph{The drift:} the share at a Node in quanta is $(\Gamma^2 / p_i^2)\,A^2\sin^2\omega_b / T$ (Eq.~(5) at the Node's paces), so with $A^2\sin\omega_b$ fixed, $e_i(t_2) / e_i(t_1) = [\sin\omega_b(t_2) / \sin\omega_b(t_1)]\,[p_i^2(t_1) / p_i^2(t_2)]$, the drift ratio of the paper's Section 4.3, the Node's share of Eq.~(5). \emph{Planck's relation:} the share $A^2\sin^2\omega / T$ at the vacuum's paces is the invariant $I = A^2\sin\omega / T$ times $\sin\omega$. \emph{The body's click:} under the declaration that one click is one quantum's change in the absorbing body (Section 2.5), the NodeDetector's own record answers at its rotation by the two-mode line of S.16, one transition per click from the mode $\omega_i$ to the mode $\omega_j$, and the share it takes is $T(\sin\omega_j - \sin\omega_i) = T\sin(\omega_j - \omega_i)\,[1 - \omega_i\omega_j / 2 + O(\omega^4)]$, by $\sin\omega_j - \sin\omega_i = 2\cos\frac{\omega_i + \omega_j}{2}\sin\frac{\omega_j - \omega_i}{2}$: one photon's share $T\sin\omega_L$ at the line $\omega_L = \omega_j - \omega_i$, to a relative $\omega_0^2 / 2 < 10^{-18}$ at nature's gap (S.27). So a body counts photons, one per quantum since $2I = 1$ ($I = 1 / 2$ per quantum under $2A^2\sin\omega = T$, C.9 (a)), and reads $\hbar\omega$ per click with $T$ for $\hbar$; the share is the energy a click carries, and through a deepening well the number of quanta, $2I$, is kept.
 \emph{Status:} derived to the first order in the slow variation, the invariant computed, the body's click under the declaration of one transition per click; lattice and clicks.
\end{derivation}

\begin{derivation}[the clock's click rate, Section 4.3]
\emph{Inputs:} Eq.~(3) at $k = 0$ with $p_0^2 = \Gamma^2(1 - 2U + 2U^2)$, the composed paces, to the second order in $U$, where they agree with the composed paces of S.2: $\cos\omega(U) = 1 - g(1 - 2U + 2U^2) = \cos\omega_0 + 2gU - 2gU^2$; S.2's light speed $N^2 / \sqrt 3$ for the light clock; the one assumption $hN = 1$ (Section 2.3); S.31's scaling for the second order. \emph{Steps:} to first order $d(\cos\omega) = 2g\,dU = -\sin\omega_0\,d\omega$, so $d\omega / \omega_0 = -[2g / (\omega_0\sin\omega_0)]\,dU = -f(\omega_0)\,dU$ with $f(\omega_0) = 2(1 - \cos\omega_0) / (\omega_0\sin\omega_0)$; $f = 1.063$ at $[2, 3]$, and as $\omega_0 \to 0$, $1 - \cos\omega_0 \to \omega_0^2 / 2$ and $f \to 1$, the rate $1 - U$ of the clock's pace alone. \emph{The light clock:} a packet between two declared faces $L$ Links apart at the content $x$ has the round trip $2L / v_g = 2L\sqrt 3 / N^2$ (Derivation S.2), so its rate falls by $2x$ to the first order; between two bodies, which stand $L / h$ Links apart under the one assumption of Section 2.3, the round trip is $2L\sqrt 3 / N$ (Derivation S.31).
 \emph{Status:} derived under the one assumption; clicks.
\end{derivation}

\begin{derivation}[the resonance at twice the rotation, Section 4.3]
\emph{Inputs:} the one-Node line with the rotation modulated by a level $b_L\cos\omega_L t$ read into the content, the law's two-mode line; light's declared read is the holder's rotation (Sections 3.5 and 4.2), under which S.41 derives the transfer line by the two-mode coupling, this derivation's content-modulation line being the pace read's own, with no analogue under the rotation (S.41's Status): $2\cos\omega(t) = 2 - 2g\,(1 - 2x + 2x^2)$ with $x = x_0 + b_L\cos(\omega_L t) / \Gamma$, so $d(2\cos\omega) / dx = 4g(1 - 2x)$ and $\anext + \abefore = [2\cos\omega_b + \epsilon\cos(\omega_L t + \vartheta)]\,\anow$ with $\epsilon = 4g(1 - 2x_0)\,b_L / \Gamma$. \emph{Steps:} write $a_t = \mathrm{Re}[z_t e^{i\omega_b t}]$ with $z$ slowly varying, $z_{t \pm 1} = z_t \pm z'$. Then $a_{t+1} + a_{t-1} = \mathrm{Re}[e^{i\omega_b t}(2\cos\omega_b\,z_t + 2i\sin\omega_b\,z')]$, and the product $\epsilon\cos(\omega_L t + \vartheta)\,\mathrm{Re}[z_t e^{i\omega_b t}]$ has, at $\omega_L = 2\omega_b$, the resonant part $(\epsilon / 2)\,\mathrm{Re}[\bar z_t e^{i\vartheta} e^{i\omega_b t}]$. Equating the slowly varying coefficients of $e^{i\omega_b t}$, $2i\sin\omega_b\,z' = (\epsilon / 2)\,e^{i\vartheta}\bar z$, so $z' = -i\epsilon e^{i\vartheta}\bar z / (4\sin\omega_b)$; this antilinear map $z \mapsto \lambda\bar z$ with $|\lambda| = \mu$ has the real eigenvalues $\pm\mu$ on two lines of the $z$ plane, for $\omega_b$ away from $\pi / 2$ and $\omega_L$ off every other resonance, where the non-resonant terms average out; at $\omega_b = \pi / 2$ the drive $\omega_L = \pi$ aliases both sidebands onto the conjugate and the two-step transfer, of trace $-2 - \epsilon^2\cos^2\vartheta$ at the phase $\vartheta$ of the kick against the mode, grows at $\operatorname{arsinh}(\epsilon|\cos\vartheta| / 2)$ per interval, $\operatorname{arsinh}(\epsilon / 2)$ at $\vartheta = 0$ and nothing at $\vartheta = \pi / 2$, twice the formula, so the amplitude grows as $e^{\mu t}$ along one line and the share, quadratic in the amplitude, as $e^{2\mu t}$. Detuned by $\Delta = \omega_L - 2\omega_b$, the substitution $z = \zeta e^{i\Delta t / 2}$ gives $\zeta' = -i(\Delta / 2)\zeta - i\mu e^{i\vartheta}\bar\zeta$, whose exponents are $\pm\sqrt{\mu^2 - \Delta^2 / 4}$: growth only for $|\Delta| < 2\mu$, a full width $\epsilon / \sin\omega_b$. \emph{Check:} the recurrence stepped $20{,}000$ intervals at $2\cos\omega_b = 1.4317$ and $\epsilon = 0.013$ grows at $0.004653$ to $0.004656$ per interval over seven modulation phases, the formula's first-order $0.004654$ to four digits. \emph{The two-mode line:} with two bound modes the same modulation multiplies $a_t = \mathrm{Re}[z_i e^{i\omega_i t} + z_j e^{i\omega_j t}]$, and its product with $\cos\omega_L t$ carries the frequencies $\omega_i \pm \omega_L$. The term resonant with the mode $j$ from the mode $i$ is at $\omega_L = \omega_j - \omega_i$, a transfer between the two modes, and the term that drives both from the conjugate of the other is at $\omega_L = \omega_i + \omega_j$, of which the one-mode case above is the degenerate case $i = j$ (at $i = j$ the difference line is $\omega_L = 0$, a static shift and no line). The transfer's rate, $\Omega_R = \epsilon_{ij} / (4\sqrt{\sin\omega_i\sin\omega_j})$, follows by the same slowly varying coefficients with the two modes' $2i\sin\omega_i z_i'$ and $2i\sin\omega_j z_j'$ coupled through $\epsilon_{ij}$, a Rabi oscillation.
 \emph{Status:} two derivations; derived; clicks.
\end{derivation}

\begin{derivation}[the cross current of two standing records, Section 4.3]
\emph{Inputs:} two standing records $a_t = A\cos(\omega_1 t + \alpha)$ and $b_t = B\cos(\omega_2 t + \beta)$ at the two ends of a Link; the static tail line of S.11, $\cosh\kappa_b = 3\cos\omega_b / \cos\omega_0 - 2$ at the bodies' rotation; one quantum's current measure $W_c\sin\omega_0$ (Section 2.5). \emph{Steps:} the current $\num\,(a_t b_{t-1} - a_{t-1} b_t)$ is $\num AB[\cos(\omega_1 t + \alpha)\cos(\omega_2 t - \omega_2 + \beta) - \cos(\omega_1 t - \omega_1 + \alpha)\cos(\omega_2 t + \beta)]$, a sum of cosines of $(\omega_1 - \omega_2)t$ and $(\omega_1 + \omega_2)t$ with the constant parts $\frac{1}{2}\cos(\alpha - \beta + \omega)$ and $-\frac{1}{2}\cos(\alpha - \beta - \omega)$ only when $\omega_1 = \omega_2$; for $\omega_1 \ne \omega_2$ its mean over $n$ beats is $0$ to the order of the fast term over the window, $1 / ((\omega_1 + \omega_2)\,n\,T_b)$ with $T_b$ the beat's length, so no quantum passes at the first order. \emph{The reach:} two bodies' tails are $a_1 e^{-\kappa_b d_1}$ and $a_2 e^{-\kappa_b d_2}$, $\kappa_b$ the bodies' own inverse reach at $\omega_b$ ($\cosh\kappa_b = 3\cos\omega_b / \cos\omega_0 - 2$), at a Link $d = d_1 + d_2$ Links from the two centres, so the current there is of the order $a_1 a_2 e^{-\kappa_b d}$ per interval and falls under one quantum's measure of a current for a gapped family, $\Wc\sin\omega_0$, at $d = \ln(\num\,a_1 a_2 / (3\,\den\,T\sin\omega_0)) / \kappa_b$.
 \emph{Status:} derived; lattice.
\end{derivation}

\section*{The forms of nature (Section 5.1)}

\begin{derivation}[the fall and the Newtonian potential the orbits measure, Eq.~(8)]
\emph{Inputs:} Eq.~(3) at a content $U$ for a reader at the weight $1$, $\cos\omega(k, U) = 1 - g(1 - 2U + 2U^2) - \frac{\num}{3\,\den}(1 - 2U)^2(1 - \cos k)$ along the motion, the quadratic paces to the first order in $U$, where they agree with the composed paces of S.2. \emph{Steps:} at small $k$ and $U$, $\omega(k, U) = \omega_0(U) + k^2 / (2m^*)$ with the inertia $m^* = 3\,\den\sin\omega_0 / \num = 3\tan\omega_0$ (from $\sin\omega_0\,d\omega = \frac{\num}{3\,\den}\,k\,dk$), so the group velocity is $v = k / m^*$. The rays of $\omega(k, x)$, $x$ the coordinate here, give $dk / dt = -\partial\omega / \partial x = -(d\omega_0 / dU)\,\nabla U = (2g / \sin\omega_0)\,|\nabla U|$ toward the well (Derivation S.15), so the acceleration is
\[
a = \frac{1}{m^*}\frac{dk}{dt} = \frac{\num}{3\,\den\sin\omega_0}\cdot\frac{2g}{\sin\omega_0}\,|\nabla U|
= \frac{2\cos\omega_0(1 - \cos\omega_0)}{3\sin^2\omega_0}\,|\nabla U| = \frac{2\cos\omega_0}{3(1 + \cos\omega_0)}\,|\nabla U|,
\]
which in the gap $g = 1 - \cos\omega_0$ is $[(1 - g) / (3(1 - g / 2))]\,|\nabla U|$: $0.2667$ at $[2, 3]$ and $1 / 3$ as the gap closes. Nature's fall is $a = c^2\,|\nabla U_K|$ with $U_K = G_K M / (c^2 r)$ and $c^2 = 1 / 3$ here, so $U_K = [2\cos\omega_0 / (1 + \cos\omega_0)]\,U$, Eq.~(8), $0.800$ at $[2, 3]$.
 \emph{Status:} derived; clicks.
\end{derivation}

\begin{derivation}[the post-Newtonian parameters against the Newtonian potential the orbits measure, Section 5.1]
\emph{Inputs:} Eq.~(8), $U = \lambda\,U_K$ with the ratio $\lambda = U / U_K = (1 + \cos\omega_0) / (2\cos\omega_0) = (\den + \num) / (2\num)$ for $0 < \num \le \den$ ($\cos\omega_0 > 0$, the map $U_K / U = 2\num / (\den + \num)$ inverted; at $\num = \den$ the parameters are the limit $1$, $1$, $1$ with no bound orbit; at $[0, 1]$ $U_K = 0$ and the column is empty; a negative numerator defines none); the composed clock's metric $N^2 = e^{-2U}$, $h^2 = e^{2U}$ (S.34); the three forms in $U$: the bending $4U_b$, the periapsis $6\pi U_1 / (1 - e^2)$ with $U_1 = \mu / a$ (the post-Newtonian formula, standard \cite{will}), the redshift $f(\omega_0)\,U$. \emph{Steps:} (the orbit step) a slow packet of matter moves by the geometric optics of its band at the composed paces, whose rays are the geodesics of the metric $N^2 = e^{-2U}$, $h^2 = e^{2U}$ in the time rescaled by $1 / \cos(\omega / 2)$, the band being Klein-Gordon in $s = 2\sin(\omega / 2)$ (S.31); at the computed pair the rescaling moves Kepler's $G_K M$ between orbits by $\tan^2(\omega_0 / 2)\,U = 0.2\,U$, so $\gamma_K$ and $\beta_K$ against $U_K$ are defined once the calibrating orbit is named; at S.54's bound, $\omega_e \le 9.6 \times 10^{-15}$, the rescaling's shift $\tan^2(\omega_e / 2)$ is $2.3 \times 10^{-29}$ and its factor $1 / \cos(\omega_e / 2)$ is $1$ to $1.2 \times 10^{-29}$ and the rays are the geodesics; its expansion, $g_{00} = -(1 - 2U + 2U^2)$ to the second order, $g_{ab} = (1 + 2U)\,\delta_{ab}$ to the first and $g_{0a} = 0$ for the static source (the odd part the holder's odd lines carry, sourced by the count's flux, S.56), the orders the periapsis reads, is the post-Newtonian metric with $\gamma = 1$ exactly and $\beta = 1$ in $U$ at the small gap, so the periapsis advances per orbit by Einstein's $6\pi U_1 / (1 - e^2)$ with $U_1 = \mu / a$, that is $6\pi\mu / (a(1 - e^2))$. With the two ends' product $p_0(i)\,p_0(j) / \Gamma$ for the Link the exterior equation is $\nabla \cdot (e^{-4U}\nabla U) = 0$, so $e^{-4U} = 1 - 4m / r$; then $U = u + 2u^2 + O(u^3)$ with $u = m / r$ and the clock's $N^2 = e^{-2U} = 1 - 2u - 2u^2$, so $\beta_K = -1$ against $-g_{00} = 1 - 2u + 2\beta u^2$, and the advance $(2 + 2\gamma - \beta) / 3 = 5 / 3$ at $\gamma = 1$ and $\beta = -1$ in $U$, the small-gap value, $72$ arcseconds per century against $43$ (the paper's Section 2.3). (The three parameters against $U_K$:) the bending is $4\lambda U_{K, b} = 2(1 + \den / \num)\,U_{K, b}$, and the post-Newtonian bending is $(1 + \gamma)\,2GM / (c^2 b) = (1 + \gamma)\,2U_{K, b}$, so $1 + \gamma_K = 1 + \den / \num$, $\gamma_K = \den / \num$. The periapsis is $\lambda$ times Einstein's $6\pi U_{K, 1} / (1 - e^2)$, and the post-Newtonian advance is $(2 - \beta + 2\gamma) / 3$ of Einstein's, so $(2 - \beta_K + 2\,\den / \num) / 3 = (\den + \num) / (2\num)$, whence $\beta_K = 2 + 2\,\den / \num - 3(\den + \num) / (2\num) = (\den + \num) / (2\num)$. The redshift is $f(\omega_0)\lambda\,U_K$, and $f\lambda = \frac{2(1 - \cos\omega_0)}{\omega_0\sin\omega_0}\cdot\frac{1 + \cos\omega_0}{2\cos\omega_0} = \frac{\sin^2\omega_0}{\omega_0\sin\omega_0\cos\omega_0} = \frac{\tan\omega_0}{\omega_0}$, so the redshift's post-Newtonian parameter $\alpha_K$, $z = (1 + \alpha_K)\,U_K$, has $1 + \alpha_K = \tan\omega_0 / \omega_0$. At $[2, 3]$: $1.5$, $1.25$ and $1.329$; all three tend to $1$ as $\omega_0 \to 0$, and $\gamma_K - 1 < 2 \times 10^{-5}$ requires $\den / \num - 1 < 2 \times 10^{-5}$. The push on a moving record, of which the bending is the $\beta = 1$ case, is derived from the same dispersion in S.51; the orbit step above holds for the quadratic band, and at the computed pair the band's $k^4$ term, of the same post-Newtonian order, adds $1 / 8$ of the advance to first order for an orbit in a coordinate plane and $1 / 12$ in the $(111)$ plane, $13$ percent in a coordinate plane, orientation dependent within a percent, and $8$ to $10$ percent in the $(111)$ plane by the ray integration on the exact band, the lattice's $\sum_a k_a^4 / 24$ part of S.21's fourth-order term, computed on the exact band twice ($\mu / p$ from $5 \times 10^{-4}$ to $1.1 \times 10^{-2}$, $p$ the orbit's semi-latus rectum in Links, $e$ from $0.25$ to $0.33$; the pure-axis coefficient's $17$ percent realised by no orbit), vanishing as $2(1 - \cos\omega_0) / \cos\omega_0$ against the advance in $U$ (two recomputations on the exact band).
 \emph{Status:} derived under the one assumption; clicks; the computed pair's values lattice, the ray integrations computed.
\end{derivation}

\begin{derivation}[the bending of light and Shapiro's delay, Section 5.1]
\emph{Inputs:} light's index $n = 1 + 2U$ with $U = \mu / r$ (Derivation S.2). \emph{Steps:} (a) \emph{The delay.} The delay along a straight path at the impact parameter $b$ from $-x_1$ to $x_2$ is $\int(n - 1)\,dl / c = 2\sqrt 3\,\mu\int_{-x_1}^{x_2} dl / \sqrt{l^2 + b^2} = 2\sqrt 3\,\mu\,[\operatorname{arsinh}(x_2 / b) + \operatorname{arsinh}(x_1 / b)]$, which for $x_1, x_2 \gg b$ is $2\sqrt 3\,\mu\ln(4x_1x_2 / b^2)$: Shapiro's delay with $\gamma = 1$ in $U$, and with $\gamma_K$ against $U_K$ by the same factor as the bending, since the bending is the transverse derivative of the same integral. (b) \emph{The bending.} With $n - 1 = 2U = 3s / (2\pi\Gamma r)$, where $s$ is the source's quanta per interval over gravity's divisor $E_g$ (S.25), $r^2 = b^2 + l^2$, the transverse gradient is $\partial n / \partial b = -3s\,b / (2\pi\Gamma\,(b^2 + l^2)^{3/2})$, and $\int_{-\infty}^{\infty} b\,dl / (b^2 + l^2)^{3/2} = 2 / b$, so the turn is $3s / (\pi\Gamma b) = 4U_b$ with $U_b = 3s / (4\pi\Gamma b)$, the index's $2U$ giving twice the Newtonian $2U_b$: Eq.~(9) of the paper.
 \emph{Status:} derived; clicks; the bending the one assumption's calibration and no test.
\end{derivation}

\begin{derivation}[the moving clock to the second order, Section 5.1]
\emph{Inputs:} the exact band of a free record of the pair $[\num, \den]$ along an axis, from Eq.~(1) at the vacuum's pace (S.1): $\cos\omega = c_s^2(2 + \cos k)$ with $c_s^2 = \num / (3\,\den)$, $\cos\omega_0 = 3c_s^2$; the rotation at the moving centre, $\Omega(k) = \omega(k) - k\,\omega'(k)$, the phase rate less $k$ times the group velocity $v = \omega'(k)$, taken as what a clock carried by the packet reads (a definition of the packet's clock, the second order a definition of $c_m$). \emph{Steps:} write $s = \sin\omega_0$, $c = \cos\omega_0$ and $u = 1 - \cos k = k^2 / 2 - k^4 / 24 + \dots$; then $\cos\omega = c - c_s^2 u$, and with $\omega = \omega_0 + \delta$, $\cos\omega = c - s\delta - c\delta^2 / 2 + \dots$, so $s\delta + c\delta^2 / 2 = c_s^2 u$ and
\[
\omega(k) = \omega_0 + \alpha k^2 - \beta k^4 + O(k^6), \qquad \alpha = \frac{c_s^2}{2s}, \qquad \beta = \frac{c_s^2}{24s} + \frac{c\,c_s^4}{8s^3}.
\]
Hence $v = \omega' = 2\alpha k - 4\beta k^3$, $\Omega = \omega - k\omega' = \omega_0 - \alpha k^2 + 3\beta k^4$, and inverting $v(k)$, $k^2 = v^2 / (4\alpha^2) + \beta v^4 / (4\alpha^5) + \dots$, so
\begin{align*}
\frac{\Omega}{\omega_0} &= 1 - \frac{v^2}{4\alpha\omega_0} - \frac{\beta\,v^4}{16\alpha^4\omega_0} + O(v^6) = 1 - \frac{v^2}{2c_m^2} - \frac{2\omega_0}{\sin 2\omega_0}\cdot\frac{v^4}{8c_m^4} + O(v^6),\\
c_m^2 &= 2\alpha\omega_0 = \frac{c_s^2\,\omega_0}{s} = \frac{\omega_0}{3\tan\omega_0},
\end{align*}
where the fourth-order coefficient is $\beta / (16\alpha^4\omega_0)$ over Lorentz's $1 / (8c_m^4)$, which reduces with $3c_s^2 = c$ to $\omega_0(s / c + c / s) = \omega_0 / (\sin\omega_0\cos\omega_0) = 2\omega_0 / \sin 2\omega_0$. So the clock of a free quantum is Lorentz's factor at the family's own light speed $c_m$ to the second order in $v$ exactly, and at the fourth order, along an axis, it carries the factor $2\omega_0 / \sin 2\omega_0$, $1.69$ at $[2, 3]$ and $1$ as the gap closes; along the unit vector $\mathbf n$ the band's $k^4$ term carries $\sum_a n_a^4$, so the factor is $\omega_0\,[(\sum_a n_a^4)\tan\omega_0 + \cot\omega_0]$, $1.2224$ along a face diagonal and $1.0657$ along the body diagonal, where the group velocity is parallel to $\mathbf k$: the full factor is reached only in the limit of a small rest rotation, the same limit in which the post-Newtonian parameters reach $1$ (S.19). \emph{Check:} with the exact band at $[2, 3]$ the residual $(f_v - 1 + v^2 / (2c_m^2)) / v^4$, $f_v = \Omega / \omega_0$, at $v = 0.02$ is $-3.37$ along an axis against $1.6926$ times Lorentz's $-1.988$, that is $-3.365$, and $-2.12$ along the body diagonal against $1.0657 \times (-1.988) = -2.119$. \emph{A body:} a bound body's centre moves on the body's own band, $\omega_b(k)$, whose rest rotation and curvature are computed from the fixed point (S.13); the same steps on that band give the same form with the body's $c_b^2 = \omega_b / m_b^*$, $m_b^*$ its band inertia and $c_b$ its kinetic scale as $m^*$ and $c_m$ are the free packet's, and the body's band is a computed input, not derived in closed form here.
 \emph{The largest group speed:} on the full band $v_g$ is largest at $\mathbf k = (\pi / 2, \pi / 2, \pi / 2)$, $(\num / \den) / \sqrt 3 = \cos\omega_0 / \sqrt 3$, light's speed times $\cos\omega_0$, $0.3849$ at $[2, 3]$.
 \emph{Status:} derived, the second order a definition; lattice and clicks.
\end{derivation}

\begin{derivation}[de Broglie's relation, Section 5.1]
\emph{Inputs:} S.18's band at the composed paces and S.21's $c_m$. \emph{Steps:} from Derivation S.18, $v = k / m^*$ with $m^* = 3\tan\omega_0$ at the vacuum's content, so $k = m^* v = \omega_0 v / c_m^2$ for a free quantum ($\omega_b$, a body's own rotation, in its place on a body's band, S.21) with $c_m^2 = \omega_0 / m^* = \omega_0 / (3\tan\omega_0)$, the family's own light speed; $0.2508$ at $[2, 3]$.
 \emph{Status:} derived; lattice and clicks.
\end{derivation}

\section*{The two slits (Section 6)}

\begin{derivation}[the fringe's spacing and the near-field minima, Section 6.2]
\emph{Inputs:} the gaps centred at the rows $18$ and $30$, $d = 12$; the screen $L = 24$ Links from the wall; $\lambda = 8$ Links; gravity's weight $E_g = 10^6$ in this universe, chosen so that the packet's own gravity does not focus it. \emph{Steps:} the far-field spacing is $\lambda L / d = 16$ Links. In the near field the minima are where the two gaps' paths to the screen differ by half a wavelength, $\sqrt{L^2 + (y - 18)^2} - \sqrt{L^2 + (y - 30)^2} = \pm 4$, whose solutions are $y = 15.28$ and $y = 32.72$: the rows $15$ and $33$. An opaque region in one gap takes that gap's inflow, so the screen reads the other gap's envelope and no fringe, by the continuity of the one inflow through the two boundaries. \emph{Check:} the two roots computed from the equation above; the law's real line on this universe gives the exact lay $1{,}810.4$ quanta by Eq.~(5) and $284.7$ units through the screen over the window (the line's $284.7$ was computed after the run; the one number that stood before the run is the blind row's $273$), and $278.9$ on the integer lay of $1{,}809$ quanta, laid at the level $60$'s paces and stepped at the level $0$'s, $6.1$ below the run's $285$ (the paper's Section 6.2); the integer lay's $278.9$ lies below the exact lay's $284.7$ because the before level's mismatch at the lay seeds an opposite-branch admixture that the receding face loses (the two roots $e^{\pm i\omega}$ of the leapfrog line, the backward branch's amplitude from the before level's difference, $18.80$ at most, and its exit through the receding face); the folded lattice returns the Node to itself through both $z$ Ports, so the band is $2\cos\omega = \tfrac23(\cos k_x + \cos k_y + 1)$ exactly at $k_z = 0$; the peak of the screen's inflow is at the interval $69.0$ on the lattice as declared, $62.0$ with the lattice refined fourfold and $62.1$ eightfold (the line evaluated); the gap NodeDetector's net inflow is $-0.58$ over the window; the run's reading is deposited with the paper (Statements and Declarations) and holds the two seeds' rows. The blind row's $273$ is light's line in real arithmetic propagated on the design's lattice from the two gaps as its sources and summed at the screen over the window, Huygens's construction with the lattice's own line as its propagator, called here the Huygens sum; the line stepped over the whole lattice from the laid packet is a second computation, $284.7$ at this universe's exact lay; the run of Rule3's integer line on this universe gives $N = 285$ at both seeds, $24$ and $25$, the seed-$24$ clicks' sum of absolute differences from the blind row $28$ of $275$ ($275$ the blind row's sum by its per-Node rounding, $273$ its $N$); the group velocity $\sin k / (3\sin\omega) = 0.547$ Link per interval at $k = \pi / 4$ and $\omega = 0.4456$, the flight over $L' = 34$ Links from the packet's centre (the column $10$) to the screen $62.2$ intervals; the rows $15$ and $33$ fall in the four-row regions $3$ and $8$ counted from $0$; the blind row's dips at the regions $4$ and $8$; the reversal check of Section 3.1, the run forward and back over $40$ intervals with the counting window opening at $45$, so that the check contains no click; the integer step's bound and walk are S.53 (i)'s.
 \emph{Status:} derived for the spacing and the minima; computed for the lays, the totals and the arrival; the run's numbers computed, a run; the universe a declaration; lattice and clicks as each sentence marks.
\end{derivation}

\section*{The forms, the bodies and the calibration (Sections 4 to 8)}

\begin{derivation}[the Link's bound from light's dispersion, Section 5.1]
\emph{Inputs:} the group speed $v_g = (1 / \sqrt 3)(1 - k^2 / 12)$ (Derivation S.1); the bound on a quadratic dependence of light's speed on its energy, MAGIC's subluminal quantum-gravity (QG) scale of quadratic Lorentz violation $E_{\mathrm{QG}, 2}$ above $6.3 \times 10^{10}$ GeV (Acciari et al.\ 2020 \cite{magic}; the lattice band is subluminal, $v_g = (1 / \sqrt 3)(1 - k^2 / 12)$, so the subluminal limit applies), that is $|\delta v / v| = (3 / 2)(E / E_{\mathrm{QG}, 2})^2 < 3.8 \times 10^{-16}$ at $E = 1$ TeV in the standard quadratic form. \emph{Steps:} $k^2 / 12 < 3.8 \times 10^{-16}$ gives $k < 6.8 \times 10^{-8}$ per Link at $1$ TeV; with $k = 2\pi\,\mathrm{Link} / \lambda$ and $\lambda = 1.24 \times 10^{-18}$ m at $1$ TeV, the Link is below $k\lambda / (2\pi) = 1.3 \times 10^{-26}$ m along an axis, where the $k^4$ coefficient of the band's $\omega^2$ is $-1 / 54$ ($-1 / 135$ averaged over directions, $2.1 \times 10^{-26}$ m; $0$ along a diagonal, no bound); over the unit of S.27, $6.69 \times 10^{-13}$ m, the axis bound gives $\omega_e < 2.0 \times 10^{-14}$ ($\omega_e$ the matter family's rest rotation, S.27's $\omega_0$) ($3.1 \times 10^{-14}$ on the average) and $\gamma_K - 1 = \omega_e^2 / 2 < 2.0 \times 10^{-28}$. \emph{The orientation:} each source's bound carries its direction's coefficient $C(\mathbf n) = (3\sum_a n_a^4 - 1) / 24$ of the group speed's fall, unknown against the lattice: $1 / 12$ on an axis, $1 / 48$ on a face diagonal, $0$ on a body diagonal and continuous toward it; the mean over directions, $\langle\sum_a n_a^4\rangle = 3 / 5$, gives $2 / 5$ of the axis's coefficient (the $-1 / 135$ against $-1 / 54$ above); the Link bound weakens tenfold only within $4$ degrees of a body diagonal, one percent of the sky, and a hundredfold within $0.4$ degrees, so the comparison of MAGIC's and LHAASO's bounds holds under the mean's premise (two derivations, the second S.54's two-way coefficient check).
 \emph{Status:} derived; lattice and clicks.
\end{derivation}

\begin{derivation}[Newton's constant as a reading, Eq.~(11)]
\emph{Inputs:} the massless holder's rest far from a static source of $M$ quanta per interval, $3M / (4\pi r)$ over the divisor $E_g$, the write $s = M / E_g$; the clock's content $U = c_\ell / \Gamma$, $c_\ell$ the content in levels. \emph{Steps:} $U = 3M / (4\pi r\,\Gamma E_g) = 3s / (4\pi\Gamma r)$, and nature's $U = GM / (c^2 r)$ with $M$ in quanta and $c^2 = 1 / 3$ gives $G / c^2 = 3 / (4\pi\Gamma E_g)$, that is $G_{\mathrm{clock}} = 1 / (4\pi\Gamma E_g)$; Kepler's $G$, read from the orbits of matter, is $[2\cos\omega_0 / (1 + \cos\omega_0)]\,G_{\mathrm{clock}}$ by Eq.~(8).
 \emph{Status:} derived; lattice, clicks.
\end{derivation}

\begin{derivation}[the fine-structure constant as a coefficient of the declared universe, Section 5.1]
\emph{Inputs:} the holder of the sign under its rotation (Section 3.5) with its declared write weight $k$, the written level $k$ times the plain one and the turn $\theta = \mathrm{level} / \Gamma$, $k = 1$ in the rule's declared universe; a circular record of a family at the vacuum's paces, $z = Ae^{i\omega_s t}$, the quantum of a family one unit of the invariant, $2A^2\sin\omega_s = T$, with the share $T\sin\omega_s$ and $|W| = A^2\sin\omega_s = T / 2$ (S.9); the holder's rest far from a static source of $s$ per interval, $3s / (4\pi r)$ over its weight $E_s$ (Derivation S.25's line); de Broglie's relation and light's $c = 1 / \sqrt 3$ (S.1, S.22). \emph{Steps:} (a) One quantum writes $s = W / (E_s T) = 1 / (2E_s)$ per interval, for every family, since the sourcing family's denominator enters no write (the count is the share over the wall $3\,\den_s T$ and the source the form over $E_s T$, the same number); so $|q_p| = |q_e|$ with nothing declared. The angle at $r$ is $\theta(r) = 3ks / (4\pi r\Gamma) = 3k / (8\pi r\Gamma E_s)$. (b) A reader of unit sense gains the angle per interval, so its wave number changes by the gradient, $dk_{\mathrm{dB}} / dt = \mp\nabla\theta$ by the sense of (c), of magnitude $3k / (8\pi r^2\Gamma E_s)$; Coulomb's law in nature's units, $dk_{\mathrm{dB}} / dt = \alpha c / r^2$, gives $\alpha_{\mathrm{law}} = \sqrt 3\,r^2\,dk_{\mathrm{dB}} / dt = (3\sqrt 3 / 8\pi)\,k / (\Gamma E_s) = 0.20675\,k / (\Gamma E_s)$; the holder's weight enters the source and the read alike as $k$ and $E_s$, and the family's gap cancels. (c) The reader's shift: the step carries $e^{i\theta_t}$ on $z_{\mathrm{next}}$, so a sense $+$ quantum goes from $\omega$ to $\omega + \theta$ at its fixed Wronskian (S.33, S.41 (b)), its energy from $T\sin\omega$ to $T\sin(\omega + \theta)$, $\delta E = +T\theta\cos\omega$ per quantum; a body of $n$ quanta and sense $\pm$ shifts by $\pm nT\theta\cos\omega = 2\theta\cos\omega\,W$ with $W = \pm nT / 2$, the force $-2W\cos\omega\,\nabla\theta = \mp nT\cos\omega\,\nabla\theta$, which is $\mp n\nabla\theta$ by the sense in the unit $T\cos\omega$, one quantum's energy shift per unit turn (its energy $T\sin\omega$), at the vacuum's rotation, Coulomb's magnitude with $q = \pm ne$ and $e = T / 2$; the angle $\theta = k_r L / \Gamma$ carries the sign of the level and the level the sign of its source's Wronskian, so two quanta of one sense raise each other's rotation, $\delta E > 0$ growing toward the source, and move apart; which sense moves apart is the declared sign of Section 3.5 (under the law's one convention like senses move apart, as nature's measurement that like charges repel declares), the magnitude untouched; the reader sources the holder it reads at the same weight, so the pair is reciprocal (the reciprocity Theorem 3 needs). (d) The number: nature's $\alpha = 1 / 137.036$ \cite{codata} fixes $k / (\Gamma E_s) = 1 / 28.33$ and nothing else; $k$ is the universe's integer read from this one click as $E_g$ is from $G$, and at $k = 1$ no integer $\Gamma E_s$ gives it ($28$ gives $1 / 135.43$ and $29$ gives $1 / 140.27$). (d$'$) The energy line. The weight is the product of the write's and the read's, $k = k_w k_r$, with the read's $k_r = 1$ on the integer lattice (a read weight would multiply a level the write already rounded away). Under the time-Link turn a plane record's form and the gauge-turned record's Wronskian (S.41 (a)) are conserved exactly per Node in the rational line for any level in time (the turn the time Link's gauge, S.33, S.41 (a)) and its rotation shifts by the turn per interval, so the energy its clicks read changes by $T[\sin(\omega + \sigma\sigma_s\Delta\theta) - \sin\omega]$ per quantum of sense $\sigma$ in a level of source sense $\sigma_s$ (S.41 (b): a quantum goes to $\omega + \sigma\sigma_s|\theta|$), $\sigma\sigma_s T\cos\omega\,\Delta\theta$ to first order in the turn (the remainder $(\Delta\theta / 2)\tan\omega_s$ relative, below $10^{-18}$ at nature's turn per interval), per quantum $T\cos\omega_s\,(k_r / \Gamma)\,\Delta\ell$ in a level of its own sense (a convention with one turn angle on both levels would lose $2W\cos\omega\,\Delta\theta$ of form per Node per step; the law's convention loses none); the holder's own form changes by $\sum_i s_i(\ell_{t+1} - \ell_{t-1})$ per step under the write, per quantum $(k_w / E_s)\,\Delta\ell$, the two factors $2$ cancelling, and for a pair of opposite senses the cross form carries $W_1 W_2 < 0$; so the energy its clicks read and the holder's form move together under the turn and the write, both up for a body in its own level and both down for a pair of opposite senses, and agree to first order in the turn if and only if $k_r E_s T\cos\omega_s = k_w\Gamma$: one energy read twice, by the clicks and by the form, an exact condition on the energies read and not an identity of any one record's form, the ratio of the two changes $\chi = k_r E_s T\cos\omega_s / (k_w\Gamma)$ at every step and for every motion, at the vacuum's paces of these inputs (in a well the body's clock enters the turn and the write alike), the ratio $+\chi$ under S.41 (b) by the algebra above (the two senses' product the sign); the law's line is a condition on the energies read and no identity of a record's form. Under it, with $k_r = 1$, $\alpha_{\mathrm{law}} = (3\sqrt 3 / 8\pi)\,T\cos\omega_s / \Gamma^2$, at nature's gap $T = \Gamma^2 / 28.33$, $E_s$ a unit of the level and the declared universe one integer fewer; the line is per family through $\cos\omega_s$, which in nature is $1$ to $10^{-18}$ for every family, so one write weight serves all; light and every real-line family read no turn and stand outside it. The rule's own universe has no plane family and keeps $T = 2^{15}$; the universe's gate condition, $E_s T\,\num = k_w\Gamma\,\den$ for a plane family reading a holder under the rotation, enforces the line. The per-step identity above is the one-field case, a level built by the body itself; under the rule in force, no record reading its own write, the reader's loss is by the other record's field and the writer's gain its own field's, and what holds per step is the two-record balance at the same coefficient per quantum, not run; the universe's gate condition is unchanged by this. The radiated share of a charged body (C.6 (c)): the time level carries half of Larmor's power in form; the odd lines' part depends on their wall and coupling, which of the two couplings the odd lines carry open; no rate and no lifetime is claimed. Two derivations on the identity and its number; the line a decision of the model (assumption). (e) $\Gamma$ is bounded by no click of the clocks (the gap bound binds the matter family's own denominator, $\den > 1.5 \times 10^{19}$ for the least pair, S.27, a pair of a universe meant for nature and not Universe24's): the content is an integer and the clock reads it as $p_0 = \Gamma(1 - 1 / \Gamma)^c$, so the finest level at one Node and one interval is $1 / \Gamma$; but the write adds $(\sum wq + r) \ddiv E$ with the remainder $r$ kept, so a sub-unit source does not vanish but writes one level once every $1 / s$ intervals, the level at a Node a sequence of $0$ and $1$ whose mean is the exact rest $3s / (4\pi r)$, a first-order sigma-delta modulation; a clock in such a level reads the level $1$ for the fraction $\Delta U\,\Gamma$ of its intervals and its mean redshift is $\Delta U$, finer than $1 / \Gamma$, and a bound body spanning many Nodes averages their independent dithers, so the gravitational clocks of nature bound the mean potential and not $\Gamma$; nothing bounds $\Gamma$ from above and no line fixes it: $\Gamma$ is free and large (two derivations). \emph{Check:} $3\sqrt 3 / (8\pi) = 0.20675$; $137.036 \times 0.20675 = 28.33$; the share's formula at $[2, 3]$ and $\Gamma = 6{,}000$ gives $|W| / (\mathrm{count}\,T) = 0.6708 = 1 / (2\sin\omega_0)$ for a record laid by share ($A^2 = T / (2\sin^2\omega_0)$ for a record laid by share) and $1 / 2$ for a record of count $1$, one quantum of the family under the record's own unit (C.9 (b)); $3\sqrt 3 / (8\pi) / 28.33 = 0.20675 / 28.33 = 1 / 137.0$.
 \emph{Status:} derived under the declared weights; lattice and clicks.
\end{derivation}

\begin{derivation}[the units, the share's ceiling and the matter pair's gap, Section 5.1]
\emph{Inputs:} the electron as the free quantum of matter, rotating at $\omega_0$ per interval, and nature's electron at $m_e c^2 / \hbar$ per second; light's band $\cos\omega = 1 - \frac{1}{3}\sum_a(1 - \cos k_a)$ (S.1); a photon of $1.4$ PeV, LHAASO's record of 2021 \cite{lhaaso} (photons above $2$ PeV were reported in 2024 \cite{lhaaso2024}); the share $T\sin\omega$ a click reads (S.14). \emph{Steps:} one interval is $\omega_0\hbar / (m_e c^2) = \omega_0 \times 1.288 \times 10^{-21}$ s, and one Link, light moving $1 / \sqrt 3$ Link per interval at $c$, is $\sqrt 3\,\omega_0\,\bar\lambda_e = \omega_0 \times 6.69 \times 10^{-13}$ m with $\bar\lambda_e = \hbar / (m_e c) = 3.862 \times 10^{-13}$ m. The band's top along the body diagonal, $k_a = \pi$ for all $a$, is $\cos\omega = -1$, $\omega = \pi$, and along an axis, $k = \pi$ alone, $\cos\omega = 1 / 3$, $\omega = 1.231$. Along an axis the group velocity vanishes at the edge; along the diagonal, $k_a = q$ on every axis, the band is $\omega = q$ exactly, so the speed stays $1 / \sqrt 3$ up to the corner, where the band has a cusp and no one gradient. That is a lattice statement of the highest rotation, which no click reads as an energy: what a click reads of a quantum of frequency $\omega$ is its share, $T\sin\omega$ (S.14), at most $T$, reached at $\omega = \pi / 2$, which the band contains along the diagonal (along an axis it ends at $1.231$, $\sin\omega = 0.943$), and an energy measurement of nature, a calorimeter or an air shower, is a sum of transitions' shares. The electron's rest energy is $T\sin\omega_0$, so a photon of $1.4$ PeV read as a share gives $\sin\omega_0 \le 0.511\,\mathrm{MeV} / 1.4\,\mathrm{PeV}$, $\omega_0 < 3.7 \times 10^{-10}$ for every orientation of the lattice (the diagonal's $3.65 \times 10^{-10}$, the axis's $3.44 \times 10^{-10}$; the diagonal's is the one quoted). Then $\den = d / (1 - \cos\omega_0) \approx 2d / \omega_0^2 > 1.5 \times 10^{19}$ with $d = \den - \num \ge 1$ (the bound at $d = 1$) and $\gamma_K - 1 = 1 / \cos\omega_0 - 1 \approx \omega_0^2 / 2 < 6.7 \times 10^{-20}$ (S.19); $\Gamma$ reads nothing of $\omega_0$, $\alpha$ fixing $k / (\Gamma E_s)$ and not $\Gamma$ (S.26). Were the frequency read as the energy, $T\omega$, the bound would be $1.1 \times 10^{-9}$ along the diagonal and $4.5 \times 10^{-10}$ along an axis: the law's resonance adds frequencies (S.41) and its share adds sines (S.14), the two agreeing to $\omega^2 / 6$ relative, below $10^{-18}$ at nature's gap, and parting at the Link scale, where which of the two a click takes the law as written does not decide and no measurement of nature reaches; named, not claimed. \emph{Check:} $\hbar / (m_e c^2) = 1.0546 \times 10^{-34} / 8.187 \times 10^{-14} = 1.288 \times 10^{-21}$ s (the CODATA values \cite{codata}); $\sqrt 3 \times 3.862 \times 10^{-13} = 6.69 \times 10^{-13}$ m; $0.511 / (1.4 \times 10^9) = 3.65 \times 10^{-10}$; $2 / (3.65 \times 10^{-10})^2 = 1.5 \times 10^{19}$; $(3.65 \times 10^{-10})^2 / 2 = 6.7 \times 10^{-20}$; $\sin 1.231 = 0.943$; $[\pi - \omega(\pi - \epsilon)] / (\sqrt 3\,\epsilon) = 0.5774$ at $\epsilon = 10^{-2}$ and $10^{-3}$ on the diagonal.
 \emph{Status:} derived; lattice and clicks.
\end{derivation}


\begin{derivation}[the content's index for every family, Section 5.1]
\emph{Inputs:} Eq.~(3) at the vacuum's pace and at a content $U$ under the composed paces, $(p_0 / \Gamma)^2 = N^2$ and $(p_a / \Gamma)^2 = N^4$ with $N = (1 - 1 / \Gamma)^c = e^{-U}$ to $O(1 / \Gamma)$ (Derivation S.2); a stationary boundary conserves the rotation $\omega$. \emph{Steps:} outside, $\cos\omega = 1 - g - \frac{\num}{3\,\den}(1 - \cos k)$ with $g = 1 - \num / \den$; inside, $\cos\omega = 1 - gN^2 - \frac{\num}{3\,\den}N^4(1 - \cos k_{\mathrm{in}})$. Equating, $\frac{\num}{3\,\den}N^4(1 - \cos k_{\mathrm{in}}) = \frac{\num}{3\,\den}(1 - \cos k) + g(1 - N^2)$, so $1 - \cos k_{\mathrm{in}} = [1 - \cos k + 3(\den - \num)(1 - N^2) / \num] / N^4$, exact in $k$ on the domain where the right side lies in $[0, 2]$. With $1 - N^2 = 2U - 2U^2 + \dots$ and $N^{-4} = 1 + 4U + \dots$ the first order is $[1 - \cos k + 6(\den - \num)\,U / \num] / (1 - 2U)^2$, the quadratic forms' line, the mass term's second order $-6(\den - \num)U^2 / \num$. For light, $g = 0$: $1 - \cos k_{\mathrm{in}} = (1 - \cos k) / N^4$, so $\sin(k_{\mathrm{in}} / 2) = \sin(k / 2) / s$ with $s = N^2$, exact in the law's variable $2\sin(k / 2)$ where $\sin(k / 2) \le s$. The phase index is $k_{\mathrm{in}} / k = (2 / k)\arcsin[\sin(k / 2) / s] = 1 / s + k^2(1 / s^3 - 1 / s) / 24 + O(k^4)$, which is $1 / s = e^{2U}$, Derivation S.2's long-wavelength index, at $k \to 0$ alone. At $k = \pi / 2$ and $s = 0.8$ the exact ratio is $1.380$ against $1 / s = 1.25$, and to first order in $U$ the index is $1 + [4\tan(k / 2) / k]\,U$, the coefficient $2$ at $k \to 0$ and $8 / \pi$ at $k = \pi / 2$: the lattice's own chromatic dispersion, absent in the ray limit in which the bending, the delay and the lens are read. For slow matter, $k, k_{\mathrm{in}} \ll 1$: $k_{\mathrm{in}}^2 = [k^2 + 6(\den - \num)(1 - N^2) / \num] / N^4$, so $n_m^2 = k_{\mathrm{in}}^2 / k^2 = [1 + 6(\den - \num)(1 - N^2) / (\num\,k^2)] / N^4$, to first order $[1 + 12(\den - \num)\,U / (\num\,k^2)] / (1 - 2U)^2$. \emph{Check:} with $k^2 = m^{*2} v^2$ and $m^* = 3\tan\omega_0$ (S.18), the $1 / k^2$ term is $12(\den - \num) U / (\num\,m^{*2} v^2) = [4(1 - \cos\omega_0) / (3\cos\omega_0\tan^2\omega_0)]\,U / v^2$, which as $\omega_0 \to 0$ is $(2 / 3)\,U / v^2 = 2(U / 3) / v^2 = -2\Phi / v^2$ with $\Phi = -c^2 U$, $c^2 = 1 / 3$ (nature's $n^2 = 1 - 2\Phi / v^2$ for a slow body in a potential): nature's index for a slow body in a potential. The factor $6(\den - \num) / \num$ over $m^{*2}$ is the matter band's own, the Kepler column of Eq.~(8), and $1 / N^4$ the Link's double read.
 \emph{Status:} two derivations; derived; lattice and clicks.
\end{derivation}


\begin{derivation}[the clusters' cosmology: the ball of dust, Section 4.4]
The paper's claim from this derivation is the first order in $z$, Friedmann's dust equations and Hubble's law; the second order is derived here and stands in the paper's Section 8.1 beside the supernovae as a finding against nature. \emph{Inputs:} Newton's pull between bodies at Kepler's $G_K$ (Section 5.1, Derivation S.25); a finite ball of dust of radius $a(t)$ and uniform density $\rho$ on the static lattice; the Doppler period of a receding source, derived in the steps, and the moving clock's factor (S.21); S.1's band for the Doppler period; S.4's covariance for the reading's kind; Section 8.1's bound for the comparison. \emph{Steps:} a shell at $a$ feels the mass inside it, $M(a) = (4\pi / 3)\rho a^3$, and nothing from outside (Newton's theorem on a sphere, which holds for the $1 / r$ level of Section 5.1), so $\ddot a = -G_K M(a) / a^2 = -(4\pi G_K / 3)\,\rho\,a$. Multiplying by $\dot a$ and integrating with $\rho a^3$ constant gives $\dot a^2 / a^2 = (8\pi G_K / 3)\,\rho - K / a^2$ with $K$ the start's declaration, positive for a bound ball: Friedmann's two equations for dust, as Milne and McCrea derived them for the Newtonian ball \cite{mccrea}. The deceleration parameter is $q_0 = -\ddot a\,a / \dot a^2 = (4\pi G_K \rho / 3) / H^2 = \Omega_m / 2$ with $\Omega_m = 8\pi G_K\rho / (3H^2)$, positive for every $\rho > 0$. \emph{The redshift:} a source receding at $\beta c$ emits two crests a period $P$ apart from points $\beta c P$ apart, so the second trails the first by $(1 + \beta) c P$ and a resting NodeDetector reads the period $(1 + \beta) P$ exactly, the wave number set at the emission and kept by the line's translation invariance (S.1), the period read $(1 + \beta)$ times longer and the source's own period is longer by the moving clock's factor $1 / \sqrt{1 - \beta^2}$ at $c_m = c$ (S.21), so $1 + z = (1 + \beta) / \sqrt{1 - \beta^2} = 1 + \beta + \beta^2 / 2 + \dots$, nature's relativistic Doppler form to the second order, the receding source's classical factor $1 + \beta$ differing at the second order ($1 / (1 - \beta) = 1 + \beta + \beta^2$ is the receding receiver's). At the computed pair $c_m = 0.867c$ the second order differs, as every number of the computed pair does. Every click's duration is stretched by the same $1 + z$, the light curves' time dilation, and a surface's brightness falls by $(1 + z)^{-4}$, Tolman's dimming, at the small gap, $(1 + z)^{-2}$ from the two rows alone, the aberration's two powers needing the band's Lorentz invariance, which holds at the small gap (S.21) and not at a finite gap (S.4); light from a source at $r$ falls into the ball's deeper centre, a blueshift $\Delta U = (\Omega_m / 4)(H_0 r / c)^2$ beside the Doppler, the second order that makes $d_L(z)$ Friedmann's with $q_0$ unchanged. \emph{Status:} two derivations; derived; clicks; the second order, $q_0 = \Omega_m / 2 > 0$ under the declared homologous start, a finding against nature in the paper's Section 8.1.
\end{derivation}


\begin{derivation}[the shadow's capture radius from light's index, Section 5.1]
\emph{Inputs:} light's index $n(r) = e^{2U}$ with $U = m / r$ (Derivation S.2 under the composed paces), $m = GM / c^2$ in Links (the $\mu$ of S.19 and S.20). \emph{Steps:} in a static isotropic index a ray at the impact parameter $b$ keeps $n(r)\,r\sin\psi = b$ along its path, so a circular ray stands where $d(n r) / dr = 0$: $d(r e^{2m / r}) / dr = e^{2m / r}(1 - 2m / r) = 0$ at $r = 2m$, the photon sphere, and the capture radius is $b = n(2m)\,2m = 2e\,m = 5.437\,m$. Schwarzschild's is $3\sqrt 3\,m = 5.196\,m$, so the shadow is larger by $2e / (3\sqrt 3) - 1 = 0.046$. \emph{Comparison:} with the quadratic index $1 / (1 - 2U)$ the same condition gives $r = 4m$ and $b = 8m$, $+54$ percent, outside the Event Horizon Telescope's fractional deviation from Schwarzschild for Sgr A$^*$, $-0.08 \pm 0.09$ \cite{eht_sgra}; the composed index's $+4.6$ percent sits $1.4$ standard deviations from it, bounded and not confirmed.
 \emph{Status:} derived under the one assumption; clicks; the shadow bounded by measurement and not confirmed.
\end{derivation}

\begin{derivation}[the two clocks under the one assumption of Section 2.3, Section 5.1]
\emph{Inputs:} the one assumption of Section 2.3, a uniform level $U$ a change of units, time by $N$ and length by $h$ with $hN = 1$; light's speed $N^2$ (S.2). \emph{Steps:} a bound body in the level is the flat body in the local units, so its size is $1 / h = N$ Links by the one assumption. Its band at the Node is $2 - 2\cos\omega = 2(1 - \nu)N^2 + (\nu / 3)\,N^4\sum_a s(k_a)^2$ with $\nu = \num / \den$ and $s(x) = 2\sin(x / 2)$, since $(p_0 / \Gamma)^2 = N^2$ and $(p_a / \Gamma)^2 = N^4$, and its lattice wave numbers are $k = k_{\mathrm{proper}} / N$, so $(\nu / 3)\,N^4\sum_a s(k_a)^2 = (\nu / 3)\,N^2\sum_a k_{\mathrm{proper}, a}^2\,[1 - k_a^2 / 12 + \dots]$: every mode of the record, the rest rotation and every level of its spectrum, scales by $N$ in $s(\omega)$, exactly at long wavelength and to the lattice's $k^2 / 12$ at finite $k$. In $\omega$ the rest line's rate is $-2\tan(\omega_0 / 2) / \omega_0 = -f(\omega_0)$ per unit $U$, Derivation S.15's number. A difference $s(\omega_2') - s(\omega_1')$ scales by $N$ exactly and is no observable. The temporal beat a NodeDetector reads, $\Delta' = \omega_2' - \omega_1'$ in $\cos(\Delta' t)$, with $\omega' = 2\arcsin[N\sin(\omega / 2)] = N\omega\,[1 - (1 - N^2)\omega^2 / 24 + O(\omega^4)]$, is $\Delta' = N\Delta\,[1 - (1 - N^2)(\omega_1^2 + \omega_1\omega_2 + \omega_2^2) / 24 + O(\omega^4)]$: at $N = 0.8$, $\omega_1 = 0.4$ and $\omega_2 = 0.9$ the exact beat is $0.3916$ against $N\Delta = 0.4$ and the series $0.3920$. The relative correction, $(1 - N^2)(\omega_1^2 + \omega_1\omega_2 + \omega_2^2) / 24$ for the beat, $(1 - N^2)\omega^2 / 8$ at $\omega_1 = \omega_2$, is $2.5 \times 10^{-38}$ at $\omega = 1.0 \times 10^{-14}$ and $U = 10^{-9}$, an illustration below the bounds ($1.0 \times 10^{-37}$ at the axis bound's $2.0 \times 10^{-14}$, S.24; $3.4 \times 10^{-29}$ at the LHAASO bound's $3.7 \times 10^{-10}$, S.27) and $1.4$ percent at $\omega_b = 0.78$ and $U = 0.1$, the lattice's arithmetic and not the law's. A cavity whose mirrors are bodies stands $L / h = LN$ Links apart and its round trip is $2LN\sqrt 3 / N^2 = 2L\sqrt 3 / N$ intervals, the rate $N$ exactly with no $\omega$ correction, the atom's beat to its correction $(1 - N^2)\omega^2 / 8$ above; a cavity of declared faces keeps its $L$ Links and ticks at $N^2$. So under the one assumption of Section 2.3 no clock type is singled out, which is what the null redshift tests measure, while a cavity of declared faces is a NodeDetector's region and no body of nature.
 \emph{Status:} derived under the one assumption; lattice and clicks.
\end{derivation}

\begin{derivation}[the conserved form's Link factor, two witnesses, Theorem 3]
\emph{Inputs:} Theorem 3's line without its division on a periodic lattice at the content $0$ with every Link's factor $q$, so $p_0 = p_i = \Gamma$, $R = 2\,\num\,q^2$, $w = 6\,\den\,\Gamma^2$ and $S = 12\,\num\,\Gamma^2 - 6R$; the form $E$ of Eq.~(4) and the same form with the plain Link term $-2\,\num\sum_i \mathrm{now}_i\,S_6(\mathrm{before})_i$; exact fractions. \emph{First witness:} three Nodes along $x$ with $y$ and $z$ folded, $\num = \den = 1$, $\Gamma = 4$, $q = 2$ ($R = 8$, $w = 96$, $S = 144$), $\mathrm{now} = (1, 0, 0)$ and $\mathrm{before} = 0$: one step gives $\mathrm{next} = (11 / 6, 1 / 12, 1 / 12)$; the plain form goes from $6$ to $-21 / 4$ and Eq.~(4) stays $6$. \emph{Second witness:} $3^3$ Nodes at $[2, 3]$, $\Gamma = 10$, $q = 5$ ($R = 100$, $w = 1{,}800$, $S = 1{,}800$), a uniform $\mathrm{now} = 3$ over $\mathrm{before} = 0$: the step gives $\mathrm{next} = 4$ at every Node with the remainder $0$, so the integer step and the line agree; the plain form per Node goes from $162$ to $-54$ and Eq.~(4) stays $162$ ($4{,}374$ over the lattice). \emph{Check:} at $q = \Gamma$ the factor is $1$ and the two forms coincide, the domain of the plain current.
 \emph{Status:} theorem (a corollary of Theorem 3); computed for the two witnesses; lattice.
\end{derivation}

\begin{derivation}[magnetism: the continuity identity and the covector, Section 5.1]
\emph{Inputs:} a plane record $z = \mathrm{re} + i\,\mathrm{im}$ stepping by Rule3 under the holder of the sign (Section 3.5): $e^{i\theta_t} z_{\mathrm{next}} + e^{-i\theta_{t-1}} z_{\mathrm{before}} = [S\,z_{\mathrm{now}} + \sum_a R_a (e^{i\theta_a} z_{+a} + e^{-i\theta_a} z_{-a})] / w$, with $\theta_t$ this interval's time turn on the level next, $\theta_{t-1}$ the previous interval's on the level before and $\theta_a$ the phase on the axis $a$'s Links, each read at its interval's start (the convention of the paper's Section 3.5, stated in S.33 and S.41 (b)); the Wronskian booked from the turned levels $u_i = e^{-i\theta_{t-1}} z_{i,\mathrm{before}}$ and $v_i = e^{i\theta_t} z_{i,\mathrm{next}}$ as the step reads them, $W_i = \mathrm{Im}(\bar z_{i,\mathrm{now}}\, u_i)$, and $W_i(t + 1) = \mathrm{Im}(\bar v_i z_{i,\mathrm{now}})$ after the step, the two differing by the Link fluxes of (a); the Link quantity $g_{ij} = \mathrm{Im}(\bar z_i\, e^{i\theta_{ij}} z_j)$, antisymmetric, $g_{ji} = -g_{ij}$, since the phase is read $+\theta_a$ through the $+a$ Port and $-\theta_a$ through the $-a$ Port; S.9's Wronskian current; S.21's packet clock and S.22's de Broglie relation for the covector's reading; S.41 (a) and (b) for the turned record's Wronskian and the time turn. \emph{Steps:} (a) \emph{The identity} (the level before under the previous interval's angle and the level next under this one's, so the Wronskian is conserved exactly for any level in time, a uniform level a gauge, S.41 (a); the algebra below is written at one angle $\theta_0$ over the step, the accumulated phase removing a varying turn). Multiply the step by $w\,\bar z_{\mathrm{now}}$ and take the imaginary part; the $S$ term is real and drops: $w\,\mathrm{Im}(\bar z_{\mathrm{now}} e^{i\theta_0} z_{\mathrm{next}}) + w\,\mathrm{Im}(\bar z_{\mathrm{now}} e^{-i\theta_0} z_{\mathrm{before}}) = \sum_a R_a\,(g_{i,i+a} + g_{i,i-a})$. The first term is $-W_i(t+1)$ and the second $W_i(t)$, and $g_{i,i-a} = -g_{i-a,i}$, so
\[
w\,[W_i(t+1) - W_i(t)] + \sum_a \bigl[R_a\,g_{i,i+a} - R_a\,g_{i-a,i}\bigr] = 0 :
\]
an exact lattice continuity equation at every Node and interval, the density $\rho_i = w\,W_i$ and the flux through the Link $(i, i+a)$ equal to $R_a\,g_{i,i+a}$, for any paces, any phases and any levels, the per-Node identity; the summed conservation needs the weight $1 / p_i^2$, $\mathbf D\mathbf M$ symmetric (S.9); the integer step's remainders add the term of S.9's domain. (b) \emph{Flux over density is the group velocity.} For $z = A\,e^{i(\mathbf k\cdot\mathbf x - \omega t)}$, $\rho = w A^2\sin\omega$ and the flux is $R_a A^2 \sin k_a$ (common sign), so flux over density is $(R_a / w)\sin k_a / \sin\omega = d\omega / dk_a$ from the band $2w\cos\omega = S + 2\sum_a R_a\cos k_a$: the current is $\rho\,v$ exactly, and the law's Node current $J_a = g_{i,i+a} + g_{i-a,i}$ obeys $J_a = 2\,(w / R_a)\,W v$ on a plane record, $(6\,\den / \num)\,W v$ at the vacuum's paces where $R_a / w = \num / (3\,\den)$ (the case (d) names), nature's $j = \rho v$ with the band's constant and no $\sin\omega_b$. (c) \emph{The covector.} Inserting the plane wave in the step with the phases (the main text's convention, Section 3.5) gives $\omega = \omega_{\mathrm{band}}(\mathbf k + \boldsymbol\theta) + \theta_0$ for the sense $+$ (the convention of S.41 (b)): the pair $(\theta_0, \theta_a)$ shifts $(\omega, k_a)$ and transforms as $(\omega, \mathbf k)$ does; at long wavelength and small gap every band is the Lorentz form $\omega^2 = \omega_0^2 N^2 + c_s^2 k^2$ with $c_s^2 = \num_s / (3\,\den_s)$ (from $2\cos\omega = S / w + (2\num / 3\den)\sum_a\cos k_a$), so $(\theta_0, \theta_a)$ is a covector under the band's boost: minimal coupling, $A_0$ on the time step and $A_a$ on the Link. (d) \emph{The halved current over the time level's wall.} The sign holder under the rotation lays the time line and three odd axis lines sourced by the sign's current; the derivation's own is one kernel over one wall and the halving: Maxwell's pair with one kernel sources $(\theta_0, \theta_a)$ from $(\rho, \text{flux} / c^2)$ over one wall. The Node's flux on the axis $a$ is the mean of its two Links', $R_a (g_{i,i+a} + g_{i-a,i}) / 2 = R_a J_a / 2$, and $R_a / w = \num / (3\den)$ at the vacuum's pace, so flux over $c_s^2\rho$ is $[(\num / 3\den)\,J_a / 2] / [(\num / 3\den)\,W] = (J_a / 2) / W$: the odd line on the axis $a$ is sourced by the halved current $J_a / 2$ over the same wall as the time level's $W$, with no number, when the pair is read at the source family's own $c_s$; read at light's $c^2 = 1 / 3$ the factor $c_s^2 / c^2 = \num_s / \den_s$ appears, $1$ for light, $1 - 10^{-28}$ at nature's matter, $2 / 3$ at $[2, 3]$ in this small-gap form ($c_s^2 = \num_s / (3\den_s)$), where the band's own $c_m^2 = \omega_0 / (3\tan\omega_0)$ of S.21 and S.22 gives $0.752$; the two agree as the gap closes. The halving is not a choice: the Node's current is the mean of its two Links as the continuity identity defines it, as the pace reads the tension on an axis as the mean of its two ends' tensions. (e) \emph{Coulomb over $\gamma$.} For a source moving at $v$ along the Link and a reader co-moving with it, the time level is the moving density's, $\gamma$ times the rest level by the current's covector ((a), (c)) at fixed total $W$, within (c)'s domain, so the electric force across the beam is $\gamma F_0$; the odd level is sourced by $J_a / 2 = (3\den / \num)\,Wv = Wv / c_s^2$, the flux over $c_s^2$ of (d), the reader at $v$ feels $-v\,\partial\theta_a / \partial y$, a magnetic force $\beta^2\gamma F_0$ toward the source; the total is $\gamma(1 - \beta^2) F_0 = F_0 / \gamma$, with $\gamma$ Lorentz's at nature's gap, where the band is Lorentz's to $9 \times 10^{-20}$ at the LHAASO bound (S.21); at the computed pair $[2, 3]$ the band's own factor differs from it, and $\beta = 0.5$ lies beyond the axis's largest group speed, so the two numbers below are the form's at nature's gap. \emph{Check:} $1 / \gamma = 0.99499$ at $\beta = 0.1$ and $0.86603$ at $\beta = 0.5$, the prediction of (e); the identity (a) holds in exact arithmetic on a chain with contents and phases at every Node (the same check as S.9's, the phases a unitary on each read). \emph{Domain:} the covector is the band's at long wavelength and breaks at the lattice's anisotropy $k^2 / 12$ ($2 \times 10^{-4}$ at $k = 0.05$, $1.0 \times 10^{-28}$ at nature's Compton wavelength); the pairing holds at the source family's own $c_s$, so a source and a reader of one family see one light speed.
 \emph{Status:} theorem for (a) and (b), the conservation; derived at long wavelength for (c) to (e), the force's form, Peierls' coupling; lattice for the identity, clicks for the force.
\end{derivation}

\begin{derivation}[the composed paces, the acts' factors and the exponential metric, Section 2.3]
\emph{Inputs:} the model's one premise beside the line, that a level is a count of units and every act is measured in the Node's own units (Section 2.3); the clock $p_0$ as Rule3's coefficient in $S$; the band at a pace, S.1; the one assumption $p_a = p_0^2 / \Gamma$.
\emph{Steps:} (a) \emph{The composition.} One unit of level is one act on the clock as it stands, so it slows the clock by one part in $\Gamma$ of the clock as it stands: $p_0(c + 1) = p_0(c)\,(1 - 1 / \Gamma)$, whence $p_0(c) = \Gamma(1 - 1 / \Gamma)^c$. The functional equation $P(U + l) = P(U)P(l)$ is Cauchy's, whose continuous solutions are the exponentials and no other, so $N = p_0 / \Gamma = e^{-U}(1 + O(U / \Gamma))$ with $U = c / \Gamma$; the quadratic form's subtraction $p \to p - 2$ is as memoryless and is excluded only by the premise that the read is an act measured in the Node's units, like the write. \emph{Check:} at $\Gamma = 6{,}000$ and $c = 300$, $(1 - 1 / \Gamma)^c = 0.951225$ against $e^{-0.05} = 0.951229$, a relative $4 \times 10^{-6}$, which is $U / (2\Gamma)$. (b) \emph{The ruler.} $p_i = p_0^2 / \Gamma = \Gamma N^2$ (the Link pace $p_a$ at the vacuum's factor $q = \Gamma$) gives the axis ruler $h = p_0 / p_a = 1 / N$: $hN = 1$ (the one assumption, fixed by the bending's measurement, the perihelion choosing the reading and testing the choice, Section 2.3). (c) \emph{The metric.} S.1 at the paces reads, with $\nu = \num / \den$, $2(1 - \cos\omega) = 2(1 - \nu)N^2 + (\nu / 3)\,N^4\sum_a s(k_a)^2$ with $s(x) = 2\sin(x / 2)$; at long wavelength and a small gap, where $c_m^2 \to \nu / 3$, $\omega^2 = N^2\,[\omega_0^2 + c_m^2\,(k / h)^2]$ with $k / h = kN$: Klein-Gordon's on a static metric with the lapse $N$ and the spatial scale $h$, $ds^2 = -c^2N^2 dt^2 + h^2 d\mathbf x^2 = -e^{-2U}c^2dt^2 + e^{2U}d\mathbf x^2$, $c^2 = 1/3$, $t$ in intervals and $\mathbf x$ in Links, the exponential metric in isotropic coordinates \cite{yilmaz}, whose post-Newtonian expansion is $g_{00} = -(1 - 2U + 2U^2)$ and $g_{ij} = (1 + 2U + 2U^2)\delta_{ij}$: $\gamma = 1$ and, for this metric, $\beta = 1$, the spatial second order $2$ against Schwarzschild's $1.5$ (S.19 for $\gamma = \beta = 1$; the second order this derivation's). (d) \emph{The acts' factors.} Rule3 measured in the Node's units: a count per Node is a count per proper volume $h^3$, a step per interval is a step per proper interval $N$, and the write is Rule3's second difference in time, so a count source carries $N^2 / h^3 = N^5 = p_x p_y p_z / (p_0\Gamma^2)$ (check: $\Gamma^3 N^6 / (\Gamma N\cdot\Gamma^2) = N^5$); the Wronskian, one time difference, carries $N / h^3 = N^4 = p_x p_y p_z / (p_0^2\Gamma)$ and its readers' turn the remaining $N$, $\theta = (\mathrm{level} / \Gamma)(p_0 / \Gamma)$, the turn per proper interval; the Link angle is plain by the law's line, the Port's phase the odd lines' level times $4\omega_0 / c^2$, $c^2 = 1/3$ light's own, read with no other factor (S.56 (a), (f)). With a tension the three rulers differ, $p_0 / p_a$ along the axis $a$, and the volume is their product. (e) \emph{The truncation.} The clock's quadratic form $(\Gamma - c)^2 + c^2 = \Gamma^2(1 - 2U + 2U^2)$ is $\Gamma^2 e^{-2U} = \Gamma^2(1 - 2U + 2U^2 - 4U^3 / 3 + \dots)$ to the second order, differing at $-4U^3 / 3$; the Link's $(\Gamma - 2c)^2 = \Gamma^2(1 - 4U + 4U^2)$ is $\Gamma^2 e^{-4U} = \Gamma^2(1 - 4U + 8U^2 - \dots)$ to the first order only. \emph{Check:} at $U = 0.05$ the clock's form gives $0.905000$ against $e^{-0.1} = 0.904837$, the difference $1.63 \times 10^{-4}$ against $4U^3 / 3 = 1.67 \times 10^{-4}$; the Link's gives $0.8100$ against $e^{-0.2} = 0.8187$. The horizon of the quadratic forms at $c = \Gamma / 2$ is the truncation's zero of $\Gamma - 2c$ and not the law's (S.3). \emph{Status:} derived under the one assumption, fixed by the bending's measurement (assumption, fitted); lattice.
\end{derivation}

\begin{derivation}[the preferred-frame parameter of a diagonal metric, Section 8.1]
\emph{Inputs:} the model's metric is diagonal in the lattice's frame by construction: the content enters the clock $p_0$ and the Link pace $p_a$, both real and even in time and in space, the content's rest is Poisson's in the lattice's frame with no speed in its kernel, and no odd real Link read exists, since a real asymmetric hopping with a curl breaks the reciprocity the conserved form needs (Theorem 3); so $g_{0a} \equiv 0$ in the lattice's frame for every motion of the source (Section 8.1; a write with no odd Link part; the law's write (Section 2.3, S.56) carries the odd part and gives S.56 (g)'s $0$ for a plane configuration of one sense, a body $\alpha_1 = 8f - 8$ at its share $f$). The parametrised post-Newtonian metric in the theory's preferred frame (Will's review \cite{will}, in the standard gauge): $g_{0j} = -\tfrac{1}{2}(4\gamma + 3 + \alpha_1 - \alpha_2 + \zeta_1 - 2\xi)\,V_j - \tfrac{1}{2}(1 + \alpha_2 - \zeta_1 + 2\xi)\,W_j$, with $V_j$ and $W_j$ the two vector potentials of the matter current, related by $V_j - W_j = \chi_{,0j}$ for the superpotential $\chi$.
\emph{Steps:} (a) The post-Newtonian gauge freedom $x^0 \to x^0 + \lambda\,\chi_{,0}$ adds $\lambda\chi_{,0j} = \lambda(V_j - W_j)$ to $g_{0j}$, moving weight between $V_j$ and $W_j$ and leaving the sum of their two coefficients, $(4\gamma + 3 + \alpha_1 - \alpha_2 + \zeta_1 - 2\xi) + (1 + \alpha_2 - \zeta_1 + 2\xi) = 4\gamma + 4 + \alpha_1$, gauge-invariant. (b) A metric whose $g_{0j}$ vanishes in its preferred frame for every source current has both coefficients $0$ in some gauge, so the invariant sum vanishes: $4\gamma + 4 + \alpha_1 = 0$, $\alpha_1 = -4(1 + \gamma)$, which is $-8$ at the model's $\gamma = 1$ (S.19 in $U$). (c) The frame-dragging precession of a gyroscope near a rotating body is $(4\gamma + 4 + \alpha_1) / 8$ of general relativity's Lense-Thirring value, $1$ in general relativity, $0$ under the write without the odd Link part and $0.50$ for quartz under the law's write ($19.6$ against $37.2 \pm 7.2$ milliarcseconds per year \cite{everitt}, S.56 (g), the paper's Section 8.1): the write without the odd part's $0$ is the same fact as its $\alpha_1 = -8$, read by two experiments of nature. (d) Had the model's read been covariant, the boost of its static metric $N^2 = e^{-2U}$, $h^2 = e^{2U}$ to the velocity $v$ would give $g_{0i} = \gamma_L^2 v_i (N^2 - h^2) = -4Uv_i$ to first order ($\gamma_L$ Lorentz's, not the post-Newtonian $\gamma$ of (b)), Einstein's, with no new coefficient; the diagonal read in the lattice's frame is what removes it, so the finding is the frame's and not a number's (the carrier's derivation, C.4). \emph{Check:} $-4.000$ for the boosted cross term at first order; nature's $\alpha_1$ from the binary pulsar PSR J1738+0333 is $|\alpha_1| < 4 \times 10^{-5}$ \cite{shao} and from lunar laser ranging $10^{-4}$ \cite{muller}, so the finding stands at $2 \times 10^5$ and $10^5$; Gravity Probe B's drag at $5$ standard deviations; the two are one finding of that write, not two. \emph{Domain:} $\alpha_1$ alone is gauge-invariant here; with $g_{0j} = 0$ in the lattice frame, $g_{00} = -e^{-2U}$ and $g_{ij}$ isotropic, the standard post-Newtonian gauge, the two $g_{0j}$ coefficients vanish separately and give $\alpha_1 = -8$ and $\alpha_2 = -1$ together ($\alpha_2 = -1$ with $\zeta_1 = 2\xi$, $\zeta_1 = 0$ under the post-Newtonian conservation laws and $\xi = 0$, the preferred-location parameter, taken; $\alpha_1 = -8$ needs no such condition), the sum $4\gamma + 4 + \alpha_1 = 0$ the gauge-invariant statement, against $|\alpha_2| < 4 \times 10^{-7}$ (the Sun's spin \cite{nordtvedt}) and $< 1.6 \times 10^{-9}$ (solitary pulsars \cite{shao2013}); under the write without the odd Link part $\alpha_1 = -8$ holds for every motion of the source; under the law's write (S.56) the count's flux sources $\mathbf V$ and the parameters are S.56 (g)'s; $\alpha_2 = -1$ holds for a static source and is not derived for a moving one, every value of order one being excluded by $|\alpha_2| < 4 \times 10^{-7}$ in any case; the condition below is $\alpha_2$'s alone, the velocity-dependent terms of $g_{00}$ depending on how a moving record's count writes, which the law does not fix: named, not claimed. \emph{Status:} derived; the frame's parameters of a write with no odd Link part ($g_{0j} \equiv 0$): $\alpha_1 = -8$, $\alpha_2 = -1$; the law's write (Section 2.3, S.56) carries the odd part and gives S.56 (g)'s $0$ for a plane configuration of one sense, a body $\alpha_1 = 8f - 8$ at its share $f$; $\alpha_1 = -8$ is named in the paper's Sections 1.2 and 9 as the write without the odd Link part's; clicks.
\end{derivation}



\begin{derivation}[the mass defect at a kept count and the binding bound, Section 8.1]
\emph{Inputs:} Theorem 3, the weighted form conserved by Rule3 at fixed paces for any records; the invariant of S.14, a body of $C$ units of the invariant at the rotation $\omega_b$ having the share $C\sin\omega_b$ (the energy per quantum $\sin\omega_b$ in units of $T$, Planck's line); the rest rotation in a content, $\omega_0(U) = N\omega_0$ to the order $\omega_0^2$ (S.15 with $f \to 1$); the massless kernel of a holder at the weight $W$: the level at $r$ Links from a count $C$ is $3CW\,G_0(r) / E$ with $G_0$ the lattice Green's function per unit source, $G_0(0) = 0.2527$ (Watson) and $1 / (4\pi r)$ far, $e^{-r / R}$ beyond a gapped holder's range (S.11); S.24's Link bound for the weight's scale; S.26's coupling for the sign holder's part; C.6 (b)'s binding at a size and S.5's share identity; the measured mass defects \cite{ame2020}. \emph{Steps:} (a) \emph{No energy from combinatorics.} Two records that approach, overlap, exchange quanta through their tails, join into one standing record or part again are records stepped by Rule3; if the paces do not change, the weighted total form before and after is the same number in the rational line (Theorem 3), the plain total where the paces are uniform, within the remainder term of the share's change, Section 3.4 (S.5), with the rounding, and a click moves share from one record to another and conserves the total in the credit's unit, the weighted total where the two records stand at one pace. So no joining, parting or crediting of clusters lowers the share at a kept count: there is no energy released, and a nucleus bound that way would have no binding energy. (b) \emph{What the defect is.} The one act that changes a body's share at a kept count is a change of its paces, a content read; a body in its own content $x_{\mathrm{self}}$ rotates at $N\omega_0$ with $N = e^{-x_{\mathrm{self}}}$, so its click mass is $C\sin(N\omega_0) = C\sin\omega_0 - C\omega_0 x_{\mathrm{self}}\cos\omega_0 + O(x^2)$: the mass defect is at most the self-read's redshift $C\omega_0 x_{\mathrm{self}}$ to the order $\omega_0^2$, the bound mode's depression $\omega_0 - \omega_b$, the potential's redshift less the bound mode's kinetic rise (half the redshift for a $1 / r$ well by the virial, C.6 (a), (b)), which is what C.6 (b) computes at the fixed point; this is the law's ``a body that binds is lighter in every holder it sources by its binding''. \emph{Check:} a well deepened from $300$ to $600$ levels at $\Gamma = 6{,}000$ lowers the share read at the vacuum's coefficients (Section 3.4) by $3.8$ percent to first order ($x\omega_0 / \tan\omega_0$ at $x = 0.05$; $3.86$ by $\sin(N\omega_0)$ from $e^{-0.05}$ to $e^{-0.10}$); the share at the Node's own paces carries the factor $p_i^2(t_1) / p_i^2(t_2)$ of the drift ratio (S.14) and rises. (c) \emph{The bound.} A body of count $C$ over $r$ Links, bound by a holder at the weight $W$ and the divisor $E$, writes on itself at most the one-Node value, the lattice kernel at the source, $x_{\mathrm{self}} \le 3G_0(0)\,CW / (\Gamma E) = 0.758\,CW / (\Gamma E)$, independent of $r$; for a uniform ball of radius $r$ the centre's self-level is $1.5$ times the whole count's level at $r$ ($0.36\,CW / (\Gamma E\,r)$ at the centre) and the count's share-weighted average $1.2$ times it, the first-order shift of a mode under a non-uniform content, a profile-dependent scaling and not a bound, so
\[
\frac{\text{defect}}{\text{rest}} \;\le\; 0.76\,\frac{CW}{\Gamma E}\,, \qquad \frac{\text{defect}}{\text{rest}} \times r \;\approx\; 0.29\,\frac{CW}{\Gamma E} \text{ for a uniform ball, } 0.36 \text{ at its centre}.
\] (d) \emph{The numbers.} The deuteron's defect is $2.224\,\mathrm{MeV} / 1875.6\,\mathrm{MeV} = 1.2 \times 10^{-3}$ and the nucleon's in iron $8.8 / 939 = 9 \times 10^{-3}$ \cite{ame2020}, so, with the ball's estimate, they need $CW / (\Gamma E r) \ge 4 \times 10^{-3}$ and $3 \times 10^{-2}$; at $r$ about $6 \times 10^{10}$ Links for a nucleon ($0.84$ fm \cite{codata} at the MAGIC Link, $1.3 \times 10^{-26}$ m) this is $W$ of order $10^{19}$ to $10^{22}$ at $\Gamma = 10^{14}$ (an illustration; the bound's form is the claim and the number is not) and $E = 1$, the lower end for a cloud of $C = 1{,}836$ records in the electron's unit and the upper for $C = 1$, at the deuteron's defect (iron's raises each by $7.9$), taken for illustration, growing with $\Gamma$, which no click bounds (S.26), under the reading of the nucleon as a cloud, which C.6 replaces (the nucleon one quantum of its own family, its binding the Yukawa well with its declared $W$); the bound's form does not move with the reading of $\Gamma$, its numbers do. The credit's quantum does not supply it: the deuteron's $2.22$ MeV is $4.35$ quanta of $0.511$ MeV and iron's nucleon's $8.8$ MeV $17.2$, not wholes, and no act credits share to nowhere. (e) \emph{The one lattice lightening without content}, the band's $\sin\omega_b$ falling to $0$ at the band's top, acts at rotations of order one and makes a body nearly massless, not $99$ percent of its quanta's mass: not a nuclear binding. \emph{Check of the kernel:} Watson's integral gives $G_0(0) = 0.252731$, $3G_0(0) = 0.7582$; $G_0(1) = G_0(0) - 1 / 6 = 0.08606$, $3G_0(1) = 0.2582$, $1.08$ of the far kernel $3 / (4\pi) = 0.2387$ at $r = 1$. \emph{Status:} derived, a bound conditional on the declared $W$, which no click bounds; the size law kept; lattice for the well, clicks for the defect.
\end{derivation}

\begin{derivation}[the gapped holder's kernel, Section 4.1]
\emph{Inputs:} a holder's static line in the vacuum, from the start's condition (Section 4.1): $6\,\den\,L - \num\,S_6(L) = 3\,\den\,\sigma$, that is $\num\,\Delta L - 6(\den - \num)\,L = -3\,\den\,\sigma$ with $\Delta = S_6 - 6$ the lattice Laplacian and $\sigma$ the source per interval; the reach $\kappa^2 = 6(\den - \num) / \num$ of S.11, for $0 < \num \le \den$ as S.11 has it (the division by $\num$ below and the positive kernel need it; at $\num = 0$ the static line has local support, at a negative numerator the alternating tail of S.11); the write and the read of every rotation holder the same act, $W$ over $E_s T$ into the level and the level over $\Gamma$ into the angle (Section 2.3), whatever the holder's pair; S.24's Link bound and S.26's far kernel; the Yukawa potential's critical coupling $0.8399$, a standard result \cite{rogers1970}, named where the threshold $0.840$ is used. \emph{Steps:} (a) Dividing by $\num$: $\Delta L - \kappa^2 L = -(3 / \nu)\,\sigma$, whose lattice Green's function is $G_\kappa(\mathbf r) = \int \frac{d^3k}{(2\pi)^3}\,\frac{e^{i\mathbf k\cdot\mathbf r}}{6 - 2\sum_a\cos k_a + \kappa^2}$, far away $e^{-\kappa r} / (4\pi r)$ per unit source: the level is $3\,e^{-r / R} / (4\pi\nu r)$ per unit $\sigma$ with $R = 1 / \kappa$, Yukawa's kernel with the massless one's constant, the range the gap's and the strength at the Link the massless one's: $G_\kappa(1) = G_0(1)\,(1 - 0.92\,\kappa + O(\kappa^2))$, so the strength at one Link is the massless one to one part in $R$ (and the factor $1 / \nu = 1 + O(\text{gap})$). \emph{Check:} the exact infinite-volume integrals, $G_0(1) = 0.08606$ (Watson) and $G_{1/20}(1) = 0.08216$ by the exact lattice Yukawa integral, the ratio $0.9546$ at $R = 20$ Links ($0.08209$ and $0.954$ at the first order in $\kappa$, $G_0(1) - \kappa / 4\pi$), the slope $1 - 0.9246\,\kappa$; $G_{1/400}(1) / G_0(1) = 0.9977$; the $400^3$ mesh's own pair at $R = 20$, $0.0857$ and $0.0822$, has the ratio $0.959$ against the exact $0.9546$. (b) \emph{The coupling.} The sign holder's force constant, $\alpha_{\mathrm{law}}$ (S.26), is the write's half, the kernel's $3 / (4\pi)$ and light's $\sqrt 3$, none of them the holder's own; a gapped holder written and read by the same acts has the same three factors and the kernel above, so it couples at $\alpha_{\mathrm{law}}\,4\pi G_\kappa(1) / \nu$ at one Link, the far kernel's $\alpha_{\mathrm{law}} / r$ times three factors, the lattice's $4\pi G_0(1) = 1.0815$, the gap's $G_\kappa(1) / G_0(1) = 1 - 0.92\,\kappa$ and the source's $1 / \nu = 1 + \kappa^2 / 6$: $1.0329\,\alpha_{\mathrm{law}}$ at $R = 20$ ($1.0815 \times 0.9546 \times 1.0004$), $1.08\,\alpha_{\mathrm{law}}$ for the massless holder, and $\alpha_{\mathrm{law}}\,e^{-r / R} / (\nu r)$ beyond: every rotation holder Rule3 steps, massless or gapped, abelian or three-line, couples at the charge's $\alpha$ at the Link. A coupling $k / (\Gamma E_s)$ itself, the universe's $1 / 28.33$, would need the factor $8\pi / (3\sqrt 3) = 4.84$ on the kernel, which no gap supplies; a contact strength $3 / \kappa^2$ appears only for $\kappa > 1$, a range below one Link, where the kernel's lattice sum departs from its continuum integral. (c) \emph{The numbers.} The W boson's range $\hbar / (m_W c) = 2.45 \times 10^{-3}$ fm and the pion's $1.41$ fm are $1.9 \times 10^8$ and $1.1 \times 10^{11}$ Links at the Link's bound $1.3 \times 10^{-26}$ m (S.24), so $1 / R$ is $5 \times 10^{-9}$ and $10^{-11}$: the strength at the Link is the massless one to those parts. A pion-gap holder at its range with the law's coupling has the well's depth $\alpha\hbar c\,e^{-1} / R = 1.44\,\mathrm{MeV\,fm} \times 0.368 / 1.41\,\mathrm{fm} = 0.38$ MeV at $r = R$ and binds no state: a screened Coulomb well binds only for $\mu g R / (\hbar c)^2 \ge 0.840$, and at $g = \alpha\hbar c$, $R = 1.414$ fm, $\mu = 469.5$ MeV (nature's deuteron binds two nucleons about their common centre, so this is nature's own reduced mass, the comparison's convention for nature's column; the law's body is one quantum about the pinned centre at $\mu = m^*$, postulate 2 of C.1, the two columns differing by the factor $2$ in $\mu$ with $W \propto 1 / \mu$ at fixed $g^2$, C.6 (a)) the number is $0.0245$, $34$ times under the threshold, so the deuteron's $2.22$ MeV stands against no binding at all; a holder with the W boson's gap with the law's $\alpha$ and nature's $m_W$ gives the Fermi constant $G_F = \pi\alpha / (\sqrt 2\,m_W^2) = 2.5 \times 10^{-6}\,\mathrm{GeV}^{-2}$ against nature's $1.166 \times 10^{-5}$ \cite{pdg}, short by $4.6$, or $m_W = 37$ GeV from $G_F$ against $80.4$ \cite{pdg}. \emph{Check of Yukawa's identity:} the range $R = \sqrt{\num / (6(\den - \num))}$ equals the Compton length $c / \omega_g$ of the holder's own travelling quanta to the gap's first order (S.11), with no input. \emph{Status:} derived, one derivation, the comparisons with nature's numbers supporting findings; lattice for the kernel, clicks for the comparisons, supporting findings.
\end{derivation}

\begin{derivation}[the Greenberger-Horne-Zeilinger (GHZ) correlation under the declared credit, Section 7.2]
\emph{Inputs:} the GHZ family, light's pair with four real lines laid as one event and never summed at a Node, every part laid alike and stepped by the same line, so the four are identical at every NodeDetector (the determinism of equal parts, Section 7.2); three NodeDetectors A, B, C, each with its setting $(p, q) = (\cos a, \sin a)$ up to the norm and its pattern, one integer pair per part: A $(p, p, q, q)$, B $(p, q, p, q)$, C $(p, -q, -q, -p)$, the $+$ port reading the part $k$ as $e_k(+) = \alpha_k p + \beta_k q$, with $\alpha_k$, $\beta_k$ here the pattern's selectors $(1, 0)$, $(0, 1)$, $(0, -1)$, $(-1, 0)$ multiplying $(p, q)$, not S.47's laid amplitudes, and the $-$ port as $e_k(-) = \alpha_k(-q) + \beta_k p$; the credit, the pair's own over three sides: $J(\pm, \pm, \pm) = \bigl(\sum_{k=1}^4 e_k(A)\,e_k(B)\,e_k(C)\bigr)^2$ on equal parts, the common factor of the parts' amplitudes leaving every ratio.
\emph{Steps:} (a) With $c_a = \cos a$, $s_a = \sin a$ and the same for $b$ and $c$, the $+$ ports read A as $(c_a, c_a, s_a, s_a)$, B as $(c_b, s_b, c_b, s_b)$ and C as $(c_c, -s_c, -s_c, -c_c)$, so $\sum_k e_k(A)e_k(B)e_k(C) = c_a c_b c_c - c_a s_b s_c - s_a c_b s_c - s_a s_b c_c = \cos(a + b + c)$. A $-$ port replaces $(p, q)$ by $(-q, p)$: A's $-$ port reads $(-s_a, -s_a, c_a, c_a)$, and the sum is $-s_a c_b c_c + s_a s_b s_c - c_a c_b s_c - c_a s_b c_c = -\sin(a + b + c)$; by the symmetry of the three patterns every port triple with an even number of $-$ ports gives $\pm\cos(a + b + c)$ and with an odd number $\pm\sin(a + b + c)$ (for instance $(-, -, +)$ gives $-\cos$, $(-, -, -)$ gives $\sin$). (b) The eight joint shares are therefore $\cos^2(a + b + c)$ at the four even triples and $\sin^2(a + b + c)$ at the four odd, summing to $4$ at every setting; the normalised shares are $\cos^2 / 4$ and $\sin^2 / 4$. (c) $E_3$, the shares signed by the product of the three signs, is $4\cdot\cos^2 / 4 - 4\cdot\sin^2 / 4 = \cos 2(a + b + c)$. A pairwise $E_{AB}$ sums over C's two ports first, which gives $(\cos^2 + \sin^2) / 4 = 1 / 4$ for each of the four sign pairs of A and B, so $E_{AB} = (1 - 1 - 1 + 1) / 4 = 0$, and the same for the other two pairs; each single marginal is $1 / 2$: no signalling and no correlation in any two, the correlation standing in the three alone. (d) Mermin's combination $M = E(a, b', c') + E(a', b, c') + E(a', b', c) - E(a, b, c)$ at $a = b = c = 0$ and $a' = b' = c' = \pi / 4$ is $\cos\pi + \cos\pi + \cos\pi - \cos 0 = -4$, $|M| = 4$, where every local assignment of signs to the three ports, and every mixture of them, gives $|M| \le 2$ \cite{mermin}. (e) With integer settings $(p, q)$ the norms $p^2 + q^2$ multiply every share alike and cancel, so $x = (1, 0)$ and $y = (1, 1)$ give the eight shares exactly, $0$ and $1 / 4$ on the three settings with one $x$ ($E_3 = -1$) and $1 / 4$ and $0$ on $(x, x, x)$ ($E_3 = 1$), $M = -4$ exactly. \emph{Check:} the trigonometric identity $\cos(a + b + c) = c_a c_b c_c - c_a s_b s_c - s_a c_b s_c - s_a s_b c_c$ and $\sin(a + b + c) = s_a c_b c_c + c_a s_b c_c + c_a c_b s_c - s_a s_b s_c$. \emph{Nature's angles:} the law's $E_3 = \cos 2(a + b + c)$ is nature's $\cos(\theta_A + \theta_B + \theta_C)$ at $\theta = 2a$, nature's $x$ and $y$ at $\theta = 0$ and $\pi / 2$ the law's bases $(1, 0)$ and $(1, 1)$, as Bell's angles are (S.10). \emph{Status:} derived under the declared credit, the credit the one assumption (Section 7.2), computed for the numbers ($M = -4$); clicks.
\end{derivation}

\begin{derivation}[the prediction, the check and the open question, Section 5.2]
\emph{(1) A clock in two places.} \emph{Inputs:} a bound body's beat between two of its modes, $\Omega = \omega_j - \omega_i$, scales by the clock $N = e^{-U}$ at the body's Node (S.31, exactly at long wavelength, to $(1 - N^2)\omega^2 / 8$ at finite $\omega$); a record laid on two paths is one record (Section 2.5.1), and the body's internal state is carried on each path by its own rotation. \emph{Steps:} on two paths at the contents $U_1$ and $U_2$ for a hold time $\tau_{\mathrm{hold}}$ the internal beats, of frequency $\Omega = \omega_j - \omega_i$, accumulate the phases $\Omega N_1 \tau_{\mathrm{hold}}$ and $\Omega N_2 \tau_{\mathrm{hold}}$, differing by $\Omega(N_1 - N_2)\tau_{\mathrm{hold}} = \Omega\,\Delta U\,\tau_{\mathrm{hold}}$ to first order; for a body prepared in an equal superposition of the two modes the internal states on the two paths overlap as $|\cos(\Omega\,\Delta U\,\tau_{\mathrm{hold}} / 2)|$, which multiplies the path fringes' visibility. \emph{Check:} strontium's clock line $429.228$ THz \cite{sr_clock}, $\Omega = 2\pi \times 4.29228 \times 10^{14} = 2.697 \times 10^{15}\,\mathrm{s}^{-1}$; $\Delta U = gh / c^2 = 9.807 \times 1 / (2.998 \times 10^8)^2 = 1.091 \times 10^{-16}$ for one metre ($g$ Earth's acceleration, $h$ the height); $\tau_{\mathrm{hold}} = 1$ s: $\Omega\,\Delta U\,\tau_{\mathrm{hold}} = 0.294$ rad, the visibility $\cos(0.147) = 0.989$; the correction $(1 - N^2)(\omega_1^2 + \omega_1\omega_2 + \omega_2^2) / 24$, $(1 - N^2)\omega^2 / 8$ at $\omega_1 = \omega_2$, is $1.0 \times 10^{-37}$ at the axis bound's $\omega = 2.0 \times 10^{-14}$, $3.4 \times 10^{-29}$ at LHAASO's ($U = 10^{-9}$; S.31). The semiclassical expectation \cite{zych} has the same magnitude; the law's statement is that the one record's beat reads the content at its Node, so the dephasing is the clock's and no new line.
\emph{(2) The gravity of light.} \emph{Inputs:} the one write: a holder of the content gains the form $D$ of every record over its wall (Section 2.2); Theorem 3 and S.5: a travelling light record's form summed over its Nodes is its share, its count $n$ (energy $nT$), and $D = A^2\sin^2\omega$ per Node for a travelling wave as for a standing one; the tension sources the axis parts alone, which enter the Link's factor and not the clock (Section 2.3). \emph{Steps:} a light packet of count $n$ writes $n$ quanta per interval into the massless holder as a static body of count $n$ does; this, the integrated write, is the derived part. The premise taken next is the static far kernel, $3n / (4\pi r)$ over $E_g$, which is the field of a source at rest: under it the clock's well is $U = 3n / (4\pi\Gamma E_g r)$, Eq.~(11)'s $G_{\mathrm{clock}}$ times $n$, and a static body falls toward the packet at $G_{\mathrm{clock}}\,n / r^2$, that is $1 \times E / c^2$ for the packet's active mass. The field of the packet itself, a source moving at the holder's own characteristic speed, is not derived here and is open: the retarded field of a sub-characteristic source of the same holder is $\phi \propto 1 / \sqrt{s_\parallel^2 + (1 - \beta^2)\,s_\perp^2}$, directional and $\gamma_L$ times the static field across the motion, singular at the characteristic speed in the continuum, which the lattice's dispersion blunts and does not remove. In general relativity a pencil of light pulls a static body at $2 \times E / c^2$ \cite{tolman}, the beam's longitudinal stress $\rho c^2$ along the beam and $0$ across entering $\rho + \sum_i p_i / c^2 = 2\rho$; in the model the stress of the wave, the Link's booking $2\,\num\,b^2\sin^2 k$ and the Node's tension $\num\,b^2\sin^2 k$ (S.7), enters the Link's factor on its axis and the clock reads the form alone, so under the static kernel the factor is $1$: a prediction for a declared universe, no measurement of nature, the moving packet's field open.
\emph{(3) The black hole's quanta.} \emph{Inputs:} the composed paces (S.34): $p_0 = \Gamma(1 - 1 / \Gamma)^c$, $p_i = p_0^2 / \Gamma$ rounded to the unit, every Link coefficient $R_a = 2\,\num\,p_i^2 q^2 / \Gamma^2$; Theorem 1, the step a bijection; a collapse ends in Nodes whose content exceeds the frozen content (Section 2.3). \emph{Steps:} $p_i$ rounds to $0$ where $\Gamma e^{-2U} < 1 / 2$, that is $U > \tfrac12\ln(2\Gamma)$, the frozen content $c_f = (\Gamma / 2)\ln(2\Gamma)$ ($4.70\,\Gamma$ levels, $28{,}206$ at $\Gamma = 6{,}000$ with $p_0$ rounded to the unit before squaring, the law's rounding); there the Node's own six coefficients are $0$, so it reads nothing of its neighbours, while its neighbours' coefficients, each the reading Node's own, still read it: the cut is one-way, the frozen Node dark by the share's declaration and not cut; what leaves a frozen region is the levels frozen in it, read by the surface's active shell at its pace, matter's a near-static source of a static well, light's a level drifting linearly that no line of the law closes; whether the law's black hole radiates, and what, is not established by the line and is open, $T_H = 0$ not claimed (a reading of the line excluded by two derivations). Hawking's temperature \cite{hawking1974}\cite{hawking1975} for one solar mass, the number the open question is set against (a comparison only, as Bekenstein's entropy \cite{bekenstein} beside it), is $\hbar c^3 / (8\pi G M k_B) = 6.17 \times 10^{-8}$ K. The region's entropy as a count of the state is the count of its kept levels, the declared bound's logarithm per frozen Node; what it means for radiation is open. \emph{Check:} $\hbar c^3 / (8\pi G M_\odot k_B)$ with $M_\odot = 1.988 \times 10^{30}$ kg gives $6.17 \times 10^{-8}$ K; $\tfrac12\ln(12{,}000) = 4.696$. \emph{Status:} (1) derived under the one assumption, S.31's scaling; untested; clicks. (2) derived, the integrated write, the coefficient under the static far kernel, the moving packet's field open; clicks, untested. (3) derived under the one assumption, what it means for radiation open; lattice.
\end{derivation}

\begin{derivation}[what a body writes into the holder of the sign, Section 4.3]
\emph{Inputs:} the continuity identity of S.33: at every Node $w[W_i(t+1) - W_i(t)]$ equals the net of the Link fluxes $R_a g$, so the weighted sum $\sum_i W_i / p_i^2$ over a body is exactly conserved for a static self-level, the plain sum its uniform-pace case (S.33 (a), S.9), and a single Node's $W_i$ changes only by a current through its Links (the credit books the conserved form's own current through a front, the Link's factor squared its one weight, the Node's $1 / p_i^2$ cancelling against the read's $p_i^2$, so the count is the conserved form's flux, exact at paces fixed in time; a gapped quantum's share in a well differs from $W_{\mathrm{rec}}$ by the rest term's shift $4U(\den - \num) / \den$, the law's physics read as the count; where the paces move the share changes by the work term, the body's; a ratio of two readings at one pace is unchanged); in particular at one isolated Node, $z_{t+1} + z_{t-1} = 2\cos\omega_t\,z_t$ with any real $\omega_t$ keeps $W = \mathrm{Im}(\bar z_t z_{t-1})$ exactly ($W_{t+1} = \mathrm{Im}(2\cos\omega_t|z_t|^2 - \bar z_{t-1}z_t) = W_t$). The holder's line at the vacuum's pace, $\ell_{t+1} - 2\ell_t + \ell_{t-1} = \tfrac13\Delta\ell + s_t$, with $s$ the write per interval (the Wronskian over $E_s T$): the wave equation at the speed $1 / \sqrt 3$ with the retarded solution $\ell(\mathbf r, t) = 3\,s(t - \sqrt 3\,r) / (4\pi r)$ for a point source; a static record has $D = 0$ and $F = 0$ and carries no quanta (Section 2.4); the share of a travelling wave of the holder's record is $A^2\sin^2\Omega / T$ quanta per Node (S.5 at $[1, 1]$) moving at $1 / \sqrt 3$.
\emph{Steps:} (a) \emph{A steady rotation.} A body at rest with the standing plane record $z_i = \phi_i e^{-i\omega_b t}$, $\phi$ real up to one phase, has $W_i = \phi_i^2\sin\omega_b$ constant at every Node and the Link quantity $g_{ij} = \mathrm{Im}(\bar z_i z_j) = 0$: the source is static, the holder's level is its rest, no travelling event is born, no click anywhere. (b) \emph{Uniform motion.} A uniformly moving source writes no travelling light in the long-wavelength band, $\omega = c|k|$, where no sub-$c$ source is phase-matched: its retarded level is the uniformly moving rest, which carries no energy flux to infinity, as for a uniformly moving charge in Maxwell's theory. On the lattice light's phase speed falls to $0.392$ Link per interval at the zone's edge along an axis, so a source faster than that is matched to zone-edge light (the forced recurrence's denominator $2\cos(k\cdot v) - 2\cos\omega(k)$ vanishing there), which a body smooth over many Links sources only by its Fourier weight at that wave number, a small number set by the body's tail ($e^{-\pi^2 R^2 / 2}$ for a Gaussian of rms $R$ Links, $\tanh^2(1 / (2R))$, of order $1 / (4R^2)$, for an exponential tail), the exclusion's work carried by the group-speed condition that follows; at the computed pair no body reaches $0.392$ (matter's largest group speed $0.385$ along the diagonal); from the pair $[9, 10]$ upward a body moving along an axis can reach it, so the exclusion is an approximation with a named domain and not an exact one (a finding of two derivations). (c) \emph{A breathing body.} Breathing is a change of the amplitudes $A_i(t)$ across the body by currents between its Nodes; the weighted total $\sum_i W_i / p_i^2$ is kept, so where the paces are uniform over the body the monopole of the source is constant and radiates nothing, exactly; where they vary across the body the source written with the plain weight departs from the conserved total at first order in the pace's variation, $2\,\Delta p / p$ of the total, below $10^{-9}$ for an atom in a laboratory well (the write's weight is the law's line, open). The far level of a spherically symmetric source distribution $\sigma_i(t)$ of constant total $S$ is, by the retarded integral with the retardation expanded, the known multipole expansion of the holder's wave equation's solution, standard mathematics of the law's line ($\sigma(t - \sqrt 3|\mathbf r - \mathbf r'|) = \sigma(t_r) + \sqrt 3(\hat{\mathbf r}\cdot\mathbf r')\dot\sigma + \tfrac32(\hat{\mathbf r}\cdot\mathbf r')^2\ddot\sigma + \dots$, the first order averaging to $0$ over the sphere and $\langle(\hat{\mathbf r}\cdot\mathbf r')^2\rangle = r'^2 / 3$),
\[
\ell(r, t) = \frac{3S}{4\pi r} + \frac{3}{4\pi r}\cdot\frac{1}{2}\,\ddot Q(t - \sqrt 3\,r), \qquad Q(t) = \sum_i \sigma_i(t)\,r_i^2 ,
\]
so a body of mean square radius $\langle r^2\rangle$ breathing with its second moment $Q(t) = Q_0(1 + 2\delta\cos\Omega t)$, $\delta$ defined by it with the monopole $S$ constant and $Q_0 = S\langle r^2\rangle$, writes the travelling level $\ell_1 = 3\,s_{\mathrm{eff}} / (4\pi r)$ at the frequency $\Omega$ with $s_{\mathrm{eff}} = \delta\,\Omega^2\,S\langle r^2\rangle$: the light born carries the breathing's frequency $\Omega$ and not the rotation $\omega_b$, as Section 4.3 says, with the amplitude $S\cdot 2\delta$ ($\delta$ the breathing's depth) multiplied by $\Omega^2\langle r^2\rangle / 2$, the second moment's factor. The quanta radiated per interval through a far sphere are the share density times the speed times the area, $4\pi r^2\cdot\ell_1^2\sin^2\Omega / (\sqrt 3\,T) = (3\sqrt 3 / 4\pi)\,s_{\mathrm{eff}}^2\sin^2\Omega / T$, of order $\delta^2\Omega^6 S^2\langle r^2\rangle^2 / T$: the scalar analogue of the quadrupole formula, and a far NodeDetector's click rate is that current's share through its front. Section 4.3's three conclusions stand: no light from a steady rotation, none from a uniform motion in the long-wavelength band, light at the beat $\Omega$ from a breathing body. \emph{Check:} $9 / (4\pi\sqrt 3) = 3\sqrt 3 / (4\pi) = 0.4135$; the static kernel $3 / (4\pi r)$ is S.25's; the retarded expansion is the standard multipole one at the speed $1 / \sqrt 3$. \emph{Status:} derived, (b) by two derivations as stated; lattice for the Wronskian and the source, clicks for the far NodeDetector's rate.
\end{derivation}

\begin{derivation}[the two-mode line under the rotation, Section 4.3]
\emph{The convention:} this derivation's time turn $\mathcal U_t = e^{-i\theta_t}$ is the main text's convention and S.33's (S.33's step multiplied by $e^{-i\theta_t}$ gives this form): the sense $+$ rotates at $\omega_b + \theta$. \emph{Inputs:} the plane record's step under the holder of the sign (Section 3.5), the time turn $\mathcal U_t$ applied to the plain step, $z_{t+1} = \mathcal U_t\bigl(w^{-1}[S z_t + \sum_a R_a(e^{i\theta_a}z_{+a} + e^{-i\theta_a}z_{-a})] - \mathcal U_{t-1}z_{t-1}\bigr)$ with $\mathcal U_t = e^{-i\theta_t}$ the time turn between $t$ and $t+1$ and $\theta_a$ the Link phases, both the holder's level over $\Gamma$ at the interval's start; light as the holder's own record, a level $b_L\cos\omega_L t$ at the body with the gradient $\nabla b_L$ across it; the body's bound modes $\phi_m$, real eigenvectors of $\mathbf M$, symmetric at the paces uniform over the body taken here (one $w$), with $(\mathbf M\phi_m)_i = 2w\cos\omega_m\,\phi_{m,i}$, orthonormal; at nonuniform paces $\mathbf D\mathbf M$ is the symmetric one (S.45), the modes are orthonormal in the weights $d_i = 1 / p_i^2$ and every Link term below carries $d_i R_{ij}$ in place of $R_a$. \emph{Steps:} (a) \emph{The uniform turn is a gauge.} Put $z_t = e^{-i\Phi_t}u_t$ with $\Phi_{t+1} = \Phi_t + \theta_t$. Then $\mathcal U_t z_t = e^{-i\Phi_{t+1}}u_t$, $\mathcal U_t \mathcal U_{t-1} z_{t-1} = e^{-i\Phi_{t+1}}u_{t-1}$ and $z_{t+1} = e^{-i\Phi_{t+1}}u_{t+1}$, so $u$ obeys the plain line $u_{t+1} = \mathbf M u_t / w - u_{t-1}$ exactly, for any $\theta_t$: a time turn uniform over the body changes the record's phase and nothing of its moduli, its form or its modes; at a single Node light read as a turn drives no resonance and no absorption. (b) \emph{A constant turn.} For $\theta_t = \theta$ the plane solution $z = e^{-i\Omega t}$ has, under this section's step, $2\cos(\Omega - \theta) = 2\cos\omega_b$, so the sense $+$ rotates at $\omega_b + \theta$ and the sense $-$ at $-\omega_b + \theta$: both shift by $+\theta$, their magnitudes split by $2\theta$, the energy shift following the record's own sense, Section 3.5's line. (c) \emph{The Link phases couple the modes.} To first order in $\theta_a$ the Link read adds $V z = i\sum_a R_a[\theta_{i,a}\,z_{i+a} - \theta_{i-a,a}\,z_{i-a}]$ at the Node $i$, an $i$ times a real antisymmetric matrix, Hermitian. Write $z_t = \sum_m c_m(t)\,\phi_m\,e^{-i\omega_m t}$ with the $c_m(t)$ slowly varying (the mode amplitudes, not the kinetic scale $c_m$); $w(z_{t+1} + z_{t-1}) - \mathbf M z_t = -2iw\sum_m\sin\omega_m\,c_m'\,\phi_m e^{-i\omega_m t}$, and projecting $V z$ on $\phi_j$ with $\theta_{i,a}(t) = \theta_L\,\epsilon_a(i)\cos\omega_L t$ ($\epsilon_a$ the light's profile across the body; $\theta_L = |\partial_a b_L| / (\Gamma\omega_L)$, the Link phase of the gauge of (a) for light's time level $b_L$ with its gradient, so that $E_0 = \omega_L\theta_L$ reproduces the force field of S.26 (c)) gives $\langle\phi_j, V\phi_m\rangle = i\,\theta_L\cos(\omega_L t)\,K_{jm}$ with
\begin{align*}
K_{jm} &= \sum_{\text{Links } (i, i+a)} R_a\,\epsilon_a(i)\,\bigl[\phi_j(i)\phi_m(i+a) - \phi_m(i)\phi_j(i+a)\bigr],\\
K_{mj} &= -K_{jm}, \quad K_{mm} = 0,
\end{align*}
the inter-mode Link current of S.33's $g$ taken between two modes. So
\[
-2iw\sin\omega_j\,c_j' = i\theta_L\cos(\omega_L t)\sum_m K_{jm}c_m e^{-i(\omega_m - \omega_j)t};
\]
the product carries the frequencies $\omega_m - \omega_j \pm \omega_L$ and resonates at $\omega_L = |\omega_j - \omega_m|$, the difference line, for drives below the sum; at $\omega_L = \omega_j + \omega_m$ the same matrix element couples $a_m\phi_m e^{-i\omega_m t}$ to the other sense $b_j\phi_j e^{+i\omega_j t}$, since a second-order recurrence has both senses as solutions and the one-sense record is initial data: $a'' = +[\theta_L K_{jm} / (4w\sqrt{\sin\omega_j\sin\omega_m})]^2 a$, a growth at the rate $\Omega_R$ (S.16's $\epsilon_{ij}$ is $\theta_L K_{jm} / w$ here, one rate in two notations) with S.16's detuned exponents $\pm\sqrt{\Omega_R^2 - \Delta^2 / 4}$, the sum line S.16 itself names; with $K$ real $\frac{d}{dt}(|a|^2\sin\omega_m - |b|^2\sin\omega_j) = 0$, the Wronskian kept while the share grows, so the sum line creates content of the opposite sense in pairs, a pair line at $\omega_j + \omega_m \approx 2\omega_0$, in nature's names $2mc^2$, named and not claimed (the Floquet multipliers at $\Omega = \pi / 2$ for $\omega = \pi / 6, \pi / 3$, $\theta_L K / w = \epsilon = 0.01$: $0.9849$ and $1.0153$, the growth halving with $\epsilon$; at the difference, all four of modulus $1$); the one-mode line $2\omega_b$ of the pace read (S.16) is the degenerate case; and $K_{mm} = 0$ says a Link phase shifts no single mode at first order, the gauge statement of (a) for the Links. (d) \emph{Rabi's rate.} At $\omega_L = \omega_j - \omega_m$ the resonant terms are $c_j' = -\theta_L K_{jm}c_m / (4w\sin\omega_j)$ and $c_m' = +\theta_L K_{jm}c_j / (4w\sin\omega_m)$, so $c_j'' = -\Omega_R^2 c_j$ with
\[
\Omega_R = \frac{\theta_L\,|K_{jm}|}{4w\sqrt{\sin\omega_j\sin\omega_m}},
\]
a full transfer between the two modes every half period at resonance, one transition per click; detuned by $\Delta$ the transfer generator $[[-i\Delta / 2, -g], [g, i\Delta / 2]]$, with $g = \Omega_R$ and $\Delta$ the detuning, has the eigenvalues $\pm i\sqrt{g^2 + \Delta^2 / 4}$, so $P(t) = [g^2 / (g^2 + \Delta^2 / 4)]\sin^2(\sqrt{g^2 + \Delta^2 / 4}\,t)$, Rabi's law at the generalised frequency, oscillatory at every detuning with no threshold, the envelope's full width $4\Omega_R$ in the detuning at half height (the parametric exponents of S.16 belong to the sum line of (c), not here; a finding of two derivations). (e) \emph{The matrix element is the dipole.} For a uniform gradient, $\epsilon_a = 1$, the identity $\sum_i x_a(i)[\phi_j(\mathbf M\phi_m) - (\mathbf M\phi_j)\phi_m]_i = 2w(\cos\omega_m - \cos\omega_j)\langle\phi_j, x_a\phi_m\rangle$, evaluated with $\mathbf M$'s hopping and $x_a(i+a) - x_a(i) = 1$, gives $K_{jm} = 2w(\cos\omega_j - \cos\omega_m)\,d_{jm} = -4w\sin\tfrac{\omega_j + \omega_m}{2}\sin\tfrac{\omega_j - \omega_m}{2}\,d_{jm}$ with $d_{jm} = \langle\phi_j, x_a\phi_m\rangle$ the dipole: the velocity form equals the length form times the line's frequency, the discrete Ehrenfest relation; for a small line $\Omega_R = \theta_L\,\omega_L\,|d_{jm}| / 2$, nature's $d\cdot E_0 / (2\hbar)$ with the field $E_0 = \omega_L A_0 = \omega_L\theta_L$ and $\theta_L = |\partial_a b_L| / (\Gamma\omega_L)$ the phase per Link of the vector potential, the rotating-wave factor $1 / 2$ included; a static holder level with a gradient across the body ($\omega_L = 0$, $\nabla b_L \ne 0$) mixes the modes without a transition, Stark's shift; a uniform static phase is a gauge. \emph{Check:} the slowly varying method is S.16's; $K_{mm} = 0$ for every real mode by antisymmetry; the identity in (e) follows from $\mathbf M$'s hopping terms alone, the $S$ term cancelling. \emph{What rests on it:} the body's click of S.14, one transition at the line $\omega_j - \omega_i$ taking the share $T(\sin\omega_j - \sin\omega_i)$, stands under the rotation read as under the pace read, the difference line being the absorption line under the rotation and the sum a pair line; the transfer rate is the law's $\Omega_R$ above with the current's matrix element; the fraction absorbed, with the body's write back, stays open as Section 4.3 says. \emph{Status:} derived, one derivation, the sum and difference assignment of (c) and (d) checked by two; the pace read's one-mode line (S.16) has no analogue under the rotation; clicks.
\end{derivation}


\begin{derivation}[the two forces between quanta and their ratio between families, Section 5.1]
\emph{Inputs:} the sign holder's coupling between two quanta, $\alpha_{\mathrm{law}} = (3\sqrt 3 / 8\pi)\,k / (\Gamma E_s)$ (S.26); gravity's level from a source of $s_g$ wells per interval, $3s_g / (4\pi r\Gamma E_g)$ far away, Eq.~(11) of the paper (S.25), and the clock's rate $1 - U$ in it (S.2); the quantum of a family, its energy $T\sin\omega_s$, its share $\sin\omega_s$ in units of $T$ and its Wronskian $T / 2$. \emph{Steps:} (a) The two holders are one kind of family, the kernel $3 / (4\pi r)$ and the same acts of write and read. The sign holder is written by a quantum's Wronskian, $1 / (2E_s)$ per interval for every family, and the content holder by its form, $\sin\omega_s / E_g$ per interval, so a quantum of the family $s$ at $r$ writes the level $U(r) = 3\sin\omega_s / (4\pi r\Gamma E_g)$. (b) A reader quantum of the family $s'$ has its rotation redshifted by $U$, its energy by $\delta E = -T\sin\omega_{s'}\,U$, so the pull is $F_g = 3T\sin\omega_s\sin\omega_{s'} / (4\pi\Gamma E_g r^2)$; the electric force between two quanta by S.26 (c) is $F_e = 3kT\cos\omega / (8\pi\Gamma E_s r^2)$, $\cos\omega \to 1$ at a small gap. In the form of a coupling, with light's $c = 1 / \sqrt 3$ as in S.26 (b), gravity's between two quanta of one family is $\alpha_g(s) = (3\sqrt 3 / 4\pi)\sin^2\omega_s / (\Gamma E_g)$, the square of the quantum's energy over $T$. (c) The ratio at one family: $F_e / F_g = \alpha_{\mathrm{law}} / \alpha_g(s) = k\,E_g / (2E_s\sin^2\omega_s)$ to the first order in the gap ((b)'s $\cos\omega_s$ and the finite-gap redshift's $2\cos\omega_s / (1 + \cos\omega_s)$, S.18, both $1$ to $10^{-21}$ at nature's gaps at the MAGIC bound ($\omega_p^2 / 4 = 1.1 \times 10^{-13}$ at LHAASO's $\omega_p = 6.7 \times 10^{-7}$) and $2 / 3$ and $0.8$ at $[2, 3]$), the unit of charge ($T / 2$) over the quantum's energy ($T\sin\omega_s$), squared, times the universe's two weights $k$ and $E_g$ over $E_s$. (d) Between two families the weights cancel: $(F_e / F_g)_e / (F_e / F_g)_p = \sin^2\omega_p / \sin^2\omega_e = (m_p / m_e)^2$ with no number, the masses being the energies $T\sin\omega_s$ (S.27), an identity in the small-gap limit, forced by charge universality and $F_g \propto m^2$, no test. \emph{Check:} nature's $e^2 / (4\pi\epsilon_0 G m_e^2) = 4.17 \times 10^{42}$ and $e^2 / (4\pi\epsilon_0 G m_p^2) = 1.24 \times 10^{36}$ \cite{codata}, the ratio $3.36 \times 10^6$ with these three digits ($3.37 \times 10^6 = 1836.15^2$ with the paper's four-digit inputs); at the MAGIC bound $\omega_e < 2.0 \times 10^{-14}$ (S.24), $1 / \sin^2\omega_e > 2.5 \times 10^{27}$ and $k\,E_g / E_s = 2 \times 4.17 \times 10^{42}\sin^2\omega_e$, about $3 \times 10^{15}$ there (about $10^{24}$ at the slack LHAASO bound, $\omega_e < 3.7 \times 10^{-10}$, S.27), its largest value, and smaller at a smaller $\omega_e$ (what grows at a smaller gap is $1 / \sin^2\omega_e$, the ratio the weights must supply per unit); on the lattice the two holders carry two integers and one kernel, and the $10^{42}$ appears only when a NodeDetector compares the two forces on one quantum, its charge read as its invariant and its mass as its energy. \emph{Status:} derived from the definitions, an identity and no test (Section 5.1's mark); clicks.
\end{derivation}




\begin{derivation}[the atom with the nucleus's angle alone, the record's own level, and the holders the electron does not read, Section 4.2]
\emph{Inputs:} the sign holder's angle of a nucleus of charge $Z$ at $r$, $\theta_0 = Z\alpha_{\mathrm{law}}\,c / r$ by the far kernel (S.26), $c = 1 / \sqrt 3$ (S.1); a bound electron record of the matter family read into Schr\"odinger's slow limit (Section 2.4) with the band's inertia $m^* = 3\tan\omega_s$ per interval (S.22) and light's $c = 1 / \sqrt 3$; the one write of every record into the one holder (Section 2.2). \emph{Steps:} (a) \emph{Bohr's form with the nucleus's angle alone.} The slow limit is the hydrogen-like problem with the coupling $Z\alpha$, the inertia $m^*$ and the speed $c$: the levels $E_n = -m^*(Z\alpha)^2 c^2 / (2n^2) = -m^*(Z\alpha)^2 / (6n^2)$ per interval, Rydberg's $1 - 1 / n^2$ between them, and the radius $a_0 = 1 / (m^* Z\alpha c) = \sqrt 3 / (m^* Z\alpha)$ Links, which at nature's gap ($m^* = 3\omega_s$) is $\bar\lambda_e / (Z\alpha)$; the lines are the beats of two levels (S.41); the electron's rotation is redshifted by the binding, $(Z\alpha)^2 / 2$ of its energy at a small gap, the mass defect as the click layer reads a frequency. (b) \emph{The orbital moment.} The kinetic term of a plane record in a plaquette flux $\Phi$ (S.33's covector) gives the orbital splitting $\Phi\,m_\ell / (2m^*)$, one Bohr magneton per unit of angular momentum, Zeeman's known form derived here under S.33's covector with the orbital $g = 1$; the spin term's $g = 2$ belongs to the spin holder, named not reached in Section 1.1 of the paper. (c) \emph{The record's own level.} The holder of the sign is one row, and a record cannot tell its own write from the nucleus's, so the bound electron reads the nucleus's angle $Z\alpha / r$ and its own, of the opposite sense and the strength $1 / Z$ of the nuclear one, spread over its own radius: the Hartree self-term, repulsive and at full strength. With $x = a_0 / a$ for a hydrogen-like trial record the mode's frequency over $Z^2$ times Rydberg's unit is $\epsilon = x^2 - 2x + (1.25 / Z)\,x$, the kinetic, the nuclear and the self-term $J$; its minimum is not a stationary state of the law's iteration. The trial mode, a one-parameter variational estimate of Hartree's fixed point and not the fixed point itself, is the minimum of $E_{\mathrm{tot}} = x^2 - 2x + (0.625 / Z)\,x$ at fixed norm, $x = 1 - 0.3125 / Z$, and its frequency, which the click reads, is $\epsilon = E_{\mathrm{tot}} + (0.625 / Z)\,x$: $E_{\mathrm{tot}} / (Z^2 R) = -(1 - 0.3125 / Z)^2$ and $\epsilon / (Z^2 R) = -(1 - 0.3125 / Z)^2 + 0.625\,(1 - 0.3125 / Z) / Z$, with the radius $a / a_0 = 1 / (1 - 0.3125 / Z)$. At $Z = 1$: $E_{\mathrm{tot}} = -0.473$, $\epsilon = -0.043$, the radius $1.45$; at $Z = 2$: $-0.712$, $-0.448$, $1.19$; at $Z = 10$: $-0.938$, $-0.878$, $1.03$; at $Z = 30$: $-0.979$, $-0.959$, $1.01$. The Rydberg ratios return as $1 - 1.25 / Z$, $95.8$ percent at $Z = 30$. So a record reading its own sign field is reflected on the lattice in Bohr's form at large $Z$ only, and not hydrogen, excluded by line (1) of the model's decision: a record reads every sign field but its own, and hydrogen's lines are S.52's. (d) \emph{The holders the electron does not read.} A nuclear holder the electron reads is excluded by the spectrum: a one-Node proton by the finite-size part of hydrogen's $1s$ level, $1.1$ MHz for nature's radius (the finite-size shift's known form, derived here as the slow limit's own perturbation; the form factor and the muonic radii the comparison), which a point nucleus would lack and the measured Lamb shift resolves at the kilohertz level \cite{bezginov}, and by the form factor, $1$ at every momentum transfer where nature's falls; a $1$ fm cloud bound by a gapped holder at the binding's content $8 / 938$ by muonic hydrogen, $5.2$ meV (the measured Lamb shift's own coefficient, $5.2275\,r_p^2$ meV \cite{antognini} with $r_p$ in fm, at $r_p = 1$ fm) against the electronic and muonic radii's agreement \cite{antognini}; an ungapped holder by the spectrum. So the nuclear holders are read by the nucleon families and not by the electron, the selective read of Online Resource 1, Section B.10 (the numbers computed by hand from the law's lines, one derivation). \emph{Check:} $\alpha^2 / 2 = 2.66 \times 10^{-5}$; $(1 - 0.3125)^2 = 0.473$, $0.473 - 0.625 \times 0.6875 = 0.043$; at $Z = 30$ the self-term is $4.1$ percent of the binding. \emph{Status:} (a) and (b) derived in form, two derivations; (c) the finding and its numbers two derivations; (d) a comparison with nature's Lamb shift and muonic radius, the selective read a declaration. Under the two lines of S.52 the self-term of (c) is absent and (a) is the law's; clicks.
\end{derivation}


\begin{derivation}[the equivalence principle between families, Section 5.1]
\emph{Inputs:} the fall of a free packet of the pair $[\num, \den]$ in a pace gradient, $a = c^2\,[2\nu / (1 + \nu)]\,\nabla U$ with $\nu = \num / \den$ (S.18), the proton and the neutron two families of the gaps $\omega_p$ and $\omega_n$ under the families' reading. \emph{Steps:} $2\nu / (1 + \nu) = 1 - \tan^2(\omega_s / 2)$ exactly, so two families fall apart by $[\tan^2(\omega_n / 2) - \tan^2(\omega_p / 2)]$, $(\omega_n^2 - \omega_p^2) / 4$ at small gaps. Two bodies differ in their fall by that times the difference of their neutron fractions: beryllium ($5 / 9$) against $^{48}$Ti ($26 / 48$) differ by $0.0139$, and with $\omega_n^2 - \omega_p^2 = (1838.68^2 - 1836.15^2)\,\omega_e^2 = 9{,}297\,\omega_e^2$ the E\"otv\"os parameter is $\eta = 0.0139 \times 2{,}324\,\omega_e^2 = 32.3\,\omega_e^2$, below $4.3 \times 10^{-18}$ at the LHAASO bound $\omega_e < 3.65 \times 10^{-10}$ ($1.3 \times 10^{-26}$ at MAGIC's $\omega_e < 2.0 \times 10^{-14}$, S.24), under the beryllium against titanium bound of $10^{-13}$ from the torsion balance \cite{schlamminger} (MICROSCOPE's titanium against platinum, $10^{-15}$, Touboul et al.\ 2022 \cite{microscope}, for which the neutron fractions at natural abundances by mass numbers are $0.5409$ and $0.6003$, the difference $0.0594$ and $\eta = 0.0594 \times 2{,}324\,\omega_e^2 = 138\,\omega_e^2$, $135.6\,\omega_e^2$ for the single isotopes $^{48}$Ti and $^{195}$Pt at $0.5417$ and $0.6000$, the bound $1.8 \times 10^{-17}$ at $\omega_e < 3.65 \times 10^{-10}$; the conclusion holds at either); the electron against the proton falls apart by $\omega_p^2 / 4 = 1.1 \times 10^{-13}$ at the same bound, which no click has read (the electron's free fall known to $10$ percent \cite{witteborn}). The caveat: S.18's fall is a free packet's; a nucleon bound in the nuclear holder falls as the body's record does, and the body's fall per family is not derived, so the row is the free-packet estimate of the composition dependence, a free-packet estimate; its relation to the bound body's response is undetermined. \emph{Check:} $(1 - \nu) / (1 + \nu) = \tan^2(\omega / 2)$ with $\nu = \cos\omega$; $(1838.68^2 - 1836.15^2) = 9{,}297$; $9{,}297 / 4 \times 0.0139 = 32.3$. \emph{Status:} derived in form, two derivations, the number a free-packet estimate; its relation to the bound body's response is undetermined; clicks.
\end{derivation}


\begin{derivation}[the proofs of the theorems of Sections 3 and 4.2]
\emph{Inputs:} Rule3's line, Eq.~(1), and the share's identity, Eq.~(4); the momentum reading $P_a$ of S.4; S.1's band; the stable body's functional of Section 4.2; the momentum flux identity of the paper's Section 3.3 for the guide's lay. \emph{Theorem 1, the direction.} Forward, $\sigma = 1$, the line is the definition of $\ddiv$ and $\dmod$. Backward, $\sigma = -1$: write $Y$ for the right side's first two terms, $\sum_j R_{ij}\,\mathrm{arr}_j + S\,\anow$. Then $u = w\,\anext - Y + r'$ and $z = -\lfloor u / w \rfloor = \lceil (Y - w\,\anext - r') / w \rceil$. From the forward line, $Y - w\,\anext - r' = w\,\abefore - r$ with $0 \le r < w$, whose ceiling divided by $w$ is $\abefore$ exactly; and $\rho' = u - w\lfloor u / w \rfloor = w\,\abefore - (Y - w\,\anext - r') = r$. The theorem holds for the lattice's step at every interval; a NodeDetector's write and a body's window, click or not, are the NodeDetectors' acts outside the reversal, as the lay is; the exact inverse holds between clicks; across a click the write is a declaration.

\emph{Theorem 2, what the acts can form.} Each act is a linear form of the levels or the floor of one, and $|\lfloor L(x)\rfloor| \le |L(x)| + 1$; so sums and compositions of the four acts, at declared coefficients, grow at most linearly in the levels, while a product $\mathrm{now}_i\,\mathrm{before}_j$ grows as their square. So nothing bilinear comes from the acts alone. The remark of the paper's Section 3.2, that on a bounded domain every integer function is a finite sum of floors of affine forms, in two lines: on a finite set $D$ of integer points choose an integer vector $\mathbf c$ with $\langle\mathbf c, \mathbf x\rangle$ injective on $D$ (a base larger than $D$'s width per coordinate does it) and $N$ larger than its spread; for each $\mathbf a \in D$ the indicator is two floors, $\mathbf 1_{\mathbf x = \mathbf a} = \lfloor(\langle\mathbf c, \mathbf x - \mathbf a\rangle + N) / N\rfloor - \lfloor(\langle\mathbf c, \mathbf x - \mathbf a\rangle + N - 1) / N\rfloor$ (the difference is $1$ at $\langle\mathbf c, \mathbf x - \mathbf a\rangle = 0$ and $0$ at every other value in $(-N, N)$); then $f = \sum_{\mathbf a \in D} f(\mathbf a)\,\mathbf 1_{\mathbf x = \mathbf a}$, a finite sum of floors of affine forms with integer coefficients. A floor of an affine form of floors with a bounded term, $\lfloor a\,(\mathbf L\mathbf x + \mathbf b(\mathbf x)) + c\rfloor$ with $\mathbf b$ bounded, is $a\,\mathbf L\mathbf x$ plus a term bounded by $|a|\sup|\mathbf b| + |c| + 1$, affine plus bounded again, so every fixed finite composition of the acts is an affine form plus a bounded term.

\emph{Theorem 3, the conserved form.} Without the rounding the line is $\mathbf a_{\mathrm{next}} + \mathbf a_{\mathrm{before}} = \bM\,\mathbf a_{\mathrm{now}}$ with $\bM$ the matrix whose diagonal is $S_i / w$ and whose entries on the six Links are $R / w$. Let $\bD$ be the diagonal matrix of the weights $1 / p_i^2$; then $\bD\bM$ is symmetric iff $d_i M_{ij} = d_j M_{ji}$ on every Link, $d_i R_a(i \to j) = d_j R_a(j \to i)$: with $R_a(i \to j) = 2\,\num\,p_i^2 q_{ij}^2 / \Gamma^2$ the weight $1 / p_i^2$ cancels the Node's factor on both sides of every Link and the Link's factor $q_{ij}^2$ is the same from both ends, so $\bD\bM$ is symmetric for any content and any tensions. For a symmetric $\bK = \bD\bM$ the quantity $Q(\mathbf x, \mathbf y) = \mathbf x \cdot \bD\mathbf x + \mathbf y \cdot \bD\mathbf y - \mathbf x \cdot \bK\mathbf y$ satisfies $Q(\mathbf a_{\mathrm{next}}, \mathbf a_{\mathrm{now}}) = Q(\mathbf a_{\mathrm{now}}, \mathbf a_{\mathrm{before}})$: with $\mathbf a_{\mathrm{next}} = \bM\mathbf a_{\mathrm{now}} - \mathbf a_{\mathrm{before}}$, $\mathbf a_{\mathrm{next}} \cdot \bD\mathbf a_{\mathrm{next}} - \mathbf a_{\mathrm{next}} \cdot \bK\mathbf a_{\mathrm{now}} = \mathbf a_{\mathrm{before}} \cdot \bD\mathbf a_{\mathrm{before}} - \mathbf a_{\mathrm{before}} \cdot \bD\bM\mathbf a_{\mathrm{now}}$, which is the other side. Multiplied by $w$ this is Eq.~(4) exactly. \emph{The form where the paces move.} With $Q(\mathbf x, \mathbf y; \mathbf M) = \mathbf x \cdot \mathbf D\mathbf x + \mathbf y \cdot \mathbf D\mathbf y - \mathbf x \cdot \mathbf D\mathbf M\mathbf y$ and the weights $\mathbf D$ of one interval held, the step $\mathbf a_{\mathrm{next}} = \mathbf M(t)\,\mathbf a_{\mathrm{now}} - \mathbf a_{\mathrm{before}}$ gives $Q(\mathbf a_{\mathrm{next}}, \mathbf a_{\mathrm{now}}; \mathbf M(t)) = Q(\mathbf a_{\mathrm{now}}, \mathbf a_{\mathrm{before}}; \mathbf M(t))$ by Theorem 3's substitution; the form of the next interval is read with $\mathbf M(t + 1)$, so $E(t + 1) - E(t) = w\,[Q(\mathbf a_{\mathrm{next}}, \mathbf a_{\mathrm{now}}; \mathbf M(t + 1)) - Q(\mathbf a_{\mathrm{next}}, \mathbf a_{\mathrm{now}}; \mathbf M(t))] = \mathbf a_{\mathrm{next}}^{\top}\mathbf D\,(w\mathbf M(t) - w\mathbf M(t + 1))\,\mathbf a_{\mathrm{now}}$, exact in the rationals, $0$ at paces fixed in time (a theorem by substitution). Its sign under a deepening well is a lattice reading: S.36 (b)'s witness, the well deepened from $300$ to $600$ levels at $\Gamma = 6{,}000$ under a standing mode, lowers the share read at the vacuum's coefficients by $3.86$ percent (lattice).

\emph{The proposition, the stable body (Section 4.2).} With the dilation $\lambda$ about the body and $K$, $V_r$ the body's own energies at $\lambda = 1$, $F(\lambda) = K / \lambda^2 - \sum_r \sigma_r V_r \lambda^{-s_r}$; $F'(1) = 0$ is $-2K + \sum_r \sigma_r s_r V_r = 0$; $F''(1) = 6K - \sum_r \sigma_r s_r (s_r + 1) V_r = \sum_r \sigma_r (3s_r - s_r^2 - s_r)\,V_r$: Derrick's scaling, Pohozaev's identity on the lattice.

\emph{The momentum flux (the momentum density $P_a$, the momentum flux identity of the paper's Section 3.3).} Substitute the line into $w\,P_a(t + 1)$ at coefficients uniform over the lattice: the $S$ term cancels there (at non-uniform static paces the momentum density gains the Link-wise differences of the coefficients, Section 3.3), the $w$ term is $w\,P_a(t)$ and the $R$ term is the difference of the fluxes $G_{ab}$ as the momentum density's form states; the integer step subtracts the remainder term $\delta_i\,\Delta_a(\mathrm{now})_i - \mathrm{now}_i\,\Delta_a(\delta)_i$ with $\delta = r - r'$ (in this orientation of $P_a$).
\emph{The guide's lay (two derivations).} (1) A source $s\,e^{-i\Omega t}$ at one Node of a width-one guide, whose band is $\cos\Omega = (\num / 3\den)\cos k$ (the transverse mode at $\cos(\pi / 2) = 0$), drives two outgoing waves $C e^{ik|x|}$ with $C\,(2\cos\Omega - (2\num / 3\den)\,e^{ik}) = s$ at the source Node, so $|C| = 3\den\,s / (2\num\sin k)$. (2) The source adds at the interval $t$ the increment $A_t\cos(\Omega t)$, $A_t$ its amplitude; the two outgoing waves have $B = 3A / (2\sin k)$ on the width-one guide ($\num = \den$, the emission's case; $(3\den / 2\num)\,A / \sin k$ for a general pair), and the form of a travelling wave is $D = B^2\sin^2\Omega$ exactly at every Node, with no time mean; each interval's emission occupies $v_g = \sin k / (3\sin\Omega)$ Links on each side, so the two trains' form is $\sum D = 2\,(9 / (4\sin^2 k))(\sum_t A_t^2)\sin^2\Omega\,\sin k / (3\sin\Omega) = (3 / 2)(\sum_t A_t^2)\sin\Omega / \sin k$, the sum over the increments' squared amplitudes. (3) One quantum on a real line is $\sum D = T\sin\Omega$ (Section 2.5's packet, $\sum a^2 = T / \sin\Omega$), so one quantum is laid at $\sum_t A_t^2 = (2 / 3)\,T\sin k$, $A_t$ the amplitude of the increment added at the interval $t$; the increments' own squares, $\sum_t A_t^2\cos^2(\Omega t)$, sum to half of it, $(1 / 3)\,T\sin k$, for an envelope narrow in band against $\Omega$, to the cycle average (a short envelope departs from the half; the convention named; two derivations).

\emph{Status:} Theorems 1 and 3 proved for every lattice, Theorem 2 on the infinite lattice at fixed coefficients, the write across a click a declaration, an assumption; the stable body's proposition derived for the continuous functional, the integer step's version open; lattice.
\end{derivation}




\begin{derivation}[the click and the NodeDetector: the passage from the lattice to the clicks, Section 2.5]
\emph{The click, algebraically.} Let $X$ be the state space: at every Node, for every part of every family, one integer level pair and the write remainders, and nothing else. The step $S_t$ is Rule3 over the interval with the paces read from the holders at its start, a bijection of $X$, exact by the division with its remainder; the line without its division is linear in the levels at fixed paces, the linearity the backward reading uses, the integer map not affine; its forward orbit is the future and its inverse orbit the past. The readings $D$, $F$ and $W$ are bilinear forms on $X$, covariant under the symmetry group of Section 2.1 and never written. A NodeDetector $R$ is a declared region (two Nodes or more with Nodes alone, one or more with a record of its own) with its front Ports and its coefficients, which define $Q_R$: for one line a quadratic form of the levels, the inflow through its front Ports accumulated over the window; for a pair, the window's sum of the squared cross-side sum over the parts, quadratic in the joint amplitude, the product of the two sides' levels, and so quartic in the levels. A click is the pair $(R, t)$, the credit of one whole quantum to the NodeDetector $R$ at the interval $t$, drawn with the probability $s_R / \sum_{R'} s_{R'}$, where $s_R = Q_R(x_t) / W_{\mathrm{rec}}$, the record's own unit ($\Wc$ for a record laid by share with its count, $\Wc\sin\Omega$ for a born quantum), is the NodeDetector's form on the forward orbit over the quantum's measure of a current. In the meeting's words, the inverse orbit from the credited state coincides with the forward orbit on the realised orbit (in the state space, the record over every path at once), and the share is the bilinear form's value at the coincidence. The click's decision is the one map in the model that is not a bijection, from the vector of shares to one chosen outcome, the same under every element of the group; its write is one act on $X$, at the one Node of the region from which the quantum reached the NodeDetector, which the NodeDetector knows and no click reports. The exchange is this: the arriving record's levels and remainder at that Node to $0$, the hole, its count down by one, and the absorbing record's realised part laid there at its count up by one, the absorption, the phase the product's and the remainder at the lay's origin, the quantum conserved between them, the absorbing record's change spreading by Rule3 as an event and the absorbed record erased from that Node outward at the causal bound, inside the click's light cone, nothing else written (\claimmark{assumption}, the NodeDetector's declaration). Every observable is $Q_R / W_{\mathrm{rec}}$ on the forward orbit, in the record's own unit ($\Wc$ for a packet laid at count-unit, $\Wc\sin\Omega$ for a born quantum), or a ratio of two such. \emph{The one act at three storeys.} Rule3 on the levels divides the six weighted arrivals with the Node's own terms and the kept remainder by the wall $w$, the quotient the next level, the remainder kept at the Node, the inverse a bijection by it. The hold's write divides the sources' forms by $E_s T$, the quotient the held level's gain, the remainder kept at the Node, exact by it. The click divides the inflow through the front over the window, with the books' remainder, by $\Wc$, the quotient the clicks and the remainder kept in the books for the next window, the two records at the Node its neighbours, the inverse crossed from the click line by the remainder and the count. The count is the level of the top storey; the cancellation lives in the count because the count is the quotient; the exchange at the Node is the writing of the quotient into the two records; the reversibility between clicks is Rule3's; the three holders are the middle storey and the conversion the top storey's one click; the one thing in the click that is not Rule3 is the draw (\claimmark{derived} for the form, two derivations; \fence{clicks}). The hole's booking identity (Section 5.2, cited in Section 2.5.2): with $wQ(p, q) = \sum_i [w\,p_i^2 + w\,q_i^2 - p_i(S q_i + \sum_j R_{ij} q_j)]$, $w$ times the two-level form $Q$ of Eq.~(4), on a pair of consecutive levels, Rule3's line without the remainder, $w(\mathrm{next}_i + \mathrm{before}_i) = S\,\mathrm{now}_i + \sum_j R_{ij}\mathrm{now}_j$, gives $wQ(\mathrm{next}, \mathrm{now}) = wQ(\mathrm{now}, \mathrm{before})$ at paces uniform over the lattice, where $R_{ij} = R_{ji}$, and then $wQ(\mathrm{now}, \mathrm{before}) = w\sum_i D_i$ (at non-uniform static paces Theorem 3's weighted form takes $Q$'s place); a face writing $x$ in place of the level $v$ the step would write at one Node, the leaving level $b$ standing, changes $wQ$ by $(x - v)[w(x + v) - (S\,\mathrm{now}_i + \sum_j R_{ij}\mathrm{now}_j)] = w\,(x - v)(x - b)$, the line giving $S\,\mathrm{now}_i + \sum_j R_{ij}\mathrm{now}_j = w(v + b)$; exact in the line at any paces as a one-Node identity with the Node's own coefficients, in the plain unit of $wQ$ (the Node's weight $1 / p_i^2$ relating it to the weighted share), and in integers within $|x - v|$ in the unit $Q$, that is $w\,|x - v|$ in the plain unit of $wQ$, the carried remainders' term $(x - v)(r - r')$, and the most one write removes is $w(v - b)^2 / 4$ at the midpoint (three independent derivations). The credit rounds the top storey half up per window; the carried remainder is the law's form with its own blind on $N$. The window is declared, the universe's draw window set explicitly to the whole passage of a packet, or the credit floors at $0$ (Section 3.4); a Node whose pace rounds to $0$, a frozen Node, stands in no NodeDetector's front and is never clicked (\claimmark{assumption} for the definition, \claimmark{derived} for the precisions; \fence{clicks}).

\emph{One form for every click.} At the credit the parts of a record are read with the NodeDetector's coefficients, each part paired with the same part through the common origin, summed and squared; the share is that density, and the draw is by the shares. The two slits are the case with nothing to pair: light is one real line, and the share is the record's own inflow. Bell's check is the case through the common origin: a pair born as one event is two trees with one root, and the backward reading from the first click passes the root and continues forward on the partner's tree (Section 7.1). The joint share is computed from the forward records alone, by the linearity of the rule, the inverse orbit existing by Theorem 1 (\claimmark{derived}, \fence{clicks}).

\emph{The NodeDetector is the one declaration kind}

\emph{Assumption (the NodeDetector's declaration).} A NodeDetector with Nodes alone is one connected region of two Nodes or more declared in the universe's declaration, never one Node, at least half the wavelength of the family it reports across the beam ($4$ Links at $\lambda = 8$ Links); a narrower region, a single Node or a region in pieces is excluded by the declaration. A region that absorbs is opaque, backed by a receding face or by the wall's own face, so that what enters it leaves the declared lattice and never returns; a bare region, backed by no face, reports the inflow through its front and absorbs nothing, its net over the window nothing (Online Resource 1, Section B.7). Its click is its report, per interval and family: the net current into the region through its front boundary Ports, the Ports leading in from the declared lattice, summed in integers, with the region's name and the family and never a Node; one quantum's measure $W_c\sin\omega_0$ for a gapped family and $W_c$ itself, $3\,\den\,T$ for a massless family, of the inflow over a passage is one unit of count, one quantum's share, $T$ for a bolometer and $T\sin\omega$ for a counter (C.9 (a), (b)), so a screen saw $N$ units where its regions' inflows sum to $N\,\Wc$. The credit: one quantum to one NodeDetector, drawn by the shares of the regions' inflows, read outside the lattice over every declared NodeDetector of the universe, one credit per quantum and never two; the NodeDetector draws with the world's declared draw (a body's NodeDetector with its own seed), the lattice draws nothing and the draw is deterministic per seed. The draw is a declared linear congruence of the full period over $2^b$ values \cite{hull_dobell}, so over the period each outcome's frequency is within $2^{-b}$ of its weight's share and not exact: the pick is $(x \times \mathrm{total}) \ddiv 2^b$ for the draw's state $x$ and the outcomes' total, and over the full period each residue appears once, so each outcome's frequency is within $2^{-b}$ of its share, $2^{-63} = 1.1 \times 10^{-19}$ at the declared $b = 63$. After a realised click the NodeDetector writes once, at the one Node from which the quantum reached it. A NodeDetector whose declared universe gives it a record of its own (parts, transitions, a lifetime) is the same kind and emits as well: its region is one Node or more, resolving the place and not the wave, held by its declaration and never bound by a well, and its one write is at the Node the draw picks by the share, the absorption's hole at the Node drawn by the arriving record's inflow over the region and the emission's quantum at the Node drawn by its own record's share. The credit is the one non-local place of the description and no line of the law.


Why a region: a one-Node record is every wave number at once, and the region's half-wavelength is the NodeDetector's declaration; so no click names a position finer than a region, and every per-Node number is a diagnostic; the emission's one-Node write is every wave number at once as well (the transition beat labels the emitting transition; the emitted excitation's frequency distribution and its frequency-selective absorption are not derived from this write; assumption for the write; lattice; S.1, S.41).

What produces a click is deterministic: the meeting, the current of the amplitudes through a region's boundary. What is not known is which of the possibilities becomes the click: when one possibility is spread over several NodeDetectors, the lattice gives the shares, no line of the law gives which NodeDetector fires, and the NodeDetector draws. Exclusivity, one click per quantum, is derived under the counting rule's one draw, the NodeDetector's declaration, from the count conserved; Born's rule, the credit's chance equal to the intensity's share, is a declaration of the NodeDetector and not a finding of the model. The uncertainty relation is not a separate rule: it is the NodeDetector's declarations, the region (its Nodes), the universe's draw window and the draw (its seed), the write exact inside at a Node and an interval that are the lattice's alone and never reported, the report uncertain outside (\claimmark{assumption}). What each declaration costs: a region of $n$ Nodes reads a fringe of period $\Lambda$ with the factor $\sin(\pi n / \Lambda) / (n\sin(\pi / \Lambda))$; a frequency is read to $\Delta\omega \ge 2\pi / n_w$ over $n_w$ intervals; and $N$ quanta drawn by the shares scatter as $\sqrt{NP_i(1 - P_i)}$, so a pattern is in the histogram and never in one click (\claimmark{derived}, \fence{clicks}; S.8).

\emph{The click decides in the credit and writes at one Node}

\emph{Assumption (the click's decision and its write).} The lattice is deterministic and local: Rule3 with no draw, Theorem 1's bijection at every interval. The click is the meeting of two quanta at the NodeDetector, the incoming record's quantum and the NodeDetector's own transition; the NodeDetector, declared in Universe24's declaration, draws with the world's declared draw (a body's NodeDetector with its own seed) whether the meeting is realised, and the randomness is the NodeDetector's, the lattice has none. Forward, to the future, one write: the exchange at the one Node from which the quantum reached the NodeDetector, which the NodeDetector knows and no click reports (the uncertainty principle: for us the click is a region and a window), the Node of the realised write drawn after the outcome by the same draw's next pick, one Node draw per click and none for a candidate not realised, for an absorption by the inflow booked through each Node over the window (the region's share their sum) and for an emission by the record's share over its Nodes, the Node a lattice diagnostic beside the click line; the exchange: the arriving record's levels and remainder at that Node to $0$, the hole, its count down by one; the absorbing record's realised part laid at that Node at its count up by one, the absorption, with the phase the product's and the remainder at the lay's origin, the part it leaves laid at its count down by one; the quantum conserved between them, and the absorbing record's change spreads by Rule3 as an event, one Link per interval; no act faster than one Link per interval: the front erases the absorbed record from that Node outward at the causal bound, one shell of Link distance per interval, the hole's act on each shell, the clicked record alone; the other paths are a cancelled future at three levels: for every click there will be, the record's count $0$; at every NodeDetector they reach where the declared universe declares a face behind it, their share leaving the lattice there with no credit; on the lattice the front erasing them inside the click's light cone, beyond it their levels remaining, still sourcing the holders they write to, until it arrives; a record's front starts when its count reaches $0$. Backward, to the past, a reading and never a write: the credit's reading through the root, which chooses which tree is the past and gives Bell at a distance, the one non-local place; it cannot be a write, the lattice being deterministic forward.

A single quantum's count stands at $0$ after its click, the count conserved and read by the credit through the root, the NodeDetectors' shared other side on the lattice, the Nodes' levels carrying the shares and never the count: antibunching and one quantum one click per record between far NodeDetectors, the counting rule's one draw over all the NodeDetectors that read one record (\claimmark{assumption} for the one draw, \claimmark{derived} from it for the rest; \fence{clicks}; the pair's $478 / 169$ and the GHZ's $-4$ with one click per record). Nothing at one NodeDetector changes the amplitudes at another except by Rule3 from the Node of the exchange; no signal passes through the amplitudes (Rule3 is local, Theorem 1) nor through a bound state's clicks (C.8 (b)); a Nodes-only region's count is coupled to the others' inflows by the credit's rounding, the declared non-locality of C.1 (3), by at most one quantum per window, $1 / N$ relative. A NodeDetector is a declared region of the lattice and stands nowhere outside it: it has two sides, the front where the incoming quantum meets its own transition and the other side where its own record connects to the rest of the lattice, the apparatus's bodies and, through the record it absorbed, that record's root, the root's event, so that two NodeDetectors reading one record are connected through the root, their common past; it can be entangled on both sides. A bound body clicks by quanta from outside meeting its own transition, never by clicks from outside, and only where it is declared a NodeDetector: the future is light's forward record at the body's boundary, the past the body raised by one quantum read backward to the source, and the meeting is the resonance of the body's rotation with the light's (\claimmark{derived}, \fence{clicks}; Section 4.3); a NodeDetector with a record of its own writes once at its window's close and only its own record: on a click, the exchange at the one Node the draw picks by the share, the arriving record's levels there to $0$ and its own part raised; on no click, the null window's re-lay of its own record at its Nodes in the complement of its outcome set with its count whole, no Node drawn and nothing of the arriving record written. The arriving record is written only at a realised click, once. A spread body's click is the exchange at its entry Node and its settling a reading, named open, and the emission lays the light quantum at the Node drawn by the NodeDetector's own record's share, at the resonance $\omega_e - \omega_g$, the share's difference booked and stated explicitly (S.14, S.41); a body in the dark is only the future and still gravitates; emission is the same meeting with the roles reversed, two local writes and one draw of the credit binding them. The Node stays generic, knowing no family and no click; the write is the NodeDetector's act (\claimmark{assumption}, the NodeDetector's declaration, the law's line of the click's write; the decision's form \claimmark{derived}). The lattice holds only the future, every possibility, and has no arrow: Rule3 is a bijection at every interval (Theorem 1). The clicks have the arrow: each is forward only, the one act outside the reversal; at each click what was future becomes past, the present being the meeting; the chain of clicks is the past advancing into the future, written once and never undone. The direction of time is the clicks' and not the lattice's (\claimmark{theorem} for the step; the clicks' arrow the NodeDetector's declaration). The two slits' run is with NodeDetectors of Nodes alone, which report their clicks and raise no configuration of their own (Section 2.5), the write at the window's end leaving the numbers unchanged.

\emph{The passage from the lattice to the clicks.} Every formula of Section 5.1 is derived in three steps: the law's line in the lattice's words; each word traded for a NodeDetector's report of it ($\omega$ for a clock's rate, $U$ for the ratio of two clocks' rates); and what remains, a relation among clicks alone, set beside nature's. The credit's bookkeeping is that the shares sum to what the screen saw, one whole per click, the whole the NodeDetector's own transition; a frequency of nature is read as the inflow's beat within the window or as the resonance of the NodeDetector's own record, and the lattice's quantities, the form, the content, the level, the rotation, are read by no click except as the share's inflow. \emph{Status:} the click's algebraic definition and the NodeDetector's declarations the law's, each declaration an assumption; the credit a declaration; lattice for the record, clicks for every number of nature.
\end{derivation}


\begin{derivation}[why there is entanglement, and its scope, Sections 2.5 and 7]
\emph{Bell: the declared universe.} The entanglement the model computes is one event with two trees. The pair is laid as the pair family, two real lines of light never summed at a Node, with equal parts. Each side's NodeDetector, a region at least half a wavelength across, reports the two parts' signed level sums over its region and knows nothing of angles; the ports enter at the credit. The NodeDetector applies the declared setting, a pair of integers $(p, q)$, with the port coefficients $e(+) = (p, q)$ and $e(-) = (-q, p)$, exactly orthogonal, so each side's marginal is one half for any setting. The settings of the Clauser-Horne-Shimony-Holt test are $(1, 0)$, $(1, 1)$, $(12, 5)$ and $(5, 12)$, at $0$, $45$, $22.62$ and $67.38$ degrees. The credit is the one form of the click: the joint share pairs the parts by their label across the two sides through the root,
\begin{equation}
J(o_A, o_B) = \sum_{\text{window}}\Bigl|\sum_k c_k\,e_k(o_A)\,\alpha_k(o_A)\,e_k(o_B)\,\beta_k(o_B)\Bigr|^2, \qquad c = (1, 1), \tag{S.47a}
\end{equation}
with $o_A$, $o_B$ the two sides' ports ($+$ or $-$), $\alpha_k$ and $\beta_k$ the part $k$'s amplitude at the port on each side and $e_k(o)$ the port's integer coefficient on the part $k$, the sum over the window's intervals standing outside the square, on both members of the level pair. The joint share is never the square of a local sum, which one NodeDetector could form alone and which gives $S$ at most $2$. One draw is made per pair, the NodeDetectors' draw with the declared seed, the lattice drawing nothing: the first side from its marginal and the partner from the conditional; each side's write at the Node of its exchange spreads by Rule3 and changes no number of the row, one click per record (S.46) (\claimmark{assumption}, \fence{clicks}).

\emph{Bell: what the model says it must give}

(1) Rule3 gives each part's amplitudes at the ports: the two parts of one beam, laid equal, are two records of one family stepped by the same line along the same path, identical exactly, integer for integer (\claimmark{theorem}, \fence{lattice}). (2) The meeting with the linearity of Rule3 gives the joint amplitude: the realised click at A's port weights the parts by A's analyser, is read back to the root where the two parts were one event and forward to B, where B's analyser weights them, and the root sums the two branches, $\sum_k c_k\,e_k(a)\,e_k(b) = \cos(a - b)$ for the normalised vectors $e(a)$ and $e(b)$ with $c$ the parts' laid weights, $(1, 1)$ for the symmetric lay (\claimmark{derived}, \fence{clicks}). (3) That the share is the density of the meeting, its square, is Born's rule for a pair: nothing in Rule3 forces it; everything after it, the fractions, $S$, GHZ's $M$, is derived in the algebra (\claimmark{assumption}, \claimmark{fitted}; \fence{clicks}). From the three, with the pair's common phase $\theta$ uniform over the pairs and A's $+$ port at $a$ and B's at $b$,
\begin{equation}
\begin{aligned}
&\text{the meeting:} && P(++) = \tfrac{1}{2}\cos^2(a - b), \quad E(a, b) = \cos 2(a - b), \quad S = 2\sqrt 2;\\
&\text{a local credit by the shares:} && P(++) = \tfrac{1}{4} + \tfrac{1}{8}\cos 2(a - b), \quad E = \tfrac{1}{2}\cos 2(a - b), \quad S = \sqrt 2;\\
&\text{a local credit by the sign:} && E = 1 - 4\,|a - b| / \pi \ (|a - b| \le \pi / 2), \quad S = 2;\\
&\text{the credit by the contrast:} && E = \cos 2(a - b) \text{ at the efficiency } 2 / \pi, \quad S = 2\sqrt 2,
\end{aligned}\tag{S.47b}
\end{equation}
The second, third and fourth lines follow by averaging over the pair's phase $\theta$ (\claimmark{derived}, \fence{clicks}; S.10). So $E(a, b) = \cos 2(a - b)$ under the meeting. For the declared pairs it is exact, Lagrange's identity of the NodeDetector on equal parts: $E = ((p_A p_B + q_A q_B)^2 - (p_A q_B - q_A p_B)^2) / ((p_A^2 + q_A^2)(p_B^2 + q_B^2))$, so at the four settings $E = 119 / 169$, $-119 / 169$, $120 / 169$ and $120 / 169$ and $S = 478 / 169 = 2.8284$ exactly, an efficiency of $1$ by construction and the marginal one half on each side whatever the partner's setting (\claimmark{computed}, \fence{clicks}; S.10). For unequal parts, $r$ their products' ratio at the ports, the same credit gives $E = \cos 2a\cos 2b + \rho\sin 2a\sin 2b$ with $\rho = 2r / (1 + r^2)$ and $S = \sqrt 2\,(1 + \rho)$ at the ideal angles (\claimmark{derived}; \fence{clicks}). The maximum of $S$ over the four settings for $E = \cos 2a\cos 2b + \rho\sin 2a\sin 2b$ is $2\sqrt{1 + \rho^2}$, Tsirelson's bound $2\sqrt 2$ at $\rho = 1$ and never the algebraic $4$, and the declared pairs' $478 / 169 = 2.82840$ sits below $2\sqrt 2 = 2.82843$ by the rational angles (\claimmark{derived}, \fence{clicks}).

What the lattice gives without the declaration is the second and third lines of Eq.~(13) of the paper, two results on local credits on the same declared universe (\claimmark{derived} under the uniform phase, \fence{clicks}): a credit of each side by the shares of its ports gives $S = \sqrt 2$, and a credit by the sign of its larger port gives the sawtooth, whose sum is exactly $2$, Bell's bound. Bell's value therefore sits in the credit rule, a declaration, and is no independent evidence that the rule accounts for entanglement. Nature beside the derivation: the two loophole-free photon tests of 2015 for the pair's row, and $S = 2.42 \pm 0.20$ in the test of that year on two spins joined at a click \cite{hensen}, the swap kind, which the law names as met by the same form and which is not derived here. The same credit on four lines gives the three-quantum GHZ correlation $E_3 = \cos 2(a + b + c)$, every pairwise $E$ equal to $0$ and Mermin's $M = -4$ where every local credit gives at most $2$ (\claimmark{derived} under the declared credit, \fence{clicks}; S.38). No signal passes between the sides, each side's marginal being one half; no settings are chosen by the lattice's own past, so a superdeterministic world is not this one.

\emph{The scope.} Entanglement in the law is a property of the lay: parts laid alike as one event (the pair family, the GHZ family with four parts for three NodeDetectors, S.38) share one root, and the credit's joint share pairs their parts through it; Rule3 preserves the lay and, by the real line's linearity within its bounded remainder (the paper's Theorem 2), creates none. What the law meets is therefore the pair laid alike as one event and two multi-part records meeting at one click, where the click is the root through which their parts are paired, the swap of entanglement; no other entanglement is named as met. \emph{Status:} the joint share a declaration; the correlations under it derived, two derivations; the swap's joint share \claimmark{hypothesis}, with the criterion that separates a transferred quantum from a coupling, the credit's bilinear form after the exchange; clicks.
\end{derivation}






\begin{derivation}[the properties of the lattice and the calibration series, in full, Section 8.2]
\emph{What the lattice fixes with no number.} Before any declared universe the lattice fixes: the geometry, a cube of Nodes with an integer state, six Ports and the $48$ symmetries with the reflections, so parity is kept and there is one preferred frame; no $V$-proportional odd read exists for a one-part real line (S.56 (d), a theorem); a plane configuration reads gravity's odd lines as the Port's phase, $g_{0j} = -4V_j$, a plane configuration of one rotation sense reads $\alpha_1 = \alpha_2 = 0$ (S.56), bodies $\alpha_1 = 8f - 8$ at the share $f$, $-4$ at one half (the misses of 8.1); $\alpha_1 = -8$ is the write without the odd Link part (S.35) (\claimmark{derived}; S.56 (d); S.35, C.4); the dynamics, one rule per family, explicit, local and bijective with the remainder kept, so reversibility exactly, integer for integer, no noise but the remainders' and no loops, hence no running of $\alpha$ and no Lamb shift (\claimmark{derived}; a finding against nature, S.50 (6)); the conservation, per family and no total across families: the weighted form and so the energy per family at static paces, the gauge-turned record's Wronskian (S.41 (a)) and so the charge, in the rational line (the integer step adding S.9's remainder term), at paces fixed in time, exactly for any charge-potential level in time under the time-Link turn, no spontaneous click and so no lifetime (\claimmark{derived}); the band, light at $1 / \sqrt 3$ Link per interval under the causal bound $1$, the anisotropy $k^2 / 12$ at short wavelength, a gapped holder Yukawa's kernel with the massless constant (\claimmark{derived}); the paces, the exponential metric with $\gamma = \beta = 1$ in $U$ exactly, $\beta_K = 1$ as the gap closes, the finest potential $1 / \Gamma$ and the frozen content $U_f = \frac12\ln(2\Gamma)$, $(\Gamma / 2)\ln(2\Gamma)$ levels (\claimmark{derived} under the one assumption); the holders, each one's field Poisson's with the kernel $3 / (4\pi r)$, its strength a declared weight and its range its gap (\claimmark{derived}); the click layer, the bilinear credit and so the angle it gives, one whole quantum per click and the draw by the shares (\claimmark{assumption}); and the scale, one scale $\Gamma$ and pure numbers of order one, every large number a power of $\Gamma$ or a declaration. In one sentence: the lattice is a machine of forms; it fixes every form and every ratio of order one, conserves, reverses and makes no noise; it fixes no strength, no scale and no large number; those are the universe's.

\emph{The calibration series.} Every measured number of an experiment is first written as a click pattern on a declared NodeDetector against a declared clock. The families' pairs and the units $\Gamma$ and $T$ would be fixed by that match and by nothing read off the lattice, Universe24's integers being the rule's own; the two ratios the model fixes without it are light's speed of $1 / \sqrt 3$ Link per interval and the causal bound of $1$ (\claimmark{derived} for the speed, \fence{lattice}; the bound is postulate 4 of C.1). The passage from the lattice to the clicks is S.46's. No coefficient of the declared universe is derived from the lattice; the lattice dictates the form in which each enters and reads none of them. Each is read from one measurement of nature, as the paper's Section 8.2 and S.50 set out: every coefficient of the declared universe is read by one click on a bound body; the lattice fixes the form of the reading; the declared universe holds the integer. Each stands in one of two places: $T$ and the NodeDetector's unit of credit are the click's, the gaps, $k$, $E_g$ and $\Gamma$ the lattice's, read through bound bodies and given by no NodeDetector. The law has no number from the lattice itself, no fixed point of the whole that selects $\Gamma$ and no ratio of the weights; if such a number exists, it is in the cosmology. One derived check beside them, the two clocks: a cavity of bodies and an atom's beat shift alike by $N$ at long wavelength, as nature's clocks do, an atom against an atom to $2.1 \times 10^{-5}$ \cite{bauch} and a cavity against an atom to $1.7 \times 10^{-2}$ \cite{turneaure}, a cavity of declared faces at $N^2$ (\claimmark{derived}; \fence{clicks}; S.31).

\emph{Status:} each item as marked; lattice for the forms, clicks where a number of nature is named, lattice for the helium and positronium comparisons of (5) (S.52).
\end{derivation}


\begin{derivation}[the conservation picture and the credit's bookkeeping, Section 2.4]
\emph{What is conserved, and where.} On the lattice the conservations are per family: the weighted form of a family while its reads are static (Theorem 3), the gauge-turned record's Wronskian (S.41 (a)) of a plane family, its total with the weights $1 / p_i^2$, exactly in the line for any level in time under the time-Link turn (Section 3.5, S.41), the quanta through the Ports, the wave number of a free record and reversibility exactly, integer for integer. Energy is conserved within a family by theorem, between a charged body and its light by the law's line (Section 5.1), and between a body and the holders of the content it is not conserved, the loss being the mass defect (Section 8.1). The lattice conserves no total: the sum of the families' forms is not conserved and not assumed, the share moving with the paces with no current to compensate (Section 2.4 of the paper). The click layer has books of its own: a click credits one whole quantum at the NodeDetector's own whole, drawn by the inflow of the record's share through its front over its window; a frequency of nature, the lattice's quantities and the credit's bookkeeping are read as S.46 states (Sections 2.5 and 4.3), the bookkeeping up to the absorptions from undepleted records, read in the books and read as the deficit, the lattice's share of such a record standing within its declared tolerance while its count falls. In the clicks there are conservations other than the lattice's, and neither layer keeps a total energy (\claimmark{theorem} for the lattice's, \claimmark{assumption} for the credit's).

\emph{The table.}
\begin{center}\small\begin{tabular}{p{1.5cm}p{2.8cm}p{3.3cm}p{3.7cm}}
\hline
Layer & Conserved & Scope & By what\\
\hline
lattice & the weighted form $E$ of a family & per family, paces fixed, in the rational line (the integer step within the remainder term of the share's change, Section 3.4 (S.5)) & Theorem 3\\
lattice & the Wronskian $W$ of a plane family, its total weighted by $1 / p_i^2$ & per family, in the rational line (the integer step adding S.9's remainder term), at paces fixed in time, exactly for any charge-potential level in time under the time-Link turn & S.9, S.33, S.41\\
lattice & the wave number $k$ of a free record & per family, paces uniform over the lattice & the band's translation invariance\\
lattice & the sum of $D$ over families & not conserved & Section 2.4 of the paper: the share moves with the paces, no current compensates\\
clicks & the number of quanta, the charge per quantum, the whole per click & per whole quantum, across families & the credit's integer unit; the energy line for a charged body and its light\\
content holders & the body's share & not conserved: the loss is the mass defect & S.36\\
\hline
\end{tabular}\end{center}
On the lattice nothing is assumed: the conservations are theorems of Rule3's line without its division, at paces fixed in time. What is not conserved and not assumed is the sum of all families' energy; the holders are open to their readers and the law has no theorem equating one family's loss with another's gain. The energy line closes it for one sector: for a charged body and its light the energy its clicks read and the holder's form move together under the turn and the write, both up for a body in its own level and both down for a pair of opposite senses, and agree to first order in the turn if and only if $k_r E_s T\cos\omega_s = k_w\Gamma$: one energy read twice, by the clicks and by the form, a condition and not an identity, at the vacuum's paces of S.26 (S.26 (d$'$)). For the content holders it is not closed and not to be: a body's count falls in a well, its potential being in its count, while the holder's waves gain a different share of the loss, the mass defect, the energy leaving the books, the law's own statement and nature's reading of bound bodies. \emph{Status:} the lattice's conservations theorems; the credit's bookkeeping an assumption (the counting rule), an assumption for the clicks' total; the energy line the law's line, two derivations.
\end{derivation}


\begin{derivation}[the explicit findings, with their numbers, Section 8.1]
The measurements of nature named below carry their citations here or in the paper's Sections 5.1 and 8.1. The paper's thirteen misses map onto these seventeen so (frame dragging for bodies, the thirteenth, is S.56's and stands outside these seventeen): (1) the $\times$ polarisation is one of the paper's nine, and the frame's (1), the write without the odd Link part's, which the paper's Sections 1.2 and 9 name, stands in its own derivation and is not counted among the thirteen. (2) one of the nine. (3) one of the three. (4) one of the nine. (5) the fine structure and the two many-body numbers are among the nine, the exclusion among the three, its $g - 2$ with (6). (6) one of the three. (7) one of the nine. (8) a control excluded by line (1), not counted; the reduced-mass and isotope-shift failure of (5) takes its place among the nine. (9) stands in its own derivation and section and is not counted. (10) one of the nine. (11) not counted. (12) not counted. (13) not counted. (14) the gate-joined ions and transmons of B.8, a finding against nature not counted among the thirteen. (15) not counted. (16) not counted. (17) one of the nine. Each finding is a place where the law as written and a measurement of nature differ, named with its number and its root; none is smoothed. (1) The lattice's frame: a metric diagonal in the preferred frame, $\alpha_1 = -8$ against $|\alpha_1| < 4 \times 10^{-5}$ and, in the standard gauge, $\alpha_2 = -1$ against $|\alpha_2| < 4 \times 10^{-7}$ from the Sun's spin and $1.6 \times 10^{-9}$ from solitary pulsars, conditional on the static read (Eq.~(1)'s second line writes the time part and the three axis parts and never an off-diagonal $T_{ab}$, so the carried perturbation has $h_{00}$ and $h_{aa}$ only, the breathing $h_{xx} + h_{yy}$ and the lattice-aligned $+$, $h_{xx} - h_{yy}$, and no $\times$, $h_{xy}$), and the Nordtvedt parameter $\eta_N = 4\beta - \gamma - 3 - \alpha_1 + 2\alpha_2 / 3 = 7.3$ ($8.0$ without $\alpha_2$'s term) against lunar laser ranging's $\eta_N = (4.4 \pm 4.5) \times 10^{-4}$ \cite{will,williams}, the $\alpha_2$ part conditional on the static read; no $\times$ polarisation; no frame dragging under the write without the odd Link part ($0$ against Gravity Probe B's $37.2 \pm 7.2$ milliarcseconds per year), $19.6$ for quartz under the law's write (S.56 (g)). (2) The nuclear binding by any derived body: the binding fractions under the declared holders miss in two directions, the deuteron above and iron below, by factors of order ten and two under the rms convention of the quoted radii ($12$ and $2$; $4$ and $6$ under a Bohr-like radius), the profile and the screening entering at the same order (C.6 (b)) (C.6); the size law of one well, $R \propto 1 / A$, against $A^{1/3}$ (C.6); the bound $0.76\,CW / (\Gamma E)$ per quantum at one Node is conditional on the declared $W$, which no click bounds, and is no finding (S.36; \fence{lattice} for the well, \fence{clicks} for the defect it bounds). (3) The neutron's $878.4 \pm 0.5$ s \cite{pdg}: no spontaneous click, no rate; the conversion a declared list of the one click act at a declared rate, the table and the rate nature's numbers and nothing of the neutron's mass claimed (Section 8.1). (4) The $Z$ and parity: no fourth rotation of a three-line record, the step commuting with the reflections. (5) The anomaly $g - 2$ and Pauli's two per orbital: no loops, no vector giving the exclusion; the law does not support the reduced mass, as S.52 (c) states once (no amplitude over two positions on the lattice; the pair's joint amplitude at the click, in the credit's pairing, S.47), so the isotope shifts and the muonic atoms are outside what it derives: with the nucleus pinned (line (2)) the H-D $1S$-$2S$ shift is $0$ against nature's measured isotope shift, and muonic hydrogen's levels lack the factor $0.899$ of the reduced mass; the fine structure: the band's $(Z\alpha)^4$ term is the spin-$0$ form $n / (l + 1 / 2)$, splitting hydrogen's $2p$ from $2s$ by $29$ GHz ($\alpha^4 m_e c^2 = 1.4490 \times 10^{-3}$ eV, its twelfth $29.2$ GHz the Klein-Gordon $2s$ to $2p$ split, its thirty-second $10.95$ GHz Dirac's $2p$ split), where nature's hydrogen has Dirac's $n / (j + 1 / 2)$ ($2p$'s splitting $10.97$ GHz measured, Dirac's $10.95$, $2s_{1/2}$ and $2p_{1/2}$ $1.06$ GHz apart \cite{bezginov}), and the model has no spin holder, so the fine structure, the hyperfine line and the Zeeman anomaly have no line; the many-body correlation is not reproduced by a product of records (the ground levels, computed by hand from the modes: helium $1.4$ percent short, positronium $43$ percent of nature's) (S.52). (6) $\alpha$'s running and the Lamb shift: no loops. (7) The mass ratios of nature not whole numbers, so, with the cluster's defect bounded at $C\omega_0 x_{\mathrm{self}}$ (S.36 (d)), $10^{-22}$ at the write weight $1$, no particle a cluster of the electron's quanta at that weight; the masses declared. (8) The record's own level: hydrogen's $1s$ at $-0.043$ of Rydberg's unit at $1.45\,a_0$, the Rydberg scaling returning at large $Z$ as $1 - 1.25 / Z$ ($E_{\mathrm{tot}}$ as $(1 - 0.3125 / Z)^2$) (S.43), excluded by line (1) of the model's decision (S.52). (9) The radiated share's coefficient: the time level carries half of Larmor's power in form, the odd lines' part depends on their wall and coupling, which of the two couplings the odd lines carry open; no lifetime number (S.26, C.6). (10) At the second order in $z$ the ball of dust alone decelerates, $q_0 = \Omega_m / 2 > 0$, where the supernovae read $q_0 \approx -0.5$ \cite{riess}\cite{perlmutter} (the flat $\Lambda$CDM model, cold dark matter with a cosmological constant); the paper's claim is at $z \ll 1$ only, and no line of the law closes it. (11) The content holders' balance is not a line of the law: a body's loss in a well is the mass defect (S.36, S.49). (12) The electromagnetic binding energy's seat, open: the holders of the content gravitate with their binding, the sign holder's cross term sources nothing as written (Section 5.1, S.36). (13) Bjorken scaling: under the dimension-$3$ family a hard click reads three fractions summing to one; their dependence on the momentum transfer is uncomputed, the lines' Fourier weights at $q$ depending on the lines' lay and the read operator not derived (C.5 (c)), where nature's measurements show logarithmic violations of order ten percent across the measured decades of $q^2$; the momentum sum rule, nature's quarks carrying about half the nucleon's momentum and the gluons the rest, is not reached, the model having no gluon field; what the law does fix is that no pair is created and no coupling runs, the root of (6). (14) Bell violations between separately prepared bodies joined by a coupling or by a transferred quantum (gate-joined ions, $S = 2.25 \pm 0.03$ \cite{rowe}; transmons with a transferred photon, $S = 2.0747 \pm 0.0033$ \cite{storz}, loophole-free) exceed the local bound the law's credits give for such a join: on the law as written nothing after an exchanged quantum is bilinear in the two bodies' amplitudes, and every later credit is local, $S \le 2$; the largest discrepancy the paper names between the law as written and measurement, two derivations. The pair laid as one event and the swap of two multi-part records at a click the law meets (S.47). (15) The shelved ion's telegraph: bimodal counts and exponential dark periods, which the law without the NodeDetectors' writes does not give (the ion's world is not computed in this work); read against their blinds by the body's own emission click at a declared rate (computed, the rate declared), the shelving at the Zeno-limited rate $\Omega_w^2 / R$, the bright periods' mean $R / \Omega_w^2$ and the dark $\tau_D$, the rate not derived. (16) The sign holder's radiation is a scalar's: a spherically symmetric breathing charge writes light through its second moment (S.40 (c)), where nature's vector field writes none (hydrogen's $2s$, metastable, decaying by two photons); the root the holder's time level written as a scalar, the odd lines' part open (S.26 (d$'$)). (17) The weak coupling: with every rotation holder at $\alpha$ and nature's $m_W$, $G_F = \pi\alpha / (\sqrt 2\,m_W^2) = 2.5 \times 10^{-6}\,\mathrm{GeV}^{-2}$ against $1.166 \times 10^{-5}$, short by $4.6$, the mixing angle's $1 / \sin^2\theta_W$ absent (S.37 (c)). \emph{Status:} explicit findings; the frame's parameters of a write with no odd Link part ($g_{0j} \equiv 0$): $\alpha_1 = -8$, $\alpha_2 = -1$; the law's write (Section 2.3, S.56) carries the odd part and gives S.56 (g)'s $0$ for a plane configuration of one sense, a body $\alpha_1 = 8f - 8$ at its share $f$; $\alpha_1 = -8$ is named in the paper's Sections 1.2 and 9 as the write without the odd Link part's; clicks where a number of nature is named.
\end{derivation}


\begin{derivation}[the push on a moving record, Section 5.1, Eq.~(9)]
\emph{Inputs:} the composed paces, $p_0 / \Gamma = e^{-U}$ and $p_a / \Gamma = e^{-2U}$ (Eq.~(2), the one assumption $hN = 1$); the band at a pace (Eq.~(3)) at a small gap and long wavelength, $E^2 = m^2 + P^2$ with $E = \omega$, $m = N\omega_0$ and $P = cN^2 k$, the exponential metric's dispersion ($g_{00} = -e^{-2U}$, $g_{aa} = e^{2U}$, $\gamma = 1$; $g_{0a} = 0$ for the static source, the odd part the holder's odd lines carry, sourced by the count's flux, S.56, which the push of a static source does not read); the force on a record as the gradient of its rotation at fixed wave number, $\mathbf F = -(\partial\omega / \partial U)_k\,\nabla U$ (the fall of S.18 is its $k \to 0$ case). \emph{Steps:} (a) $\partial(E^2) / \partial U = \partial(m^2 + P^2) / \partial U = -2m^2 - 4P^2$, so $-(\partial\omega / \partial U)_k = (m^2 + 2P^2) / E = (E^2 + P^2) / E$; (b) with $\beta = P / E = v / (cN^2)$, $v = \partial\omega / \partial k = c^2 N^4 k / \omega$ the group velocity and $cN^2$ the local light speed, $\mathbf F = E\,(1 + \beta^2)\,\nabla U$; (c) at $\beta = 0$ the pull $E\,\nabla U$, over the band's inertia $m^* = 3\tan\omega_0$ the fall $a = f(\omega_0)\,\omega_0 / m^*\,\nabla U = c^2 (U_K / U)\,\nabla U$, with $U_K$ Kepler's potential of S.18 and $f(\omega_0) = 2(1 - \cos\omega_0) / (\omega_0\sin\omega_0)$ (S.15), $0.2667\,\nabla U$ at $[2, 3]$ (S.18, the same number); at $\beta = 1$ twice the pull, the bending $4U_b$ (S.20); (d) the units: the paper's momentum is the wave number, so the lattice form is $(E^2 + c^2\,\mathbf p\cdot\mathbf p) / E$ and the covariant one $(E^2 + P^2) / E$; the form $E^2 + 3\gamma\,\boldsymbol\Pi\cdot\boldsymbol\Pi$ holds for the covariant variable $\boldsymbol\Pi = E\mathbf v$ alone (the paper's $\mathbf p = E\mathbf v / c^2$ is the wave number, and the boxed momentum at the composed paces is $P = cN^2 k = \Pi / (cN^2)$, one letter per quantity), where $3\,\boldsymbol\Pi\cdot\boldsymbol\Pi = E^2\beta^2$ with $c^2 = 1 / 3$ at $N = 1$. \emph{Check:} on the exact band at $[2, 3]$ the rest push over $E$ is $f(\omega_0) = 1.0634$, the redshift's factor ($\alpha_U = f(\omega_0) - 1 = 0.063$), and the coefficient of $\beta^2$ in the push over its rest value $f(\omega_0)\,E$ is $1.3353 = \gamma_K - \alpha_K / 2$ (two derivations), and with the factor $\omega_0 / \tan\omega_0$ the paper's coefficient, $(\gamma_K - \alpha_K / 2)\,\omega_0 / \tan\omega_0 = 1.0045$; at $[999, 1000]$ the three are $1.00017$, $1.00067$ and $1.00000003$ ($1 + \omega_0^2 / 12$, $1 + \omega_0^2 / 3$ and $1 + \omega_0^4 / 120$), all $\to 1$; the angle's ratio to Newton's at $\beta = 0.38$ is $1 + \beta^2 = 1.144$, the force's $\gamma_L(1 + \beta^2) = 1.237$. \emph{Status:} derived from Eqs.~(2) and (3) under the one assumption (assumption, fitted), two derivations; as physics the standard post-Newtonian factor $1 + \gamma\beta^2$, where $\gamma = 1$ follows from $hN = 1$, which the bending fixes, and the finite-gap coefficient is the law's own; the chain to Einstein's weak-field forms is the one dispersion read four ways (the bending at $\beta = 1$, Newton at $\beta = 0$, the redshift and Shapiro's delay by its time part and index, the perihelion by the orbit step of S.19 with its $k^4$ caveat at the computed pair); clicks, the deflection angle's ratio $1 + \beta^2$ to Newton's.
\end{derivation}


\begin{derivation}[Sections 4.2 and 8.1: the atom under two lines: no record reads its own write of the sign; the nucleon a compact body]
\emph{Inputs:} the hydrogen-like problem of S.43 (a) with the nucleus's angle alone; the self-read of S.43 (c); two lines, the model's decision: (1) every record that reads the holder of the sign is one quantum (count $1$), writes its own field of the sign (its Wronskian over $E_s T$, Rule3 at the paces of its reads) and reads every sign field but its own, the light born by its write living in its field; (2) the holders of the content keep the self-read, the nucleon families read the nuclear holders' fields including their own, so a nucleon is a compact self-bound body at the declared weight, its centre one classical variable in an outside well. \emph{Steps:} (a) \emph{The chain, by the law's lines and nothing else.} The nucleus, one quantum of its family, writes its Wronskian $T / 2$ into the sign holder's field every interval (the one write, Section 2.2; S.9), and the holder's rest far from it is the kernel's, the angle $\theta_0(r) = 3kZ / (8\pi r\Gamma E_s) = Z\alpha_{\mathrm{law}}\,c / r$ with $c = 1 / \sqrt 3$ and $\alpha_{\mathrm{law}} = (3\sqrt 3 / 8\pi)\,k / (\Gamma E_s)$ (S.25, S.26 (a), (b)). Under line (1) the electron's record reads that field alone, so its Rule3 line is the plane record's step under the time turn $\theta_0(r)$ (S.33's inputs; the sign holder rotates the two-part record), $e^{i\theta_t} z_{\mathrm{next}} + e^{-i\theta_{t-1}} z_{\mathrm{before}} = [S z_{\mathrm{now}} + \sum_a R_a(z_{+a} + z_{-a})] / w$ (the convention of S.41 (b)), and nothing else acts on it. Its stationary modes $z = \phi\,e^{-i\omega t}$ obey, by the band at small wave number (S.1, S.22: $\omega = \omega_s + k^2 / (2m^*)$ with $m^* = 3\tan\omega_s$) and the turn's shift of the rotation (S.26 (c), S.41 (b): $\omega \to \omega + \theta$ for the sense $+$, the electron's $\theta_0(r)$ of the sign that lowers its rotation in the nucleus's level, the attraction's sign the declared sign; $2\cos(\omega + \theta_0)$ equal to the band's $2\cos\omega(k)$), the law's own eigenvalue line on the lattice, $-\Delta\phi / (2m^*) - \theta_0(r)\,\phi = (\omega - \omega_s)\,\phi$ with $\Delta$ the six-Port second difference; its radial form is solved on $2{,}000$ radial points with no formula put in: the bound modes sit at $\omega_n - \omega_s = -m^*(Z\alpha_{\mathrm{law}})^2 c^2 / (2n^2) = -m^*(Z\alpha_{\mathrm{law}})^2 / (6n^2)$ per interval, $Z^2$ and $1 / n^2$ exact in the small-gap limit, the radius $a_0 = \sqrt 3 / (m^* Z\alpha_{\mathrm{law}})$ Links; the levels are a lattice reading, the diagnostic behind the click, and the lines nature measures are the beats of two modes read by a click (S.14, S.41), the table below. Into nature's units: the interval of S.27 ($\omega_e\hbar / (m_e c^2)$) and the law's $\alpha$ read from Coulomb's force (S.26 (b)) turn the mode's depth into $m_e c^2\alpha^2 / 2$ at $n = 1$, $-1 / 2$ hartree exactly in the slow limit ($-0.49995$ on the mesh) and $-1 / 8$ at $n = 2$. What the law computes is the coefficient $1 / 2$ and the dependence on $n$, $Z$ and the family's gap $m^*$; the two inputs, $\alpha$ and the electron's unit, are the universe's calibrations from clicks other than the spectrum (Coulomb's force, the electron's mass), so the chain is not circular: no formula of the atom enters, the absolute hartree is a comparison through the unit and the coefficients and the ratios are the derived content. (b) Without (2) the nucleus is a free record of inertia $m_p^* = 1{,}836\,m_e^*$ reading the electron's field, a well $27$ eV deep and $a_0$ wide: its stationary spread is $0.21\,a_0$ on $2{,}000$ radial points, the nucleus's kinetic share its own reading (the $1s$ cloud's curvature $k = 4 / 3$ in atomic units (a.u.), $\omega_{\mathrm{zp}} = \sqrt{k / m_p} = 0.027$ a.u.), and the self-consistent product of the two records sits at $-0.4671$ Ha (hartree), $6.58$ percent above Bohr, the $2s$ at $-0.1147$, the $1s$--$2s$ at $9.59$ eV against $10.20$; the spread scales as $m^{-1/4}$, so deuterium's $1s$--$2s$ differs from hydrogen's by $172$ meV against nature's $2.8$ meV, the measured H--D 1S--2S isotope shift \cite{parthey2010}. (c) Under (2) the nucleus does not spread, (a) holds, and the law does not support the reduced mass: it holds no amplitude over two positions, no wave of two records on pairs of Nodes (every record's levels stand at its own Nodes, jointly with the other records there and coupled to them by the read, and the pair's joint amplitude enters at the click, in the credit's pairing, S.47; postulates 1 and 2's locality), so a product of two bound records about one centre has no relative motion of the pair, and under (2) the nucleon's centre is one classical variable, so the isotope shifts and the muonic atoms are outside what the law derives; helium in the mean field without exchange $-2.8617$ against $-2.9037$ Ha \cite{pekeris}; positronium, two free records of one gap about one centre, $-0.217\,R$ ($-0.1085$ Ha, the minimum of Pekar's functional, Pekar's constant $0.108513$ \cite{lieb}) against nature's $-0.50\,R$ ($-0.25$ Ha, $6.80$ eV): a product about a centre has none of the pair's relative motion, which carries the whole binding of two equal records; a named finding. \emph{Check:} the pinned nucleus reproduces $-1 / (2n^2)$ to $10^{-4}$ on the $2{,}000$ points; the free-nucleus numbers by the self-consistent two-record Hartree iteration ($60$ iterations, mixing $0.5$); the zero-point estimate $\tfrac32\omega_{\mathrm{zp}} = 0.040$ Ha against the self-consistent $0.0136$ plus the smearing's $0.019$. \emph{Status:} derived under the two lines, the model's decision (line (1) \claimmark{assumption}, \claimmark{fitted}; line (2) \claimmark{assumption}); two derivations, with two costs named: under (1) the fields at a Node number the universe's charged records, so the bound of postulate 2 of C.1 is the universe's declaration; under (2) the nucleon's classical centre is the pinned nucleus by declaration, which gives the levels, the reduced mass unsupported by construction (no amplitude over two positions), the isotope shifts with it. \emph{What rests on it:} S.43's finding is replaced in Section 8.1 of the paper by the statement that the law does not support the reduced mass and that the many-body correlation is not reproduced by a product of records, Pauli's exclusion having no line (S.50 (5)); \fence{lattice} for helium and positronium (a comparison and no click claimed); \fence{clicks} for the one-electron lines.


\noindent\emph{The atoms' table: the clicks the law derives, with how they were derived; lattice for helium and positronium (a comparison and no click claimed), clicks for the one-electron lines.} Nature's measurements give lines and thresholds, the differences of levels, and each row is the beat of two modes read by a click (S.14, S.41), computed by hand from the modes of steps (a) to (c); a level alone is a lattice reading and stands in the text as the diagnostic behind the click. Every row is derived from the law's lines alone, no equation outside the law entering but the self-consistent iteration's own loop, which solves the law's two lines together (Bohr's, Schr\"odinger's and Hartree's forms in no column but nature's); nature's column is the measured value with its source, the comparison, entering no derivation. The inputs from nature, named once: the electron's unit of S.27 ($m_e$ and $\alpha$, read from the electron's gap and from Coulomb's force, clicks other than the spectrum; the hartree $m_e c^2\alpha^2$), so an absolute eV is a comparison through the unit and a coefficient or a ratio is the derived content, and the proton's inertia $1{,}836$ as the family's declared pair. The rows are a supporting finding and no more, derived under the two lines, the model's decision, two independent derivations (the chains and the circularity check, row by row), \fence{clicks}. Not in it: the slow limit of the band only, so the $(Z\alpha)^2$ order, the Lamb shift and the nuclear size are not in it; the reduced mass unsupported by construction (step (c)); Pauli's exclusion without a line (two records in one mode need none).
\begin{center}\small\begin{tabular}{p{2.0cm}p{2.4cm}p{1.8cm}p{5.1cm}}
\hline
The click & Nature measured, with its source & The law's click under the two lines & How derived, by the law's lines\\
\hline
Hydrogen's $1s$--$2s$ line & $10.20$ eV, $2{,}466{,}061{,}413{,}187$ kHz \cite{parthey} & $10.20$ eV & the beat of the $n = 1$ and $n = 2$ modes of step (a) read by a click, the share $T(\sin\omega_j - \sin\omega_i)$ (S.14, S.41): $3 / 4$ of the $1s$ depth derived, the eV through the unit\\
The ionisation threshold & $13.598$ eV \cite{nist} & $13.6$ eV & the $1s$ mode's depth of step (a) read as a click's share (S.14): the coefficient $1 / 2$ derived, the eV through the unit\\
He$^+$'s $1s$--$2s$ line & $40.81$ eV (the He\,II spectrum \cite{nist}) & $40.82$ eV ($Z^2 \cdot 3 / 8$ hartree at $Z = 2$) & step (a) at $Z = 2$, the angle $Z\alpha_{\mathrm{law}}c / r$ scaling the two modes by $Z^2$, the beat $3 / 4$ of the $1s$: the $Z^2$ ladder derived as a click, the $(Z\alpha)^2$ order not in the slow limit\\
Helium's ionisation threshold & $24.587$ eV \cite{nist} & $23.45$ eV, $4.6$ percent & two electron records in one mode about the pinned nucleus, each reading the nucleus's field and the other's and never its own (line (1)), iterated by the law's write and read to the fixed point, $-2.8617$ Ha (step (c)), against He$^+$'s one record at $-2$ Ha exactly: the threshold their difference, a comparison and no click claimed (lattice); nature's total energy two thresholds summed and no click; the correlation the product misses sits wholly in this threshold\\
Positronium's $1s$--$2s$ line & $5.10$ eV, $1{,}233{,}607$ GHz \cite{fee} & $1.9$ eV & two free records of one gap about one centre, each reading the other's field, iterated to the fixed point (step (c)), the beat of the product's two modes ($-0.1085$ and $-0.039$ Ha, the diagnostic behind the click); a product about a centre has none of the pair's relative motion, the ground level at $43$ percent of nature's, the line at $37$\\
Control: hydrogen's $1s$--$2s$ with the nucleus a free record & --- & $9.59$ eV & without line (2) the nucleus's record of count $1$ reads the electron's field and spreads by $0.21\,a_0$ at the inertia $1{,}836$, the two records iterated to the fixed point (step (b)): the control that separates line (2) from no line (2)\\
\hline
\end{tabular}\end{center}
\end{derivation}





\section*{The slow limit (Section 2.2)}

\begin{derivation}[Schr\"odinger's equation as the slow limit of the line, Section 2.2]
\emph{Inputs:} Rule3's line at the vacuum's paces without the remainder, $\anext + \abefore = (\num / 3\,\den)\sum_{\text{Ports}} arr$ (S.1); $\cos\omega_0 = \num / \den$; $\Delta$ the six-Port second difference, $\sum_{\text{Ports}} a = 6a + \Delta a$; the rest rotation at a content, $\omega_0(U) = \omega_0 - \omega_0 f(\omega_0)\,U + O(U^2)$ (S.15). \emph{Steps:} write $a_t = \mathrm{Re}[\psi_t e^{-i\omega_0 t}]$ with $\psi$ slowly varying, $\psi_{t \pm 1} = \psi \pm \psi' + \psi'' / 2 + \dots$; then $\anext + \abefore = \mathrm{Re}[e^{-i\omega_0 t}(2\cos\omega_0\,\psi - 2i\sin\omega_0\,\psi' + \cos\omega_0\,\psi'')]$, and the right side is $\mathrm{Re}[e^{-i\omega_0 t}(2\cos\omega_0\,\psi + (\num / 3\,\den)\,\Delta\psi)]$, since $(\num / 3\,\den)\,6 = 2\cos\omega_0$. The $2\cos\omega_0\,\psi$ cancels, the band at $k = 0$; dropping $\psi''$,
\[
i\psi' = -\frac{\num}{6\,\den\sin\omega_0}\,\Delta\psi = -\frac{\Delta\psi}{2m^*}, \qquad m^* = \frac{3\,\den\sin\omega_0}{\num} = 3\tan\omega_0,
\]
Schr\"odinger's free equation on the lattice, $\hbar = 1$ per interval, the inertia the band's (S.18, S.22: the band's own expansion $\omega - \omega_0 = k^2 / (2m^*)$ is the same statement). At a content $U$ the same steps with $\omega_0(U)$ add $V\psi$ with $V = -\omega_0 f(\omega_0)\,U$, the well lowering the rotation. The three conditions, each with its order: (i) the remainder dropped, below one level per Node per interval (the budget clause), the hard bound per laid mode $(n - 1) / \cos(\omega / 2) + 1 / \sin\omega + 1$ (the carried remainders a difference sequence, summed by parts; about $n$ at long wavelength, $1.73\,n + 0.33$ at the plane's corner for the massless pair, $1.60\,n + 0.42$ at $[2, 3]$; the band's two ends excluded by the division act's two sums and by the even periodic box excluded by the declaration) and the walk $\sigma\sqrt n$, per band mode $\sqrt{n / (12(1 + \cos\omega_k))}$, $2.4$ levels at the carrier and $3.0$ per Node over the lattice's zone at $n = 130$, the bound per laid mode $135.6$ at the carrier; (ii) long wavelength, the spatial second difference's $k^4$ term dropped, of relative size $k^2 / 12$ along an axis (the band's own $k^4$ coefficient relative to its $k^2$ is $1 / (12\sin^2\omega_0)$, $0.15$ at $[2, 3]$, the rest of it the envelope's term of (iii)); (iii) slow variation, $\psi''$ dropped against $2\sin\omega_0\,\psi'$, of relative size $(\omega - \omega_0) / (2\tan\omega_0)$, below $(\omega - \omega_0) / \omega_0 = k^2 / (2m^*\omega_0)$, the Klein-Gordon correction. Within the three the statement is exact algebra of the line; in nature's names it is the slow limit of the discrete Klein-Gordon equation, and the model's own result is $m^* = 3\tan\omega_0$, the potential's coefficient $f(\omega_0) = 2(1 - \cos\omega_0) / (\omega_0\sin\omega_0)$ and the lattice's departures named. \emph{Check} (the line evaluated numerically): a Gaussian packet of amplitude width $40$ Links at $k_0 = 0.08$ on a chain of $4{,}000$ Nodes with the two other axes folded, laid with the exact band's before level and stepped $3{,}000$ intervals by the line at $[2, 3]$ ($m^* = 3.3541$): the envelope's centre at $71.4$ Links against Schr\"odinger's $k_0 t / m^* = 71.6$ at $t = 3{,}000$ intervals (the exact band's group velocity times $t$, $71.4$); the density's rms width $32.38$ against Schr\"odinger's $a\sqrt{1 + (t / (2m^* a^2))^2} = 32.40$ with $a = 40 / \sqrt 2$. \emph{Status:} derived, the slow limit with its three conditions named, one derivation (S.18 and S.22 the same statement), the Check the line's numerical evaluation; lattice.
\end{derivation}
\begin{derivation}[the lattice's scale from light's anisotropy and nature's photon bounds, Section 8.1]
\emph{Inputs:} the massless band at the vacuum, $\cos\omega(\mathbf q) = (\cos q_x + \cos q_y + \cos q_z) / 3$ (S.1); nature's quadratic Lorentz-violation form $v / c = 1 - \tfrac32 (E / E_{\mathrm{QG},2})^2$ with the Fermi Large Area Telescope's (Fermi-LAT) bound $E_{\mathrm{QG},2} > 1.3 \times 10^{11}$ GeV for the subluminal case (the gamma-ray burst GRB 090510; Vasileiou et al.\ 2013 \cite{vasileiou}, after Abdo et al.\ 2009 \cite{fermi_lat}); the electron's rest energy as $\hbar\omega_0 / \tau_{\mathrm{int}}$ at the pair's rotation. \emph{Steps:} (1) with $\mathbf n$ the unit direction and $S_4 = \sum_a n_a^4$, $\cos\omega = 1 - k^2 / 6 + (k^4 / 72)\,S_4 + O(k^6)$ and $\cos\omega = 1 - \omega^2 / 2 + \omega^4 / 24$, so $\omega^2 = k^2 / 3 - (k^4 / 108)(3S_4 - 1) + O(k^6)$: the coefficient of $k^4$ is $-1 / 54$ along an axis, $-1 / 216$ along a face diagonal and $0$ along a body diagonal, where $\cos\omega = \cos(k / \sqrt 3)$ exactly; the group velocity $v_g / c = 1 - \tfrac32\,\epsilon_4 k^2$ with $\epsilon_4 = (3S_4 - 1) / 36$, subluminal everywhere and $0$ on the body diagonals. (2) With $E = \hbar\omega / \tau_{\mathrm{int}}$ and $k = E\ell / (\hbar c)$, $E_{\mathrm{QG},2} = \hbar c / (\ell\sqrt\epsilon_4)$, $4.243\,\hbar c / \ell$ along an axis; Fermi's bound \cite{vasileiou} gives $\ell < 6.4 \times 10^{-27}$ m (an axis; $1.0 \times 10^{-26}$ m at the angular mean) and $\tau_{\mathrm{int}} = \ell / (\sqrt 3\,c) < 1.2 \times 10^{-35}$ s. (3) The electron's pair fixes $\tau_{\mathrm{int}} = \hbar\omega_0 / (m_e c^2)$ by substitution and nothing of the law fixes the pair: at Fermi's bound $\omega_0 \le 9.6 \times 10^{-15}$, the gap $1 - \num / \den \le 4.6 \times 10^{-29}$, $\den \ge 2.2 \times 10^{28} = 2^{94.1}$ (Cassini's bound alone, at $2\sigma$, gives $\omega_0 < 0.0116$; LHAASO's highest photon $\den \ge 1.5 \times 10^{19}$); the proton's gap is $1836^2$ times the electron's, so one denominator serves both only at $\den \ge 2^{94}$. \emph{Check:} the quartic coefficient by the series and by the exact arccosine at $k = 0.01$ in the three directions, $-0.018519$, $-0.004630$, $0$, and $d\omega / dk$ at $k = 0.01$ along an axis, $v_g\sqrt 3 - 1 = -8.333 \times 10^{-6}$ both ways; two methods. \emph{Status:} derived, a constraint on any declared universe meant for nature, the law fixing no number; the anisotropy of light's time of flight, $0$ on the body diagonals to $1 / 54$ on the axes, a prediction of the law's own readable only at a Link near the bound; lattice.
\end{derivation}


\section*{The click's algebra beyond the law (Section 2.5.2)}

\begin{derivation}[the click's algebra beyond the law, Section 2.5.2]
\emph{Inputs:} Rule3 at the vacuum's paces and its band $\cos\omega = (\num / 3\den)\sum_a \cos k_a$ (S.1); the readings at a Node, the Wronskian $W = \mathrm{re}_{\mathrm{now}}\,\mathrm{im}_{\mathrm{before}} - \mathrm{im}_{\mathrm{now}}\,\mathrm{re}_{\mathrm{before}}$ of a plane configuration and the Link current $F_{ij} = \num\,(\mathrm{now}_i\,\mathrm{before}_j - \mathrm{before}_i\,\mathrm{now}_j)$ of a real line (Eq. (1), Section 2.2); the invariant $2A^2\sin\omega = T$ per quantum and the energy $T\sin\omega$ (Sections 4.3 and 5.1, the factor $2$ between a real line's $I = A^2\sin\omega / T$ and the invariant the convention of (b)); light's speed $c = 1 / \sqrt 3$; and the proposition of Section 2.5.2, a line beyond the law's write under its own name, that the write carries the arriving wave vector $\mathbf k$ as a phase twist of the taker's record, the twist $\delta k_a = \mathrm{sgn}(k_a)\,(|k_a|\,\mathrm{div}\,M)$ per Link along each axis for a record of count $M$ and the remainder $r_a = \mathrm{sgn}(k_a)\,(|k_a|\,\mathrm{mod}\,M)$ carried at the body, so that $k_a = M\,\delta k_a + r_a$ with $|r_a| < M$ and $r_a$ of $k_a$'s sign, the rule odd under $k_a \to -k_a$ exactly (the law's $\mathrm{div}$, the floor division of Section 2.2 whose remainder is non-negative, gives $(-3)\,\mathrm{div}\,7 = -1$ with the remainder $4$ against $3\,\mathrm{div}\,7 = 0$ with the remainder $3$; the twist, a line beyond the law, takes the magnitude's division with the sign carried so that it commutes with the cube's 48, reflections included).
\emph{Steps:} (a) \emph{Lemma A.} For a plane configuration $a_i = A\,e^{i(\mathbf k\cdot\mathbf x_i - \omega t)}$, with $\mathrm{now} = A e^{i\phi}$ and $\mathrm{before} = A e^{i(\phi + \omega)}$ one interval earlier, $W = A^2[\cos\phi\sin(\phi + \omega) - \sin\phi\cos(\phi + \omega)] = A^2\sin\omega$ at every Node, free of $\mathbf k$; for a real line along the axis $a$ with $\phi_j = \phi_i + k_a$, $\mathrm{now}_i\mathrm{before}_j - \mathrm{before}_i\mathrm{now}_j = A^2[\cos\phi\cos(\phi + k_a + \omega) - \cos(\phi + \omega)\cos(\phi + k_a)] = \tfrac12 A^2[\cos(k_a + \omega) - \cos(\omega - k_a)] = -A^2\sin\omega\sin k_a$, so $F_{ij} = -\num\,A^2\sin\omega\sin k_a$, free of $\phi$. A phase common to the two levels at a Node leaves $W$; a phase step across a Link moves $F$. (b) \emph{The norm factor.} The one-Node recurrence $x_{t+1} + x_{t-1} = 2\cos\omega\,x_t$ conserves $Q = x_t^2 + x_{t-1}^2 - 2\cos\omega\,x_t x_{t-1}$: for a real line $B\cos(\phi - \omega t)$, $Q = B^2\sin^2\omega$; for a plane configuration of amplitude $A$, $Q = 2A^2\sin^2\omega$, so $Q / \sin\omega = 2A^2\sin\omega = T$ and $W = A^2\sin\omega = T / 2$ per plane quantum. The identity $2A\cos\phi = Ae^{i\phi} + Ae^{-i\phi}$ is therefore exact in the count, $4A^2\sin^2\omega$ on both sides: a real line of amplitude $2A$ is two plane quanta of opposite sense, neutral, $W = +A^2\sin\omega - A^2\sin\omega = 0$; one real quantum has $B_q^2 = 2A_q^2$, and a real line of amplitude $B$ carries $B^2 / (2A_q^2)$ quanta. (c) \emph{The twist.} Under the proposition the twist $e^{i\,\delta\mathbf k\cdot\mathbf x}$, $\delta k_a = \mathrm{sgn}(k_a)\,(|k_a|\,\mathrm{div}\,M)$, is a phase common to the two levels at each Node, so by (a) it leaves $W$ at every Node, the charge kept, and moves the Link current; it keeps the count whole only with the write's own scaling of the amplitude, Section 2.5's $A^2 = (N + 1)\,T / (2\sin\Omega)$, $\Omega$ the transition rotation of Section 2.5, taken at the twisted rotation, $A^2 = (N + 1)\,T / (2\sin\omega(k / M))$, the invariant $2A^2\sin\omega = T$ per quantum holding at every $k$ and the rotation's difference $M[\omega(k / M) - \omega_0]$ being the recoil of (e), the beat's balance, booked by the write, in the click's unit $T\sin\omega$ it is $\cos\omega_0$ times that to leading order, the contrast of (f); the integer wave vector, in the wave-vector unit $2\pi / L_{\mathrm{box}}$ per Link, $L_{\mathrm{box}}$ the declared box's side in Links, is conserved across the click modulo $2\pi$ per Link with the remainder carried: for a taker whose books hold $r_a$ from earlier clicks the division is of $r_a + k_a$, $\delta k_a = \mathrm{sgn}(r_a + k_a)\,(|r_a + k_a|\,\mathrm{div}\,M)$ and $r'_a = r_a + k_a - M\,\delta k_a$, so $r_a + k_a = M\,\delta k_a + r'_a$ at every click, the form stated above is the first click's, with $r_a = 0$, the remainder in the taker's books (i), while the momentum reading $F \propto \sin k_a$ is conserved only to the order $k^3$ and folds at the zone's edge. The wave vector the twist carries is read at the taker's Ports and nowhere else: by (a) the arriving light's Link current along the axis $a$ is $F_{ij} = -\num\,B^2\sin\omega\sin k_a$ with $B^2\sin\omega = T$ per real quantum (b), so $\sin k_a$ per axis is in the NodeDetector's own reads over its boundary, the integer $k_a$ its nearest whole in the box's units to the order $k^3$ of the reading, and no number beyond the reads enters. (d) \emph{The band's inertia and kinetic scale.} With $\cos\omega_0 = \num / \den = \nu$ and $\sum_a(1 - \cos k_a) = k^2 / 2 + O(k^4)$, $\cos\omega - \cos\omega_0 = -\nu k^2 / 6$, so $\omega - \omega_0 = \nu k^2 / (6\sin\omega_0) = k^2 / (2m^*)$ with $m^* = 3\sin\omega_0 / \nu = 3\tan\omega_0$, and $\omega^2 = \omega_0^2 + c_m^2 k^2 + O(k^4)$ with $c_m^2 = \omega_0 / (3\tan\omega_0) = c^2\,\omega_0 / \tan\omega_0$ (S.18, S.21; $c_m / c = 0.867$ at $[2, 3]$). (e) \emph{The recoil.} A body of count $M$ taking the twist $k / M$ per Link gains the recoil, in rotation, $M[\omega(k / M) - \omega_0] = k^2 / (2Mm^*) + O(M^{-3})$, its energy in the click's unit $T\cos\omega_0\,k^2 / (2Mm^*)$ to the same order, vanishing for a large body; the recoil is even in $k$ and the twist odd, so after an absorption and an emission at $\mathbf k_{\mathrm{out}}$ the body keeps $\mathbf k_{\mathrm{in}} - \mathbf k_{\mathrm{out}}$ exactly, under the inversion of any axis as well. (f) \emph{Compton.} The click balances the rotations by the beat rule of Section 4.3, the law's own (the born light at $\omega_e - \omega_g$, the energy in the click's unit $T\sin\omega$ following to the order $\omega_i\omega_j$ while no layer keeps a total, Section 2.4), and under the proposition the wave vector; balanced in the rotations the non-relativistic factor below is $\omega_0 / \tan\omega_0$, where a balance of the energies $T\sin\omega$ would give $\omega_0\cos^2\omega_0 / \sin\omega_0 = 0.5015$ at $[2, 3]$, the two one at nature's gap; with light's band $\omega_L = c|\mathbf k_L|$ and the matter band of (d) in the Klein-Gordon form, a quantum at rest absorbing $\omega_{L,\mathrm{in}}$ and emitting $\omega_{L,\mathrm{out}}$ at the angle $\theta$ obeys $(1 - c_m^2 / c^2)(\omega_{L,\mathrm{in}}^2 + \omega_{L,\mathrm{out}}^2) - 2\omega_{L,\mathrm{in}}\omega_{L,\mathrm{out}}(1 - (c_m^2 / c^2)\cos\theta) + 2\omega_0(\omega_{L,\mathrm{in}} - \omega_{L,\mathrm{out}}) = 0$ with $c_m^2 / c^2 = \omega_0 / \tan\omega_0$. At $c_m^2 = c^2$, nature's gap, this is $1 / \omega_{L,\mathrm{out}} - 1 / \omega_{L,\mathrm{in}} = (1 - \cos\theta) / \omega_0$, Compton's formula, $\lambda_{\mathrm{out}} - \lambda_{\mathrm{in}} = (2\pi c / \omega_0)(1 - \cos\theta)$ with $\lambda = 2\pi c / \omega_L$, the Compton wavelength $2\pi / (\sqrt 3\,\omega_0)$ Links. At a finite gap the outgoing quantum gains, to first order in $1 - c_m^2 / c^2$, $\delta\omega_{L,\mathrm{out}} = (1 - c_m^2 / c^2)\,c^2 q^2\,\omega_{L,C} / (2\omega_0\omega_{L,\mathrm{in}})$, $q$ the recoil wave vector and $\omega_{L,C}$ Compton's value; in the non-relativistic regime $\omega_{L,\mathrm{in}} \ll \omega_0$ this is one factor on the whole shift, $\lambda_{\mathrm{out}} - \lambda_{\mathrm{in}} = (2\pi / (m^* c))(1 - \cos\theta) = (2\pi / (\sqrt 3\,\omega_0))\,(\omega_0 / \tan\omega_0)\,(1 - \cos\theta)$, exact at every gap in that regime: the inertia $3\tan\omega_0$ against the rest rotation $\omega_0$, $\omega_0 / \tan\omega_0 = 1 - \omega_0^2 / 3 + O(\omega_0^4) = 1 - \tfrac23(1 - \nu) + O((1 - \nu)^2)$ ($0.778$ at $[2, 3]$ to first order), $0.752$ at $[2, 3]$ and $1$ at nature's gap; the factor $0.752$ is conditional on the beat rule, the law's balance of Section 4.3, and not on the click's energy $T\sin\omega$, under which it would be $0.5015$. (g) \emph{Absorption alone.} A free quantum at rest absorbing $\omega_L = c k_L$ with no emission would need $c k_L = \sqrt{\omega_0^2 + c_m^2 k_L^2} - \omega_0$, that is $k_L(c^2 - c_m^2) = -2c\,\omega_0 < 0$ for $c_m \le c$: forbidden at every gap, the free quantum's one window holding two writes. (h) \emph{The zone's edge.} A single real quantum making a plane pair needs $k_L^2(c^2 - c_m^2) \ge 4\omega_0^2$: impossible at nature's gap, and at $[2, 3]$ only at $k_L \ge 2\omega_0 / \sqrt{c^2 - c_m^2} = 5.85$ radians per Link, outside the zone's $|k_L| \le \pi\sqrt 3 = 5.44$.
(i) \emph{The integer recoilless condition.} The twist is an integer division of the magnitude with the sign carried, so $\delta k_a = \mathrm{sgn}(k_a)\,(|k_a|\,\mathrm{div}\,M) = 0$ whenever $|k_a| < M$ in the box's unit $2\pi / L_{\mathrm{box}}$ per Link: a body whose count exceeds the quantum's wave number takes no twist from one click; the remainder $\mathrm{sgn}(k_a)\,(|k_a|\,\mathrm{mod}\,M) = k_a$ accumulates at the body, and one whole unit per Link, of $k_a$'s sign, is written at every $M / |k_a|$ clicks on average, the first at the $\lceil M / |k_a|\rceil$-th, the same clicks for $-k_a$ as for $k_a$ with the kicks opposite, so that opposite remainders cancel, and $M \sum \delta k_a + r_a = \sum k_a$ exactly over any run of clicks, the booked total conserved at every click and the kick never per click below the threshold; at $|k_a| = M$ the twist is one unit at every click and the remainder $0$. The threshold $|k_a| < M$ compares a wave number in the declared box's units with a count, so it is a property of the computed world and no formula shared with nature; the box-free statement is the recoil in rotation of (e). (j) \emph{Light's wave vector in a well.} At static paces Rule3 keeps a wave's frequency, and light's local speed in a well of potential level $U$ is $c / n$ with the index $n = e^{2U}$ (S.2), so its wave vector is $k = \omega n / c$ at long wavelength (the paper's Eq. for the lens, the chromatic term $k^2(e^{6U} - e^{2U}) / 24$ beyond it; S.28): under the proposition the wave vector passed per click, the twist, is $n$ times the vacuum's at the same frequency, Minkowski's form of the momentum of light in a medium, and the recoil in rotation $k^2 / (2Mm^*)$ is $n^2$ times the vacuum's.
\emph{Check:} (a) and (b) evaluated in real arithmetic to $10^{-15}$, (d) at $k = 0.05$, (f) against the exact quadratic at $\omega_{L,\mathrm{in}} = 0.02$ and $\theta = \pi / 2$ (the frequency drop's ratio $0.757$, the factor $0.752$ and the relativistic remainder), (g) and (h) at $[2, 3]$, the accumulation of (i) for $k_a = \pm 3$ and $M = 7$ (the kicks at the third, fifth and seventh clicks, $+1$ for $+3$ and $-1$ for $-3$) and at the boundary $k_a = \pm 7$ (one kick at every click), the books' recurrence stepped for one repeated $k_a$ (its kicks) and for the mixed sequence $3, -5, 10, 3, 3, -7, 2$ against $M = 7$ (the kicks $0, 0, 1, 0, 1, -1, 0$, the books $2$, $7 \cdot 1 + 2 = 9$ the sum) with the identity $r_a + k_a = M\,\delta k_a + r'_a$ stated at every click, the rule odd under $k_a \to -k_a$ at $\pm 3$, $\pm 7$ and $\pm 10$ against $M = 7$, and the two recoils of (e) at $[2, 3]$, $k = 0.03$ and $M = 1000$, in rotation $1.34164 \times 10^{-7}$ and in the click's unit $8.9443 \times 10^{-8}$ at $T = 1$, their ratio $\cos\omega_0 = 2 / 3$, with the two factors of (f), $0.752$ under the beat and $0.5015$ under the energy, distinct at $[2, 3]$ and one at the closing gap; one Node and one Link, no run.
\emph{Status:} (a), (b), (d) derived; (c), (e), (f), (g), (h), (i), (j) derived under the proposition of Section 2.5.2 (hypothesis, fitted); lattice, with S.2 and S.28 for (j), Compton's factor lattice at the computed pair, nature's by construction, no test; (i)'s threshold a property of the declared box and no formula shared with nature. The angular shares of the emitted direction are not derived here.
\end{derivation}

\section*{The clock family on the Ports (Sections 2.3 and 5.1)}

\begin{derivation}[gravity's potential field on the Ports, Sections 2.3 and 5.1]
\emph{Inputs:} Rule3 at the paces, $z_{\mathrm{next}} = \mathbf M z_{\mathrm{now}} - z_{\mathrm{before}}$ for a plane configuration written as one complex level $z$ per Node, with the Link coefficients $M_{ij} = R_{ij}$, $R_{ij}$ here the normalised hop $R_a / w = \num / (3\,\den)$ at the vacuum's paces and the self term $S_i / w$, and the weights $d_i = 1 / p_i^2$ of Theorem 3, $d_i R_{ij} = d_j R_{ji}$ (Section 3.3; S.45); the band $\cos\omega = (\num / 3\den)\sum_a\cos k_a$ at the vacuum's paces (S.1), its group velocity $v_g = (\num / 3\den)\sin k / \sin\omega$ along an axis (S.1) and its inertia $m^* = 3\tan\omega_0$ (S.18, S.53); the clock $p_0 = \Gamma e^{-U}$ and the static metric $ds^2 = -e^{-2U}c^2dt^2 + e^{2U}d\mathbf x^2$, $c^2 = 1/3$, under the one assumption $hN = 1$ (Section 2.3; S.34); the booked count and its change by the currents (Section 3.4, Eq. (5)); the Peierls angle of a space Link written from the holder's odd lines, as magnetism's (S.41); Will's parametrised post-Newtonian form of $g_{0j}$, $-\tfrac12(4\gamma + 3 + \alpha_1 - \alpha_2)V_j - \tfrac12(1 + \alpha_2)W_j$ at $\zeta_i = \xi = 0$, and $\eta_N = 4\beta - \gamma - 3 - \alpha_1 + 2\alpha_2 / 3$ \cite{will}; and the law's write of the odd Link part (Section 2.3): the three odd Link parts are three odd lines of gravity's holder, sourced each interval at the matter's Links by the count's flux $f(F(\mathrm{next}, \mathrm{now}))$ in the count's unit (the flux per interval and not its sum, as the time level is sourced by the count $N$), carried by the holder's own line with its wall, so that for a rigidly moving source their static solution is $\mathbf V = U\mathbf v_s$ to first order as $U$ is Poisson's for $N$ (no number: one wall for the count and its flux, the one assumption for a moving source); the Port's phase reads the odd line's level along the Link, $V_{ij} = -V_{ji}$, and turns the arriving plane record by $\theta_{ij} = -\theta_{ji} = 4V_{ij}\,\omega_0 / c^2 = 12\,V_{ij}\,\omega_0$ of the record's family, $c^2 = 1/3$ light's own, the three axes' $3$ and no number of nature, as the hop's coefficient $M_{ij} = R_{ij}e^{+i\theta_{ij}}$.
\emph{Steps:} (a) \emph{The conserved form.} The Port's phase reads the odd line's level along the Link, the mean of its two ends' levels, as the sign's odd lines are read (S.33 (d)), so that $\theta_{ij} = -\theta_{ji}$ exactly. With $\theta_{ji} = -\theta_{ij}$ the weighted matrix $\mathbf K = \mathbf D\mathbf M$ has $K_{ji} = d_j R_{ji}e^{-i\theta_{ij}} = \overline{K_{ij}}$, Hermitian, by Theorem 3's own condition $d_iR_{ij} = d_jR_{ji}$ and nothing of $|e^{i\theta}| = 1$. For $Q(z, z') = \mathrm{Re}[z^\dagger\mathbf D z + z'^\dagger\mathbf D z' - z^\dagger\mathbf K z']$ and $z_{\mathrm{next}} = \mathbf M z_{\mathrm{now}} - z_{\mathrm{before}}$: $z_{\mathrm{next}}^\dagger\mathbf D z_{\mathrm{next}} - z_{\mathrm{before}}^\dagger\mathbf D z_{\mathrm{before}} = z_{\mathrm{now}}^\dagger\mathbf M^\dagger\mathbf D\mathbf M z_{\mathrm{now}} - 2\,\mathrm{Re}(z_{\mathrm{now}}^\dagger\mathbf M^\dagger\mathbf D z_{\mathrm{before}})$ and $\mathrm{Re}(z_{\mathrm{next}}^\dagger\mathbf K z_{\mathrm{now}}) = z_{\mathrm{now}}^\dagger\mathbf M^\dagger\mathbf D\mathbf M z_{\mathrm{now}} - \mathrm{Re}(z_{\mathrm{before}}^\dagger\mathbf K z_{\mathrm{now}})$, and with $\mathbf M^\dagger\mathbf D = \mathbf K^\dagger = \mathbf K$ the two cancel in $Q(z_{\mathrm{next}}, z_{\mathrm{now}}) - Q(z_{\mathrm{now}}, z_{\mathrm{before}})$ exactly: Theorem 3 with the phase, $\theta = 0$ its real case. The same holds for the first-order coefficient $M_{ij} = R_{ij}(1 + i\theta_{ij})$, an exact rotation by $\arctan\theta_{ij}$ times the even magnitude $\sqrt{1 + \theta_{ij}^2}$, the same for $ij$ and $ji$; on the lattice its two integer numerators $(\Gamma_\theta\,\mathrm{re}_j - V'_{ij}\,\mathrm{im}_j,\ \Gamma_\theta\,\mathrm{im}_j + V'_{ij}\,\mathrm{re}_j)$, $V'_{ij} = \Gamma_\theta\theta_{ij}$ an integer over the common divisor $\Gamma_\theta$, go over $w\Gamma_\theta$ in Rule3's one division per Node and per part, as the law's line has it, with S.33's remainder term; the magnitude scales the hop by $1 + \theta^2 / 2$, an even term of the hop's magnitude with no drag in it. The convention, once: the hop from $j$ to $i$ carries the phase $+\theta_{ij}$ with $\theta_{ij} = 4V_{ij}\omega_0 / c^2$ along the current, the symbol $\cos(k + \theta)$, so a plane record's wave number shifts by $+\theta$ and $\delta\omega = \omega(k + \theta) - \omega(k) = +v_g\theta$ to first order, Einstein's sign ($H = p^2 / 2m + 4\,\mathbf V\cdot\mathbf p$, the Lense-Thirring $2G\,\mathbf J\cdot\mathbf L / c^2r^3$); with the opposite assignment $-v_g\theta$. (b) \emph{Reversibility.} The backward step is $z_{\mathrm{before}} = \mathbf M(V_t)\,z_{\mathrm{now}} - z_{\mathrm{next}}$ with the same matrix, no inverse rotation; the order that keeps it exact is the count's: read the odd lines' levels $V_t$, step, source the odd lines by the flux $f(F(z_{\mathrm{next}}, z_{\mathrm{now}}))$ of the state after the step, $f$ the holder's floor division with its remainder kept. Theorem 1 holds as for the time level (the holder's line reversible, the source recomputed from the two held records backward); no count register is needed. A flux read directly from the two levels the step reads, $F(\mathrm{now}, \mathrm{before})$, would put $z_{\mathrm{before}}$ into its own coefficient, an implicit step, and Theorem 1 would fail; the flux of the state after the step is the law's. (c) \emph{Linearity and locality.} The step is linear in the record at the written current (the current quadratic in the record, the paces' own status), one Link's reach, explicit, the arriving now and before only; the clock $p_0$ and the rulers $p_a$ are untouched. (d) \emph{The odd part cannot be a pace.} For a one-part real line $a_{\mathrm{next},i} = \sum_j c_{ij}a_{\mathrm{now},j} - \sum_j d_{ij}a_{\mathrm{before},j}$ the symbol at the wave number $k$ is $\lambda^2 - C_1(k)\,\lambda + C_0(k) = 0$, $C_1$ and $C_0$ the step symbol's coefficients, both roots unimodular at every $k$ only if $|C_0| = 1$ and $C_1 / \sqrt{C_0}$ is real with $|C_1 / \sqrt{C_0}| \le 2$. An odd real coefficient $\epsilon$ on the neighbour's level gives $C_0 = 1$ and $C_1 = C_1^{(0)} + 2i\epsilon R\sin k \notin \mathbb R$ ($C_1^{(0)}$ its value at $\epsilon = 0$), so the roots cannot both be unimodular ($\lambda_2 = 1 / \lambda_1 = \overline{\lambda_1}$ would make $C_1 = 2\,\mathrm{Re}\,\lambda_1$ real): to first order $|\lambda| = 1 + \epsilon R\sin k / \sin\omega = 1 + \epsilon v_g$ and no frequency shift, the number the phase delivers as the shift $+\theta v_g$ turned by $i$. On the neighbour's change $\mathrm{now} - \mathrm{before}$ it gives $C_0 = 1 + 2i\epsilon R\sin k$, $|C_0| > 1$, one root growing, $\delta\omega = -\epsilon R\sin k$ and $|\lambda| - 1 = \epsilon R\sin k\tan(\omega / 2)$; on the neighbour's next (time-centred) the roots are unimodular but the step implicit. So local, explicit and exact are two of three for a one-part line whose odd part is to be proportional to $V$: a theorem (the one exception the staggered step, a shift of the before level by whole Nodes, whose drag is quantised and not $V$'s), the five forms of C.4 its instances; the two-part record escapes it because its coefficient is Hermitian and not real, the symbol $2R\cos(k + \theta) + S$ real and $|\lambda| = 1$. (e) \emph{The number.} A uniform $\theta$ along an axis shifts the wave number, $\delta\omega = \omega(k + \theta) - \omega(k)$ exactly under the convention of (a), $+v_g\theta$ to first order; on a ring the two senses split by $2v_g\theta$ (the Sagnac form). For slow matter $v_g = k / m^*$, so under the delay form $\theta = 4V\omega_0 / c^2$, $\delta\omega = +(k / m^*)\,4V\omega_0 / c^2 = +4\,(c_m^2 / c^2)\,\mathbf V\cdot\mathbf k = 4\,(\omega_0 / \tan\omega_0)\,\mathbf V\cdot\mathbf k$ with $c_m^2 = \omega_0 / m^*$ (S.21), $3.01\,\mathbf V\cdot\mathbf k$ at $[2, 3]$ and Einstein's $4\,\mathbf V\cdot\mathbf k$ at nature's gap, that is $\delta\omega / \omega_0 = +4\,\mathbf V\cdot\mathbf v / c^2 = 12\,\mathbf V\cdot\mathbf v$ with $\mathbf v = \mathbf k / m^*$, for every family, to first order in $V$ and $k$, the quotient $c_m^2 / c^2 = 0.752$ at the computed pair against Einstein's $1$; the share form $\theta = 4V\sin\omega_0 / c^2$ would carry $\sin\omega_0 / \omega_0$ ($0.886$ at $[2, 3]$) and is not the law's. Einstein's shift of a slow mass's energy by $+4\,\mathbf V\cdot\mathbf p$ has this form at nature's gap. (f) \emph{The coefficient 4 from the law's own exponents.} A source moving at $\mathbf v_s$ carries the static write boosted: with $c^2$ in $g_{00}$, $ds^2 = -e^{-2U}c^2dt^2 + e^{2U}d\mathbf x^2$, the boost at $c$, $dt_r = dt - \mathbf v_s\cdot d\mathbf x / c^2$ and $d\mathbf x_r = d\mathbf x - \mathbf v_s\,dt$ to first order (the band's own boost at long wavelength and small gap, S.33 (c)), the cross term of $-e^{-2U}c^2dt_r^2 + e^{2U}d\mathbf x_r^2$ is $2(e^{-2U} - e^{2U})\,\mathbf v_s\cdot d\mathbf x\,dt = -8U\,\mathbf v_s\cdot d\mathbf x\,dt + O(U^2, v_s^2)$, the $c^2$ cancelling, so $g_{0j} = -4Uv_{s,j} = -4V_j$ with $\mathbf V = U\mathbf v_s$; the flat part cancels, $2(1 - 1) = 0$ (derived under the one assumption and the band's own Lorentz form at long wavelength and small gap (S.21), the boost a property of the band and not of the lattice's group). The one-way time of light along a Link, from $-c^2dt^2 - 8V\,dt + 1 = 0$, is $dt = 1 / c - 4V / c^2$ for one sense and $1 / c + 4V / c^2$ for the other, so the delay per Link is $4V / c^2 = 12V$ intervals and the Port's phase is the family's $\omega_0$ times it. The $4$ is $2 + 2\gamma$, the clock's exponent and the ruler's, equal at $\gamma = 1$ ($hN = 1$), the exponent of the Link's coefficient relative to the clock's ($R_{ij} \propto p_a^2 = p_0^4q^2 / \Gamma^4 \propto e^{-4U}$), Will's $2(1 + \gamma)$. The odd lines' source $U\mathbf v_s$ carries no number because the count flows at the group velocity: a packet's count through a Link per interval is the Node's count times $v_g = v_s$ to first order (the continuity of the conserved share, Theorem 4, Section 3.4); read as the raw Wronskian, $F_{ij} / W = -\num\sin k_a \approx -\num\,m^*v_s$ would bring $\num$ and $m^*$ in. (g) \emph{The parameters.} The harmonic form $g_{0j} = -4V_j$ is the standard post-Newtonian form $-\tfrac72 V_j - \tfrac12 W_j$ under the relabelling $x^0 \to x^0 + \tfrac12\partial_t\chi$, $\chi$ the superpotential, a relabelling of the reading and not of the lattice; so, for a plane record of one rotation sense, $\alpha_1 = \alpha_2 = 0$, and at nature's gap, $\beta = \gamma = 1$, $\eta_N = 0$ ($\alpha_2$ and $\eta_N$ under the relabelling of S.56 (g), $g_{00}$'s velocity terms not fixed, S.35); for bodies the signed share (the paper's Section 8.1); at the computed pair $\beta_K = 1.25$ and $\gamma_K = 1.5$ give $\eta_N = 0.5$, the finite gap's own departure. The write without the odd Link part has $g_{0j} \equiv 0$, so $7 + \alpha_1 - \alpha_2 = 0$ and $1 + \alpha_2 = 0$: $\alpha_2 = -1$, $\alpha_1 = -8$ (S.35), the pair of the write without the odd Link part, of which the paper's Sections 1.2 and 9 name $\alpha_1 = -8$. (h) \emph{The two predictions.} Light is a one-part real line with $\omega_0 = 0$: nothing to rotate, by (d) no $V$-proportional real odd coefficient allowed, so light is not dragged, the Lense-Thirring deflection of light near a rotating body absent, where frame dragging has been read on matter alone (gyroscopes, satellites, pulsars); and a Port cannot read a record's own frequency without reading the record (c), so the phase is the family's $\omega_0$ times the delay and the coupling of a fast record falls as $\omega_0 / \omega = 1 / \gamma_L$, $\gamma_L$ the Lorentz factor, against Einstein's $+4\,\mathbf V\cdot\mathbf p$ at every speed; every measured drag is slow. A free one-part matter line (the scalar entry, dimension $1$), which nature's metric would drag, is likewise undragged, untested; bound bodies are dragged by their plane records' signed share of inertia, the neutron's three real lines not at all, which is where every measurement sits.
\emph{Check:} algebra and a nine-Node numerical check of the symbol's roots: (a) by the three identities above; (d) by the roots of the quadratic symbol at $\epsilon = \theta = 0.02$ and $k = 0.5$, light $[\Gamma, \Gamma]$ and matter $[2, 3]$, against the closed forms: on the level $|\lambda| = 1.01136$ and $1.00278$ against $1 + \epsilon v_g = 1.01130$ and $1.00277$; on the change $\delta\omega = -3.18 \times 10^{-3}$ and $-2.131 \times 10^{-3}$ against $-\epsilon R\sin k = -3.20 \times 10^{-3}$ and $-2.131 \times 10^{-3}$, $|\lambda| - 1 = 4.7 \times 10^{-4}$ and $1.00 \times 10^{-3}$ against $\epsilon R\sin k\tan(\omega / 2) = 4.6 \times 10^{-4}$ and $1.00 \times 10^{-3}$; time-centred and the phase $|\lambda| = 1.000000$, the phase's $\delta\omega = 1.129 \times 10^{-2}$ and $2.82 \times 10^{-3}$ against $v_g\theta = 1.130 \times 10^{-2}$ and $2.77 \times 10^{-3}$; (e) the ring's splitting at $[2, 3]$, $\theta = 0.01$: $2v_g\theta = 0.001742$, $0.003224$, $0.004531$ at $k = 0.3$, $0.6$, $1.0$, equal to $\omega(k + \theta) - \omega(-k + \theta)$ to six digits; the centred difference $(\omega(k + \theta) - \omega(k - \theta)) / (2 \cdot 4Vk)$ at the Inputs' own $\theta = 12V\omega_0$, $V = 10^{-3}$, $k = 0.1$, is $0.7500$ against $c_m^2 / c^2 = 0.7523$, the $k^3$ term of $v_g$ lowering it by $0.3$ percent (the forward and backward differences $0.788$ and $0.712$, the second order in $\theta$); (f) $72$ arcseconds (the perihelion at the small-gap $\beta = -1$, S.19, S.34) and $39.2$ milliarcseconds per year (Gravity Probe B, S.35 (c)) are not this derivation's; (g) the gyroscope's precession in the harmonic form, Gravity Probe B's $39.2$ milliarcseconds per year for a plane configuration of one rotation sense, $19.6$ for its quartz rotors ($Z/A = 0.50$) against its measured $37.2 \pm 7.2$ \cite{everitt}, $2.4$ standard deviations, the coefficient S.35 (c)'s; the Hermitian form on a periodic chain of $9$ Nodes with random paces in $[0.8, 1.0]$ and angles in $[-0.3, 0.3]$ over $2{,}000$ steps conserved to a relative $3 \times 10^{-15}$ for both the exact and the first-order hop (a check, not a result).
\emph{Domain:} every plane record (the charged families and the bound bodies made of them) at any pace; a real line (light; the scalar entry of matter) carries no phase; first order in the source's velocity and in $V$; the count's flux booked in the count's unit as the content is, the odd lines its carrier. \emph{Status:} derived under the one assumption, the law's write of the odd Link part; (a) to (c) theorems, (d) a theorem for a $V$-proportional odd part; (e) to (g) derived under the one assumption for a plane record of one rotation sense, the drag odd in the sense (the symbol $\cos(k \mp \theta)$, the two senses' shifts opposite at one velocity; the paper's Section 8.1), clicks at nature's gap for the three parameters, the paper's Sections 5.1 and 8.1, which the signed share of a body's inertia, $Z / A$ under the paper's proton hypothesis (its Section 4.1), turns into findings against nature; (h) derived under the one assumption, lattice, untested for the two predictions and for antimatter's anti-drag and positronium's zero.
\end{derivation}


\section*{C. The derivations and passages the paper's statuses cite beyond S.1 to S.56}\label{sec:carried}

\noindent The blocks below hold the derivations and the passages the paper's and the supplement's statuses cite beyond Derivations S.1 to S.56, numbered C.1 to C.9 and cited as (C.$n$); their pointers are written as the paper's Section, Eq. and Theorem numbers and this supplement's S.$n$ and C.$n$, and the derivations among them keep their form, Inputs, Steps, Check and Status.

\subsection*{C.1 The postulates}

The model rests on the law's seven postulates, with the NodeDetector's write added by decision. (1) The model's world is Nodes and Events. An Event is a level stepped by Rule3 with its remainder. A click is a NodeDetector's report and, by that decision, its one write at the Node the quantum reached it from. Nothing else happens on the lattice; the click layer's state is the books of Section 2.4 of the paper. (2) A Node holds bounded local information, read from itself and its six neighbours, and nothing of the past is kept beyond the one level before. (3) Consistency is local and causal, one Link per interval for every reading of the interval (Section 2.4 of the paper), with one exception: the credit of a click by the shares over every declared NodeDetector, the one non-local element, not derived. (4) The causal speed is one Link per interval. (5) Every physical calculation is on bounded integers, bounded for the run's length by the load's bound (the width's room, the amplitude bound and the register's room checked when the universe is declared, the uniform light state $a_t = A + Bt$ the case the bound names), with no number but integers, no square root and no random draw in the lattice's step; the NodeDetector draws with the world's one draw, a body's NodeDetector with its own seed. (6) A measurement is a NodeDetector's click: the net current into a declared region through its front boundary over one interval, accumulated over the window, and never a Node. (7) The lattice, where no one measures, is deterministic and reversible to the bit; a NodeDetector's write is the NodeDetector's act, forward only and outside the reversal; and the NodeDetectors' clicks are the only thing claimed to represent nature. The line and its clock are fixed: every mend of the model is a family, a holder, a declaration or a read, never a change of Rule3 (\claimmark{assumption}).

\subsection*{C.2 The functional of the bound state}

Let a body of $C$ quanta have the record $\phi$ over its region, with $\sum\phi_i^2 = 1$ and the counts $C_i = C\phi_i^2$. Let every holder $r$ it sources be a hollow ($\sigma_r = +1$) or a hill ($\sigma_r = -1$), with the static kernel $\mathbf G_r$ of its line and the divisor $E_r$. The well at a Node, in levels, is then $\ell_i = \sum_r \sigma_r (C / E_r) \sum_j G_{r,ij}\,\phi_j^2$ ($\ell_i = \Gamma U_i$; Newton's potential is $\Phi = -c^2 U$, S.28). To first order in $\ell_i / \Gamma$ the well raises the one-Node rotation $2\cos\omega$ by $4g\ell_i / \Gamma$, where $g = 1 - \cos\omega_0$ is half the rest term $m^2$ of Section 2.2 of the paper. The body's fixed point is then a critical point on the sphere of
\begin{equation}\label{eq:functional}
F[\phi] = \frac{\num}{3\,\den}\sum_{\text{Links}} (\phi_i - \phi_j)^2 - \frac{2g}{\Gamma}\sum_r \sigma_r\,\frac{C}{E_r}\sum_{ij}\phi_i^2\,G_{r,ij}\,\phi_j^2,
\end{equation}
the band's pressure less the wells' energy, each well counted once because the record is its own source (Hartree's form). The record is the top mode of Rule3's symmetric form at the paces, and $dF / d\phi = 2\lambda\phi$ on the sphere is the fixed point's own equation, with $2\cos\omega_b = 2\cos\omega_0 - \lambda$ (\claimmark{theorem}, the functional's variational derivative, S.13's first step; \fence{lattice}). Let $\phi_\lambda(x) = \lambda^{-d / 2}\phi(x / \lambda)$ be the body dilated about itself, with $d$ the lattice's dimension. The pressure is then $K / \lambda^2$ and each holder's well energy $V_r / \lambda^{s_r}$, with $K$ and $V_r$ the body's own pressure and well energies at $\lambda = 1$, where $s_r$ is the holder's scaling power at the body's size: $s_r = d$ beyond the holder's reach, and $s_r = d - 2$ within it or for the massless holder.

\subsection*{C.3 The two branches, and the status of a standing body}

The binding holder's fixed point has two branches, a wide cloud and a compact bound state one Node wide, and the start's width alone chooses between them (Section 4.2 of the paper; \claimmark{theorem} for the equation; the branches a lattice reading, named; S.13; \fence{lattice}). The integer two-branch computation is not reproduced here; no computed two-branch result is claimed. No standing body is claimed from a run. The computed body is coarse at the Node's width by construction (a few Links, about $10^3$ levels per Node). So its drift and breathing are the lattice's budget and the lay's, and not the law's: $(\mathrm{Link} / R)^2$, the Node's width against the body's size, and $\sqrt n / A$, the integer step's walk over $n$ intervals against the amplitude $A$. The first is $2 \times 10^{-23}$ at a radius of $1.4$ fm and Fermi-LAT's Link (S.54); the law's derivations are of the line without the division. The body's rows carry only what is computed: its definition (Section 4.2 of the paper), the static fixed point and the modes by hand (S.43, S.52), the stability conditions as derived (Section 4.2 of the paper) and its clicks (Section 4.3 of the paper), and no sentence that a body stands under Rule3. Whether a laid body stands under Rule3 within that walk is not this paper's claim.

\paragraph{The atom: known from clicks, laid on the lattice.} The atom's electron record is a plane record of the holder of the sign; the rule's own universe declares no plane family (Section 4.1 of the paper), so the atom here and the body that writes light (Section 4.3 of the paper) are the law's lines for a row no declared universe of this paper declares, and the declared matter is a real line. Under line (1) every charged record carries its own field of the sign, one field per record (S.52), the one exception to the reciprocity of the paper's Section 2.3. The atom is known from nature's measurements: its spectra, its lifetimes and its scattering. Nothing of it is seen on the lattice. What is tested is whether the way from the clicks to the lattice and back carries the atom. This is a supporting finding, with its state named (S.43). In the model, an atom is a bound body of the holder of the sign; a hydrogen nucleus, one quantum of the proton's family, writes the time angle $\theta_0 = Z\alpha c / r$, with $c = 1 / \sqrt 3$, into the sign holder's field by its Wronskian. Electrons, which are quanta of the matter family, read that angle into their rotation. The standing records that Rule3 makes in that angle are the atom's levels, Schr\"odinger's in the slow limit. With the nucleus's angle alone, the levels have Bohr's form, $E_n \propto 1 / n^2$ with $a_0 = \sqrt 3 / (m^* Z\alpha)$ (\claimmark{derived} in the slow limit, \fence{clicks}; S.43). A record that reads its own sign field (S.43 (c)) returns Rydberg's form only at large $Z$, as $1 - 1.25 / Z$, and does not give hydrogen. Under line (1) of S.52 (\claimmark{assumption}, \claimmark{fitted}) a record reads every sign field but its own, and hydrogen's lines are those of S.52. The holders of the content keep the self-read (line (2) of S.52, \claimmark{assumption}). The atom of this paragraph is derived and not run, its derivations standing.

\paragraph{What the atom's clicks give, and what they do not.} What a click at the atom gives is derived: the born light is at the resonance $\omega_e - \omega_g$ and not at the energy-conserving frequency, with the share's difference booked; absorption erases the photon's empty wave inside the click's causal cone; and there is no recoil, consistently, since the write conserves no momentum and the nucleus pinned by line (2) of S.52 has infinite inertia, so nature's recoil shift is outside the claim. Bohr's levels, the threshold, the lines and the two-mode binding are unchanged by it, and the cause of a click is outside the paper's claims. The lines are beats of two levels, and the mass defect is the click layer's reading of a frequency. The orbital Zeeman splitting has Land\'e's $g_L = 1$ (\claimmark{derived}; supporting findings with their numbers in S.43; \fence{clicks} for the levels' ratios, \fence{lattice} for the records). The electron and the nucleons are quanta of their families. An electron is one quantum, of energy $T\sin\omega_e$ with $\omega_e$ here the electron family's rest rotation, ten orders below the compact threshold, and nothing smaller than a quantum clicks. In the model, mass is the family's gap, $T\sin\omega_s$ ($s$ the family), bound or free; what decays is a body, from which a click at its front parts a quantum; a quantum itself never breaks, since it has no part that the click credits (\claimmark{derived}; \fence{clicks}). The criterion is one line, a rule of reading: what nature splits with energy is read as a body on the lattice, and what no energy splits is read as a quantum of a family. A nucleus of several nucleons, the atom and the molecule are bodies by the one mechanism the law has for a body, a record standing in the well of a holder: the sign holder for the atom and the nuclear holder for the nucleus (\claimmark{assumption}). A nuclear holder that the electron reads is set against hydrogen's $1s$ and muonic hydrogen, a comparison that supports excluding it (\claimmark{derived} for the shift's form, the exclusion a comparison; one derivation; \fence{clicks}; S.43). So the nuclear holders of Online Resource 1, Section B.10 are read by the nucleon families and not by the electron, a selective read that is declared. The table of S.52 sets the atom's clicks, the lines and the thresholds, against nature's under the two lines decided for the law, with how each was derived; this is a supporting finding and no more. What the atom of the model does not give is listed with its roots in Section 8.1 of the paper.

\subsection*{C.4 The carrier: the five forms tried}

\renewcommand{\thederivation}{C.\arabic{derivation}}\renewcommand{\theHderivation}{C.\arabic{derivation}}\setcounter{derivation}{3}
\begin{derivation}[the carrier: the five forms tried, Section 8.1]
\emph{Inputs:} $g_{0a}$ is the term of the wave operator odd in time and odd in space together, $2g^{0a}\partial_t\partial_a$; Rule3's step is explicit (next from now and before at the Node and its six neighbours), of one Link's radius, bijective (Theorem 1), and conserves the form (Theorem 3), whose proof needs the Link coefficients reciprocal; a real record has no phase; the dispersion of a term is read by inserting the plane wave $a = A\cos(kx - \omega t)$.
\emph{Steps, the four ways on the synchronous clock:} (i) \emph{The one-sided real mixed term} $C_a[(a_{+a} - a_{-a})_{\mathrm{now}} - (a_{+a} - a_{-a})_{\mathrm{before}}]$ added to the line gives $w(e^{-i\omega} + e^{i\omega}) = S + 2\sum_a R_a\cos k_a + C_a(2i\sin k_a)(1 - e^{i\omega})$, and $1 - e^{i\omega} = -2i\sin(\omega / 2)\,e^{i\omega / 2}$, so the term is $4C_a\sin k_a\sin(\omega / 2)[\cos(\omega / 2) + i\sin(\omega / 2)] = 2C_a\sin k_a\sin\omega + 4iC_a\sin k_a\sin^2(\omega / 2)$: the $g_{0a}$ part beside an imaginary part, a growth or decay of $\omega$ that breaks the form and the exact inverse. \emph{Check:} at $\Gamma = 6{,}000$, light's family, $k = 0.3$, $C / w = 10^{-2}$: $\omega = 0.16984 - 2.47 \times 10^{-4}\,i$ against the band's $0.17277$ (the first-order prediction $-2.56 \times 10^{-4}\,i$). (ii) \emph{The time-centred mixed term}, with the neighbours' next, is conservative but implicit: it needs the neighbours' next, outside the one-Link dependency radius and the explicit inverse. (iii) \emph{A real asymmetric hopping} with a curl, $R_{i\to j} / R_{j\to i}$ not a ratio of Node weights, has no symmetric step kernel, so Theorem 3 fails. (iv) \emph{A Peierls phase} on the Link carries $g_{0a}$ for a plane record only, a real record having no phase, and it shifts $k$ by the same angle for every record, a force independent of $\omega_0$: it couples to the charge and not to the mass; light and neutral matter are real lines. So none of the four ways tried, real, local, explicit and bijective lines on the synchronous clock, carries $g_{0a}$.
\emph{The fifth way, and why it fails:} clock the cube's two sublattices alternately, half an interval apart, each Node reading its neighbours in the two slices that bracket its own time, their mean into the hopping and their difference into an odd Link level $C_{ij}$: $w(a_i(t+1) + a_i(t-1)) = S a_i(t) + \sum_j R_{ij}[a_j(t + \tfrac12) + a_j(t - \tfrac12)] / 2 + \sum_j s_{ij} C_{ij}[a_j(t + \tfrac12) - a_j(t - \tfrac12)]$. Its dispersion, $2w\cos\omega = S + 2\cos(\omega / 2)\sum_a R_a\cos k_a + 4\sum_a C_a\sin(\omega / 2)\sin k_a$, is real for every $k$, the step's eigenvalues lie on the unit circle (to $10^{-15}$ in real arithmetic on a ring of $12$ Nodes, $C$ uniform or random Link factors), the mixed term is exactly $g_{0a}$, and it is real, local, explicit and bijective. But a uniform record sampled at the neighbours' half-slices is $a(t)\cos(\omega / 2)$, so the hopping does not cancel at $k = 0$ and the rest rotation in a content becomes $1 - \cos\omega = (1 - \nu)N^2 + \nu N^4(1 - \cos(\omega / 2))$ in place of $(1 - \nu)N^2$; at small $\omega$, $\omega^2 = 8(1 - \nu)N^2 / (4 - \nu N^4)$, and for $\nu \to 1$, $\omega / \omega_0 = N\sqrt{3 / (4 - N^4)} = N(1 - 2U / 3 + \dots)$: a clock in a well runs at $\omega_0(1 - 5U / 3)$, a redshift $5 / 3$ of general relativity's. \emph{Check:} $\sqrt{3 / (4 - 0.9^4)} = 0.9472$, the three-dimensional staggered band's $0.948$ at $N = 0.9$, $\nu = 0.99$ on a $13^3$ lattice, its light speed $c^2 = 4 / 9$ and its top $2\pi / 3$; Pound and Snider's $0.9990 \pm 0.0076$ \cite{pound_snider} excludes the $67$ percent at $88$ standard deviations, Gravity Probe A's $1.4 \times 10^{-4}$ at two standard deviations (Vessot et al.\ 1980 \cite{vessot}) at $5 \times 10^3$. Demanding the exact redshift for every small-gap family, $P_0(4 - \nu) = N^2(4 - \nu P_a)$ for all $\nu$ ($P_0 = (p_0 / \Gamma)^2$ and $P_a = (p_a / \Gamma)^2$ the clock's and the Link's pace factors), forces $P_0 = N^2$ and $P_a = 1$: the Link may not read the content. So the fifth way fails the clock as the first three fail the inverse, the form or the radius and the fourth the real record, and the result stands on Rule3's clock for the five forms tried, a derivation and not a theorem over every form: for a one-part real line $\alpha_1 = -8$ and the drag $0$ stand on the five forms tried and on S.56 (d)'s theorem; for a plane configuration the Port's phase of S.56 carries $g_{0a}$. \emph{Status:} theorem for a one-part real line (S.56 (d)), the staggered step's whole-Node shift aside, whose drag is quantised and not $V$'s; derived for the five forms tried, one derivation; clicks.
\end{derivation}

\subsection*{C.5 Confinement as the indivisibility of a record}

\renewcommand{\thederivation}{C.\arabic{derivation}}\setcounter{derivation}{4}
\begin{derivation}[confinement as the indivisibility of a record, Section 4.1]
\emph{Inputs:} a family of dimension $3$: three planes (three complex lines) at every Node, each two real lines with a rotation in its plane and one Wronskian, stepped by the one line, with a rotation among the three planes on each Link (the colour holder's phases); the NodeDetector's definition (Section 2.5): a click is one whole quantum of a record credited to a declared region over its window, drawn by the share $Q_R / W_c$ of the inflow through the region's front; the share of a record of several lines is the sum of the lines' forms (Theorem 3 line by line, the Link rotation a unitary on the triple, so the step kernel stays symmetric in the paces' weights and the weighted summed form is conserved at static paces); the readings of Eq.~(1) are the record's, never one line's: the plane's $\mathrm{re}$ and $\mathrm{im}$ are booked together in $D$, $F$ and $W$, and so are the three planes.
\emph{Steps:} (a) \emph{No click reports one line.} The click's object is the record's quantum, the share summed over the lines; no reading of the law credits a line apart, and the credit draws regions, not lines (Section 2.5, ``never a Node'', and the one form of the click). So a record of dimension $3$ clicks whole or not at all: no NodeDetector ever reports one colour plane, as none reports a plane's real part alone. This is confinement with no coupling and no number: the quark model's ``no free quark'' is the dimension's indivisibility under the click. (b) \emph{The partons.} A hard click, a scattering read at the momentum transfer $q$, reads the record's current at the wave number $q$ through the region's front; the lines' forms sum to the record's, so what it reads is three fractions summing to one, the three planes' shares of the record's form (the sum to one the dimension's), not of the momentum: the quark model's three is the dimension's three, and the fractions are the record's lay and its Link rotations; a line's share is read as a fraction over many whole clicks, never as one line's click, which (a) forbids. (c) \emph{Scaling.} The levels are deterministic and the count is conserved, so no pair is created and no coupling runs (Section 8.1; S.50 (6) and (13), the one root); what a read at the wave number $q$ gives is the lines' Fourier weights at $q$, which depend on the lines' lay (the fixed profiles $\delta_0 + \delta_1$, $\delta_0$, $\delta_0$ give $2 / 3, 1 / 6, 1 / 6$ at $q = 0$ and $1 / 2, 1 / 4, 1 / 4$ at $q = \pi / 2$) and are $q$-independent only for lines of one shape; the hard click's read operator is not derived here, so the fractions' dependence on $q$, Bjorken's scaling or its violation, is uncomputed, not exact, where nature's measurements show logarithmic violations of order ten percent across the measured decades of $q^2$ (from the SLAC experiments to HERA). (d) \emph{The colourless bodies.} The holder's sources are the three lines' Wronskian-like currents; a body whose three senses sum to $0$ writes no monopole into the colour holder, its far level falling as the kernel's gradient and not the kernel (for lines laid alike at every Node the exterior level is exactly $0$): the colourless states, and a pair of opposite sense or three coplanar axes are the two ways to $0$. (e) \emph{What it is not.} The binding between two nucleons is between separable records (a nucleus parts into them) and needs the total form to fall at a kept count, which Theorem 3 forbids at fixed paces (S.36); so the word ``strong force'' parts in two, the confinement the click's under the dimension-$3$ hypothesis and the nuclear binding open. \emph{Check:} the conserved form under the Link rotations was verified in exact arithmetic on a chain at the vacuum's paces with random rotations (the kernel symmetric there, the form's drift $0$); the credit's unit is the check's own definition, read in every checked declared universe. \emph{Status:} derived from the click's one form under the dimension-$3$ family, an explicitly named hypothesis, two independent derivations; clicks (no free quark, the partons' fractions); the scaling uncomputed, the read operator undeclared (S.50 (13)).
\end{derivation}


\subsection*{C.6 The nuclear holder's well: the threshold, the binding at a size, the capture and its reverse}

\renewcommand{\thederivation}{C.\arabic{derivation}}\setcounter{derivation}{5}
\begin{derivation}[the nuclear holder's well: the threshold, the binding at a size, the capture and its reverse, Section 8.1]
\emph{Inputs:} a gapped holder of the content with the reach $1 / \kappa$ (S.11) and a declared write weight $W$, sourced by the nucleon families alone at that weight and read by them at the read weight $1$, as every holder's level is read (S.26 (d$'$)); the far level of a source with the reach, $3W\sin\omega_p\,e^{-\kappa r} / (4\pi r)$ per quantum (S.25, S.37); the clock's read of a level, $\omega \to \omega e^{-x}$ with $x = L / \Gamma$ (S.2); the matter band's inertia $m^* = 3\tan\omega_p$, $3\omega_p$ at a small gap (S.22); the second moment's light (S.40) and the two-mode line under the rotation (S.41). \emph{Steps:} (a) \emph{The well and its threshold.} Each quantum reads the other's level into its clock, a potential $-\omega_p x$ per quantum in rotation units, so the pair's potential is $V(r) = -g^2 e^{-\kappa r} / r$ with $g^2 = 3W\omega_p\sin\omega_p / (4\pi\Gamma)$, $3W\omega_p^2 / (4\pi\Gamma)$ at a small gap (one quantum writes the holder at the write weight $W$, the far level with the reach is $3W\sin\omega_p e^{-\kappa r} / (4\pi r)$, and the reading quantum's clock reads it at the weight $1$, $\omega \to \omega e^{-x}$ with $x = L / \Gamma$, a potential $-\omega_p x$ in rotation units; the weight enters once, and a negative write weight writes a hill and is repulsive by the same line), and each record moves at its own inertia $\mu = m^* = 3\tan\omega_p$ about the pinned centre, two records being two fields with no joint amplitude (postulate 2 of C.1). The slow limit's own equation in Yukawa's well binds if and only if $2\mu g^2 / \kappa \ge 1.6798$, the known threshold of that equation (mathematics), that is $W \ge 1.17\,\Gamma\kappa / \omega_p^3$: $591$ (the small-gap formula at $[2, 3]$; $502$ at the exact inertia) in the rule's own universe ($[2, 3]$ as the nucleon, the binding holder $[2400, 2401]$ at the range $20$ Links, $\Gamma = 6{,}000$), and $6.5 \times 10^{11}\,\Gamma$ at nature's deuteron range $1.41$ fm and the LHAASO Link (the Link $2.5 \times 10^{-22}$ m at $\omega_e = 3.7 \times 10^{-10}$, S.27, $\kappa = 1.76 \times 10^{-7}$ per Link, $\omega_p = 6.8 \times 10^{-7}$). Under the pair hypothesis alone (the pair of two bound records, the law's line under its own name) the relative part carries $\mu = m^* / 2$ and the threshold doubles, $W \ge 2.35\,\Gamma\kappa / \omega_p^3$, $1{,}183$ at $[2, 3]$ ($1{,}004$ at the exact inertia), nature's number doubling with it. (b) \emph{The binding at a body's size.} A body of radius $R$ Links of quanta of the gap $\omega_s$ standing in any well of the $1 / r$ kind that fits its size has, by the virial theorem of the slow limit's own equation (a known form, mathematics of that equation), the binding fraction $b = 1 / (6\omega_s^2 R^2) = (\lambda_s / R)^2 / 2$ with $\lambda_s = 1 / (\sqrt 3\,\omega_s)$ its Compton length, the kinetic rotation over the rest rotation, whatever the holder's weight; the factor is the convention's: by the virial theorem the binding is the kinetic rotation $\langle k^2\rangle / (2m^*)$ with $m^* = 3\omega_s$, $1 / (6\omega_s^2 R^2)$ under $kR = 1$ with $R$ the Bohr-like radius, and $1 / (2\omega_s^2 R^2)$ under $R$ the form-weighted root-mean-square radius, where the hydrogenic ground state has $\langle k^2\rangle\langle r^2\rangle = 3$; one convention, its factor, the order of the estimate, a finding, the profile entering through $\langle k^2\rangle\langle r^2\rangle$ at the next order: the deuteron at $2.13$ fm $4.9 \times 10^{-3}$ under the Bohr-like radius and $1.46 \times 10^{-2}$ under the rms radius against nature's $1.19 \times 10^{-3}$, iron at $3.75$ fm $1.6 \times 10^{-3}$ and $4.70 \times 10^{-3}$ against nature's $9.4 \times 10^{-3}$ (the proton, one quantum in the model, has no binding of nature's to set beside it; its $3.1 \times 10^{-2}$ at $0.84$ fm is the one quantum's own well depth, a diagnostic, no binding). (c) \emph{The capture.} Two slow packets entering the well are a superposition of the bound mode and the continuum; the relative coordinate breathes at the beat $\Omega$ between them, the proton's charge rides it, so the body writes a dipole $d = S r_p$ into the holder of the sign, $S = k_w / (2E_s)$ and $r_p = r / 2$, the first order of S.40's expansion since the charge is not symmetric about the centre, and radiates at the beat in the dipole's form, the time level carrying half of Larmor's power in form, the odd lines' part depending on their wall and coupling, which of the two couplings the odd lines carry open; the body's own paces move with the level it writes, so it settles into the bound mode deterministically, the binding leaving as light at the beat; no rate and no settling time is claimed. (d) \emph{The reverse.} A light record of rotation $\omega_L \ge B$ ($B$ the bound mode's binding in rotation units) at the body drives the bound mode into the continuum by the two-mode line under the rotation (S.41), the matrix element the inter-mode current, the dipole; the bound mode's tail $e^{-\sqrt{2\mu B}\,r} / r$ and the continuum's density give the shape $(\omega_L - B)^{3/2} / \omega_L^3$, Bethe and Peierls's \cite{bethe_peierls}, peaking at $2B$. (e) \emph{The size law.} A cloud of $n$ quanta in one well has $R^* \propto 1 / n$ and $b \propto n^2$, iron $28$ times smaller than the deuteron and bound $800$ times harder, against nature's $A^{1/3}$ and its saturation; the second holder, of shorter range and negative weight, is the saturating declaration. \emph{Check:} $1.6798 \times 4\pi / 9 = 2.345$; at $[2, 3]$, $\omega_p = 0.8411$ and $\kappa = 1 / 20$: $2.345 \times 6{,}000 \times 0.05 / 0.595 = 1{,}183$; $\lambda_p = 0.210$ fm, $(0.210 / 2.13)^2 / 2 = 4.9 \times 10^{-3}$. \emph{Status:} one derivation under the hypothesis, the normalisation of the radiated share for a second; nothing run; clicks.
\end{derivation}

\subsection*{C.7 The masses as families}

\renewcommand{\thederivation}{C.\arabic{derivation}}\setcounter{derivation}{6}
\begin{derivation}[the masses as families, Section 4.1]
\emph{Lattice readings, not findings.} The rule's own universe at $[2, 3]$ is no finding against nature, which measures at its own pair (\fence{lattice}; Section 5.1). The masses are families. \emph{What the click's structure derives of the masses} (two derivations), from Rule3's vacuum dispersion $\cos\omega = (\num / (3\den)) \sum_a \cos k_a$ alone: the rest rotation $\cos\omega_0 = \num / \den$, the inertial mass from the band's curvature $m^* = 3\tan\omega_0$ and the family's own light speed $c_m^2 = \omega_0 / (3\tan\omega_0)$, the band's kinetic scale (Section 5.1; S.18, S.21, S.22; $0.2508$ at $[2, 3]$, $c_m / c = 0.867$), so $m^* c_m^2 = \omega_0$ exactly, the rest rotation, an identity of the dispersion, the share's energy per quantum $T\sin\omega_0$ a second reading of the same quantum, the two one as the gap closes; the mass ratio $\sin\omega_p / \sin\omega_e$, two declared rotations, and at one denominator $m^2$ proportional to the gap $\den - \num$ to first order; the ceiling $m_p / m_e \le 1 / \sin\omega_e$, above $1.4$ PeV at LHAASO's bound ($1.34$ in the rule's universe); at a conversion the rotations add at the Node with each above its rest, so the excess is kinetic by construction and one constraint on three wave numbers leaves the electron's spectrum continuous with its endpoint at the excess, at the level of the law's words; dimension $3$ adds no mass, one quantum being one unit of the invariant over all its lines. Not derived: the number $1{,}836$; nothing in the click, the draw or the dimension fixes a family's gap, and the equivalence between families holds as the gaps close: the free packet's fall coefficient $2\nu / (3(1 + \nu))$ (S.18) is $4 / 15$ at $[2, 3]$ and $2 / 7$ at $[3, 4]$, $2 / 105$ apart, while charge universality is exact under its convention. A cloud's click mass is its count at first order in $C / \Gamma$ (\claimmark{derived} from the functional; \fence{lattice}), and nature's mass ratios are not whole numbers, $m_p / m_e = 1836.153$, $m_\mu / m_e = 206.768$ and $m_\pi / m_e = 273.13$, so, with the cluster's defect bounded at $C\omega_0 x_{\mathrm{self}}$ (S.36 (d)), $10^{-22}$ at the write weight $1$, no particle of nature can be read as a cluster of the electron's quanta at that weight (\claimmark{derived} under the declared weight; \fence{clicks}). So the masses are the families' declared integers, and the proton is one quantum of its own family (\claimmark{assumption}); hydrogen's clicks read it as one unit of $W$ spread over $0.84$ fm, four Compton lengths of its family (\claimmark{derived}, one derivation); charge universality follows from the writes exactly, and the equivalence between families as the gaps close, the free packet's fall differing between families at finite gaps by the band's coefficient $2\nu / (3(1 + \nu))$ (\claimmark{derived}); if the families share one denominator, a declaration, the masses are the square roots of the gaps and the electron's gap is a declared integer and not $1$ (\claimmark{hypothesis}). The model's particle physics is families with declared masses and conversions by click, the conversion an explicitly named hypothesis.

\emph{Status:} the mass ratios' reading derived under the hypothesis of Section 9 and B.12 (one denominator); the families declared; clicks.
\end{derivation}


\subsection*{C.8 One photon on two bodies: the anticoincidence}

\renewcommand{\thederivation}{C.\arabic{derivation}}\setcounter{derivation}{7}
\begin{derivation}[one photon on two bodies: the anticoincidence, Section 2.5]
\emph{Inputs:} a light record $L$ of one whole quantum at the born unit $\Wc\sin\Omega$ (Section 2.5), laid between two bound bodies $A$ and $B$ at equal distances, each with the modes $g$ and $e$ and the transition $\omega_e - \omega_g = \Omega$, a NodeDetector on each body reading its $e$-mode, the whole the transition (Section 2.5); the share's conservation through the Ports (Theorem 3); the two-mode line's transfer at resonance (S.41); the credit's one draw per record over all the NodeDetectors that read it (Section 2.5). \emph{Steps:} (a) the inflow of $L$ into the body $X$'s region over the window is $\Phi_X$, and by the share's conservation at the vacuum's paces of the bodies' regions (the weights $1 / p_i^2$ entering in a well) $\Phi_A + \Phi_B \le \Wc\sin\Omega$, one whole quantum; the absorbed share is the two-mode line's transfer, $s_X = \eta\,\Phi_X / (\Wc\sin\Omega)$ with $\eta = \sin^2(\Omega_R\tau_w)$, $\Omega_R = |\epsilon_{ge}| / (4\sqrt{\sin\omega_g\sin\omega_e})$ and $\tau_w$ the passage, so by symmetry $s_A = s_B = s \le 1 / 2$. (b) The bodies' rule: a body's absorption (Section 2.5) weighs each transition by its own transfer share and the rest of the unit by the outcome that none absorbs, so for two bodies $P(A\ \text{only}) = s_A$, $P(B\ \text{only}) = s_B$, $P(\text{both}) = 0$ and $P(\text{neither}) = 1 - s_A - s_B$, at the symmetric $s$ the row $s$, $s$, $1 - 2s$; a body's marginal is its own share whatever the other's, while $s_A + s_B$ is at most the unit (a region with Nodes alone carries the other rule, the rounding of the total, Section 2.5); the anticoincidence parameter $\alpha = P(\text{both}) / (P(A)\,P(B))$ is $0$, where independent draws, classical light, give $\alpha = 1$. (c) The control of two quanta in one record: $P(\text{both}) = 2s^2$, $P(A\ \text{twice}) = P(B\ \text{twice}) = s^2$, and $\alpha$ depends on its definition, both forms the independent draw's algebra and derived: by the mean counts, $\langle N_A\rangle = 2s$, $\alpha = 2s^2 / (4s^2) = 1 / 2$; by the probabilities per trial, Grangier's, $P(A) = 1 - (1 - s)^2 = 2s - s^2$ and $\alpha = 2s^2 / (2s - s^2)^2 = 2 / (2 - s)^2$, the two meeting as $s \to 0$ and differing at $s = 1 / 2$ ($8 / 9$ against $1 / 2$); the absorbed fraction of S.41 stays open, and the run's reading names which definition it is read under; the control is two records of count $1$ absorbed by two bodies, two independent absorptions, $P(\text{both}) = s_A s_B$ and $\alpha = 1$ in expectation, the $1 / 2$ of the ordered pairs being the one record's and not the control's; three bodies: $\sum_X s_X \le 1$ and $P(\text{two or more}) = 0$ for one quantum. \emph{Check:} Grangier, Roger and Aspect's anticoincidence \cite{grangier} on a beam splitter, $\alpha = 0.18 \pm 0.06$ with a source-limited single photon, and modern single-photon sources below $0.01$. \emph{Status:} derived from the credit's declaration, as Born's rule is, two derivations; clicks. Nothing of the click's write enters it: the anticoincidence rests on the one draw under the law as it stands; its readings are no claim of this paper.
\end{derivation}

\subsection*{C.9 The quantum's unit and the hazard}

(a) \emph{The quantum of a field.} One unit of the configuration's invariant, $2A^2\sin\omega_s = T$: the share $T\sin\omega_s$, which a click reads, $\hbar\omega_s$ its small-gap name, the Wronskian $T / 2$ for a charged circular quantum ($0$ for a real configuration, Section 3.5 of the paper); what a counter clicks and what stands for a particle of nature, one photon, one electron, one muon, one proton.

(b) \emph{The counting unit.} A configuration of count $1$ is one quantum of the invariant, $2A^2\sin\omega = T$ for a plane configuration, $A$ the amplitude of each of its two lines, and $B^2\sin\omega = T$ for one real line (S.55 (b)); its form is $T\sin\omega$, the unit of energy; a bolometer region credits it at $W_c\sin\omega_0$ of inflow, the integer square root of $9T^2(\den^2 - \num^2)$, $3\,\den\,T$ for light; a configuration of a field with a gap is laid by $A^2$ the integer square root of $\text{count}^2 T^2 \den^2 / (4\,n_{\mathrm{laid}}^2 (\den^2 - \num^2))$, $n_{\mathrm{laid}}$ the number of the configuration's lines laid alike ($1$ for one part); over $n$ Nodes each Node's $A^2$ is that divided by $n$, its own root rounded down, that is $\text{count}\,T / (2\,n_{\mathrm{laid}}\sin\omega)$ to the unit, the root taken on the whole product and rounded down (count-unit), and a massless configuration's quantum is read at its wave's own frequency; not a particle.

(c) \emph{The hazard.} The rate $1 / \tau$ per proper interval at which a bound state in the dark emits, drawn at the now, which makes the emission's survival geometric in the grain, $(1 - h)^n$, $h$ the hazard per window, $n_w / \tau$ for a window of $n_w$ intervals and $n$ the windows, the exponential's discrete form; $\tau$ the law's number per transition, declared (Section 2.5 of the paper; S.50).

\begin{thebibliography}{99}\setlength{\itemsep}{0pt}\setlength{\parsep}{0pt}
\bibitem{watson} G. N. Watson, Three triple integrals, The Quarterly Journal of Mathematics os-10, 266 (1939), \href{https://doi.org/10.1093/qmath/os-10.1.266}{\nolinkurl{https://doi.org/10.1093/qmath/os-10.1.266}}.
\bibitem{will} C. M. Will, The confrontation between general relativity and experiment, Living Reviews in Relativity 17, 4 (2014), \href{https://doi.org/10.12942/lrr-2014-4}{\nolinkurl{https://doi.org/10.12942/lrr-2014-4}}.
\bibitem{cfl} R. Courant, K. Friedrichs and H. Lewy, \"Uber die partiellen Differenzengleichungen der mathematischen Physik, Mathematische Annalen 100, 32 (1928), \href{https://doi.org/10.1007/BF01448839}{\nolinkurl{https://doi.org/10.1007/BF01448839}}.
\bibitem{wolfram} S. Wolfram, A New Kind of Science (Wolfram Media, Champaign, 2002).
\bibitem{fredkin} E. Fredkin, Digital mechanics: an informational process based on reversible universal cellular automata, Physica D 45, 254 (1990), \href{https://doi.org/10.1016/0167-2789(90)90186-S}{\nolinkurl{https://doi.org/10.1016/0167-2789(90)90186-S}}.
\bibitem{arrighi} P. Arrighi, An overview of quantum cellular automata, Natural Computing 18, 885 (2019), \href{https://doi.org/10.1007/s11047-019-09762-6}{\nolinkurl{https://doi.org/10.1007/s11047-019-09762-6}}.
\bibitem{deraedt} H. De Raedt, K. De Raedt and K. Michielsen, Event-based simulation of single-photon beam splitters and Mach-Zehnder interferometers, Europhysics Letters 69, 861 (2005), \href{https://doi.org/10.1209/epl/i2004-10443-7}{\nolinkurl{https://doi.org/10.1209/epl/i2004-10443-7}}.
\bibitem{thooft} G. 't Hooft, The Cellular Automaton Interpretation of Quantum Mechanics, Fundamental Theories of Physics 185 (Springer, Cham, 2016), \href{https://doi.org/10.1007/978-3-319-41285-6}{\nolinkurl{https://doi.org/10.1007/978-3-319-41285-6}}.
\bibitem{wheeler} J. A. Wheeler and R. P. Feynman, Interaction with the absorber as the mechanism of radiation, Reviews of Modern Physics 17, 157 (1945), \href{https://doi.org/10.1103/RevModPhys.17.157}{\nolinkurl{https://doi.org/10.1103/RevModPhys.17.157}}.
\bibitem{cramer} J. G. Cramer, The transactional interpretation of quantum mechanics, Reviews of Modern Physics 58, 647 (1986), \href{https://doi.org/10.1103/RevModPhys.58.647}{\nolinkurl{https://doi.org/10.1103/RevModPhys.58.647}}.
\bibitem{vaidman} Y. Aharonov and L. Vaidman, The two-state vector formalism: an updated review, Lecture Notes in Physics 734, 399 (Springer, Berlin, 2008), \href{https://doi.org/10.1007/978-3-540-73473-4_13}{\nolinkurl{https://doi.org/10.1007/978-3-540-73473-4_13}}.
\bibitem{wharton} K. B. Wharton and N. Argaman, Colloquium: Bell's theorem and locally mediated reformulations of quantum mechanics, Reviews of Modern Physics 92, 021002 (2020), \href{https://doi.org/10.1103/RevModPhys.92.021002}{\nolinkurl{https://doi.org/10.1103/RevModPhys.92.021002}}.
\bibitem{yilmaz} H. Yilmaz, New approach to general relativity, Physical Review 111, 1417 (1958), \href{https://doi.org/10.1103/PhysRev.111.1417}{\nolinkurl{https://doi.org/10.1103/PhysRev.111.1417}}.
\bibitem{dicke} R. H. Dicke, Gravitation without a principle of equivalence, Reviews of Modern Physics 29, 363 (1957), \href{https://doi.org/10.1103/RevModPhys.29.363}{\nolinkurl{https://doi.org/10.1103/RevModPhys.29.363}}.
\bibitem{puthoff} H. E. Puthoff, Polarizable-vacuum (PV) approach to general relativity, Foundations of Physics 32, 927 (2002), \href{https://doi.org/10.1023/A:1016011413407}{\nolinkurl{https://doi.org/10.1023/A:1016011413407}}.
\bibitem{boonserm} P. Boonserm, T. Ngampitipan, A. Simpson and M. Visser, Exponential metric represents a traversable wormhole, Physical Review D 98, 084048 (2018), \href{https://doi.org/10.1103/PhysRevD.98.084048}{\nolinkurl{https://doi.org/10.1103/PhysRevD.98.084048}}.
\bibitem{wilson} K. G. Wilson, Confinement of quarks, Physical Review D 10, 2445 (1974), \href{https://doi.org/10.1103/PhysRevD.10.2445}{\nolinkurl{https://doi.org/10.1103/PhysRevD.10.2445}}.
\bibitem{sciarretta} A. Sciarretta, A local-realistic model of quantum mechanics based on a discrete spacetime, Foundations of Physics 48, 60 (2018), \href{https://doi.org/10.1007/s10701-017-0129-9}{\nolinkurl{https://doi.org/10.1007/s10701-017-0129-9}}.
\bibitem{bohm} D. Bohm, A suggested interpretation of the quantum theory in terms of ``hidden'' variables. I, Physical Review 85, 166 (1952), \href{https://doi.org/10.1103/PhysRev.85.166}{\nolinkurl{https://doi.org/10.1103/PhysRev.85.166}}.
\bibitem{philippidis} C. Philippidis, C. Dewdney and B. J. Hiley, Quantum interference and the quantum potential, Il Nuovo Cimento B 52, 15 (1979), \href{https://doi.org/10.1007/BF02743566}{\nolinkurl{https://doi.org/10.1007/BF02743566}}.
\bibitem{gleason} A. M. Gleason, Measures on the closed subspaces of a Hilbert space, Journal of Mathematics and Mechanics 6, 885 (1957), \href{https://doi.org/10.1512/iumj.1957.6.56050}{\nolinkurl{https://doi.org/10.1512/iumj.1957.6.56050}}.
\bibitem{deutsch} D. Deutsch, Quantum theory of probability and decisions, Proceedings of the Royal Society A 455, 3129 (1999), \href{https://doi.org/10.1098/rspa.1999.0443}{\nolinkurl{https://doi.org/10.1098/rspa.1999.0443}}.
\bibitem{zurek} W. H. Zurek, Probabilities from entanglement, Born's rule $p_k = |\psi_k|^2$ from envariance, Physical Review A 71, 052105 (2005), \href{https://doi.org/10.1103/PhysRevA.71.052105}{\nolinkurl{https://doi.org/10.1103/PhysRevA.71.052105}}.
\bibitem{meyer1997} D. A. Meyer, Quantum mechanics of lattice gas automata: one-particle plane waves and potentials, Physical Review E 55, 5261 (1997), \href{https://doi.org/10.1103/PhysRevE.55.5261}{\nolinkurl{https://doi.org/10.1103/PhysRevE.55.5261}}.
\bibitem{eht_sgra} The Event Horizon Telescope Collaboration, First Sagittarius A* Event Horizon Telescope results. VI. Testing the black hole metric, The Astrophysical Journal Letters 930, L17 (2022), \href{https://doi.org/10.3847/2041-8213/ac6756}{\nolinkurl{https://doi.org/10.3847/2041-8213/ac6756}}.
\bibitem{psaltis} D. Psaltis et al. (The Event Horizon Telescope Collaboration), Gravitational test beyond the first post-Newtonian order with the shadow of the M87 black hole, Physical Review Letters 125, 141104 (2020), \href{https://doi.org/10.1103/PhysRevLett.125.141104}{\nolinkurl{https://doi.org/10.1103/PhysRevLett.125.141104}}.
\bibitem{tolman} R. C. Tolman, P. Ehrenfest and B. Podolsky, On the gravitational field produced by light, Physical Review 37, 602 (1931), \href{https://doi.org/10.1103/PhysRev.37.602}{\nolinkurl{https://doi.org/10.1103/PhysRev.37.602}}.
\bibitem{zych} M. Zych, F. Costa, I. Pikovski and C. Brukner, Quantum interferometric visibility as a witness of general relativistic proper time, Nature Communications 2, 505 (2011), \href{https://doi.org/10.1038/ncomms1498}{\nolinkurl{https://doi.org/10.1038/ncomms1498}}.
\bibitem{born} M. Born, Zur Quantenmechanik der Sto\ss{}vorg\"ange, Zeitschrift f\"ur Physik 37, 863 (1926), \href{https://doi.org/10.1007/BF01397477}{\nolinkurl{https://doi.org/10.1007/BF01397477}}.
\bibitem{cirelson} B. S. Cirel'son, Quantum generalizations of Bell's inequality, Letters in Mathematical Physics 4, 93 (1980), \href{https://doi.org/10.1007/BF00417500}{\nolinkurl{https://doi.org/10.1007/BF00417500}}.
\bibitem{pearle} P. M. Pearle, Hidden-variable example based upon data rejection, Physical Review D 2, 1418 (1970), \href{https://doi.org/10.1103/PhysRevD.2.1418}{\nolinkurl{https://doi.org/10.1103/PhysRevD.2.1418}}.
\bibitem{eberhard} P. H. Eberhard, Background level and counter efficiencies required for a loophole-free Einstein-Podolsky-Rosen experiment, Physical Review A 47, R747 (1993), \href{https://doi.org/10.1103/PhysRevA.47.R747}{\nolinkurl{https://doi.org/10.1103/PhysRevA.47.R747}}.
\bibitem{gargmermin} A. Garg and N. D. Mermin, Detector inefficiencies in the Einstein-Podolsky-Rosen experiment, Physical Review D 35, 3831 (1987), \href{https://doi.org/10.1103/PhysRevD.35.3831}{\nolinkurl{https://doi.org/10.1103/PhysRevD.35.3831}}.
\bibitem{deraedt2007} H. De Raedt, K. De Raedt, K. Michielsen, K. Keimpema and S. Miyashita, Event-based computer simulation model of Aspect-type experiments strictly satisfying Einstein's locality conditions, Journal of the Physical Society of Japan 76, 104005 (2007), \href{https://doi.org/10.1143/JPSJ.76.104005}{\nolinkurl{https://doi.org/10.1143/JPSJ.76.104005}}.
\bibitem{michielsen2014} K. Michielsen and H. De Raedt, Event-based simulation of quantum physics experiments, International Journal of Modern Physics C 25, 1430003 (2014), \href{https://doi.org/10.1142/S0129183114300036}{\nolinkurl{https://doi.org/10.1142/S0129183114300036}}.
\bibitem{hensen} B. Hensen et al., Loophole-free Bell inequality violation using electron spins separated by 1.3 kilometres, Nature 526, 682 (2015), \href{https://doi.org/10.1038/nature15759}{\nolinkurl{https://doi.org/10.1038/nature15759}}.
\bibitem{giustina} M. Giustina et al., Significant-loophole-free test of Bell's theorem with entangled photons, Physical Review Letters 115, 250401 (2015), \href{https://doi.org/10.1103/PhysRevLett.115.250401}{\nolinkurl{https://doi.org/10.1103/PhysRevLett.115.250401}}.
\bibitem{shalm} L. K. Shalm et al., Strong loophole-free test of local realism, Physical Review Letters 115, 250402 (2015), \href{https://doi.org/10.1103/PhysRevLett.115.250402}{\nolinkurl{https://doi.org/10.1103/PhysRevLett.115.250402}}.
\bibitem{hossenfelder} S. Hossenfelder and T. Palmer, Rethinking superdeterminism, Frontiers in Physics 8, 139 (2020), \href{https://doi.org/10.3389/fphy.2020.00139}{\nolinkurl{https://doi.org/10.3389/fphy.2020.00139}}.
\bibitem{gw170814} B. P. Abbott et al. (LIGO Scientific Collaboration and Virgo Collaboration), GW170814: a three-detector observation of gravitational waves from a binary black hole coalescence, Physical Review Letters 119, 141101 (2017), \href{https://doi.org/10.1103/PhysRevLett.119.141101}{\nolinkurl{https://doi.org/10.1103/PhysRevLett.119.141101}}.
\bibitem{riess} A. G. Riess et al., Observational evidence from supernovae for an accelerating universe and a cosmological constant, The Astronomical Journal 116, 1009 (1998), \href{https://doi.org/10.1086/300499}{\nolinkurl{https://doi.org/10.1086/300499}}.
\bibitem{perlmutter} S. Perlmutter et al. (The Supernova Cosmology Project), Measurements of $\Omega$ and $\Lambda$ from 42 high-redshift supernovae, The Astrophysical Journal 517, 565 (1999), \href{https://doi.org/10.1086/307221}{\nolinkurl{https://doi.org/10.1086/307221}}.
\bibitem{pantheon} D. Brout et al., The Pantheon+ analysis: cosmological constraints, The Astrophysical Journal 938, 110 (2022), \href{https://doi.org/10.3847/1538-4357/ac8e04}{\nolinkurl{https://doi.org/10.3847/1538-4357/ac8e04}}.
\bibitem{pdg} S. Navas et al. (Particle Data Group), Review of particle physics, Physical Review D 110, 030001 (2024), \href{https://doi.org/10.1103/PhysRevD.110.030001}{\nolinkurl{https://doi.org/10.1103/PhysRevD.110.030001}}.
\bibitem{hasert} F. J. Hasert et al., Observation of neutrino-like interactions without muon or electron in the Gargamelle neutrino experiment, Physics Letters B 46, 138 (1973), \href{https://doi.org/10.1016/0370-2693(73)90499-1}{\nolinkurl{https://doi.org/10.1016/0370-2693(73)90499-1}}.
\bibitem{lep} The ALEPH, DELPHI, L3, OPAL and SLD Collaborations, Precision electroweak measurements on the Z resonance, Physics Reports 427, 257 (2006), \href{https://doi.org/10.1016/j.physrep.2005.12.006}{\nolinkurl{https://doi.org/10.1016/j.physrep.2005.12.006}}.
\bibitem{wu} C. S. Wu, E. Ambler, R. W. Hayward, D. D. Hoppes and R. P. Hudson, Experimental test of parity conservation in beta decay, Physical Review 105, 1413 (1957), \href{https://doi.org/10.1103/PhysRev.105.1413}{\nolinkurl{https://doi.org/10.1103/PhysRev.105.1413}}.
\bibitem{sommerfeld} A. Sommerfeld, Zur Quantentheorie der Spektrallinien, Annalen der Physik 356, 1 (1916), \href{https://doi.org/10.1002/andp.19163561702}{\nolinkurl{https://doi.org/10.1002/andp.19163561702}}.
\bibitem{roothaan} C. C. J. Roothaan, L. M. Sachs and A. W. Weiss, Analytical self-consistent field functions for the atomic configurations $1s^2$, $1s^2 2s$ and $1s^2 2s^2$, Reviews of Modern Physics 32, 186 (1960), \href{https://doi.org/10.1103/RevModPhys.32.186}{\nolinkurl{https://doi.org/10.1103/RevModPhys.32.186}}.
\bibitem{codata} P. J. Mohr, D. B. Newell, B. N. Taylor and E. Tiesinga, CODATA recommended values of the fundamental physical constants: 2022, Reviews of Modern Physics 97, 025002 (2025), \href{https://doi.org/10.1103/RevModPhys.97.025002}{\nolinkurl{https://doi.org/10.1103/RevModPhys.97.025002}}.
\bibitem{rowe} M. A. Rowe, D. Kielpinski, V. Meyer, C. A. Sackett, W. M. Itano, C. Monroe and D. J. Wineland, Experimental violation of a Bell's inequality with efficient detection, Nature 409, 791 (2001), \href{https://doi.org/10.1038/35057215}{\nolinkurl{https://doi.org/10.1038/35057215}}.
\bibitem{storz} S. Storz et al., Loophole-free Bell inequality violation with superconducting circuits, Nature 617, 265 (2023), \href{https://doi.org/10.1038/s41586-023-05885-0}{\nolinkurl{https://doi.org/10.1038/s41586-023-05885-0}}.
\bibitem{vasileiou} V. Vasileiou et al., Constraints on Lorentz invariance violation from Fermi-Large Area Telescope observations of gamma-ray bursts, Physical Review D 87, 122001 (2013), \href{https://doi.org/10.1103/PhysRevD.87.122001}{\nolinkurl{https://doi.org/10.1103/PhysRevD.87.122001}}.
\bibitem{mossbauer} R. L. M\"ossbauer, Kernresonanzfluoreszenz von Gammastrahlung in Ir$^{191}$, Zeitschrift f\"ur Physik 151, 124 (1958), \href{https://doi.org/10.1007/BF01344210}{\nolinkurl{https://doi.org/10.1007/BF01344210}}.
\bibitem{cook} R. J. Cook and H. J. Kimble, Possibility of direct observation of quantum jumps, Physical Review Letters 54, 1023 (1985), \href{https://doi.org/10.1103/PhysRevLett.54.1023}{\nolinkurl{https://doi.org/10.1103/PhysRevLett.54.1023}}.
\bibitem{bergquist} J. C. Bergquist, R. G. Hulet, W. M. Itano and D. J. Wineland, Observation of quantum jumps in a single atom, Physical Review Letters 57, 1699 (1986), \href{https://doi.org/10.1103/PhysRevLett.57.1699}{\nolinkurl{https://doi.org/10.1103/PhysRevLett.57.1699}}.
\bibitem{gw170817} B. P. Abbott et al. (LIGO Scientific Collaboration, Virgo Collaboration, Fermi Gamma-ray Burst Monitor and INTEGRAL), Gravitational waves and gamma-rays from a binary neutron star merger: GW170817 and GRB 170817A, The Astrophysical Journal Letters 848, L13 (2017), \href{https://doi.org/10.3847/2041-8213/aa920c}{\nolinkurl{https://doi.org/10.3847/2041-8213/aa920c}}.
\bibitem{pekar} S. I. Pekar, Untersuchungen \"uber die Elektronentheorie der Kristalle (Akademie-Verlag, Berlin, 1954).
\bibitem{lieb} E. H. Lieb, Existence and uniqueness of the minimizing solution of Choquard's nonlinear equation, Studies in Applied Mathematics 57, 93 (1977), \href{https://doi.org/10.1002/sapm197757293}{\nolinkurl{https://doi.org/10.1002/sapm197757293}}.
\bibitem{magic} V. A. Acciari et al. (MAGIC Collaboration), Bounds on Lorentz invariance violation from MAGIC observation of GRB 190114C, Physical Review Letters 125, 021301 (2020), \href{https://doi.org/10.1103/PhysRevLett.125.021301}{\nolinkurl{https://doi.org/10.1103/PhysRevLett.125.021301}}.
\bibitem{lhaaso} Z. Cao et al. (LHAASO Collaboration), Ultrahigh-energy photons up to 1.4 petaelectronvolts from 12 gamma-ray Galactic sources, Nature 594, 33 (2021), \href{https://doi.org/10.1038/s41586-021-03498-z}{\nolinkurl{https://doi.org/10.1038/s41586-021-03498-z}}.
\bibitem{lhaaso2024} Z. Cao et al. (LHAASO Collaboration), An ultrahigh-energy gamma-ray bubble powered by a super PeVatron, Science Bulletin 69, 449 (2024), \href{https://doi.org/10.1016/j.scib.2023.12.040}{\nolinkurl{https://doi.org/10.1016/j.scib.2023.12.040}}.
\bibitem{mccrea} W. H. McCrea and E. A. Milne, Newtonian universes and the curvature of space, The Quarterly Journal of Mathematics os-5, 73 (1934), \href{https://doi.org/10.1093/qmath/os-5.1.73}{\nolinkurl{https://doi.org/10.1093/qmath/os-5.1.73}}.
\bibitem{everitt} C. W. F. Everitt et al., Gravity Probe B: final results of a space experiment to test general relativity, Physical Review Letters 106, 221101 (2011), \href{https://doi.org/10.1103/PhysRevLett.106.221101}{\nolinkurl{https://doi.org/10.1103/PhysRevLett.106.221101}}.
\bibitem{shao} L. Shao and N. Wex, New tests of local Lorentz invariance of gravity with small-eccentricity binary pulsars, Classical and Quantum Gravity 29, 215018 (2012), \href{https://doi.org/10.1088/0264-9381/29/21/215018}{\nolinkurl{https://doi.org/10.1088/0264-9381/29/21/215018}}.
\bibitem{muller} J. M\"uller, K. Nordtvedt and D. Vokrouhlick\'y, Improved constraint on the $\alpha_1$ PPN parameter from lunar motion, Physical Review D 54, R5927 (1996), \href{https://doi.org/10.1103/PhysRevD.54.R5927}{\nolinkurl{https://doi.org/10.1103/PhysRevD.54.R5927}}.
\bibitem{nordtvedt} K. Nordtvedt, Probing gravity to the second post-Newtonian order and to one part in $10^7$ using the spin axis of the Sun, The Astrophysical Journal 320, 871 (1987), \href{https://doi.org/10.1086/165603}{\nolinkurl{https://doi.org/10.1086/165603}}.
\bibitem{shao2013} L. Shao, R. N. Caballero, M. Kramer, N. Wex, D. J. Champion and A. Jessner, A new limit on local Lorentz invariance violation of gravity from solitary pulsars, Classical and Quantum Gravity 30, 165019 (2013), \href{https://doi.org/10.1088/0264-9381/30/16/165019}{\nolinkurl{https://doi.org/10.1088/0264-9381/30/16/165019}}.
\bibitem{ame2020} M. Wang, W. J. Huang, F. G. Kondev, G. Audi and S. Naimi, The AME 2020 atomic mass evaluation (II): tables, graphs and references, Chinese Physics C 45, 030003 (2021), \href{https://doi.org/10.1088/1674-1137/abddaf}{\nolinkurl{https://doi.org/10.1088/1674-1137/abddaf}}.
\bibitem{rogers1970} F. J. Rogers, H. C. Graboske, Jr., and D. J. Harwood, Bound eigenstates of the static screened Coulomb potential, Physical Review A 1, 1577 (1970), \href{https://doi.org/10.1103/PhysRevA.1.1577}{\nolinkurl{https://doi.org/10.1103/PhysRevA.1.1577}}.
\bibitem{mermin} N. D. Mermin, Extreme quantum entanglement in a superposition of macroscopically distinct states, Physical Review Letters 65, 1838 (1990), \href{https://doi.org/10.1103/PhysRevLett.65.1838}{\nolinkurl{https://doi.org/10.1103/PhysRevLett.65.1838}}.
\bibitem{sr_clock} International Bureau of Weights and Measures (BIPM), Recommended values of standard frequencies for applications including the practical realization of the metre and secondary representations of the definition of the second: strontium 87 atom ($f \approx 429$ THz), $429\,228\,004\,229\,872.992$ Hz with the relative standard uncertainty $1.7 \times 10^{-16}$, approved by the CCTF in September 2025, \href{https://www.bipm.org/en/publications/mises-en-pratique/standard-frequencies}{\nolinkurl{https://www.bipm.org/en/publications/mises-en-pratique/standard-frequencies}}.
\bibitem{hawking1974} S. W. Hawking, Black hole explosions?, Nature 248, 30 (1974), \href{https://doi.org/10.1038/248030a0}{\nolinkurl{https://doi.org/10.1038/248030a0}}.
\bibitem{hawking1975} S. W. Hawking, Particle creation by black holes, Communications in Mathematical Physics 43, 199 (1975), \href{https://doi.org/10.1007/BF02345020}{\nolinkurl{https://doi.org/10.1007/BF02345020}}.
\bibitem{bekenstein} J. D. Bekenstein, Black holes and entropy, Physical Review D 7, 2333 (1973), \href{https://doi.org/10.1103/PhysRevD.7.2333}{\nolinkurl{https://doi.org/10.1103/PhysRevD.7.2333}}.
\bibitem{bezginov} N. Bezginov, T. Valdez, M. Horbatsch, A. Marsman, A. C. Vutha and E. A. Hessels, A measurement of the atomic hydrogen Lamb shift and the proton charge radius, Science 365, 1007 (2019), \href{https://doi.org/10.1126/science.aau7807}{\nolinkurl{https://doi.org/10.1126/science.aau7807}}.
\bibitem{antognini} A. Antognini et al., Proton structure from the measurement of 2S-2P transition frequencies of muonic hydrogen, Science 339, 417 (2013), \href{https://doi.org/10.1126/science.1230016}{\nolinkurl{https://doi.org/10.1126/science.1230016}}.
\bibitem{schlamminger} S. Schlamminger, K.-Y. Choi, T. A. Wagner, J. H. Gundlach and E. G. Adelberger, Test of the equivalence principle using a rotating torsion balance, Physical Review Letters 100, 041101 (2008), \href{https://doi.org/10.1103/PhysRevLett.100.041101}{\nolinkurl{https://doi.org/10.1103/PhysRevLett.100.041101}}.
\bibitem{microscope} P. Touboul et al. (MICROSCOPE Collaboration), MICROSCOPE mission: final results of the test of the equivalence principle, Physical Review Letters 129, 121102 (2022), \href{https://doi.org/10.1103/PhysRevLett.129.121102}{\nolinkurl{https://doi.org/10.1103/PhysRevLett.129.121102}}.
\bibitem{witteborn} F. C. Witteborn and W. M. Fairbank, Experimental comparison of the gravitational force on freely falling electrons and metallic electrons, Physical Review Letters 19, 1049 (1967), \href{https://doi.org/10.1103/PhysRevLett.19.1049}{\nolinkurl{https://doi.org/10.1103/PhysRevLett.19.1049}}.
\bibitem{hull_dobell} T. E. Hull and A. R. Dobell, Random number generators, SIAM Review 4, 230 (1962), \href{https://doi.org/10.1137/1004061}{\nolinkurl{https://doi.org/10.1137/1004061}}.
\bibitem{bauch} A. Bauch and S. Weyers, New experimental limit on the validity of local position invariance, Physical Review D 65, 081101 (2002), \href{https://doi.org/10.1103/PhysRevD.65.081101}{\nolinkurl{https://doi.org/10.1103/PhysRevD.65.081101}}.
\bibitem{turneaure} J. P. Turneaure, C. M. Will, B. F. Farrell, E. M. Mattison and R. F. C. Vessot, Test of the principle of equivalence by a null gravitational red-shift experiment, Physical Review D 27, 1705 (1983), \href{https://doi.org/10.1103/PhysRevD.27.1705}{\nolinkurl{https://doi.org/10.1103/PhysRevD.27.1705}}.
\bibitem{williams} J. G. Williams, S. G. Turyshev and D. H. Boggs, Progress in lunar laser ranging tests of relativistic gravity, Physical Review Letters 93, 261101 (2004), \href{https://doi.org/10.1103/PhysRevLett.93.261101}{\nolinkurl{https://doi.org/10.1103/PhysRevLett.93.261101}}.
\bibitem{parthey2010} C. G. Parthey, A. Matveev, J. Alnis, R. Pohl, T. Udem, U. D. Jentschura, N. Kolachevsky and T. W. H\"ansch, Precision measurement of the hydrogen--deuterium $1S$--$2S$ isotope shift, Physical Review Letters 104, 233001 (2010), \href{https://doi.org/10.1103/PhysRevLett.104.233001}{\nolinkurl{https://doi.org/10.1103/PhysRevLett.104.233001}}.
\bibitem{pekeris} C. L. Pekeris, Ground state of two-electron atoms, Physical Review 112, 1649 (1958), \href{https://doi.org/10.1103/PhysRev.112.1649}{\nolinkurl{https://doi.org/10.1103/PhysRev.112.1649}}.
\bibitem{parthey} C. G. Parthey et al., Improved measurement of the hydrogen 1S--2S transition frequency, Physical Review Letters 107, 203001 (2011), \href{https://doi.org/10.1103/PhysRevLett.107.203001}{\nolinkurl{https://doi.org/10.1103/PhysRevLett.107.203001}}.
\bibitem{nist} A. Kramida, Yu. Ralchenko, J. Reader and the NIST ASD Team, NIST Atomic Spectra Database, version 5.11, National Institute of Standards and Technology, Gaithersburg (2023), \url{https://physics.nist.gov/asd}.
\bibitem{fee} M. S. Fee, A. P. Mills, Jr., S. Chu, E. D. Shaw, K. Danzmann, R. J. Chichester and D. M. Zuckerman, Measurement of the positronium $1^3S_1$--$2^3S_1$ interval by continuous-wave two-photon excitation, Physical Review Letters 70, 1397 (1993), \href{https://doi.org/10.1103/PhysRevLett.70.1397}{\nolinkurl{https://doi.org/10.1103/PhysRevLett.70.1397}}.
\bibitem{fermi_lat} A. A. Abdo et al. (Fermi LAT and Fermi GBM Collaborations), A limit on the variation of the speed of light arising from quantum gravity effects, Nature 462, 331 (2009), \href{https://doi.org/10.1038/nature08574}{\nolinkurl{https://doi.org/10.1038/nature08574}}.
\bibitem{pound_snider} R. V. Pound and J. L. Snider, Effect of gravity on gamma radiation, Physical Review 140, B788 (1965), \href{https://doi.org/10.1103/PhysRev.140.B788}{doi:\nolinkurl{10.1103/PhysRev.140.B788}}.
\bibitem{vessot} R. F. C. Vessot et al., Test of relativistic gravitation with a space-borne hydrogen maser, Physical Review Letters 45, 2081 (1980), \href{https://doi.org/10.1103/PhysRevLett.45.2081}{doi:\nolinkurl{10.1103/PhysRevLett.45.2081}}.
\bibitem{bethe_peierls} H. A. Bethe and R. Peierls, Quantum theory of the diplon, Proceedings of the Royal Society A 148, 146 (1935), \href{https://doi.org/10.1098/rspa.1935.0010}{doi:\nolinkurl{10.1098/rspa.1935.0010}}.
\bibitem{grangier} P. Grangier, G. Roger and A. Aspect, Experimental evidence for a photon anticorrelation effect on a beam splitter: a new light on single-photon interferences, Europhysics Letters 1, 173 (1986), \href{https://doi.org/10.1209/0295-5075/1/4/004}{doi:\nolinkurl{10.1209/0295-5075/1/4/004}}.
\end{thebibliography}

\end{document}
